\documentclass[11pt,letterpaper,twoside,openright]{book}
\usepackage[letterpaper,top=1in,bottom=1in,inner=1.25in,outer=1in]{geometry}
\usepackage[T1]{fontenc}
\usepackage{lmodern}
\usepackage{microtype}
\usepackage{booktabs}
\usepackage{longtable}
\usepackage{array}
\usepackage{enumitem}
\usepackage{fancyhdr}
\usepackage{makeidx}
\usepackage{xcolor}
\usepackage{listings}
\usepackage{upquote}
\usepackage[hidelinks]{hyperref}

% ---- House macros (frozen; chapters use only these) ----
% Inline code: a Dylan name, form or value.  The argument is typeset in
% typewriter type; escape TeX's specials in it (\#, \_, \{, \}, \%, \&,
% \textbackslash{}, \textasciitilde{}, \textasciicircum{}).
\newcommand{\code}[1]{\texttt{#1}}
% A file or directory name, or a command typed at the shell.
\newcommand{\file}[1]{\texttt{#1}}
\newcommand{\cmd}[1]{\texttt{#1}}
% A term defined where it appears, and indexed.
\newcommand{\term}[1]{\emph{#1}\index{#1}}
% A metavariable in a syntax description.
\newcommand{\meta}[1]{\textit{#1}}
% A reference to a chapter or section of the 1992 manual.
\newcommand{\dylanbook}[1]{\textsl{DIRM}~#1}
% The result of evaluating an example, in running text.
\newcommand{\yields}{\,$\Rightarrow$\,}

\newenvironment{tipbox}{\par\medskip\noindent\textbf{Tip:}\ }{\par\medskip}
\newenvironment{warnbox}{\par\medskip\noindent\textbf{Important:}\ }{\par\medskip}
\newenvironment{notebox}{\par\medskip\noindent\textbf{Note:}\ }{\par\medskip}

% Code blocks.  dylan: Dylan source; results are shown as comments,
% ; => value.  shell: a terminal session, $ before what is typed.
% repl: a session at the Common Lisp implementation's ? prompt.
% ccode: C source.  lisp: Common Lisp source.
\lstset{basicstyle=\ttfamily\small,columns=fullflexible,keepspaces=true,
  showstringspaces=false,upquote=true,xleftmargin=1.5em,
  aboveskip=0.6\baselineskip,belowskip=0.6\baselineskip,
  commentstyle=\color{black!60},literate={~}{{\textasciitilde}}1}
\lstdefinelanguage{Dylan}{morecomment=[l]{;},morestring=[b]",
  sensitive=false}
\lstnewenvironment{dylan}{\lstset{language=Dylan}}{}
\lstnewenvironment{shell}{\lstset{}}{}
\lstnewenvironment{repl}{\lstset{}}{}
\lstnewenvironment{ccode}{\lstset{language=C}}{}
\lstnewenvironment{lisp}{\lstset{language=Lisp}}{}

\pagestyle{fancy}
\fancyhf{}
\fancyhead[LE]{\slshape\nouppercase{\leftmark}}
\fancyhead[RO]{\slshape\nouppercase{\rightmark}}
\fancyfoot[LE,RO]{\thepage}
\renewcommand{\headrulewidth}{0.4pt}
\setlength{\headheight}{14pt}

\makeindex

\title{\Huge Dylan\\[0.5em]\LARGE The Language, the Interpreter and the Compiler}
\author{Blake McBride}
% Manual revision date: when the manual itself was last revised and
% rebuilt, not a release of the software.  It appears identically on the
% title page and the copyright page.
\newcommand{\revisiondate}{\today}
\newcommand{\revisionline}{Manual revised \revisiondate}

\begin{document}
\frontmatter

% ---- Title page ----
\thispagestyle{empty}
\vspace*{2in}
\begin{center}
  {\Huge Dylan\par}
  \vspace{1em}
  {\LARGE The Language, the Interpreter\par}
  \vspace{0.3em}
  {\LARGE and the Compiler\par}
  \vspace{1.5em}
  {\large A User Manual\par}
\end{center}
\vfill
\begin{center}
  {\large Blake McBride\par}
  \vspace{0.5em}
  {\large \revisionline\par}
\end{center}
\vspace{0.5in}
\clearpage

% ---- Copyright page ----
\thispagestyle{empty}
\vspace*{4in}
\noindent This document was typeset with \LaTeX{} on Linux.\par
\bigskip
\noindent Copyright \copyright{} 2026 Blake McBride\par
\noindent All rights reserved.\par
\noindent \url{https://github.com/blakemcbride/Dylan}\par
\noindent \revisionline\par
\bigskip
\noindent The language this manual describes is the one defined in
\emph{Dylan(TM): an object-oriented dynamic language}, Apple Computer
Eastern Research and Technology, April 1992, cited here as
\emph{DIRM}. Dylan is a trademark of Apple Computer, Inc.\par
\clearpage

\tableofcontents

\mainmatter
\part{Getting Started}
\chapter{Introduction}\label{ch:intro}

This book is a user manual for Dylan, the object-oriented dynamic language
that Apple Computer described in 1992, and for the two implementations of
it that this project provides: an interactive implementation written in
Common Lisp and a compiler that translates Dylan programs into C. It
teaches the language from the beginning, explains how to use each
implementation, and documents where the implementations depart from the
1992 definition.

This chapter describes the language in broad terms, introduces the two
implementations and when to choose each, gives a short history of the
project, and explains how the rest of the book is organized and what its
typographic conventions mean.

\section{What Dylan is}\label{sec:intro-what}

\index{Dylan!1992 definition}
Dylan was designed at Apple Computer's Eastern Research and Technology
laboratory and published in April 1992 as \emph{Dylan: an object-oriented
dynamic language}. That book, which this manual calls the \term{DIRM}, is
the standard both implementations follow. The language it defines is the
original Dylan, written in a fully parenthesized \term{prefix syntax} like
that of Scheme and Common Lisp. Later versions of Dylan moved to an infix
syntax that looks more like Pascal or C; that later language is not the
subject of this book. Everything here is the 1992 prefix language.

A short piece of Dylan shows its flavor:

\begin{dylan}
(define-class <point> (<object>)
  (x init-value: 0 init-keyword: x:)
  (y init-value: 0 init-keyword: y:))

(define-method norm ((p <point>))
  (sqrt (+ (* (x p) (x p)) (* (y p) (y p)))))

(norm (make <point> x: 3 y: 4))   ; => 5.0
\end{dylan}

The language's main ideas are few, and they fit together closely.

\paragraph{Objects all the way down.} Every value in Dylan is an object,
and every object is an instance of a class. Integers, characters,
strings, lists, functions and classes themselves are all objects, with
classes such as \code{<integer>}, \code{<string>} and \code{<class>}.
There are no ``primitive'' values that stand outside the object system,
so any value can be passed to any function, stored in any collection, and
used to select a method.

\paragraph{Generic functions and multiple dispatch.}\index{generic function}
In most object-oriented languages a method belongs to a class and is
chosen by the class of a single receiver object. In Dylan, methods belong
to \term{generic functions}. A generic function is an ordinary function
that holds a set of methods; when it is called, it looks at the classes of
\emph{all} of its required arguments and runs the most specific method that
applies. This is called \term{multiple dispatch}. A method may also be
specialized on one particular object (a \term{singleton}) rather than on a
class.

\paragraph{Classes with multiple inheritance.} A class names one or more
superclasses and describes \term{slots}, the named fields each instance
holds. Dylan defines the getter and setter functions for each slot
automatically, and those functions are themselves generic functions, so
slot access and method calls look exactly alike. A class may inherit from
several superclasses at once.

\paragraph{First-class functions.} Functions are objects. They can be
created anonymously with \code{method}, passed as arguments, returned as
results, and stored in data structures. A function created inside another
function remembers the variables around it (it is a \term{closure}). The
library includes a set of functional operations, such as \code{curry},
\code{compose} and \code{complement}, that build new functions from old
ones.

\paragraph{Conditions.} Errors and other exceptional situations are
represented by \term{condition} objects, which are instances of classes
such as \code{<error>} and \code{<warning>}. A program signals a condition
and a dynamically established handler decides what to do with it: return
a value, take a non-local exit, or decline and let an outer handler try.
Because conditions are objects, handlers select them by class, just as
methods do.

\paragraph{Collections.} Dylan has a rich, uniform library of collections:
lists, vectors, strings, stretchy vectors, deques, ranges, arrays and
tables. All of them share one iteration protocol, so the same operations
(\code{map}, \code{reduce}, \code{do}, \code{choose}, \code{member?},
\code{for-each} and the rest) work on every kind of collection, and new
collection classes can join in by implementing the protocol.

\paragraph{A dynamic language that can be compiled well.} Dylan programs
need no type declarations, and in the 1992 definition classes and methods
can be added and changed while a program runs. The language also offers
ways to say that parts of a program will not change, such as sealing a
class, so that a compiler can produce fast code. The two implementations
described here sit at the two ends of this range.

\section{The two implementations}\label{sec:intro-implementations}

\index{implementations!choosing}
This project provides two separate implementations of the same language.
They share the reader (the code that turns program text into data), and
they are tested against the same examples from the DIRM, but otherwise
they are independent programs.

\paragraph{The interactive implementation.}\index{interactive implementation}
The first is written in Common Lisp and runs inside SBCL. Its source is in
\file{src/}, and this book calls it \emph{the interactive
implementation} or \emph{the Common Lisp implementation}. It is not an
interpreter: it compiles each Dylan form into Common Lisp as the form is
entered and lets the Lisp compiler run it. You use it in three ways: at an
interactive read-eval-print loop, the command \cmd{dylan}; by
compiling a file of Dylan into a file of Common Lisp with
\code{dylan:compile-file}; and by loading such compiled files into a Lisp
that runs them. It implements the whole language of the DIRM, including
the parts that let a program change its own classes and generic functions
as it runs (\code{add-method}, \code{remove-method}, \code{add-slot} and
their relatives), and it has arbitrarily large integers, because Common
Lisp does.

\paragraph{The compiler to C.}\index{dylanc}
The second is a compiler, \cmd{bin/dylanc}, that translates a complete
Dylan program into C and then runs a C compiler to produce a native
executable. Its source is in \file{compiler/}, and its run-time support is
in \file{runtime/}: a small runtime written in C and the Dylan library,
itself written in Dylan. The compiled programs have proper tail calls,
generic dispatch compiled to decision trees, type inference, inlining,
unboxed floating-point arithmetic, a generational garbage collector,
and real operating-system threads. The system also adds two things the
1992 language does not have: macros, in both implementations, and
threads, in the compiler.

The price of this speed is that the compiler treats a program as
\emph{static}. Everything the program defines must be visible when it is
compiled, and the program cannot change its classes or generic functions
while it runs, cannot evaluate or load new code, and has no bignums or
complex numbers: its integers are 63 bits wide, and overflow is an error.
Chapter~\ref{ch:dylanc} and Appendix~\ref{app:differences} describe these
restrictions in full.

\paragraph{Which to use.} The two are meant to be used together.

\begin{itemize}
\item Use the \textbf{interactive implementation} to learn the language, to
  try out ideas, to explore the library, and to develop code a piece at a
  time. You type a form and see its value at once. It is also the right
  choice for programs that really do need to redefine classes or add
  methods while they run, or that need integers larger than 63 bits.
\item Use the \textbf{compiler} when you want a standalone program: a fast
  executable that starts instantly, needs no Lisp to run, and can use
  threads. A program developed interactively can normally be compiled
  unchanged, provided it stays within the static-program restrictions.
\end{itemize}

Throughout this book, examples work in both implementations unless the
text says otherwise. Where the two differ, for instance in how they print
an object that has no readable printed form, the text points it out.

\section{A short history}\label{sec:intro-history}

\index{Thomas}
In 1992, shortly after the DIRM appeared, Matt Birkholz, Jim Miller and
Ron Weiss of Digital Equipment Corporation's Cambridge Research Laboratory
wrote \term{Thomas}, an implementation of the language in Scheme. They
wrote it in four weeks, and their stated goals were to learn the language
by implementing it from the book, to help others learn from readable
source, and to run small programs; performance was explicitly not a goal.
They named it Thomas rather than Dylan because they had not received
permission to use the trademark.

This project began as a direct translation of Thomas from Scheme into
Common Lisp. Since then it has moved a long way from that starting point.
The Common Lisp implementation was checked against the DIRM chapter by
chapter and corrected wherever Thomas and the book disagreed (for example,
\code{\#f} and the empty list are now distinct objects, as the book
requires), and it gained modules, keyword checking, sealing and slot
options that Thomas lacked. It was then rebuilt for speed: generic
dispatch goes through caches, known calls are made directly, and common
arithmetic, element access, iteration and slot access are compiled in line.

The compiler to C is new. It was written for this project and shares only
the reader with the Common Lisp implementation (which also runs its macros
while compiling). Thomas's own sources are no longer part of the tree;
each file in \file{src/} still records which Thomas file it grew from, and
Thomas's original 1992 README and release notes are kept as
\file{README.thomas} and \file{RELEASE}.

\section{How this book is organized}\label{sec:intro-organization}

The book has three parts and three appendices.

\paragraph{Part I, Getting Started,} is for reading in order.
This chapter is the first. Chapter~\ref{ch:start} explains what you need
to install, how to start the interactive implementation, and how to
compile and run a first program with \cmd{dylanc}.
Chapter~\ref{ch:tutorial} is a hands-on tour of the language, from
arithmetic to classes, generic functions and conditions, ending with a
small complete program built step by step and then compiled.

\paragraph{Part II, The Language,} is a reference to the language, in the
order a reader usually needs it: syntax (Chapter~\ref{ch:syntax}), control
flow (Chapter~\ref{ch:control}), functions (Chapter~\ref{ch:functions}),
classes (Chapter~\ref{ch:classes}), numbers (Chapter~\ref{ch:numbers}),
collections (Chapter~\ref{ch:collections}), functional operations
(Chapter~\ref{ch:functional}), conditions (Chapter~\ref{ch:conditions}),
modules (Chapter~\ref{ch:modules}), and printing and output
(Chapter~\ref{ch:printing}). These chapters follow the DIRM and note,
section by section, where either implementation differs from it.

\paragraph{Part III, The Implementations,} describes each implementation
as a tool: the Common Lisp implementation (Chapter~\ref{ch:lisp}), the
\cmd{dylanc} compiler and its options (Chapter~\ref{ch:dylanc}), the
extensions---macros and threads
(Chapter~\ref{ch:extensions}), how the compiler works inside
(Chapter~\ref{ch:internals}), and how to write fast programs and measure
them (Chapter~\ref{ch:performance}).

\paragraph{The appendices} collect the differences between this system and
the 1992 manual (Appendix~\ref{app:differences}), a summary of the library
(Appendix~\ref{app:library}), and how to run the test suites and
benchmarks (Appendix~\ref{app:testing}).

\section{Conventions}\label{sec:intro-conventions}

\paragraph{Code.} Dylan code, file names, commands and other text you type
or the computer prints are set in a typewriter face, as in
\code{(+ 1 2)} or \file{src/load-dylan.lisp}. In a syntax description,
italic words such as \meta{name} or \meta{body} stand for something you
supply.

\paragraph{Examples and their results.} Most examples are written as
Dylan source. When an example shows what an expression evaluates to, the
value follows the expression in a Dylan comment that begins with
\code{; =>}:

\begin{dylan}
(+ 1 2)          ; => 3
(list 1 'a "b")  ; => (1 a "b")
\end{dylan}

The comment is not part of the value; you can type the whole line,
comment and all, and Dylan will ignore everything from the semicolon to
the end of the line. When an expression returns more than one value, the
values are separated by commas:

\begin{dylan}
(truncate/ 17 5)  ; => 3, 2
\end{dylan}

When an example writes output, the text it writes is shown in a comment
beginning \code{; prints}:

\begin{dylan}
(display "Hello")  ; prints: Hello
\end{dylan}

\paragraph{Terminal sessions.} A session at the shell shows a \code{\$}
prompt before each command you type, followed by the command's output:

\begin{shell}
$ bin/dylanc hello.dyl
$ ./hello
Hello, world!
\end{shell}

\paragraph{Interactive sessions.} A session at the interactive
implementation's read-eval-print loop shows its \code{?} prompt before each
form you type, and the \code{;Value:} line it prints in reply:

\begin{repl}
? (* 6 7)

;Value: 42
\end{repl}

\paragraph{References to the 1992 manual.} A citation such as
\dylanbook{chapter 12} refers to a chapter of the 1992 Dylan manual. An
annotated copy of it, in HTML, is in the repository under
\file{manual/Dylan-1992/}; open \file{START-HERE.html} there. When this
book and the DIRM seem to disagree about the language, the DIRM is the
standard, and any deliberate difference is listed in
Appendix~\ref{app:differences}.

\paragraph{Notes, tips and warnings.} Three kinds of aside appear in the
text. A \textbf{Note} adds a detail or a cross-reference that you can skip
on a first reading. A \textbf{Tip} suggests an easier or better way to do
something. An \textbf{Important} box warns about a mistake that is easy to
make or a difference between the implementations that can surprise you.

\section{License}\label{sec:intro-license}

\index{license}
This software is \emph{source-available}, not open source. The original
Thomas system is Copyright 1992 Digital Equipment Corporation, and the
Common Lisp implementation, which began as a translation of it, is a
derivative work governed by Digital's 1992 notice. That notice allows you
to use, copy, modify and redistribute the software at no charge to
recipients who accept the same terms, but it forbids selling it or
transferring it for compensation. Because of that restriction, the terms
meet neither the Open Source Definition nor the Free Software Definition
and are not compatible with the GPL, MIT, BSD or Apache licenses.

The notice must appear in all copies, derivative works and modified
versions. It is reproduced in full in the file \file{LICENSE} at the top of
the repository, together with an explanation of what it permits; read that
file before redistributing any part of this system.

The home of the project is \url{https://github.com/blakemcbride/Dylan}.

\chapter{Getting Started}\label{ch:start}

This chapter takes you from an empty directory to running Dylan in both
implementations. It lists what you need to install, shows how to get the
source, walks through a first session with the interactive
implementation, and compiles a first program with \cmd{dylanc}. By the
end of it you will have everything you need for the tutorial in
Chapter~\ref{ch:tutorial}.

\section{Requirements}\label{sec:start-requirements}

\index{requirements}
Both implementations are written in Common Lisp and need \textbf{SBCL},
Steel Bank Common Lisp. The system was developed on SBCL~2.6.3, and any
recent release should work. The interactive implementation needs nothing
else but \cmd{make} to build it, and even that can be done without.

The compiler to C needs two more things, and can use a third:

\begin{itemize}
\item \textbf{A C compiler that supports the \code{musttail}
  attribute}, which guarantees that a call in tail position is compiled as
  a jump. That means \textbf{GCC 15 or later} or \textbf{Clang 13 or
  later}. Any other C11 compiler can be used with the \cmd{-{}-trampoline}
  option, which implements tail calls with trampolines instead, at some
  cost in speed. \cmd{dylanc} uses the compiler named by the \code{CC}
  environment variable, or \cmd{cc} if \code{CC} is not set; the
  \cmd{-{}-cc} option names one explicitly.
\item \textbf{\cmd{make}} and a POSIX shell. \cmd{bin/dylanc} is a shell
  script, and it builds the C runtime library with \cmd{make} the first
  time it is needed.
\item Optionally, \textbf{the Boehm-Demers-Weiser garbage collector},
  \file{libgc}, \emph{built with thread support}, together with its
  header files, for programs built with \cmd{-{}-gc boehm}
  (Section~\ref{sec:dylanc-gc}); the runtime then defines
  \code{GC\_THREADS} and links with \cmd{-lgc -lm -lpthread}. The
  default collector, the Memory Pool System, is part of the source tree
  and needs nothing installed.
\end{itemize}

\paragraph{Checking what you have.} These commands report the versions
installed:

\begin{shell}
$ sbcl --version
SBCL 2.6.3
$ gcc --version | head -1
gcc (GCC) 16.2.1 20260819 (Red Hat 16.2.1-2)
$ clang --version | head -1
clang version 22.1.8 (Fedora 22.1.8-4.fc44)
\end{shell}

Your version numbers will differ; what matters is that SBCL runs and that
one of the C compilers is new enough (or that you plan to use
\cmd{-{}-trampoline}). For the Boehm collector, the header \file{gc.h}
must be present too (\cmd{ls /usr/include/gc.h}). On Fedora and related
systems its development package is called \file{gc-devel}; on Debian and
Ubuntu it is \file{libgc-dev}. The packaged collectors on these systems
are normally built with thread support. If you build the collector
yourself, configure it with threads enabled.

\begin{notebox}
The C compiler is needed only to \emph{build} programs. An executable made
by \cmd{dylanc} needs only the C library at run time (and the shared
\file{libgc}, if built with \cmd{-{}-gc boehm}); it does not need SBCL.
\end{notebox}

\section{Getting the source}\label{sec:start-source}

The project lives at \url{https://github.com/blakemcbride/Dylan}. Clone
it with Git:

\begin{shell}
$ git clone https://github.com/blakemcbride/Dylan.git
$ cd Dylan
\end{shell}

There is nothing to configure and nothing to install. Every command in
this book is run from the top directory of the clone unless the text says
otherwise. The directories you will meet first are these:

\begin{center}
\begin{tabular}{@{}ll@{}}
\toprule
\file{src/} & the interactive implementation, in Common Lisp \\
\file{Makefile} & builds \file{bin/dylan} \\
\file{bin/dylan} & the interactive implementation, once built \\
\file{bin/dylanc} & the command that runs the compiler to C \\
\file{compiler/} & the compiler to C, in Common Lisp \\
\file{runtime/} & the compiler's runtime: \file{runtime/c} in C,
  \file{runtime/dylan} in Dylan \\
\file{manual/Dylan-1992/} & the annotated 1992 Dylan manual (the DIRM) \\
\file{dylan-examples.dyl} & the examples from the 1992 manual \\
\file{tests/}, \file{bench/} & test suites and benchmarks \\
\bottomrule
\end{tabular}
\end{center}

You can run \cmd{dylan} and \cmd{dylanc} from anywhere by putting
\file{bin} on your \code{PATH}, or by making symbolic links to them in a
directory that is already on it. \cmd{dylan} is self-contained, and the
\cmd{dylanc} script finds the rest of the system relative to its own real
location, so a link works for both.

\section{The interactive implementation}\label{sec:start-rep}

\index{read-eval-print loop}\index{dylan@\cmd{dylan}}
Build the interactive implementation once, in the top directory of the
clone:

\begin{shell}
$ make
sbcl --script src/build-dylan.lisp
;Writing /home/you/Dylan/bin/dylan
\end{shell}

\noindent
This loads the whole system into SBCL and saves it as one executable,
\file{bin/dylan}. (On Windows, without \cmd{make}, run
\cmd{sbcl -{}-script src/build-dylan.lisp} yourself; it writes
\file{bin\textbackslash dylan.exe}.) Run \cmd{make} again whenever
\file{src/} changes; it does nothing when \file{bin/dylan} is up to date.

\index{fasl files}
The first build compiles every Lisp file of the implementation and
writes the compiled files (\emph{fasl} files) into a directory under
\file{src/fasl/} named for the Lisp and its version, such as
\file{src/fasl/SBCL-2-6-3/}. This takes a little while. Later builds
recompile only a changed source file and every file loaded after it,
since later files may depend on its macros; you never need to delete the
fasl directory by hand. Because the fasl files are kept apart for each
Lisp version, upgrading SBCL simply makes a new directory.

\paragraph{Starting the loop.} Run \cmd{dylan}. It starts a
read-eval-print loop that reads Dylan rather than Lisp. Its prompt is a
question mark:

\begin{repl}
$ bin/dylan

Entering Dylan read-eval-print-loop.
Leave by typing "(exit)"

? (define greeting "Hello")

;Value: greeting

? (concatenate greeting ", world")

;Value: "Hello, world"
\end{repl}

Type a form and press Enter. The loop reads the form, compiles it into
Lisp, runs it, and prints its value after \code{;Value:}. A form may
span several lines; the loop keeps reading until the parentheses balance.
A definition such as \code{define} answers the name it defined. A form
that returns several values shows each one on its own line, numbered:

\begin{repl}
? (truncate/ 17 5)

;Value[1]: 3
;Value[2]: 2
\end{repl}

\begin{notebox}
In a real terminal the loop prints an extra blank line after each form you
type. The transcripts in this book leave it out.
\end{notebox}

\paragraph{Leaving the loop.} Type \code{dylan:done} or \code{(exit)},
or end the input with Ctrl-D, to leave \cmd{dylan}.

\paragraph{Loading it into SBCL yourself.} \index{dylan:rep}
\index{load-dylan.lisp@\file{load-dylan.lisp}}
Instead of \cmd{dylan}, you can start a plain SBCL in the top directory
of the clone, load \file{src/load-dylan.lisp}, and call the function
\code{dylan:rep}, which starts the same loop. This needs no build step,
since the load compiles what has changed, and it leaves you in a Lisp
where you can mix Lisp and Dylan:

\begin{repl}
$ sbcl
* (load "src/load-dylan.lisp")

Call DYLAN:REP to start a Dylan read-eval-print loop.
;Dylan loaded.  (dylan:rep) starts a read-eval-print loop.
T
* (dylan:rep)

Entering Dylan read-eval-print-loop.
Leave by typing "(exit)"

?
\end{repl}

\noindent
There \code{dylan:done} returns to the Lisp prompt. What you defined is
not lost: if you call \code{(dylan:rep)} again, the variables, classes and
methods you defined earlier are still there. To leave SBCL altogether,
type \code{(exit)} at the \code{?} prompt, or \code{(sb-ext:exit)} at the
Lisp prompt.

\paragraph{When something goes wrong.} An error that your Dylan code does
not handle stops the loop and enters SBCL's debugger. The debugger
offers restarts that return you to the Dylan prompt; choose one by typing
its number:

\begin{repl}
? (error "Account ~a is overdrawn" 42)

debugger invoked on a DYLAN::DYLAN-CONDITION in thread
#<THREAD tid=701926 "main thread" RUNNING {1200030003}>:
  Error: Account 42 is overdrawn

Type HELP for debugger help, or (SB-EXT:EXIT) to exit from SBCL.

restarts (invokable by number or by possibly-abbreviated name):
  0: [DYLAN-RESTART] Return to the read-eval-print loop.
  1: [DYLAN-REP    ] Return to the Dylan read-eval-print loop.
  2: [ABORT        ] Return to the Dylan read-eval-print loop.
...
0] 0
;Aborted

? (+ 3 4)

;Value: 7
\end{repl}

\begin{notebox}
A few errors are detected by Common Lisp itself rather than by Dylan; for
example, \code{(car 5)} is reported as a Lisp \code{TYPE-ERROR}. For such
an error the debugger still offers \code{DYLAN-REP}, ``Return to the Dylan
read-eval-print loop'', which gives you a new \code{?} prompt. Dylan's own
handlers see these errors as instances of \code{<error>}, so
\code{handler-case} catches them normally. In a Lisp you loaded the
system into, the debugger's \code{ABORT} leaves the loop for the Lisp
\code{*} prompt; nothing you defined is lost; type \code{(dylan:rep)} to
carry on.
\end{notebox}

\paragraph{Running a file at the loop.} The loop reads from standard
input, so you can feed it a whole file of Dylan from the shell. This runs
every form in \file{program.dyl} and then leaves:

\begin{shell}
$ bin/dylan < program.dyl
\end{shell}

Each form's value is printed after \code{;Value:}, as if you had typed it.
An unhandled error prints the debugger's report and ends the run, since
the debugger finds no more input.

\section{Compiling a file to Lisp}\label{sec:start-compile-file}

\index{dylan:compile-file}
The interactive implementation can also compile a whole file of Dylan
into a file of Common Lisp, which can then be loaded and run without
reading or compiling any Dylan. Suppose \file{greet.dyl} contains this:

\begin{dylan}
(define-method greet ((name <string>))
  (display "Hello, ")
  (display name)
  (display "!")
  (newline))

(greet "world")
\end{dylan}

Load \file{src/load-compiler.lisp} and call \code{dylan:compile-file}
with the names of the input and output files:

\begin{shell}
$ sbcl --non-interactive --eval '(load "src/load-compiler.lisp")' \
       --eval '(dylan:compile-file "greet.dyl" "greet.lisp")'
\end{shell}

Compiling does not run the program; it only writes \file{greet.lisp}.
That file is ordinary Common Lisp. To run it, load
\file{src/load-runtime.lisp} and then load the compiled file:

\begin{shell}
$ sbcl --non-interactive --eval '(load "src/load-runtime.lisp")' \
       --eval '(load "greet.lisp")'
...
Hello, world!
\end{shell}

You can also load a compiled file into a session and then start the
loop; the file's definitions are then available at the \code{?} prompt,
and you can call them, redefine them, or add methods to them.

\begin{notebox}
In this version, \file{load-compiler.lisp} and \file{load-runtime.lisp}
both load the whole system, exactly as \file{load-dylan.lisp} does; they
are separate entry points so that a program can say what it needs.
Chapter~\ref{ch:lisp} describes \code{dylan:compile-file},
\code{dylan:compile-to-file} and modules in detail.
\end{notebox}

\section{A first program with dylanc}\label{sec:start-dylanc}

\index{dylanc!first program}
The compiler to C takes one or more files of Dylan and produces a native
executable. Write this into \file{hello.dyl}:

\begin{dylan}
(define-method greet ((name <string>))
  (format #t "Hello, ~a!~%" name))

(greet "world")
\end{dylan}

Then compile it and run the result:

\begin{shell}
$ bin/dylanc hello.dyl
$ ./hello
Hello, world!
\end{shell}

A program is simply a sequence of top-level forms, run in order from the
top of the file when the executable starts. There is no \code{main}
function; here the last form calls \code{greet}.

\begin{notebox}
The 1992 manual uses \code{format} in examples but does not define it;
both implementations provide it (Chapter~\ref{ch:printing}).
\end{notebox}

\paragraph{The first build takes longer.} The first time you run
\cmd{dylanc}, it compiles the compiler itself into fasl files, builds the
C runtime library with \cmd{make}, and compiles the Dylan library's C
code, all inside the clone (under \file{compiler/fasl/} and
\file{runtime/build/}). Each of these is kept and reused, so later
builds of a small program take only a second or two. The clone therefore
needs to be writable by the user who runs \cmd{dylanc}.

\paragraph{The executable's name.} By default the executable is named
after the first file, without the \file{.dyl}. The \cmd{-o} option names it
yourself:

\begin{shell}
$ bin/dylanc -o greeter hello.dyl
$ ./greeter
Hello, world!
\end{shell}

\paragraph{Several files.} A program can be split across several files.
Name them all on the command line; they are compiled, in the order given,
into one program, exactly as if they were one long file. Put definitions
before the code that uses them at top level. With \file{shapes.dyl}

\begin{dylan}
(define-class <square> (<object>)
  (side init-keyword: side:))

(define-method area ((s <square>))
  (* (side s) (side s)))
\end{dylan}

\noindent and \file{main.dyl}

\begin{dylan}
(print (area (make <square> side: 7)))
\end{dylan}

\noindent the program is built and run like this:

\begin{shell}
$ bin/dylanc -o app shapes.dyl main.dyl
$ ./app
49
\end{shell}

\paragraph{What files appear.} Normally the only file \cmd{dylanc} leaves
behind is the executable. The C it generates is written beside the
executable and deleted once the C compiler has finished with it. The \cmd{-k} (\cmd{-{}-keep-c}) option keeps it:

\begin{shell}
$ bin/dylanc -k hello.dyl
$ ls
hello  hello.c  hello.dyl  hello-dylan.c  hello-dylan-user.c  hello.h
\end{shell}

\file{hello.c} holds the program's data, its generic-function dispatch
and the C \code{main} function; \file{hello.h} is a header the other files
share; and there is one file for each module's functions:
\file{hello-dylan-user.c} for your code, which is in the module
\code{dylan-user}, and \file{hello-dylan.c} for the library's module
\code{dylan}. The generated C contains \code{\#line} directives that point
back to your Dylan source, so a C debugger or a C compiler's message
names the line of Dylan it came from. Reading this C is a good way to see
what the compiler did with your program; Chapter~\ref{ch:internals}
explains it.

\paragraph{Errors when compiling.} A syntax error stops the compilation
with a message giving the file, line and column, and a non-zero exit
status:

\begin{shell}
$ bin/dylanc bad.dyl
bad.dyl:1:15: error: End of file inside a Dylan list.
\end{shell}

A name that is used but never defined is only a warning; the program is
still built, and the error comes
when the program runs and reaches that name:

\begin{shell}
$ bin/dylanc -o bad main.dyl
main.dyl:1:8: warning: area is not defined (module dylan-user)
main.dyl:1:14: warning: <square> is not defined (module dylan-user)
$ ./bad
Dylan error: the variable area@dylan-user has no value
\end{shell}

Take these warnings seriously: they almost always mean a misspelling or a
missing file on the command line. A program that ends with an unhandled
error prints \code{Dylan error:} and a message to standard error and exits
with status~1.

\paragraph{Other options.} \cmd{bin/dylanc -h} lists every option. The
ones you are most likely to want soon are \cmd{-O0} to \cmd{-O3} (how hard
the C compiler optimizes; the default is \cmd{-O2}), \cmd{-Os} for a small
executable, \cmd{-g} for a debugging build that checks every inline
operation, and \cmd{-{}-trampoline} and \cmd{-{}-cc} for other C
compilers. Chapter~\ref{ch:dylanc} describes them all.

\section{The examples and the test suites}\label{sec:start-tests}

The repository includes the examples from the 1992 manual and a large
set of tests. You do not need to run any of them to use the system, but
they are a useful check that an installation works, and the examples are
good reading.

The interactive implementation runs the manual's examples, in
\file{dylan-examples.dyl}, and compares the results with the transcripts
in \file{src/reference/}. It must be run from the \file{src} directory:

\begin{shell}
$ cd src
$ sbcl --eval '(load "load-dylan.lisp")' \
       --eval '(load "full-test.lisp")' \
       --eval '(dylan:test-examples)'
\end{shell}

The compiler to C has its own suites, among them the manual's examples
chapter by chapter, tests of the code generator, stress tests, tests of
the \cmd{dylanc} command, and tests of the C runtime:

\begin{shell}
$ sbcl --script tests/conformance/run.lisp
$ sbcl --script tests/codegen/run.lisp
$ sbcl --script tests/stress/run.lisp
$ sh tests/driver/run.sh
$ make -C runtime test
\end{shell}

\begin{warnbox}
By default each \file{run.lisp} builds every test with both GCC and Clang,
with \code{musttail} and with trampolines, which takes a long time. Give
the mode \cmd{quick} as an argument, as in
\cmd{sbcl -{}-script tests/conformance/run.lisp quick}, for a faster run
with one configuration. Appendix~\ref{app:testing} describes the modes and
the benchmarks.
\end{warnbox}

\section{Where to go next}\label{sec:start-next}

Chapter~\ref{ch:tutorial} is a hands-on tour of the language. Start the
interactive loop, keep this book open beside it, and type the examples as
you read; at the end of the tutorial you will build a small program and
compile it with \cmd{dylanc}.

If you already know Lisp or Scheme and want the details first, the
language reference begins with Chapter~\ref{ch:syntax}. The interactive
implementation is described fully in Chapter~\ref{ch:lisp}, and the
compiler in Chapter~\ref{ch:dylanc}.

\chapter{A Tutorial}\label{ch:tutorial}

This chapter is a hands-on tour of Dylan. It starts with arithmetic and
works up, a step at a time, through variables, functions, control flow,
collections, classes, generic functions and conditions, and it ends by
building a small complete program and compiling it into an executable.
Each section introduces a few ideas with short examples that you can type
and run, and each points to the chapter in Part~II that covers the same
ground in full.

The tutorial assumes you know some other programming language, but not
Lisp or Scheme. If you do know a Lisp, you will find much that is
familiar, and you can move quickly through the first sections.

\section{How to follow along}\label{sec:tut-follow}

The easiest way to follow the tutorial is at the interactive
implementation's read-eval-print loop. Start SBCL in the top directory of
the clone and start the loop, as Chapter~\ref{ch:start} explains:

\begin{repl}
$ sbcl
* (load "src/load-dylan.lisp")
* (dylan:rep)

Entering Dylan read-eval-print-loop.
Exit by typing "dylan:done"

?
\end{repl}

Then type the examples as you come to them. The examples are written as
Dylan source, with the value of each expression shown in a comment
beginning \code{; =>}. At the loop, the same value appears after
\code{;Value:}:

\begin{dylan}
(+ 1 2)   ; => 3
\end{dylan}

When an example prints something, the text it prints is shown in a
comment beginning \code{; prints}. Every example in this chapter has been
run in both implementations: at the loop, and compiled with \cmd{dylanc}.
You can also put the examples of a section into a file and compile it, as
Section~\ref{sec:tut-compile} shows; in a compiled program, nothing is
printed unless the program prints it, so wrap an expression in
\code{print} to see its value.

\begin{notebox}
The examples build on one another. A function defined in one section is
often used in the next, so type them in order, in one session.
\end{notebox}

\section{Expressions and prefix syntax}\label{sec:tut-syntax}

\index{prefix syntax}
A Dylan program is made of \term{expressions}, also called
\term{forms}. The simplest are literal constants, which evaluate to
themselves:

\begin{dylan}
42          ; => 42
-7          ; => -7
3.5         ; => 3.5
"a string"  ; => "a string"
#\a         ; => #\a
#t          ; => #t
#f          ; => #f
\end{dylan}

\code{\#t} and \code{\#f} are the true and false values, and
\code{\#\textbackslash{}a} is the character \code{a}. Everything else is
written as a parenthesized list. A list whose first element names a
function is a \term{call}: the function comes first, and its arguments
follow, separated by white space.

\begin{dylan}
(+ 1 2)                ; => 3
(max 3 9 4)            ; => 9
(* (+ 1 2) (- 10 4))   ; => 18
\end{dylan}

This is \term{prefix syntax}: the operator comes before its operands, and
the parentheses mark exactly where each expression begins and ends. There
is no operator precedence to remember, because nesting says everything.
The last example multiplies the value of \code{(+ 1 2)} by the value of
\code{(- 10 4)}. Operators such as \code{+} and \code{max} accept any
number of arguments:

\begin{dylan}
(+ 1 2 3 4)   ; => 10
(< 1 2 3)     ; => #t
\end{dylan}

Arithmetic operators are ordinary functions, not special syntax, which is
why they are written the same way as every other call. A few lists are
not calls but \term{special forms} with rules of their own, such as
\code{if}, \code{define} and \code{method}; you will meet them in the
sections that follow.

\paragraph{Names.} Names in Dylan may contain letters, digits and many
punctuation characters, so \code{even?}, \code{set!}, \code{<point>},
\code{*count*} and \code{remove-key!} are all single names. Dylan uses
punctuation by convention to suggest what a name is for:

\begin{itemize}
\item a name in angle brackets, such as \code{<integer>}, names a class;
\item a name ending in \code{?}, such as \code{even?}, names a predicate,
  a function that answers true or false;
\item a name ending in \code{!}, such as \code{add!}, names a function
  that may modify one of its arguments;
\item a name between asterisks, such as \code{*count*}, names a module
  variable whose value changes.
\end{itemize}

Names are not case-sensitive: \code{Max}, \code{MAX} and \code{max} are the
same name. Because spaces separate the parts of a form, you must put
spaces around operators: \code{(- x 1)} subtracts, but \code{x-1} is a
single name.

\paragraph{Symbols, keywords and quoting.} A name preceded by a quote
mark is a \term{symbol}, a value in its own right rather than a reference
to a variable. A name ending in a colon is a \term{keyword}, which is used
to label arguments. Both evaluate to themselves:

\begin{dylan}
'apple       ; => apple
size:        ; => size:
\end{dylan}

The quote mark also stops a list from being evaluated, which is one way to
write a constant list. A hash sign before a list makes a constant vector:

\begin{dylan}
'(1 2 3)     ; => (1 2 3)
#(1 2 3)     ; => #(1 2 3)
\end{dylan}

\paragraph{Comments.} A semicolon starts a comment that runs to the end of
the line. Text between \code{\#|} and \code{|\#} is a comment that can span
several lines.

\smallskip
Chapter~\ref{ch:syntax} describes the syntax completely (\dylanbook{chapter
3}).

\section{Numbers and arithmetic}\label{sec:tut-numbers}

\index{numbers}
Dylan has integers, exact ratios and floating-point numbers. Division of
integers that do not divide evenly gives an exact ratio, not a rounded
result:

\begin{dylan}
(/ 10 5)      ; => 2
(/ 10 4)      ; => 5/2
(+ 1/3 1/6)   ; => 1/2
(/ 10.0 4)    ; => 2.5
(as <float> 1/4)  ; => 0.25
\end{dylan}

Mixing a float with an exact number gives a float, and \code{as} converts a
number to another class. Integer division that yields an integer is done
by the functions whose names end in a slash. Each returns two values, the
quotient and the remainder:

\begin{dylan}
(truncate/ 17 5)   ; => 3, 2
(floor/ -7 2)      ; => -4, 1
(modulo -7 2)      ; => 1
(remainder -7 2)   ; => -1
\end{dylan}

\code{truncate/} rounds the quotient toward zero and \code{floor/} rounds
it down; \code{modulo} and \code{remainder} return just the remainder that
\code{floor/} and \code{truncate/} would. Functions that return several
values are common in Dylan, and Section~\ref{sec:tut-variables} shows how
to receive them.

There are the usual comparisons and predicates:

\begin{dylan}
(= 2 2.0)     ; => #t
(< 1 2 3)     ; => #t
(even? 4)     ; => #t
(zero? 0)     ; => #t
(abs -3)      ; => 3
(min 3 9 4)   ; => 3
\end{dylan}

\code{=} compares numbers by value, whatever their class; Dylan also has
\code{id?}, which answers whether two values are the very same object.
Powers, roots and the rest of the mathematical library are there too:

\begin{dylan}
(expt 2 10)   ; => 1024
(sqrt 16)     ; => 4.0
(sqrt 2)      ; => 1.4142135623730951
\end{dylan}

Notice that \code{(sqrt 16)} answers the float \code{4.0}, not the integer
\code{4}. The 1992 manual defines \code{sqrt}, \code{exp} and the
trigonometric functions as in Common Lisp, and its own examples show them
answering floats, so both implementations do the same.

\begin{warnbox}
The two implementations differ in the size of integers. The interactive
implementation has integers of any size, as Common Lisp does. A program
compiled with \cmd{dylanc} has 63-bit integers, and an operation whose
result does not fit signals an error:
\begin{dylan}
(expt 2 62)   ; => 4611686018427387904 in the interactive implementation
              ; error: integer overflow (integers are 63 bits) when compiled
\end{dylan}
Neither implementation has complex numbers.
\end{warnbox}

Chapter~\ref{ch:numbers} covers numbers in full (\dylanbook{chapter 14}).

\section{Variables}\label{sec:tut-variables}

\index{variables}
\code{define} creates a \term{module variable}, one that the whole
program can see, and gives it a value. \code{set!} changes the value of a
variable that already exists:

\begin{dylan}
(define *count* 0)            ; => *count*
(set! *count* (+ *count* 1))  ; => 1
*count*                       ; => 1
\end{dylan}

At the loop, a \code{define} answers the name it defined, and
\code{set!} answers the new value. Writing a variable's name by itself
evaluates to its value.

\index{bind}
\code{bind} creates \term{local variables} that exist only inside its
body. Its first argument is a list of bindings, each a name and an
initial value; the remaining forms are the body, and the value of the
last one is the value of the whole \code{bind}:

\begin{dylan}
(bind ((x 3)
       (y 4))
  (sqrt (+ (* x x) (* y y))))   ; => 5.0
\end{dylan}

The bindings are made in order, so a later initial value can use an
earlier variable:

\begin{dylan}
(bind ((a 1)
       (b (+ a 1)))
  (list a b))                   ; => (1 2)
\end{dylan}

A binding with more than one name receives multiple values, one in each
name. This is how you use the quotient and remainder from
\code{truncate/}:

\begin{dylan}
(bind ((quotient rem (truncate/ 17 5)))
  (list quotient rem))          ; => (3 2)
\end{dylan}

Local variables are the usual way to name intermediate results. Module
variables are best kept for things that really belong to the whole
program; by convention, a module variable whose value changes is named
between asterisks, as \code{*count*} is. Chapter~\ref{ch:syntax} covers
variables, and Chapter~\ref{ch:modules} covers modules, the namespaces
module variables live in (\dylanbook{chapters 4, 5 and 7}).

\section{Defining functions}\label{sec:tut-functions}

\index{define-method}
The usual way to define a function is \code{define-method}. It takes the
function's name, a list of parameters, and a body:

\begin{dylan}
(define-method square ((n <number>))
  (* n n))                      ; => square

(square 12)                     ; => 144
(square 1/2)                    ; => 1/4
\end{dylan}

The parameter \code{(n <number>)} says that \code{n} must be a number. The
class is called a \term{specializer}; if you leave it out and write just
\code{n}, the parameter accepts any object. A function returns the value of
the last form in its body; there is no \code{return} statement. Functions
can of course call other functions:

\begin{dylan}
(define-method hypotenuse ((a <real>) (b <real>))
  (sqrt (+ (square a) (square b))))

(hypotenuse 5 12)               ; => 13.0
\end{dylan}

\code{define-method} does a little more than define a function: it adds a
method to a \emph{generic function}, which can hold several methods for
different classes of argument. Section~\ref{sec:tut-generic} returns to
this, and it is the heart of Dylan. For now, think of
\code{define-method} as ``define a function''.

\paragraph{Anonymous functions.}\index{method}
The special form \code{method} makes a function without naming it. Such a
function is a value like any other: it can be called directly, stored in
a variable, or passed to another function.

\begin{dylan}
((method (x y) (* x y)) 6 7)    ; => 42

(define add1 (method (x) (+ x 1)))
(add1 41)                       ; => 42
\end{dylan}

\paragraph{Keyword and rest parameters.}\index{\#key}\index{\#rest}
After its required parameters, a parameter list may name \term{keyword
parameters} after \code{\#key}. The caller passes them by name, in any
order, and each may have a default value. A parameter after \code{\#rest}
collects any remaining arguments in a sequence.

\begin{dylan}
(define-method greet (name #key (greeting "Hello"))
  (concatenate greeting ", " name "!"))

(greet "Ada")                       ; => "Hello, Ada!"
(greet "Ada" greeting: "Goodbye")   ; => "Goodbye, Ada!"

(define-method sum-all (#rest numbers)
  (reduce + 0 numbers))

(sum-all 1 2 3 4 5)                 ; => 15
\end{dylan}

\paragraph{Returning several values.}\index{values}
\code{values} returns more than one value at once, and \code{bind}
receives them:

\begin{dylan}
(define-method smallest-and-largest ((s <sequence>))
  (values (reduce min (first s) s)
          (reduce max (first s) s)))

(smallest-and-largest '(3 1 4 1 5 9 2 6))   ; => 1, 9

(bind ((lo hi (smallest-and-largest #(10 -2 7))))
  (- hi lo))                                ; => 12
\end{dylan}

\paragraph{Closures.}\index{closure}
A function made inside another function can use the local variables
around it, and it keeps them for as long as it exists, even after the
function that made it has returned. Here \code{make-counter} returns a new
function each time it is called, and each of those functions has its own
private count:

\begin{dylan}
(define-method make-counter ()
  (bind ((n 0))
    (method ()
      (set! n (+ n 1))
      n)))

(define tick (make-counter))
(tick)                          ; => 1
(tick)                          ; => 2
(tick)                          ; => 3
\end{dylan}

Chapter~\ref{ch:functions} describes functions and parameter lists in
detail (\dylanbook{chapters 6 and 10}).

\section{Making decisions}\label{sec:tut-conditionals}

\index{if}
\code{if} takes a test, a form to evaluate when the test is true, and a
form to evaluate when it is false:

\begin{dylan}
(define-method sign-of ((n <real>))
  (if (negative? n)
      "negative"
      (if (zero? n) "zero" "positive")))

(sign-of -4)                    ; => "negative"
(sign-of 0)                     ; => "zero"
\end{dylan}

\begin{warnbox}
In Dylan, \code{\#f} is the \emph{only} false value. Zero, the empty
list and the empty string all count as true:
\begin{dylan}
(if 0 'true 'false)             ; => true
(if '() 'true 'false)           ; => true
\end{dylan}
\end{warnbox}

Each branch of an \code{if} is a single form. To do several things in a
branch, group them with \code{begin}, which evaluates its forms in order
and returns the value of the last.

\index{when}\index{unless}
\code{when} and \code{unless} are handy when there is only one branch.
\code{when} evaluates its body if the test is true, \code{unless} if it is
false; otherwise each returns \code{\#f}. Their bodies may contain several
forms.

\begin{dylan}
(when (> 3 2) "yes")            ; => "yes"
(when (> 2 3) "yes")            ; => #f
(unless (> 2 3) "no")           ; => "no"
\end{dylan}

\index{cond}
\code{cond} chooses among several cases. Each clause starts with a test;
the first clause whose test is true is evaluated, and the rest are skipped.
A final clause whose test is \code{else:} catches everything else.

\begin{dylan}
(define-method classify ((n <integer>))
  (cond ((< n 0)  "negative")
        ((= n 0)  "zero")
        ((< n 10) "small")
        (else:    "large")))

(classify 7)                    ; => "small"
(classify 42)                   ; => "large"
\end{dylan}

\index{case}
\code{case} compares one value against lists of constants, using
\code{id?}. The constants are not evaluated, so symbols in them need no
quote marks.

\begin{dylan}
(define-method day-kind (day)
  (case day
    ((saturday sunday) "weekend")
    ((monday tuesday wednesday thursday friday) "weekday")
    (else: "unknown")))

(day-kind 'sunday)              ; => "weekend"
(day-kind 'friday)              ; => "weekday"
(day-kind 'someday)             ; => "unknown"
\end{dylan}

Without an \code{else:} clause, a \code{case} with no matching clause
signals an error. Its relative \code{select} lets you choose the test that
compares the values; with \code{instance?} as the test it chooses by
class. Chapter~\ref{ch:control} describes it.

\index{and}\index{or}\index{not}
\code{and} and \code{or} evaluate their arguments from left to right and
stop as soon as the answer is known. \code{and} returns \code{\#f} or the
value of its last argument; \code{or} returns the first true value.
\code{not} turns \code{\#f} into \code{\#t} and every other value into
\code{\#f}.

\begin{dylan}
(and 1 2 3)                     ; => 3
(or #f 2 3)                     ; => 2
(not 0)                         ; => #f
\end{dylan}

Chapter~\ref{ch:control} covers all of these (\dylanbook{chapter 8}).

\section{Output and loops}\label{sec:tut-loops}

\paragraph{Printing.}\index{display}\index{print}\index{newline}
Before writing loops it helps to be able to print. Both implementations
have \code{display}, which writes a string or character as it is, and any
other object in its printed form; \code{newline}, which ends the line; and
\code{print}, which writes an object the way you would type it, strings in
quotes, followed by a newline.

\begin{dylan}
(begin
  (display "The answer is ")
  (display 42)
  (newline))
; prints: The answer is 42
\end{dylan}

\noindent
\code{format} does the same in one call, \code{(format \#t "The answer
is \textasciitilde a\textasciitilde\%" 42)}; Chapter~\ref{ch:printing}
describes it and the other printing functions.

\paragraph{for.}\index{for}
\code{for} is the general loop. Its first argument lists the loop
variables, each with an initial value and a form that computes its next
value. The second is a list of an end test and the forms whose value the
loop returns when the test becomes true. Any further forms are a body,
evaluated on each pass. This loop computes a factorial without any body
at all:

\begin{dylan}
(define-method factorial ((n <integer>))
  (for ((i n (- i 1))
        (product 1 (* product i)))
       ((<= i 0) product)))

(factorial 10)                  ; => 3628800
\end{dylan}

On each pass, all the variables are stepped together, using the values
from the previous pass.

\paragraph{dotimes.}\index{dotimes}
\code{dotimes} runs its body a fixed number of times, with a variable
counting from zero:

\begin{dylan}
(define-method sum-to ((n <integer>))
  (bind ((total 0))
    (dotimes (i (+ n 1))
      (set! total (+ total i)))
    total))

(sum-to 100)                    ; => 5050
\end{dylan}

\paragraph{for-each.}\index{for-each}
\code{for-each} steps through the elements of one or more collections.
Like \code{for}, it takes a list of variables, an end test with result
forms, and a body. The loop also ends, returning \code{\#f}, when any of
the collections runs out.

\begin{dylan}
(define-method first-even ((s <sequence>))
  (for-each ((x s))
            ((even? x) x)))

(first-even '(1 3 8 5 6))       ; => 8
(first-even '(1 3 5))           ; => #f
\end{dylan}

When there is no end test, write an empty list in its place. Here two
collections are stepped together, one of them an unbounded range of
integers:

\begin{dylan}
(for-each ((name '(ann bob cy))
           (i (range from: 1)))
          ()
  (display i) (display ". ") (display name) (newline))
; prints: 1. ann
;         2. bob
;         3. cy
\end{dylan}

\paragraph{while and until.}\index{while}\index{until}
\code{while} repeats its body as long as a test is true, and \code{until}
repeats it until a test becomes true. Both return \code{\#f}, so they are
used for their effects:

\begin{dylan}
(define-method collatz-steps ((n <integer>))
  (bind ((steps 0))
    (until (= n 1)
      (set! n (if (even? n) (/ n 2) (+ (* 3 n) 1)))
      (set! steps (+ steps 1)))
    steps))

(collatz-steps 27)              ; => 111
\end{dylan}

A method's parameters are ordinary local variables, so \code{set!} can
change them, as this one changes \code{n}.

In practice, many loops in Dylan are written not with these forms but
with the collection functions of Section~\ref{sec:tut-functional}.
Chapter~\ref{ch:control} covers the loops (\dylanbook{chapter 8}).

\section{Lists, vectors and strings}\label{sec:tut-sequences}

\index{sequences}
Lists, vectors and strings are all \term{sequences}: collections whose
elements are numbered from zero. They share most of their operations.

\paragraph{Lists.}\index{lists}
\code{list} makes a list of its arguments. \code{first}, \code{second}
and \code{third} take elements from the front; \code{element} takes any
element by its index; \code{size} counts them.

\begin{dylan}
(define primes (list 2 3 5 7 11))

(first primes)                  ; => 2
(second primes)                 ; => 3
(element primes 4)              ; => 11
(size primes)                   ; => 5
(last primes)                   ; => 11
\end{dylan}

A list is a chain of \term{pairs}. \code{cons} makes a new pair, putting
an element on the front of a list; \code{car} is a list's first element
and \code{cdr} the rest of it.

\begin{dylan}
(cons 1 (list 2 3))             ; => (1 2 3)
(car primes)                    ; => 2
(cdr primes)                    ; => (3 5 7 11)
(empty? '())                    ; => #t
\end{dylan}

Many operations make new sequences and leave their arguments alone:

\begin{dylan}
(reverse primes)                ; => (11 7 5 3 2)
(append primes '(13 17))        ; => (2 3 5 7 11 13 17)
(add primes 1)                  ; => (1 2 3 5 7 11)
primes                          ; => (2 3 5 7 11)
\end{dylan}

\code{add} adds an element wherever is most natural for the kind of
sequence, which for a list is the front.

\paragraph{Vectors.}\index{vectors}
A vector holds its elements in a block, so any element can be reached in
the same time. \code{vector} makes one from its arguments; \code{make}
makes one of a given size. \code{set!} with \code{element} changes an
element.

\begin{dylan}
(define scores (vector 10 20 30))

(element scores 1)              ; => 20
(set! (element scores 1) 99)    ; => 99
scores                          ; => #(10 99 30)

(make <vector> size: 3 fill: 0) ; => #(0 0 0)
\end{dylan}

The form \code{(set! (element scores 1) 99)} works because \code{set!} can
change more than variables: given a call such as \code{(element scores
1)}, it calls the matching \term{setter} function instead. You will see
the same thing with the slots of objects in Section~\ref{sec:tut-classes}.

A vector's size is fixed. A \code{<stretchy-vector>} grows as you add to
it with \code{add!}:

\begin{dylan}
(define names (make <stretchy-vector>))
(add! names 'ada)
(add! names 'grace)
(size names)                    ; => 2
(element names 1)               ; => grace
\end{dylan}

\paragraph{Strings.}\index{strings}
A string is a sequence of characters, so the sequence operations work on
strings too:

\begin{dylan}
(define lang "Dylan")

(size lang)                     ; => 5
(element lang 0)                ; => #\D
(concatenate lang " 1992")      ; => "Dylan 1992"
(as-uppercase lang)             ; => "DYLAN"
(reverse "abc")                 ; => "cba"
(copy-sequence "Hello, world" start: 7)   ; => "world"
\end{dylan}

\code{=} compares sequences element by element, while \code{id?} asks
whether two values are the same object. \code{as} converts between kinds
of sequence:

\begin{dylan}
(= "abc" "abc")                     ; => #t
(id? "abc" (copy-sequence "abc"))   ; => #f
(as <list> "abc")                   ; => (#\a #\b #\c)
(as <string> '(#\x #\y))            ; => "xy"
\end{dylan}

\paragraph{Ranges.}\index{ranges}
A \code{<range>} is a sequence of numbers in arithmetic progression. It
does not store its elements but computes them, so a range can be very
large or even unbounded, as the one in the \code{for-each} example of
Section~\ref{sec:tut-loops} was. \code{range} makes one: \code{from:} is
the first number, \code{up-to:} an upper bound that is not included,
\code{through:} a bound that is included, \code{by:} the step, and
\code{size:} the number of elements.

\begin{dylan}
(as <list> (range from: 1 up-to: 5))            ; => (1 2 3 4)
(as <list> (range from: 1 through: 5))          ; => (1 2 3 4 5)
(as <list> (range from: 10 through: 0 by: -5))  ; => (10 5 0)
(as <list> (range from: 0 size: 4))             ; => (0 1 2 3)
(reduce + 0 (range from: 1 through: 100))       ; => 5050
(member? 50 (range from: 0 by: 10))             ; => #t
\end{dylan}

\paragraph{Deques.}\index{deques}
A \code{<deque>}, a double-ended queue, grows and shrinks efficiently at
both ends. \code{push} and \code{pop} work at the front, and
\code{push-last} and \code{pop-last} at the back:

\begin{dylan}
(define line-up (make <deque>))
(push line-up 1)
(push-last line-up 2)
(push-last line-up 3)
(as <list> line-up)             ; => (1 2 3)
(pop line-up)                   ; => 1
(pop-last line-up)              ; => 3
(as <list> line-up)             ; => (2)
\end{dylan}

\paragraph{Sorting.}\index{sort}
\code{sort} returns a sorted copy of a sequence. Its \code{test:} keyword
argument is the function that compares two elements:

\begin{dylan}
(sort '(3 1 2) test: <)                     ; => (1 2 3)
(sort (vector "pear" "apple" "fig") test: <) ; => #("apple" "fig" "pear")
(sort '(3 1 2) test: >)                     ; => (3 2 1)
\end{dylan}

Chapter~\ref{ch:collections} describes every kind of sequence and its
operations (\dylanbook{chapter 12}).

\section{Tables}\label{sec:tut-tables}

\index{tables}
A \code{<table>} maps keys to values; other languages call it a hash
table, dictionary or map. Make one with \code{make}, store values with
\code{set!} and \code{element}, and read them back with \code{element}:

\begin{dylan}
(define ages (make <table>))

(set! (element ages 'alice) 34)
(set! (element ages 'bob) 27)
(set! (element ages 'carol) 41)

(element ages 'bob)                     ; => 27
(size ages)                             ; => 3
(element ages 'dave default: 'unknown)  ; => unknown
\end{dylan}

Asking for a key that is not in the table is an error unless you supply
\code{default:}, whose value is returned instead. \code{remove-key!}
removes a key. A table is a collection whose elements are its values, so
the collection operations apply to it; \code{key-sequence} gives its keys.

\begin{dylan}
(reduce + 0 ages)                       ; => 102
(remove-key! ages 'bob)
(size ages)                             ; => 2
\end{dylan}

Here is a common use: counting how often each word occurs.

\begin{dylan}
(define word-counts (make <table>))

(for-each ((w '(the cat saw the dog and the bird)))
          ()
  (set! (element word-counts w)
        (+ 1 (element word-counts w default: 0))))

(element word-counts 'the)              ; => 3
(element word-counts 'cat)              ; => 1
\end{dylan}

The order of a table's elements and keys is not defined, so sort them when
the order matters. Tables are described with the other collections in
Chapter~\ref{ch:collections}.

\section{Classes and objects}\label{sec:tut-classes}

\index{classes}\index{define-class}
\code{define-class} defines a class. It takes the class's name, a list of
superclasses, and a description of each \term{slot}. Every class
ultimately inherits from \code{<object>}.

\begin{dylan}
(define-class <point> (<object>)
  (x init-value: 0 init-keyword: x:)
  (y init-value: 0 init-keyword: y:))
\end{dylan}

Each slot here has a name, a default value given by \code{init-value:},
and a keyword given by \code{init-keyword:} that sets the slot when an
object is made. \code{make} makes an \term{instance} of a class:

\index{make}
\begin{dylan}
(define pt (make <point> x: 3 y: 4))

(x pt)                          ; => 3
(y pt)                          ; => 4
(y (make <point>))              ; => 0
\end{dylan}

Defining the class also defined a \term{getter} function for each slot,
named after the slot, which is why \code{(x pt)} reads the slot \code{x}.
It defined a setter too, which \code{set!} uses:

\begin{dylan}
(set! (x pt) 10)                ; => 10
(x pt)                          ; => 10
\end{dylan}

\code{instance?} asks whether an object belongs to a class, and
\code{object-class} answers the object's class:

\begin{dylan}
(instance? pt <point>)          ; => #t
(instance? 5 <point>)           ; => #f
(object-class pt)
\end{dylan}

\code{(object-class pt)} prints as \code{\{the class <point>\}}, and
\code{pt} itself, like any object without a readable printed form, as
\code{\{instance of <point>\}}. Chapter~\ref{ch:printing} explains how
to give your own classes a printed form.

\paragraph{Inheritance.} A class inherits the slots of its superclasses
and adds its own. Here is a small family of shapes, which the following
sections use:

\begin{dylan}
(define-class <shape> (<object>)
  (name init-value: "shape" init-keyword: name:))

(define-class <circle> (<shape>)
  (radius required-init-keyword: radius:))

(define-class <rect> (<shape>)
  (width  init-keyword: width:)
  (height init-keyword: height:))
\end{dylan}

\code{required-init-keyword:} means that \code{make} must be given that
keyword; making a \code{<circle>} without a \code{radius:} is an error.
Instances of the subclasses have the inherited \code{name} slot as well
as their own, and they are instances of \code{<shape>} as well as of their
own classes:

\begin{dylan}
(define wheel (make <circle> radius: 2 name: "wheel"))

(name wheel)                    ; => "wheel"
(radius wheel)                  ; => 2
(instance? wheel <shape>)       ; => #t
(subclass? <circle> <shape>)    ; => #t
\end{dylan}

\paragraph{Multiple inheritance.} A class may have more than one
superclass, and it then inherits from all of them:

\begin{dylan}
(define-class <colored> (<object>)
  (color init-value: 'black init-keyword: color:))

(define-class <colored-circle> (<circle> <colored>))

(define dot (make <colored-circle> radius: 1 color: 'red))
(list (color dot) (radius dot) (name dot))   ; => (red 1 "shape")
\end{dylan}

\paragraph{Initializing objects.}\index{initialize}
A slot's initial value need not be a constant. \code{init-function:}
gives a function of no arguments that is called to compute the value each
time an object is made. This matters when each object needs its own
mutable value, such as a fresh collection, or a value computed afresh,
such as a serial number:

\begin{dylan}
(define *next-id* 0)

(define-class <ticket> (<object>)
  (id    init-function: (method ()
                          (set! *next-id* (+ *next-id* 1))
                          *next-id*))
  (title init-value: "untitled" init-keyword: title:)
  (notes init-function: (method () (make <stretchy-vector>))))

(id (make <ticket>))                         ; => 1
(id (make <ticket> title: "printer jam"))    ; => 2
\end{dylan}

Had \code{notes} used \code{init-value:} instead, every ticket would share
one and the same stretchy vector.

For initialization that slot options cannot express, add a method to the
generic function \code{initialize}, which \code{make} calls on each new
object with the keyword arguments it was given. Call \code{next-method}
first, so that the slots are filled in the usual way, and then do the
rest. Here a rectangle made without a height becomes a square;
\code{slot-initialized?} asks whether a slot has a value yet:

\begin{dylan}
(define-class <rectangle> (<object>)
  (width  init-keyword: width:)
  (height init-keyword: height:))

(define-method initialize ((r <rectangle>) #key)
  (next-method)
  (unless (slot-initialized? r height)
    (set! (height r) (width r))))

(height (make <rectangle> width: 4))             ; => 4
(height (make <rectangle> width: 4 height: 3))   ; => 3
\end{dylan}

\begin{warnbox}
A slot's getter is an ordinary function in the module, so a local
variable or parameter with the same name hides it. Inside a method with a
parameter called \code{name}, the call \code{(name s)} would try to call
the parameter rather than the getter. Give parameters names that differ
from your slots.
\end{warnbox}

Chapter~\ref{ch:classes} describes classes, slot options and the
initialization protocol (\dylanbook{chapters 6 and 11}).

\section{Generic functions and multiple dispatch}\label{sec:tut-generic}

\index{generic function}\index{multiple dispatch}
Every \code{define-method} you have written so far has added a
\term{method} to a \term{generic function}. A generic function can hold
many methods, each for a different class of argument. When the generic
function is called, it chooses the most specific method whose
specializers match the arguments and runs it.

\code{define-generic-function} declares a generic function and its
parameters. It is optional---the first \code{define-method} creates the
generic function if there is none---but it documents the protocol in one
place:

\begin{dylan}
(define-generic-function area (shape))

(define-method area ((c <circle>))
  (* 3.14159 (radius c) (radius c)))

(define-method area ((r <rect>))
  (* (width r) (height r)))

(area (make <rect> width: 3 height: 5))     ; => 15
(area wheel)                                ; => 12.56636
\end{dylan}

Because \code{area} works on any shape, functions that use it work on any
shape too:

\begin{dylan}
(define some-shapes
  (list (make <circle> radius: 1) (make <rect> width: 2 height: 2)))

(map area some-shapes)          ; => (3.14159 4)
\end{dylan}

\paragraph{The most specific method wins.} Methods may overlap. When more
than one applies, the one whose specializers are more specific is chosen.
A method on \code{<shape>} serves every shape that has no method of its
own:

\begin{dylan}
(define-method describe ((s <shape>))
  (concatenate "a shape called " (name s)))

(describe (make <rect> width: 1 height: 1))   ; => "a shape called shape"
\end{dylan}

\paragraph{Dispatch on several arguments.} A method can be specialized on
any or all of its required arguments, and the choice takes all of them
into account. This is \term{multiple dispatch}, and it is what makes
generic functions more than methods attached to a single class:

\begin{dylan}
(define-method collide ((a <circle>) (b <circle>))
  "two circles")

(define-method collide ((a <circle>) (b <rect>))
  "a circle and a rectangle")

(define-method collide ((a <shape>) (b <shape>))
  "two shapes")

(collide wheel (make <rect> width: 1 height: 1))  ; => "a circle and a rectangle"
(collide (make <rect> width: 1 height: 1) wheel)  ; => "two shapes"
(collide wheel wheel)                             ; => "two circles"
\end{dylan}

\paragraph{Specializing on a single object.}\index{singleton}
A specializer can also be one particular object, written with
\code{singleton}. Such a method applies only when the argument is that
object, and it is more specific than any method on a class. This gives a
neat way to write the base case of a recursion:

\begin{dylan}
(define-method fact ((n <integer>))
  (* n (fact (- n 1))))

(define-method fact ((n (singleton 0)))
  1)

(fact 5)                        ; => 120
\end{dylan}

Singletons also suit symbols that act as commands or names:

\begin{dylan}
(define-method speak ((animal (singleton 'dog))) "woof")
(define-method speak ((animal (singleton 'cat))) "meow")
(define-method speak ((animal <symbol>)) "...")

(speak 'dog)                    ; => "woof"
(speak 'fish)                   ; => "..."
\end{dylan}

Built-in classes such as \code{<integer>}, \code{<string>} and
\code{<symbol>} are classes like any other, so you can add methods for them
to your own generic functions. Chapter~\ref{ch:functions} describes
generic functions and how methods are ordered (\dylanbook{chapters 6 and
10}).

\section{Calling the next method}\label{sec:tut-next-method}

\index{next-method}
Often a method for a subclass wants to do what the method for the
superclass does, and a little more. Inside a method, \code{next-method}
calls the next most specific method, the one that would have run if this
one did not exist. Called with no arguments, it passes along the original
arguments:

\begin{dylan}
(define-method describe ((c <circle>))
  (concatenate "a circle, " (next-method)))

(describe wheel)                ; => "a circle, a shape called wheel"
\end{dylan}

Here the method on \code{<circle>} runs first and calls the method on
\code{<shape>} for the rest of the description. \code{next-method} may
also be given arguments, which replace the original ones; the bank
example in Section~\ref{sec:tut-program} uses this to change an amount
before passing it on. Chapter~\ref{ch:functions} has the details.

\section{Collections and functional operations}\label{sec:tut-functional}

\index{map}\index{reduce}\index{choose}
Because functions are values, much of Dylan's power comes from functions
that take other functions as arguments. These work on every kind of
collection.

\code{map} applies a function to each element and collects the results in
a new collection of the same kind as the first argument. Given several
collections, it takes one element from each. \code{map-as} names the kind
of collection to return.

\begin{dylan}
(map (method (x) (* x x)) '(1 2 3 4))   ; => (1 4 9 16)
(map + '(1 2 3) '(10 20 30))            ; => (11 22 33)
(map-as <vector> square '(1 2 3))       ; => #(1 4 9)
\end{dylan}

\code{reduce} combines the elements with a function of two arguments,
starting from an initial value; \code{reduce1} uses the first element as
the initial value.

\begin{dylan}
(reduce + 0 '(1 2 3 4 5))       ; => 15
(reduce1 max '(3 9 2))          ; => 9
\end{dylan}

\code{choose} keeps the elements that satisfy a predicate.
\code{any?} and \code{every?} ask whether some or all of them do.
\code{find-key} answers the key, or index, of the first element that
satisfies a predicate.

\begin{dylan}
(choose even? '(1 2 3 4 5 6))   ; => (2 4 6)
(any? negative? '(3 -1 4))      ; => #t
(every? positive? '(3 -1 4))    ; => #f
(find-key #(5 8 13) even?)      ; => 1
\end{dylan}

\code{do} calls a function on each element just for its effect, and
returns \code{\#f}:

\begin{dylan}
(do (method (s) (display (area s)) (newline)) some-shapes)
; prints: 3.14159
;         4
\end{dylan}

\paragraph{Making functions from functions.}\index{curry}\index{compose}
\index{complement}
\code{curry} fixes the first arguments of a function. \code{compose}
chains functions, the rightmost applied first. \code{complement} turns a
predicate into its opposite. \code{apply} calls a function with arguments
taken from a sequence.

\begin{dylan}
(define add10 (curry + 10))
(add10 5)                                ; => 15
(map (curry * 2) '(1 2 3))               ; => (2 4 6)

(define sqrt-of-abs (compose sqrt abs))
(sqrt-of-abs -16)                        ; => 4.0

(choose (complement even?) '(1 2 3 4 5)) ; => (1 3 5)

(apply max '(4 7 1))                     ; => 7
\end{dylan}

Together these replace many explicit loops. The total area of a list of
shapes, for instance, is one line:

\begin{dylan}
(reduce + 0 (map area some-shapes))      ; => 7.14159
\end{dylan}

Chapter~\ref{ch:collections} lists the collection operations and
Chapter~\ref{ch:functional} the functional operations
(\dylanbook{chapters 12 and 17}).

\section{Leaving early with bind-exit}\label{sec:tut-exits}

\index{bind-exit}
Sometimes you want to stop a computation partway through and return a
value at once, for example when a search finds what it was looking for.
\code{bind-exit} binds a name to an \term{exit procedure}. Calling that
procedure, from anywhere inside the \code{bind-exit}, however deeply
nested, ends the \code{bind-exit} immediately and makes its argument the
value of the whole form.

\begin{dylan}
(define-method find-first ((pred <function>) (s <sequence>))
  (bind-exit (return)
    (do (method (x)
          (when (pred x) (return x)))
        s)
    #f))

(find-first even? '(1 3 4 5 6))     ; => 4
(find-first even? '(1 3 5))         ; => #f
\end{dylan}

The exit procedure is called from inside the function passed to
\code{do}, yet it leaves \code{find-first} entirely. If it is never called,
\code{bind-exit} returns the value of its last form, here \code{\#f}.

\index{unwind-protect}
\code{unwind-protect} guarantees that some cleanup forms run, however its
first form is left, whether normally or by an exit:

\begin{dylan}
(bind-exit (leave)
  (unwind-protect
    (leave 'escaped)
    (display "cleanup runs anyway")
    (newline)))
; prints: cleanup runs anyway
; => escaped
\end{dylan}

Chapter~\ref{ch:control} describes non-local exits
(\dylanbook{chapter 8}).

\section{Conditions and errors}\label{sec:tut-conditions}

\index{conditions}\index{error}
When something goes wrong, a Dylan program \term{signals} a
\term{condition}, an object that describes what happened. \code{error}
signals an error. In its simplest form it takes a message, which may
contain \code{\textasciitilde{}a} where each further argument is to be
inserted:

\begin{dylan}
(define-method safe-divide ((a <number>) (b <number>))
  (if (zero? b)
      (error "cannot divide ~a by zero" a)
      (/ a b)))

(safe-divide 10 4)              ; => 5/2
\end{dylan}

If nothing handles the error, the interactive implementation enters the
debugger, as Section~\ref{sec:start-rep} showed, and a compiled program
prints the message and exits.

\paragraph{Handling errors.}\index{handler-case}
\code{handler-case} evaluates a form and, if a condition of a given class
is signaled while it runs, abandons the form and evaluates the matching
clause instead. Each clause begins with a list of the class and, after
\code{condition:}, an optional name for the condition object:

\begin{dylan}
(handler-case (safe-divide 1 0)
  ((<error> condition: c)
   (list 'caught (condition-format-string c))))
; => (caught "cannot divide ~a by zero")
\end{dylan}

An error signaled by \code{error} with a message is an instance of
\code{<simple-error>}. The function \code{condition-format-string}
returns its message, and \code{condition-format-arguments} returns the
other arguments in a list. Errors that the system itself signals are
conditions too, and \code{handler-case} catches them the same way:

\begin{dylan}
(handler-case (element (vector 1 2) 5)
  ((<error>) "no such element"))    ; => "no such element"

(handler-case (+ 1 'a)
  ((<error>) "not a number"))       ; => "not a number"
\end{dylan}

\paragraph{Your own conditions.} Conditions are objects, and their classes
are ordinary classes. To signal something more specific than a message,
define a subclass of \code{<error>} with slots that describe the problem,
and pass an instance of it to \code{error}:

\begin{dylan}
(define-class <out-of-stock> (<error>)
  (item  required-init-keyword: item:)
  (short required-init-keyword: short:))

(define-method take-from-stock ((on-hand <integer>) (wanted <integer>))
  (if (> wanted on-hand)
      (error (make <out-of-stock> item: 'widget short: (- wanted on-hand)))
      (- on-hand wanted)))

(take-from-stock 10 4)          ; => 6

(handler-case (take-from-stock 10 15)
  ((<out-of-stock> condition: e)
   (list 'short-by (short e))))  ; => (short-by 5)
\end{dylan}

A \code{handler-case} may have several clauses; they are tried in order,
so put the more specific classes first.

\paragraph{Signals that are not errors.}\index{signal}\index{warnings}
\code{signal} signals a condition that a program may ignore. Given a
message, it signals a \code{<simple-warning>}. If no handler takes it,
\code{signal} prints the warning and returns \code{\#f}, and the program
carries on:

\begin{dylan}
(handler-case (signal "low on ~a" 'paper)
  ((<warning> condition: w)
   (condition-format-arguments w)))   ; => (paper)

(signal "nobody handles this")    ; => #f
; prints: WARNING: nobody handles this          (interactive)
;         Dylan warning: nobody handles this    (compiled)
\end{dylan}

Dylan's condition system can do much more: \code{handler-bind}
establishes a handler that runs \emph{without} first leaving the
signaling code, so it can return a value to it, decline, or choose among
\term{restarts}. Chapter~\ref{ch:conditions} explains it all
(\dylanbook{chapter 13}).

\section{A complete program}\label{sec:tut-program}

This section builds a small but complete program, a model of a bank, one
piece at a time. It brings together classes, generic functions,
\code{next-method}, tables, the collection functions and conditions. Type
the pieces at the loop as you go, or put them all in a file called
\file{bank.dyl}; the next section compiles that file.

\paragraph{Accounts.} Every account has an owner and a balance. A savings
account also earns interest at some rate, and a checking account charges
a fee on each withdrawal. Balances are whole numbers of dollars, and the
interest rate is an exact ratio, so no money is lost to rounding until we
choose to round.

\begin{dylan}
(define-class <account> (<object>)
  (owner   required-init-keyword: owner:)
  (balance init-value: 0 init-keyword: balance:))

(define-class <savings-account> (<account>)
  (rate init-value: 1/50 init-keyword: rate:))

(define-class <checking-account> (<account>)
  (fee init-value: 1 init-keyword: fee:))
\end{dylan}

\paragraph{An error class.} Withdrawing more than the balance is an
error, and the condition records which account and how much was asked
for. The slot names are chosen so that they do not clash with the names of
parameters used elsewhere (see the warning in
Section~\ref{sec:tut-classes}).

\begin{dylan}
(define-class <insufficient-funds> (<error>)
  (refused-account required-init-keyword: account:)
  (refused-amount  required-init-keyword: amount:))
\end{dylan}

\paragraph{The operations.} Three generic functions make up the protocol
of an account. The basic methods are on \code{<account>}:

\begin{dylan}
(define-generic-function deposit! (account amount))
(define-generic-function withdraw! (account amount))
(define-generic-function month-end! (account))

(define-method deposit! ((a <account>) (amount <integer>))
  (set! (balance a) (+ (balance a) amount)))

(define-method withdraw! ((a <account>) (amount <integer>))
  (when (> amount (balance a))
    (error (make <insufficient-funds> account: a amount: amount)))
  (set! (balance a) (- (balance a) amount)))

(define-method month-end! ((a <account>))
  (balance a))
\end{dylan}

Try them on a plain account:

\begin{dylan}
(define acct (make <account> owner: 'dave balance: 100))

(deposit! acct 25)              ; => 125
(withdraw! acct 40)             ; => 85
(balance acct)                  ; => 85
\end{dylan}

\paragraph{Specializing with next-method.} The other kinds of account
differ in small ways, which is exactly what \code{next-method} is for. A
checking account's withdrawal is an ordinary withdrawal of the amount
\emph{plus the fee}, so its method calls the next method with a changed
argument. A savings account adds its interest at the end of the month and
then does whatever any account does.

\begin{dylan}
(define-method withdraw! ((a <checking-account>) (amount <integer>))
  (next-method a (+ amount (fee a))))

(define-method month-end! ((a <savings-account>))
  (deposit! a (round (* (balance a) (rate a))))
  (next-method))
\end{dylan}

\code{round} returns two values, the nearest integer and the remainder,
but \code{deposit!} uses only the first.

\begin{dylan}
(define chk (make <checking-account> owner: 'erin balance: 100 fee: 2))
(withdraw! chk 10)              ; => 88

(define sav (make <savings-account> owner: 'fred balance: 500))
(month-end! sav)                ; => 510
\end{dylan}

Because the withdrawal for a checking account goes through the method on
\code{<account>}, it also gets that method's check for insufficient
funds, fee included:

\begin{dylan}
(handler-case (withdraw! chk 87)
  ((<insufficient-funds> condition: c)
   (refused-amount c)))         ; => 89
\end{dylan}

\paragraph{The bank.} The bank keeps its accounts in a table, keyed by
the owner's name. \code{open-account} takes the class of account to open,
so one function opens every kind. It passes any further keyword arguments
on to \code{make} with \code{apply}.

\begin{dylan}
(define *accounts* (make <table>))

(define-method open-account ((class <class>) (who <symbol>) #rest options)
  (bind ((a (apply make class owner: who options)))
    (set! (element *accounts* who) a)
    a))

(define-method lookup ((who <symbol>))
  (element *accounts* who default: #f))
\end{dylan}

\paragraph{Handling a refusal.} \code{try-withdraw} attempts a withdrawal
and reports a refusal instead of letting the error escape. It answers
whether the withdrawal was made.

\begin{dylan}
(define-method try-withdraw ((who <symbol>) (amount <integer>))
  (handler-case (begin (withdraw! (lookup who) amount) #t)
    ((<insufficient-funds> condition: c)
     (display "refused: ")
     (display (owner (refused-account c)))
     (display " cannot withdraw ")
     (display (refused-amount c))
     (newline)
     #f)))
\end{dylan}

\paragraph{Reports.} The total of all deposits maps \code{balance} over
the table and adds up the results. The report lists the accounts in
alphabetical order of owner; since the table's keys come in no particular
order, it sorts them, comparing the symbols as strings.

\begin{dylan}
(define-method total-deposits ()
  (reduce + 0 (map balance *accounts*)))

(define-method report ()
  (for-each ((who (sort (key-sequence *accounts*)
                        test: (method (a b)
                                (< (as <string> a) (as <string> b))))))
            ()
    (display who)
    (display ": ")
    (display (balance (lookup who)))
    (newline)))
\end{dylan}

\paragraph{A month at the bank.} Finally, the program itself: open three
accounts, make some transactions, close the month, and report.
\code{(do month-end! *accounts*)} runs the month's end for every account,
each with its own method.

\begin{dylan}
(open-account <savings-account> 'alice balance: 1000)
(open-account <checking-account> 'bob balance: 200 fee: 2)
(open-account <account> 'carol)

(deposit! (lookup 'carol) 50)
(withdraw! (lookup 'bob) 48)
(try-withdraw 'carol 80)
(do month-end! *accounts*)
(report)
(display "total: ")
(display (total-deposits))
(newline)
\end{dylan}

At the loop, each form's value appears after it as usual, mixed with the
output. The output itself is this:

\begin{shell}
refused: carol cannot withdraw 80
alice: 1020
bob: 150
carol: 50
total: 1220
\end{shell}

Alice earned two percent interest on 1000; Bob's withdrawal of 48 cost
him 50 with the fee; and Carol's attempt to withdraw 80 from her 50 was
refused and reported, without stopping the program.

\section{Compiling the program}\label{sec:tut-compile}

\index{dylanc}
The bank is a complete program, and \cmd{dylanc} can turn it into an
executable. Put the definitions and the final forms of the previous
section---everything except the small experiments with \code{acct},
\code{chk} and \code{sav}---into \file{bank.dyl}, in the order they were
given, and compile it:

\begin{shell}
$ bin/dylanc bank.dyl
$ ./bank
refused: carol cannot withdraw 80
alice: 1020
bob: 150
carol: 50
total: 1220
\end{shell}

The output is the same as at the loop, because the program uses only
\code{display} and \code{newline} for its output, which both
implementations have. Nothing in the program changes a class or adds a
method while it runs, so it is a static program, and the compiler accepts
it unchanged.

You can also run the same file in the interactive implementation without
typing it, either by feeding it to the loop or by compiling it to Lisp, as
Chapter~\ref{ch:start} showed:

\begin{shell}
$ (cat bank.dyl; echo dylan:done) | \
    sbcl --non-interactive --eval '(load "src/load-dylan.lisp")' \
         --eval '(dylan:rep)'
\end{shell}

This is the usual way of working with the two implementations: develop
and test a piece at a time at the loop, where every definition can be
changed and tried at once, then compile the finished program into a fast
executable that needs neither SBCL nor Dylan to run.

\begin{tipbox}
The compiler warns about any name a program uses but never defines.
Building with \cmd{dylanc} from time to time, even while you are still
working at the loop, is a quick way to catch misspelled names.
\end{tipbox}

\section{Where to go from here}\label{sec:tut-next}

You have now seen most of the language at least once. Part~II covers each
topic thoroughly, in the order of this tutorial: syntax and variables
(Chapter~\ref{ch:syntax}), control flow (Chapter~\ref{ch:control}),
functions and generic functions (Chapter~\ref{ch:functions}), classes
(Chapter~\ref{ch:classes}), numbers (Chapter~\ref{ch:numbers}),
collections (Chapter~\ref{ch:collections}), functional operations
(Chapter~\ref{ch:functional}), conditions (Chapter~\ref{ch:conditions}),
modules (Chapter~\ref{ch:modules}), and printing
(Chapter~\ref{ch:printing}).

Part~III describes the two implementations as tools. Read
Chapter~\ref{ch:dylanc} before building larger programs with the compiler,
and Chapter~\ref{ch:extensions} for its macros and threads.
Appendix~\ref{app:differences} lists everything in which this system
differs from the 1992 manual.

The 1992 manual itself is worth reading alongside this book. It is short,
clearly written, and full of examples, and its annotated copy in
\file{manual/Dylan-1992/} links each chapter to the design notes that
later changed it.


\part{The Language}
\chapter{Syntax and Program Structure}\label{ch:syntax}

Dylan, as the 1992 manual defines it, is written in a parenthesized prefix
notation of the kind Lisp and Scheme use. A program is a sequence of
\emph{forms}; each form is either an atom, such as a number, a string or a
name, or a list whose first element says what the form does. There are no
statements, no operator precedence and no semicolons ending lines: a
semicolon starts a comment.

This chapter describes the notation itself: how the reader turns
characters into Dylan objects (Section~\ref{sec:syntax-reader}), how those objects
are evaluated as expressions (Section~\ref{sec:syntax-eval}), and how a program is
put together from top-level forms and source files
(Section~\ref{sec:syntax-program}). The standard is \dylanbook{chapters 2 and 3};
symbols and keywords are also covered in \dylanbook{chapter 15}. Both
implementations read programs with the same reader, so everything in
Section~\ref{sec:syntax-reader} holds for both.

\section{How the examples are shown}\label{sec:syntax-conventions}

The examples in this part of the book are Dylan source. Each has been run
with both implementations, compiled with \cmd{bin/dylanc} and typed at the
read-eval-print loop of the Common Lisp implementation, except where the
text says that an example applies to only one of them. The value of an
expression is shown in a comment after it, introduced by \code{=>}:

%! skip
\begin{dylan}
(+ 3 4)                 ; => 7
\end{dylan}

\noindent When an expression returns several values, the comment lists them
in order, separated by commas; \code{(values 98 102)} is shown as
\code{; => 98, 102}. When it returns no values at all, the comment says
\code{(no values)}. A long value continues on the comment lines below. Output that a form prints is shown in comments
introduced by \code{prints:}, one line of output to each comment:

%! skip
\begin{dylan}
(print "hello")         ; prints: "hello"
\end{dylan}

\noindent A form whose evaluation signals an error is marked
\code{; => error}, usually followed by a short description of the error.
The exact wording of error messages differs between the two
implementations, and the comment does not try to reproduce either.

Values that cannot be read back in, such as classes and functions, are
printed inside braces: \code{\{the class <integer>\}}, as the 1992
manual writes them.

\section{The reader}\label{sec:syntax-reader}

The \term{reader} turns the characters of a program into Dylan objects:
numbers, strings, symbols, lists and so on. A program is just those objects;
evaluation (Section~\ref{sec:syntax-eval}) gives them meaning. The reader is the
same in both implementations: the compiler to C uses the Common Lisp
implementation's reader, so a file that one reads, the other reads
identically.

\subsection{White space and delimiters}

Spaces, tabs, newlines, returns and page breaks are \term{white space}. White
space separates tokens, and the amount of it does not matter, so a program
can be laid out freely. A token also ends at a parenthesis, a double quote,
a semicolon, a quote mark, a backquote or a comma. Everything else is
part of a token, which is why names such as \code{set!}, \code{even?},
\code{<integer>} and \code{*standard-output*} are single names.

\subsection{Comments}\label{sec:syntax-comments}

A semicolon\index{comment} starts a comment that runs to the end of the
line. The characters \code{\#|} start a block comment, which runs to the
matching \code{|\#}. Block comments may span lines and may be nested, so a
block comment can comment out code that already contains one:

\begin{dylan}
(+ 1 #| an inline comment |# 2)   ; => 3
#| (print "never run")
   #| a nested comment |#
   (print "still inside the outer comment") |#
\end{dylan}

\noindent A comment is white space as far as the reader is concerned. Note
that \code{\#;}, which comments out one form in some Lisps, does not exist
in Dylan.

\subsection{Numbers}\label{sec:syntax-numbers}

Numbers are written in the usual way. Chapter~\ref{ch:numbers} describes the
number classes and arithmetic; here is only how numbers are written.

\begin{itemize}
\item An \term{integer} is an optional sign followed by decimal digits:
  \code{0}, \code{42}, \code{-17}, \code{+57}. A trailing decimal point is
  allowed and still denotes an integer: \code{1.} is the integer \code{1}.
\item A \term{ratio} is two integers separated by a slash, with no spaces:
  \code{1/3}, \code{-6/4}. Ratios are always reduced, and a ratio that
  reduces to a whole number is an integer: \code{4/2} reads as \code{2}.
\item A \term{floating-point number} has a decimal point with digits after
  it, an exponent, or both: \code{1.5}, \code{-4.0}, \code{.5},
  \code{1e10}, \code{2.5d2}. An exponent is introduced by \code{e} or
  \code{d} for a double float, the default, or by \code{f} or \code{s} for
  a single float; \code{1.5s0} is a single float.
\item An integer or ratio in another base is written \code{\#x} (hexadecimal),
  \code{\#o} (octal) or \code{\#b} (binary) followed by the digits, which
  may have a sign: \code{\#x1f4e}, \code{\#xFF}, \code{\#o-17},
  \code{\#b101}. Letters may be in either case.
\end{itemize}

\begin{dylan}
(list 123 +57 -4.0 1.5 #x1f4e)     ; => (123 57 -4.0 1.5 8014)
(list 1. 6/4 4/2 .5)               ; => (1 3/2 2 0.5)
(list #o-17 #b101 #xff/2 1.0e-3)   ; => (-15 5 255/2 0.001)
\end{dylan}

\noindent Both implementations print double floats the same way, but they
print a single float and a large or small double float differently: the
compiler to C prints \code{1.5s0} and \code{10000000000.0} where the Common
Lisp implementation prints \code{1.5f0} and \code{1.0e10}.

\begin{warnbox}
The compiler to C has no bignums: an integer is 63 bits, and an integer
literal that does not fit is an error when compiling. It has no complex
numbers either. The Common Lisp implementation has both, and reads a
complex number in Common Lisp's notation, \code{\#c(1 2)}. See
Chapter~\ref{ch:numbers}.
\end{warnbox}

\subsection{Characters}\label{sec:syntax-characters}

A \term{character} is written \code{\#\textbackslash} followed by the
character itself, or by the character's name. The names of the 1992 manual
are \code{newline}, \code{space}, \code{rubout}, \code{page}, \code{tab},
\code{backspace}, \code{return} and \code{linefeed}; their case does not
matter, so \code{\#\textbackslash Space} and \code{\#\textbackslash space}
are the same character. \code{\#\textbackslash linefeed} is the same
character as \code{\#\textbackslash newline}. A single letter after the
backslash is that letter, in its case: \code{\#\textbackslash a} and
\code{\#\textbackslash A} are different characters. Any other single
character may follow the backslash too, parentheses, semicolons and quotes
included.

\begin{dylan}
(list #\a #\A #\( #\; #\")   ; => (#\a #\A #\( #\; #\")
(list #\Space #\TAB)         ; => (#\space #\tab)
(map (method (c) (as <integer> c))
     (list #\newline #\linefeed #\space #\tab #\rubout))
                             ; => (10 10 32 9 127)
\end{dylan}

\noindent Because the reader is built on Common Lisp's, it also accepts the
other character names Common Lisp knows, such as \code{\#\textbackslash
Nul} and \code{\#\textbackslash Escape}. A program that is to be portable to
other Dylan implementations should keep to the manual's names.

\subsection{Strings}\label{sec:syntax-strings}

A \term{string} is written between double quotes. Every character between
the quotes is part of the string, white space and newlines included. A
backslash quotes the character after it, which is how a string includes a
double quote or a backslash:

\begin{dylan}
"abc"              ; => "abc"
(size "a\"b")      ; => 3
(size "a\\b")      ; => 3
(element "a\"b" 1) ; => #\"
\end{dylan}

\noindent A backslash does nothing more than quote the next character, as
\dylanbook{chapter 3} says. In particular there are no escape sequences:
\code{"a\textbackslash nb"} is the three characters \code{a}, \code{n} and
\code{b}, not a string with a newline in it.

\begin{dylan}
(size "a\nb")      ; => 3
(element "a\nb" 1) ; => #\n
\end{dylan}

\noindent To put a newline into a string, write the newline itself between
the quotes, or build the string from characters at run time:

\begin{dylan}
(size (as <string> (list #\a #\newline #\b)))   ; => 3
\end{dylan}

\subsection{Symbols}\label{sec:syntax-symbols}

Any token that is not a number, and does not end in a colon, is a
\term{symbol}. Symbols name variables, and they are also objects in their
own right, of class \code{<symbol>}\index{<symbol>@\code{<symbol>}}. A
symbol may contain letters, digits and most punctuation: \code{+},
\code{<=}, \code{set!}, \code{even?}, \code{*x*}, \code{\$pi}, \code{1+},
\code{-} and \code{...} are all symbols.

\begin{dylan}
(list '+ '<= 'set! 'even? '*x* '$pi '1+ '- '...)
                       ; => (+ <= set! even? *x* $pi 1+ - ...)
(instance? 'foo <symbol>)   ; => #t
\end{dylan}

\noindent Symbols are not case-sensitive\index{case of symbols}:
\code{Foo}, \code{foo} and \code{FOO} are the same symbol. But a symbol
remembers the case it was first written in, and that is how it prints
(\dylanbook{chapter 15}). In the example below the symbol is written
\code{Hello} first, so it prints that way however it is written later:

\begin{dylan}
(id? 'Hello 'hello)        ; => #t
'HELLO                     ; => Hello
(as <string> 'hello)       ; => "Hello"
\end{dylan}

\noindent ``First written'' means first read. The compiler to C reads a
whole program before running any of it, so a symbol's spelling is the one
in the source text that comes first. The Common Lisp implementation reads a
form at a time, so a symbol first made at run time by
\code{(as <symbol> "Abc")} keeps the spelling \code{Abc} if no form read
earlier wrote it.

Dylan programs follow naming conventions (\dylanbook{chapter 3}). Except
for setter variables, described below, the conventions mean nothing to the
language; they only help a reader.

\begin{itemize}
\item Names of classes begin and end with angle brackets: \code{<window>},
  \code{<integer>}.
\item Module variables meant to change while the program runs begin and
  end with asterisks: \code{*window-count*}.
\item Constants begin with a dollar sign: \code{\$pi}.
\item Predicates, functions that answer true or false, end with a question
  mark: \code{even?}, \code{instance?}.
\item Operations that destructively modify their arguments end with an
  exclamation point: \code{reverse!}, \code{sort!}.
\end{itemize}

\subsection{Setter names}\label{sec:syntax-setter-names}

A variable name is either a symbol or a list of the form
\code{(setter \meta{symbol})}. The second kind names a \term{setter
variable}, the variable that holds the function which stores into the
place the function in \meta{symbol} reads. The value of \code{element} reads
an element of a collection; the value of \code{(setter element)} stores
one. A setter name can be used wherever a symbol can name a variable:

\begin{dylan}
(define (setter my-variable) 20)
(setter my-variable)            ; => 20
\end{dylan}

\noindent Because \code{setter} has this special meaning, the symbol
\code{setter} itself cannot be the name of a variable. The design note that
later replaced \code{(setter foo)} with \code{foo-setter} is not part of
this language. Setters are described in Section~\ref{sec:control-assignment} and
Section~\ref{sec:functions-setters}.

\subsection{Keywords}\label{sec:syntax-keywords}

A token ending in a colon is a \term{keyword}: \code{foo:}, \code{size:},
\code{init-value:}. Keywords look like symbols but are a separate class,
\code{<keyword>}\index{<keyword>@\code{<keyword>}}, and are used as
symbolic constants, chiefly to name keyword arguments. The colon is not
part of a keyword's name. Keywords, like symbols, are not case-sensitive and
remember the case they were first written in.

\begin{dylan}
foo:                         ; => foo:
(id? hello: HELLO:)          ; => #t
(instance? foo: <keyword>)   ; => #t
(instance? foo: <symbol>)    ; => #f
(as <keyword> "size")        ; => size:
(as <string> size:)          ; => "size"
\end{dylan}

\noindent Two keywords, or two symbols, with the same name are always the
same object: \code{id?} is true of them.

\subsection{True, false and the empty list}\label{sec:syntax-booleans}

\code{\#t}\index{t@\code{\#t}} is the canonical true value and
\code{\#f}\index{f@\code{\#f}} is the false value; either may be written in
upper case. \code{\#f} is the \emph{only} false value: every other object
counts as true when tested, including \code{0}, the empty string and the
empty list.

The empty list, written \code{()}, is an object of its own, of class
\code{<empty-list>}. Unlike in older Lisps it is not the same object as
\code{\#f}, and it is true:

\begin{dylan}
(id? '() #f)                    ; => #f
(if '() 'true 'false)           ; => true
(instance? '() <list>)          ; => #t
\end{dylan}

\subsection{Lists and vectors}\label{sec:syntax-lists}

A list is written as its elements between parentheses. A dot between the
last two elements writes a dotted list, whose final tail is not the empty
list. A vector, a one-dimensional array, is written like a list with a
\code{\#} before the opening parenthesis.

\begin{dylan}
'(a (b c) ())        ; => (a (b c) ())
'(a . b)             ; => (a . b)
'#(1 2 3)            ; => #(1 2 3)
(instance? '#(1 2) <vector>)   ; => #t
\end{dylan}

\noindent A list in a program is normally a form to be evaluated, so a list
meant as data must be quoted. A vector need not be: \code{\#(1 2 3)}
evaluates to itself. The elements of a literal list or vector are not
evaluated, and literal data should not be modified.

\subsection{Quote and backquote}\label{sec:syntax-quote}

A single quote before a form, \code{'\meta{form}}, is an abbreviation for
\code{(quote \meta{form})}\index{quote@\code{quote}}, which evaluates to
the form itself, unevaluated. Quoting is needed for symbols and lists used
as data; other literals evaluate to themselves and need no quote, though
quoting them does no harm.

\begin{dylan}
(quote foo)          ; => foo
'foo                 ; => foo
'(+ 10 10)           ; => (+ 10 10)
(+ 10 10)            ; => 20
'"s"                 ; => "s"
\end{dylan}

A \term{backquote} quotes a template in which parts are filled in when it is
evaluated. Within a backquoted form, a comma, \code{,\meta{form}},
evaluates \meta{form} and puts its value in that place; \code{,@\meta{form}}
evaluates \meta{form}, which must answer a list, and splices its elements
in. \dylanbook{chapter 3} defers backquote to a section the manual does not
contain, so both implementations follow the usual Lisp rules, including
backquotes nested inside backquotes and backquoted vectors.

\begin{dylan}
(define x 5)
(define xs '(a b))
`(x ,x)              ; => (x 5)
`(1 ,@xs 2)          ; => (1 a b 2)
`(1 . ,x)            ; => (1 . 5)
`#(1 ,x)             ; => #(1 5)
\end{dylan}

\noindent The reader turns \code{'x}, \code{`x}, \code{,x} and \code{,@x}
into the lists \code{(quote x)}, \code{(quasiquote x)}, \code{(unquote x)}
and \code{(unquote-splicing x)}. The compiler to C prints these lists in
that long form; the Common Lisp implementation prints them abbreviated,
as \code{'x} and so on. A comma outside any backquote is an error. Backquote
is most useful in macros, an extension of the compiler to C described in
Chapter~\ref{ch:extensions}.

\subsection{\#key, \#rest and \#next}\label{sec:syntax-parameter-tokens}

The tokens \code{\#key}\index{key@\code{\#key}},
\code{\#rest}\index{rest@\code{\#rest}} and
\code{\#next}\index{next@\code{\#next}} mark the parts of a parameter
list (Section~\ref{sec:functions-parameters}); \code{\#rest} also appears in
\code{bind} (Section~\ref{sec:control-bind}). Like \code{\#t} and \code{\#f} they are
case-insensitive.

Both implementations also read \code{!key}, \code{!rest} and \code{!next}
as these same tokens, because Thomas, from which the Common Lisp
implementation grew, wrote them that way, and so do its programs, the
manual's examples among them. As a consequence those three names cannot
name variables. Any other name may.

\subsection{Unreadable objects}

Some objects, such as functions and classes, have no notation that reads
back as the same object. They print inside braces, \code{\{the generic
function print\}}, and there is no notation that reads them back. Any other
\code{\#} syntax than those above (\code{\#t}, \code{\#f},
\code{\#\textbackslash}, \code{\#(}, \code{\#x}, \code{\#o}, \code{\#b},
\code{\#|}, \code{\#key}, \code{\#rest}, \code{\#next}, and \code{\#c} in the
Common Lisp implementation) is an error.

\section{Expressions and evaluation}\label{sec:syntax-eval}

A program is made of expressions. Evaluating an expression yields zero or
more values, may have side effects, and may evaluate other expressions
inside it. There are four kinds of expression (\dylanbook{chapter 3}):
literal constants, variable references, function calls and syntax forms.

\subsection{Literal constants}

Numbers, characters, strings, keywords, vectors, \code{\#t}, \code{\#f}
and the empty list \code{()} evaluate to themselves. Symbols and non-empty
lists must be quoted to be used as constants, because unquoted they are
variable references and calls.

\begin{dylan}
"abc"      ; => "abc"
123        ; => 123
foo:       ; => foo:
#\a        ; => #\a
#t         ; => #t
()         ; => ()
'(1 2 3)   ; => (1 2 3)
\end{dylan}

\subsection{Variable references}

A symbol, or a setter name \code{(setter \meta{symbol})}, evaluates to the
value of the variable it names. There are two kinds of variable.
\term{Lexical variables} are created by \code{bind}, by a method's parameter
list and by several other forms, and are visible only in the text of the
form that creates them. \term{Module variables} are created by the defining
forms (Section~\ref{sec:syntax-program}) and are visible throughout their module
(Chapter~\ref{ch:modules}). An inner variable \emph{shadows} any outer one of the
same name, lexical or module.

\begin{dylan}
(define my-variable 25)
my-variable                     ; => 25
(bind ((my-variable 50))
  (+ my-variable my-variable))  ; => 100
\end{dylan}

\noindent Referring to a variable that has no value is an error. The
compiler to C also warns, when compiling, of a module variable that is used
but defined nowhere in the program.

\subsection{Function calls}\label{sec:syntax-calls}

A list whose first element is not the name of a syntax operator is a
\term{function call}. Its first element, the \term{operator}, is evaluated
to give the function; the rest are evaluated to give the arguments, and the
function is called with them. The values of the call are the values the
function returns.

The operator may be any expression that evaluates to a function, not only a
name. In this Dylan is like Scheme and unlike Common Lisp:

\begin{dylan}
(+ 3 4)              ; => 7
(* (+ 3 4) 5)        ; => 35
((if #t + *) 4 5)    ; => 9
((method (x) (* x x)) 6)   ; => 36
\end{dylan}

\noindent The arguments are evaluated from left to right. The manual allows
the operator to be evaluated at any time before, during or after the
arguments, and the two implementations differ: the compiler to C evaluates
the operator first, and the Common Lisp implementation evaluates it last. A
program should not depend on either order.

%! pre: (define trail '())
%! pre: (define-method note (x) (set! trail (cons x trail)) x)
\begin{dylan}
(list (note 1) (note 2) (note 3))   ; => (1 2 3)
(reverse trail)                     ; => (1 2 3)
\end{dylan}

\noindent (Here \code{note} records each argument in a list as it is
evaluated.)

\subsection{Syntax forms}\label{sec:syntax-forms}

A list whose first element names a \term{syntax operator} is a
\term{syntax form}. A syntax form looks like a call, but it is not one:
each syntax operator has its own rule for which of its parts are evaluated,
when, and what the others mean. In \code{(bind ((x 10)) (+ x 1))}, for
example, \code{((x 10))} is not evaluated at all; \code{bind} reads it as a
list of bindings.

There are three kinds of syntax operator.

\begin{description}
\item[Special forms] are built into the language:
  \code{begin}, \code{bind}, \code{bind-exit}, \code{bind-methods},
  \code{if}, \code{method}, \code{quote}, \code{set!}, \code{setter} and
  \code{unwind-protect}.
\item[Macros] are syntax forms defined in terms of others. The language's
  macros are \code{and}, \code{or}, \code{when}, \code{unless}, \code{cond}, \code{case}, \code{select}, \code{for},
  \code{for-each}, \code{dotimes}, \code{while}, \code{until}, and the
  condition forms \code{handler-bind} and \code{handler-case}
  (Chapter~\ref{ch:conditions}). Backquote, \code{quasiquote}, is one too.
\item[Definitions] create module variables: \code{define},
  \code{define-class}, \code{define-generic-function}, \code{define-method},
  \code{define-slot} and \code{define-module}.
\end{description}

\noindent \code{not}, by contrast, is an ordinary function. A syntax operator is
not a function and has no value: \code{(map and ...)} is meaningless, and
so is passing \code{if} as an argument.

The 1992 manual has no way for a program to define its own syntax forms.
The compiler to C adds Common Lisp--style macros, defined with
\code{define-macro} and expanded when compiling (Chapter~\ref{ch:extensions}); the
Common Lisp implementation does not have them.

\subsection{Multiple values}

An expression can return one value, several, or none. \code{values} returns
its arguments as multiple values, and \code{bind} receives them
(Section~\ref{sec:control-values}). Where only one value is wanted, as an argument
to a call, for instance, the first value is used, or \code{\#f} if there are
none.

\begin{dylan}
(values 1 2 3)             ; => 1, 2, 3
(print (list (values 1 2) 3))   ; prints: (1 3)
(print (list (values)))         ; prints: (#f)
\end{dylan}

\section{Program structure}\label{sec:syntax-program}

\subsection{Top-level forms and definitions}

A form that is not part of any other form is a \term{top-level form}; so are
the forms directly inside a \code{begin} that is itself at top level.
A program is a sequence of top-level forms. Some are definitions, which
create module variables; the rest are expressions, evaluated for their
effects.

\begin{center}
\begin{tabular}{@{}lp{0.55\textwidth}@{}}
\toprule
Definition & Creates \\
\midrule
\code{define} & a variable with an initial value (Section~\ref{sec:control-define}) \\
\code{define-method} & a method, added to a generic function (Section~\ref{sec:functions-define-method}) \\
\code{define-generic-function} & a generic function (Section~\ref{sec:functions-define-generic}) \\
\code{define-class} & a class, its slots and their accessors (Section~\ref{sec:classes-define-class}) \\
\code{define-slot} & a slot of an existing class or object (Section~\ref{sec:classes-define-slot}) \\
\code{define-module} & a module (Chapter~\ref{ch:modules}) \\
\bottomrule
\end{tabular}
\end{center}

\noindent The compiler to C adds three more definitions, which are
described in Chapter~\ref{ch:extensions}: \code{define-macro},
\code{define-for-macros} and \code{define-thread-variable}.

The 1992 manual describes definitions as if they were ordinary forms that
could appear anywhere, and both implementations accept \code{define} inside
a method, where it still defines a module variable. Later design notes
(``Definitions are Declarative'') restrict definitions to top level, and
that is the only place they belong: a definition inside a method happens
again each time the method runs, and in the compiler to C a
\code{define-method} that refers to local variables cannot be compiled:
it stops with \code{define-method: the method refers to local variables
(k)}, naming the variables.
Put definitions at top level.

\subsection{When things happen}\label{sec:syntax-when}

The two implementations agree on what a program means, but not on when a
definition takes effect, because one runs a static program and the other
is interactive.

The \textbf{Common Lisp implementation} evaluates each top-level form when
it reaches it, at the prompt or in a file. A definition takes effect then,
and not before; a later form may redefine a variable, which only draws a
warning, or redefine a class or add methods to a generic function that is
already in use.

The \textbf{compiler to C} compiles a whole program at once. Every class,
generic function and method the program defines exists before any of its
expressions run, wherever in the program it is defined. The top-level
expressions, and the initial values of \code{define}, are then evaluated in
order, from the first file to the last. A method may therefore be called
by an expression that comes before its definition in the source, but a
variable defined with \code{define} has no value until its definition has
been evaluated.

%! c-only
\begin{dylan}
(later 2)                   ; => 20
(define-method later (x) (* x 10))
later-var                   ; => error: later-var has no value yet
(define later-var 5)
\end{dylan}

\noindent In the Common Lisp implementation the first expression is an error
too, since \code{later} is not yet defined when it is evaluated.

A static program is also stricter about definitions. In the compiler to C a
class or generic function may be defined only once, a name cannot be
defined both by \code{define} and as a generic function, and a variable
that holds a class or a generic function cannot be assigned. These are
errors when compiling. Chapter~\ref{ch:dylanc} lists the full set of restrictions.

\subsection{Source files and the module: header}\label{sec:syntax-files}

A Dylan source file conventionally has the extension \file{.dyl}. It may
begin with a \term{header} (design note 33): lines of the form
\meta{keyword}\code{:} \meta{value}, each starting at the beginning of a
line, ended by a blank line. The keyword \code{module:} names the module
the file's forms belong to; any other keywords, such as \code{author:},
\code{copyright:} and \code{version:}, are accepted and ignored. A file
without a header belongs to the module \code{dylan-user}.

%! skip
\begin{dylan}
module: shapes
author: A. Programmer
version: 1.0

(define-method greet ()
  (print "hello from shapes"))
(greet)
\end{dylan}

\noindent The header is not Dylan syntax: it is read before the reader
starts, and design note 33 does not allow comments in it. The module it names must be
defined by a \code{define-module} form, which can be in an earlier file of
the program; Chapter~\ref{ch:modules} describes modules. Both implementations
honor the header when they compile a file: the compiler to C for each file
on its command line, and the Common Lisp implementation in
\code{dylan:compile-file}. At the interactive prompt, \code{(dylan:in-module
\meta{name})} sets the module instead (Chapter~\ref{ch:lisp}).

A program for the compiler to C may be spread over several files, given on
the command line in order: \cmd{bin/dylanc -o app a.dyl b.dyl}. They form
one program, and their top-level expressions run in that order
(Chapter~\ref{ch:dylanc}).

\chapter{Variables and Control}\label{ch:control}

This chapter covers the forms a program uses to name values and to direct
the flow of evaluation: defining and binding variables, assigning them,
choosing between alternatives, repeating, returning several values, and
leaving a computation early. They are the subject of
\dylanbook{chapters 4, 7 and 8}. Most of them are syntax forms, special
forms or macros (Section~\ref{sec:syntax-forms}); for each, this chapter says which
of its parts are evaluated, in what order, and what it returns.

\section{Module variables: define}\label{sec:control-define}

\index{define@\code{define}}
\begin{quote}
\code{(define \meta{variable-name} \meta{initial-value})}
\end{quote}

\noindent \code{define} creates a module variable named
\meta{variable-name} in the current module and stores in it the value of
\meta{initial-value}. The name is not evaluated; it must be a symbol or a
setter name \code{(setter \meta{symbol})}. The initial value is evaluated,
and if it returns several values, the first is used. At the interactive
prompt \code{define} returns the variable's name.

\begin{dylan}
(define foo 10)
(define bar foo)
(+ foo bar)               ; => 20
(define (setter my-place) 0)
(setter my-place)         ; => 0
(define v (values 1 2))
v                         ; => 1
\end{dylan}

\noindent A module variable can be read from anywhere in its module, by
code before or after its definition, though it has no value until the
definition has been evaluated (Section~\ref{sec:syntax-when}). A variable created by
\code{define} may be assigned with \code{set!}
(Section~\ref{sec:control-assignment}).

Defining a variable that already exists in the module is allowed by the
Common Lisp implementation, which warns \code{redefining variable} and
stores the new value; in the compiler to C, defining a variable twice with
\code{define} simply stores twice, but defining a name both with
\code{define} and as a class or generic function is an error when
compiling.

\begin{notebox}
Later design notes changed \code{define}: one renamed it
\code{define-constant} and added \code{define-variable}, another let it
declare types and receive multiple values, as \code{bind} does. None of
these is part of this language. There is no \code{define-variable} or
\code{define-constant}, and \code{define} takes exactly one plain name and
one expression; \code{(define (x <integer>) 3)} and \code{(define a b
(values 1 2))} are errors in both implementations. The compiler to C adds
\code{define-thread-variable}, which defines a variable each thread has its
own value of (Chapter~\ref{ch:extensions}).
\end{notebox}

\section{Lexical variables: bind}\label{sec:control-bind}

\index{bind@\code{bind}}
\begin{quote}
\code{(bind (\meta{binding} \ldots) \meta{form} \ldots)}
\end{quote}

\noindent \code{bind} creates lexical variables and evaluates the
\meta{form}s, its \term{body}, with them in scope. It returns the values of
the last form, or \code{\#f} if there are no forms. Each \meta{binding} is
a list of one or more variables followed by an expression, the
\meta{init}:

\begin{quote}
\code{(\meta{variable} \ldots{} [\#rest \meta{rest-variable}] \meta{init})}
\end{quote}

\noindent where each \meta{variable} is a name, or a list \code{(\meta{name}
\meta{type})} that declares the variable's type.

The bindings are made one after another. Each init is evaluated in an
environment that includes the variables of the bindings before it, so a
later init can use an earlier variable:

\begin{dylan}
(bind ((x 20)
       (y (+ x x)))
  (+ y y))                ; => 80
(bind ((x 1)))            ; => #f
\end{dylan}

\noindent A variable made by \code{bind} shadows any module variable or
outer lexical variable of the same name, for the extent of the body:

\begin{dylan}
(define foo 10)
(bind ((foo 35))
  (+ foo foo))            ; => 70
(bind ((foo 20))
  (bind ((foo 50))
    (+ foo foo)))         ; => 100
foo                       ; => 10
\end{dylan}

\subsection{Receiving multiple values}

When a binding has several variables, they receive the values of its init
in order. Missing values are \code{\#f}, and values left over are ignored,
unless the variables end with \code{\#rest} and a name, in which case that
variable receives a sequence of the remaining values.

%! print
\begin{dylan}
(bind ((a b c (values 1 2 3)))
  (list a b c))           ; => (1 2 3)
(bind ((a b (values 1)))
  (list a b))             ; => (1 #f)
(bind ((first #rest others (values 1 2 3)))
  others)                 ; => (2 3)
(bind ((#rest all (values)))
  (empty? all))           ; => #t
\end{dylan}

\noindent Multiple values also give parallel binding, like \code{let} in
Lisp, since the init is evaluated before any of its variables exist:

%! print
\begin{dylan}
(bind ((x 10) (y 20))
  (bind ((x y (values y x)))
    (list x y)))          ; => (20 10)
\end{dylan}

\subsection{Type declarations}\label{sec:control-bind-types}

A variable written \code{(\meta{name} \meta{type})} must receive an
instance of \meta{type}, or an error is signaled. The type is an
expression, evaluated before the binding's init; it must evaluate to a class
or a singleton (Section~\ref{sec:functions-singletons}).

\begin{dylan}
(bind (((n <integer>) (+ 1 2)))
  n)                      ; => 3
(bind (((x <integer>) (sqrt 2)))
  x)                      ; => error: 1.4142135623730951 is not
                          ;    an instance of <integer>
\end{dylan}

\noindent The type is checked when the variable is bound. It is not checked
again when the variable is assigned: in both implementations \code{(bind
(((x <integer>) 1)) (set! x 'a) x)} returns \code{a}. (A later design note
made the restriction apply to assignment too; the 1992 manual does not.)
The compiler to C uses declared types to choose faster code
(Chapter~\ref{ch:performance}).

\section{Assignment: set!}\label{sec:control-assignment}

\index{set"!@\code{set"!}}
\begin{quote}
\code{(set! \meta{place} \meta{new-value})}
\end{quote}

\noindent \code{set!} stores a new value in a place and returns the new
value (\dylanbook{chapter 7}). The place is either a variable name or a
call-like list.

When \meta{place} is a variable name, lexical or module, the variable is
given the value of \meta{new-value}. It is an error if there is no such
variable.

\begin{dylan}
(define counter 0)
(set! counter (+ counter 1))   ; => 1
counter                        ; => 1
(bind ((x 1))
  (set! x (* x 10))
  x)                           ; => 10
\end{dylan}

\noindent Some module variables cannot be assigned. A variable made by
\code{define-class} is read-only in both implementations. In the compiler
to C a variable holding a generic function cannot be assigned either, nor
can any variable of the \code{dylan} module; these are errors when
compiling. The Common Lisp implementation allows assigning a generic
function's variable, as the 1992 manual does.

\subsection{The extended form}\label{sec:control-extended-set}

When \meta{place} is a list \code{(\meta{name} \meta{arg} \ldots)}, it
stands for a location that the call \code{(\meta{name} \meta{arg} \ldots)}
would read, and \code{set!} stores into it by calling the setter function:

\begin{quote}
\code{(set! (\meta{name} \meta{arg} \ldots) \meta{new-value})}\\
means\\
\code{((setter \meta{name}) \meta{arg} \ldots{} \meta{new-value})}
\end{quote}

\noindent The new value comes \emph{last} among the setter's arguments.
Whatever the setter itself returns, \code{set!} returns the new value.

\begin{dylan}
(define v (vector 'a 'b 'c))
(set! (element v 1) 'sea)       ; => sea
v                               ; => #(a sea c)
((setter element) v 0 'x)       ; => x
v                               ; => #(x sea c)
\end{dylan}

\noindent The arguments and the new value are each evaluated once. The
compiler to C evaluates the arguments first, left to right, and the new
value last, the order of the call \code{set!} stands for; the Common Lisp
implementation evaluates the new value first. Slot accessors
(Chapter~\ref{ch:classes}) and \code{element} (Chapter~\ref{ch:collections}) are the
setters used most, and a program can define its own
(Section~\ref{sec:functions-setters}).

\section{Truth and if}\label{sec:control-if}

\code{\#f} is the only false value; every other object is true
(Section~\ref{sec:syntax-booleans}). A function described as ``returning true'' may
return any object other than \code{\#f}, not necessarily \code{\#t}.

\index{if@\code{if}}
\begin{quote}
\code{(if \meta{test} \meta{consequent} \meta{alternate})}
\end{quote}

\noindent \code{if} evaluates \meta{test}. If its value is true, \code{if}
evaluates \meta{consequent} and returns its values; otherwise it evaluates
\meta{alternate} and returns its values. Exactly three operands are
required: an \code{if} without an alternate, or with more than one form in
either branch, is an error in both implementations. Use \code{when} or
\code{unless} for a one-armed test, and \code{begin} to put several forms in
a branch.

\begin{dylan}
(if (< -11 0) 'negative 'non-negative)   ; => negative
(if 0 'yes 'no)                          ; => yes
(if '() 'yes 'no)                        ; => yes
(if #f 1 (values 4 5))                   ; => 4, 5
\end{dylan}

\index{not@\code{not}}
\code{(not \meta{thing})} is a function: it returns \code{\#t} if
\meta{thing} is \code{\#f}, and \code{\#f} otherwise.

\begin{dylan}
(not #f)                  ; => #t
(not 3)                   ; => #f
\end{dylan}

\section{Sequencing: begin}

\index{begin@\code{begin}}
\code{(begin \meta{form} \ldots)} evaluates the forms in order and returns
the values of the last one, or \code{\#f} if there are none. Many forms,
such as \code{bind}, \code{method} and \code{when}, already contain an
implicit \code{begin} around their bodies; an explicit one is needed where
a single form is expected, as in a branch of \code{if}.

\begin{dylan}
(if #t
    (begin (print "it was true")
           #t)
    #f)                  ; prints: "it was true"
                         ; => #t
(begin)                  ; => #f
\end{dylan}

\noindent A \code{begin} at top level makes its forms top-level forms too
(Section~\ref{sec:syntax-program}).

\section{Conditionals}\label{sec:control-conditionals}

\subsection{when and unless}

\index{when@\code{when}}\index{unless@\code{unless}}
\begin{quote}
\code{(when \meta{test} \meta{form} \ldots)}\\
\code{(unless \meta{test} \meta{form} \ldots)}
\end{quote}

\noindent \code{when} evaluates \meta{test}; if it is true, it evaluates
the forms and returns the values of the last. \code{unless} evaluates the
forms if \meta{test} is false. When the forms are not evaluated, or there
are none, the result is \code{\#f}.

\begin{dylan}
(when (> 3 2) 'big)      ; => big
(when (> 2 3) 'big)      ; => #f
(unless (> 2 3) 'a 'b)   ; => b
(when #t)                ; => #f
\end{dylan}

\subsection{and and or}

\index{and@\code{and}}\index{or@\code{or}}
\begin{quote}
\code{(and \meta{form} \ldots)}\\
\code{(or \meta{form} \ldots)}
\end{quote}

\noindent \code{and} evaluates the forms from left to right until one
returns \code{\#f} as its first value; then it returns \code{\#f} and
evaluates no more. If none does, it returns all the values of the last
form. \code{or} evaluates the forms from left to right until one returns a
true first value, and returns that value alone; if every form but the last
is false, it returns all the values of the last. With no forms, \code{and}
returns \code{\#t} and \code{or} returns \code{\#f}.

\begin{dylan}
(and 1 2 3)              ; => 3
(and 1 #f (error "never evaluated"))   ; => #f
(and 1 (values 2 3))     ; => 2, 3
(or #f 'x 'y)            ; => x
(or (values #f 3) 4)     ; => 4
(or #f (values 1 2))     ; => 1, 2
(and)                    ; => #t
(or)                     ; => #f
\end{dylan}

\subsection{cond}

\index{cond@\code{cond}}
\begin{quote}
\code{(cond (\meta{test} \meta{consequent} \ldots) \ldots)}
\end{quote}

\noindent \code{cond} evaluates the tests in order. At the first that is
true, it evaluates that clause's consequents and returns the values of the
last; later tests are not evaluated. A clause with no consequents returns
the value of its test. If no test is true, \code{cond} returns \code{\#f}.
By convention the last clause's test is \code{else:} or \code{\#t}; both are
simply true values.

\begin{dylan}
(define-method compare (new old)
  (cond ((< new old) "the new position is less")
        ((= new old) "the positions are equal")
        (else: "the new position is greater")))
(compare 1 2)            ; => "the new position is less"
(compare 3 2)            ; => "the new position is greater"
(cond (#f 'no) ((+ 1 2)))     ; => 3
(cond (#f 1))            ; => #f
\end{dylan}

\subsection{case}

\index{case@\code{case}}
\begin{quote}
\code{(case \meta{target} (\meta{match-list} \meta{consequent} \ldots) \ldots)}
\end{quote}

\noindent \code{case} evaluates \meta{target} and looks for its value in the
match lists, comparing with \code{id?}. The match lists are lists of
objects and are not evaluated. At the first clause whose list contains the
value, \code{case} evaluates that clause's consequents and returns the
values of the last, or \code{\#f} if it has none. A clause whose match list
is \code{else:} or \code{\#t} matches when no other does; it must be the last
clause. If nothing matches and there is no such clause, an error is
signaled.

\begin{dylan}
(define-method size-name (n)
  (case n
    ((0) 'none)
    ((1 2 3) 'few)
    (else: 'many)))
(size-name 2)            ; => few
(size-name 10)           ; => many
(case 'Red ((red green) 'warm-or-cool) (else: 'other))   ; => warm-or-cool
(case 9 ((3) 'three))    ; => error: no clause matches 9
\end{dylan}

\noindent Because the comparison is \code{id?}, matches are normally
numbers, characters, symbols, keywords, \code{\#t}, \code{\#f} or
\code{()}. Strings and lists are not suitable: two strings with the same
characters need not be the same object. The same match may not appear in
two clauses; the compiler to C reports that when compiling, and the Common
Lisp implementation uses the first.

\subsection{select}

\index{select@\code{select}}
\begin{quote}
\code{(select \meta{target} \meta{test} (\meta{match-list} \meta{consequent} \ldots) \ldots)}
\end{quote}

\noindent \code{select} is like \code{case}, but the matches are expressions,
evaluated, and the comparison is the function \meta{test}, called as
\code{(\meta{test} \meta{target-value} \meta{match})}. The clauses are tried
in order, and within a clause the matches from left to right; the matches
are evaluated only as they are reached, and once one matches, no more are
evaluated or tested (design note 34). \code{else:} and \code{\#t} work as in
\code{case}, and no match without them is an error.

\begin{dylan}
(define-method kind-of (x)
  (select x instance?
    ((<integer> <ratio>) "exact")
    ((<float>) "inexact")
    (else: "not a number")))
(kind-of 3)              ; => "exact"
(kind-of 2.5)            ; => "inexact"
(kind-of 'x)             ; => "not a number"
(select "b" = (("a") 1) (("b") 2))   ; => 2
(select 5 < ((3 4) 'small) ((6) 'medium) (else: 'large))   ; => medium
\end{dylan}

\noindent With a test of \code{instance?}, as in \code{kind-of},
\code{select} does what \code{typecase} does in Common Lisp.

\section{Iteration}\label{sec:control-loops}

The loops below are the language's iteration forms. Collections can also
be walked with functions such as \code{do} and \code{map}
(Chapters~\ref{ch:collections} and~\ref{ch:functional}), and any loop can be
written as a recursive method, since tail calls do not use stack
(Section~\ref{sec:functions-tail-calls}).

\subsection{for}

\index{for@\code{for}}
\begin{quote}
\code{(for ((\meta{var} \meta{init} \meta{step}) \ldots)\\
\phantom{(for }(\meta{end-test} \meta{result} \ldots)\\
\phantom{(for }\meta{body-form} \ldots)}
\end{quote}

\noindent \code{for} is the general loop. Each \code{(\meta{var}
\meta{init} \meta{step})} is a \term{variable clause}, and it is the only
kind of clause: each must have all three parts. The loop runs like this:

\begin{enumerate}
\item The inits are evaluated, in a scope that does not include any of the
  loop variables, and the variables are bound to their values.
\item The end test is evaluated. If it is true, the results are evaluated
  in order and \code{for} returns the values of the last, or \code{\#f} if
  there are none.
\item Otherwise the body forms are evaluated, for their effects.
\item The steps are evaluated, in a scope that includes the variables'
  current values, and the variables are bound afresh to the new values.
  The loop continues at step 2.
\end{enumerate}

\noindent The end-test clause may be the empty list \code{()}, meaning no
end test: the loop then runs until something exits from it
(Section~\ref{sec:control-exits}). Because all the steps are evaluated before any
variable is rebound, the variables are stepped in parallel.

\begin{dylan}
(define-method factorial ((n <integer>))
  (for ((i n (- i 1))
        (v 1 (* v i)))
       ((<= i 0) v)))
(factorial 5)            ; => 120
(for ((i 0 (+ i 1))
      (acc '() (cons i acc)))
     ((= i 4) acc))      ; => (3 2 1 0)
(for ((a 0 b) (b 1 (+ a b)) (i 0 (+ i 1)))
     ((= i 10) a))       ; => 55
\end{dylan}

\noindent Each pass binds the variables anew; it does not assign them. A
method made in the body therefore keeps the values of that pass:

\begin{dylan}
(define thunks
  (for ((i 0 (+ i 1))
        (acc '() (cons (method () i) acc)))
       ((= i 3) acc)))
(map (method (f) (f)) thunks)   ; => (2 1 0)
\end{dylan}

\subsection{for-each}

\index{for-each@\code{for-each}}
\begin{quote}
\code{(for-each ((\meta{var} \meta{collection}) \ldots)\\
\phantom{(for-each }(\meta{end-test} \meta{result} \ldots)\\
\phantom{(for-each }\meta{body-form} \ldots)}
\end{quote}

\noindent \code{for-each} iterates over the elements of one or more
collections at once. The collection expressions are evaluated first. Then,
on each pass, each \meta{var} is bound to the next element of its
collection; if any collection has no more elements, the loop ends and
returns \code{\#f}. Otherwise the end test is evaluated, with the variables
in scope; if it is true, the results are evaluated and the values of the
last returned (\code{\#f} if none). Otherwise the body forms are evaluated
and the next pass begins. As with \code{for}, the end-test clause may be
\code{()}, and each pass binds the variables anew.

\begin{dylan}
(define-method first-even ((s <sequence>))
  (for-each ((number s))
            ((even? number) number)))
(first-even '(1 3 5 4 7 9 10))   ; => 4
(first-even '(1 3))              ; => #f
(for-each ((c "abc")
           (i (range from: 0)))
          ()
  (print (list c i)))            ; prints: (#\a 0)
                                 ; prints: (#\b 1)
                                 ; prints: (#\c 2)
                                 ; => #f
\end{dylan}

\noindent The second example stops after three passes because the string
runs out, although the range \code{(range from: 0)} is unbounded. Any
collection may be used: lists, vectors, strings, ranges, tables
(whose elements are its values) and the program's own collection classes.
\code{for-each} uses the iteration protocol of Chapter~\ref{ch:collections}.

\subsection{dotimes}

\index{dotimes@\code{dotimes}}
\begin{quote}
\code{(dotimes (\meta{var} \meta{count} [\meta{result}]) \meta{form} \ldots)}
\end{quote}

\noindent \code{dotimes} evaluates \meta{count}, which should be an
integer, then evaluates the forms once for each integer from 0 up to but not
including the count, with \meta{var} bound to it. If the count is zero or
negative, the forms are not evaluated. Then \meta{result} is evaluated and
its values returned; without a result, \code{dotimes} returns \code{\#f}.
In both implementations the result form can see \meta{var}, bound to the
number of passes made (as in Common Lisp).

\begin{dylan}
(dotimes (i 3)
  (print i))             ; prints: 0
                         ; prints: 1
                         ; prints: 2
                         ; => #f
(dotimes (i 4 i))        ; => 4
(dotimes (i -2 'done))   ; => done
\end{dylan}

\subsection{while and until}

\index{while@\code{while}}\index{until@\code{until}}
\begin{quote}
\code{(while \meta{test} \meta{form} \ldots)}\\
\code{(until \meta{test} \meta{form} \ldots)}
\end{quote}

\noindent \code{while} evaluates \meta{test} before each pass and evaluates
the forms as long as it is true. \code{until} evaluates the forms as long as
\meta{test} is false. Both always return \code{\#f}.

\begin{dylan}
(define n 0)
(while (< n 3)
  (set! n (+ n 1)))      ; => #f
n                        ; => 3
(until (= n 0)
  (set! n (- n 1)))      ; => #f
n                        ; => 0
\end{dylan}

\section{Local methods: bind-methods}\label{sec:control-bind-methods}

\index{bind-methods@\code{bind-methods}}
\begin{quote}
\code{(bind-methods ((\meta{name} \meta{parameter-list} \meta{form} \ldots) \ldots)\\
\phantom{(bind-methods }\meta{body-form} \ldots)}
\end{quote}

\noindent \code{bind-methods} binds each \meta{name} to a method, made as
\code{method} makes one (Section~\ref{sec:functions-method}), and evaluates the
body forms, returning the values of the last. The scope of the names
includes the methods themselves, so the methods can call themselves and each
other. (\code{bind} cannot do that, since an init cannot see its own
variable.)

\begin{dylan}
(define-method root-mean-square ((s <sequence>))
  (bind-methods ((average (numbers)
                   (/ (reduce1 + numbers) (size numbers)))
                 (square (n) (* n n)))
    (sqrt (average (map square s)))))
(root-mean-square '(5 6 6 7 4))   ; => 5.692099788303083
\end{dylan}

\noindent The next example, Newton's method for square roots, has a method
that calls itself and two helpers:

\begin{dylan}
(define-method newtons-sqrt (x)
  (bind-methods ((sqrt1 (guess)
                   (if (close? guess)
                       guess
                       (sqrt1 (improve guess))))
                 (close? (guess)
                   (< (abs (- (* guess guess) x)) .0001))
                 (improve (guess)
                   (/ (+ guess (/ x guess)) 2)))
    (sqrt1 1)))
(newtons-sqrt 25.0)      ; => 5.000000000053722
\end{dylan}

\noindent The 1992 manual calls \code{(newtons-sqrt 25)}. With an integer
argument the arithmetic stays exact, and the guesses become ratios with
ever larger numerators and denominators: the Common Lisp implementation
returns an enormous ratio, and in the compiler to C, which has no bignums,
the integers overflow and an error is signaled. A float argument keeps the
computation in floating point.

Later versions of Dylan have a \code{local} form for this purpose; this
language has only \code{bind-methods}.

\section{Multiple values}\label{sec:control-values}

\index{values@\code{values}}\index{multiple values}
\begin{quote}
\code{(values \meta{object} \ldots)}
\end{quote}

\noindent An expression may return any number of values. The function
\code{values} returns its arguments as multiple values; \code{(values)}
returns none. \code{bind} (Section~\ref{sec:control-bind}) receives them.

\begin{dylan}
(define-method edges ((center <number>) (radius <number>))
  (values (- center radius) (+ center radius)))
(edges 100 2)            ; => 98, 102
(define-method opposite-edges ((center <number>) (radius <number>))
  (bind ((min max (edges center radius)))
    (values max min)))
(opposite-edges 100 2)   ; => 102, 98
(bind ((#rest nums (edges 100 2)))
  nums)                  ; => (98 102)
\end{dylan}

\noindent Multiple values pass through the forms that return the values
of a subform: \code{begin}, \code{bind}, \code{if}, \code{when},
\code{unless}, \code{cond}, \code{case}, \code{select}, the last form of
\code{and} and \code{or} as described above, the results of \code{for},
\code{for-each} and \code{dotimes}, \code{bind-exit}, \code{unwind-protect}'s
protected form, and calls of methods whose last form returns them. Anywhere
else, as an argument to a function, as the init of a one-variable binding,
as the value in a \code{set!} or a \code{define}, only the first value is
used, and an expression that returns no values supplies \code{\#f}.

\begin{dylan}
(begin 1 (values 2 3))   ; => 2, 3
(print (list (values 1 2) 3))   ; prints: (1 3)
(bind ((x (values))) x)  ; => #f
(values)                 ; => (no values)
\end{dylan}

\section{Non-local exits}\label{sec:control-exits}

\subsection{bind-exit}

\index{bind-exit@\code{bind-exit}}
\begin{quote}
\code{(bind-exit (\meta{exit-variable}) \meta{form} \ldots)}
\end{quote}

\noindent \code{bind-exit} binds \meta{exit-variable} to an \term{exit
procedure} and evaluates the forms. If the exit procedure is not called,
\code{bind-exit} returns the values of the last form, or \code{\#f} if
there are none. Calling the exit procedure, with any number of arguments,
leaves the \code{bind-exit} at once, and the arguments become its values.
This is a \term{non-local exit}.

\begin{dylan}
(bind-exit (return)
  1
  (return 2 3)
  4)                     ; => 2, 3
(bind-exit (return))     ; => #f
\end{dylan}

\noindent The exit procedure is an ordinary function: it can be passed to
other functions and stored. But it is valid only during the
\term{dynamic extent} of the \code{bind-exit} that made it, that is, until
that \code{bind-exit} returns or is itself exited. Calling it afterwards is
an error.

\begin{dylan}
(define-method first-even ((seq <sequence>))
  (bind-exit (exit)
    (do (method (item)
          (when (even? item)
            (exit item)))
        seq)
    #f))
(first-even '(1 3 5 4 7 9 10))   ; => 4
\end{dylan}

\noindent In the compiler to C a non-local exit is a jump and costs little;
\code{bind-exit} is the normal way to return early from a loop or a method.

\subsection{unwind-protect}

\index{unwind-protect@\code{unwind-protect}}
\begin{quote}
\code{(unwind-protect \meta{protected-form} \meta{cleanup-form} \ldots)}
\end{quote}

\noindent \code{unwind-protect} evaluates \meta{protected-form}, then the
cleanup forms. The cleanup forms are evaluated however the protected form
is left: normally, by a non-local exit through it, or by an error that a
handler outside it handles by exiting (Chapter~\ref{ch:conditions}). When the
protected form returns normally, \code{unwind-protect} returns its values;
the cleanup forms' values are discarded.

\begin{dylan}
(unwind-protect (values 1 2)
  (print 'cleanup))      ; prints: cleanup
                         ; => 1, 2
(bind-exit (k)
  (unwind-protect (k 5)
    (print 'cleanup)))   ; prints: cleanup
                         ; => 5
(handler-case
    (unwind-protect (error "boom")
      (print 'cleanup))
  ((<error>) 'caught))   ; prints: cleanup
                         ; => caught
\end{dylan}

\noindent A cleanup form may itself exit. If it does, the new exit replaces
the one in progress, and the remaining cleanup forms of that
\code{unwind-protect} are skipped:

\begin{dylan}
(bind-exit (k)
  (unwind-protect (k 1)
    (k 2)))              ; => 2
(bind-exit (one)
  (bind-exit (two)
    (unwind-protect (two 2)
      (one 1)))
  3)                     ; => 1
\end{dylan}

\noindent These are the rules design note 25 later gave for its
\code{block} form, which replaced \code{bind-exit} and
\code{unwind-protect} in later Dylan; this language keeps the two forms of
the 1992 manual.

\section{quote}

\code{(quote \meta{object})}, usually written \code{'\meta{object}}, returns
\meta{object} without evaluating it (Section~\ref{sec:syntax-quote}). It is a
special form of its own, listed here for completeness with the other
evaluation forms of \dylanbook{chapter 8}. \code{apply}, the last function
of that chapter, calls a function on a list of arguments and is described
with functions (Section~\ref{sec:functions-apply}).

\begin{dylan}
'+                       ; => +
(quote (+ 10 10))        ; => (+ 10 10)
\end{dylan}

\chapter{Functions and Methods}\label{ch:functions}

Functions are Dylan's means of doing anything. This chapter describes the
two kinds of function, methods and generic functions; how a method's
parameter list says what arguments it takes; how a generic function
chooses which of its methods to run; and the operations that call
functions and examine them. The standard is \dylanbook{chapters 6
and 10}.

\section{Methods and generic functions}\label{sec:functions-overview}

Every function is an instance of the class
\code{<function>}\index{<function>@\code{<function>}}, and there are two
kinds.

A \term{method} (class \code{<method>}) is a piece of code with a parameter
list. The parameter list may give a \term{specializer} for each required
parameter, a class or a singleton that the corresponding argument must
match. A method can be called directly, like a procedure in Scheme or a
closure in Lisp; called with arguments that do not match its specializers,
it signals an error.

A \term{generic function} (class \code{<generic-function>}) holds a set of
methods. When it is called, it finds the methods whose specializers match
the arguments, the \term{applicable} methods, and runs the most specific of
them. This choice, \term{method dispatch}, is how Dylan does
object-oriented programming: behavior is attached to classes by defining
methods on generic functions, not by putting functions inside classes.

\begin{dylan}
(define-method double ((thing <number>))
  (+ thing thing))
(define-method double ((thing <sequence>))
  (concatenate thing thing))
(double 10)                   ; => 20
(double "the rain ")          ; => "the rain the rain "
(double '(a b c))             ; => (a b c a b c)
(double 'cup)                 ; => error: no applicable method
(instance? double <generic-function>)          ; => #t
(instance? (method (x) x) <method>)            ; => #t
(instance? double <function>)                  ; => #t
\end{dylan}

\noindent Functions are ordinary objects. They can be stored in variables
and data structures, passed as arguments and returned as results.

\section{Making methods: method}\label{sec:functions-method}

\index{method@\code{method}}
\begin{quote}
\code{(method \meta{parameter-list} \meta{form} \ldots)}
\end{quote}

\noindent The special form \code{method} makes and returns a method. When
the method is called, its parameters are bound to the arguments as lexical
variables and the forms, its \term{body}, are evaluated in order; the
method returns the values of the last form, or \code{\#f} if there are none.

\begin{dylan}
((method (a b) (+ a b)) 3 4)             ; => 7
(bind ((square (method (n) (* n n))))
  (square (square 3)))                   ; => 81
(sort '(3 1 2) test: (method (a b) (> a b)))   ; => (3 2 1)
\end{dylan}

\noindent A method made by \code{method} is an \term{anonymous method}. It
can be given to a function, as the test to \code{sort} above, or kept in a
variable. To give a group of local methods names that can call one another,
use \code{bind-methods} (Section~\ref{sec:control-bind-methods}).

\subsection{Closures}\label{sec:functions-closures}

A method can refer to the variables in scope where it was made, and it
keeps them for as long as it exists, even after the form that bound them
has returned. Such a method is a \term{closure}. The variables are shared,
not copied: assignments made through one closure are seen by every other
closure over the same variable.

\begin{dylan}
(define-method make-counter ()
  (bind ((count 0))
    (values (method () (set! count (+ count 1)))
            (method () count))))
(define-method demo ()
  (bind ((next current (make-counter)))
    (next)
    (next)
    (list (next) (current))))
(demo)                         ; => (3 3)
\end{dylan}

\noindent A method that returns a method is a common idiom. The method below
takes a function and returns one that calls it twice:

\begin{dylan}
(define-method twice ((f <function>))
  (method (#rest args)
    (apply f args)
    (apply f args)))
(define print-twice (twice (method (n) (print n))))
(print-twice 55)                 ; prints: 55
                               ; prints: 55
                               ; => 55
\end{dylan}

\section{Parameter lists}\label{sec:functions-parameters}

The complete syntax of a parameter list (\dylanbook{chapter 6}) is

\begin{quote}
\code{(\{\meta{required} | (\meta{required} \meta{specializer})\}*\\
\phantom{(}[\#next \meta{next}]\\
\phantom{(}[\#rest \meta{rest}]\\
\phantom{(}[\#key \{\meta{keyparam} | ([\meta{keyword}] \meta{keyparam} [\meta{default}])\}*])}
\end{quote}

\noindent The parts must appear in this order. \code{\#next}, \code{\#rest}
and \code{\#key} may also be written \code{!next}, \code{!rest} and
\code{!key} (Section~\ref{sec:syntax-parameter-tokens}).

\begin{center}
\begin{tabular}{@{}ll@{}}
\toprule
Parameter list & Accepts \\
\midrule
\code{(a b c)} & three required arguments \\
\code{((a <list>) b \#rest c)} & two required (the first a list), then any number \\
\code{(\#rest a)} & any number of arguments \\
\code{(a b \#key foo bar)} & two required, then keywords \code{foo:} and \code{bar:} \\
\code{(\#key bar baz bim)} & keywords only \\
\code{(\#rest all \#key fee fi)} & keyword arguments, all kept in \code{all} \\
\bottomrule
\end{tabular}
\end{center}

\subsection{Required parameters}

Required parameters come first, and a call must supply an argument for
each, in order. A required parameter is a name, or a list of a name and a
specializer: \code{(x <integer>)}. A parameter without a specializer
accepts any object, as if specialized on \code{<object>}. Specializers are
described in Section~\ref{sec:functions-singletons}.

Calling a method with too few or too many arguments, or with an argument
that does not match its specializer, is an error:

\begin{dylan}
((method ((s <string>)) s) "x")   ; => "x"
((method ((s <string>)) s) 3)     ; => error: 3 does not match <string>
((method (a #rest r) a))          ; => error: too few arguments
\end{dylan}

\subsection{Rest parameters}

After \code{\#rest} comes one name. When the method is called, the
arguments left after the required ones are put into a sequence, which
becomes the value of the rest parameter; it is empty if there are none.

\begin{dylan}
(define-method show-rest (a #rest b)
  (list a b))
(show-rest 10 20 30 40)        ; => (10 (20 30 40))
(show-rest 10)                 ; => (10 ())
\end{dylan}

\noindent In both implementations the sequence is a list, but a program
should rely only on its being a sequence.

\subsection{Keyword parameters}\label{sec:functions-keywords}

After \code{\#key} come the \term{keyword parameters}. Keyword arguments are
optional, are given after the required arguments, and may come in any
order; each is a keyword followed by a value. A keyword parameter is
written in one of four ways:

\begin{center}
\begin{tabular}{@{}ll@{}}
\toprule
Written & Meaning \\
\midrule
\code{\meta{name}} & keyword \code{\meta{name}:}, parameter \meta{name}, default \code{\#f} \\
\code{(\meta{name} \meta{default})} & keyword \code{\meta{name}:}, parameter \meta{name} \\
\code{(\meta{keyword} \meta{param})} & the keyword and the parameter named separately \\
\code{(\meta{keyword} \meta{param} \meta{default})} & the same, with a default \\
\bottomrule
\end{tabular}
\end{center}

\noindent The \meta{default} is an expression evaluated each time the method
is called without that keyword. It is evaluated in a scope that includes all
the parameters before it: the required parameters, the next-method and
rest parameters, and the earlier keyword parameters. Without a default, an
unsupplied keyword parameter is \code{\#f}.

\begin{dylan}
(define-method percolate (#key (brand 'maxwell-house)
                               (cups 4)
                               (strength (if (> cups 6) 'strong 'mild)))
  (list brand cups strength))
(percolate)                          ; => (maxwell-house 4 mild)
(percolate cups: 10 brand: 'folgers) ; => (folgers 10 strong)
\end{dylan}

\noindent Naming the keyword and the parameter separately lets a method take
a keyword such as \code{position:} without calling its parameter
\code{position}, which would shadow the module variable of that name, often
a generic function the method needs:

\begin{dylan}
(define-method place (widget #key (position: the-pos 0) (size: the-size))
  (list widget the-pos the-size))
(place 'w size: 3)                   ; => (w 0 3)
\end{dylan}

\noindent All the required arguments are taken before any keyword argument.
If a call supplies too few, keywords and values are taken as required
arguments:

\begin{dylan}
(define-method show-keys (req1 req2 #key foo)
  (list req1 req2 foo))
(show-keys 'one 'two foo: 'three)    ; => (one two three)
(show-keys foo: 'three)              ; => (foo: three #f)
\end{dylan}

\noindent The arguments after the required ones must be keyword/value pairs;
anything else is an error. If the same keyword is given more than once, the
leftmost occurrence is used.

\begin{dylan}
(define-method grow ((x <number>) #key (by 1)) (+ x by))
(grow 5 by: 2 by: 100)               ; => 7
(grow 5 'by 2)                       ; => error: not keyword/value pairs
\end{dylan}

\subsection{Combining \#rest and \#key}

When a parameter list has both, \code{\#rest} comes first. The rest
parameter then receives all the keyword arguments, keywords and values, as a
sequence, and the method accepts any keyword
(Section~\ref{sec:functions-keyword-checking}); the arguments must still be
keyword/value pairs.

\begin{dylan}
(define-method test (the-req #rest the-rest #key a b)
  (list the-req the-rest a b))
(test 1 a: 2 b: 3 c: 4)              ; => (1 (a: 2 b: 3 c: 4) 2 3)
\end{dylan}

\subsection{The next-method parameter}

After the required parameters may come \code{\#next} and a name. That
parameter receives the \term{next method}, used to call less specific
methods (Section~\ref{sec:functions-next-method}). \code{define-method} gives every
method a next-method parameter named \code{next-method} unless the parameter
list names one with \code{\#next}, so it is normally written only to choose
another name, or in a method made with \code{method} that will be added to a
generic function.

\subsection{Keyword argument checking}\label{sec:functions-keyword-checking}

A method accepts only the keywords its parameter list names after
\code{\#key}, unless it has a rest parameter, in which case it accepts any
keyword. A generic function accepts, for a given call, every keyword that
any of the methods applicable to that call accepts; so if any applicable
method has \code{\#rest}, any keyword is accepted. The check is made by the
generic function before it runs any method, and a method called through a
generic function or through \code{next-method} does not check again.
Calling a method or generic function with a keyword it does not accept is an
error.

\begin{dylan}
(define-method label ((x <object>) #key price)
  (list price x))
(define-method label ((x <sequence>) #key unit-price)
  (add x (* unit-price (size x))))
(define-method label ((x <list>) #rest info #key calories)
  (add x calories))
(label 'grape price: 189)                       ; => (189 grape)
(label 'grape price: 189 unit-price: 2)         ; => error
(label (vector 3 4 5) price: 189 unit-price: 2) ; => #(6 3 4 5)
(label (vector 3 4 5) calories: 9)              ; => error
(label (list 3 4 5) protein: 7 calories: 9)     ; => (9 3 4 5)
\end{dylan}

\noindent The first \code{label} method is the only one applicable to a
symbol, so only \code{price:} is accepted. A vector is a sequence, so both
\code{price:} and \code{unit-price:} are. For a list the third method
applies, and it has \code{\#rest}, so anything goes.

The two implementations differ on one point. A method whose parameter list
has \code{\#key} with no keywords after it accepts no keywords in the
compiler to C, as the rule above says; in the Common Lisp implementation it
accepts any.

\section{Defining methods: define-method}\label{sec:functions-define-method}

\index{define-method@\code{define-method}}
\begin{quote}
\code{(define-method \meta{name} \meta{parameter-list} \meta{form} \ldots)}
\end{quote}

\noindent \code{define-method} makes a method, as \code{method} would, and
adds it to the generic function in the module variable \meta{name}. If the
variable does not exist, it is created, holding a new generic function. If
the generic function already has a method with exactly the same
specializers, the new method replaces it. \meta{name} may be a setter name,
\code{(setter \meta{symbol})} (Section~\ref{sec:functions-setters}). At the
interactive prompt \code{define-method} returns the name.

%! lisp-only
\begin{dylan}
(define-method area ((r <integer>)) (* r r))
(area 3)                       ; => 9
(define-method area ((r <integer>)) (* 3 r r))
(area 3)                       ; => 27
\end{dylan}

\noindent That is the Common Lisp implementation, where each definition
takes effect when it is evaluated. In the compiler to C all of a program's
methods exist before it starts (Section~\ref{sec:syntax-when}), so the later
definition simply wins: both calls of \code{area} above return 27.

\noindent It is an error for \meta{name} to hold something other than a
generic function. The compiler to C also refuses to add methods to the
functions the manual calls methods rather than generic functions, such as
\code{id?}, \code{not}, \code{list} and \code{apply}, and to the
reflective functions it implements in C.

\subsection{Congruent parameter lists}

All the methods of a generic function, and its
\code{define-generic-function} if it has one, must have \term{congruent}
parameter lists: the same number of required parameters, and either all of
them have \code{\#rest} or \code{\#key} (or both), or none has either.
Adding a method that is not congruent is an error, reported when compiling
by the compiler to C.

\begin{center}
\begin{tabular}{@{}l@{}}
\toprule
Groups of congruent parameter lists \\
\midrule
\code{(a b)} \quad \code{(x y)} \\
\code{(a b \#key)} \quad \code{(x y \#key foo)} \quad
  \code{(a b \#rest l)} \quad \code{(a b \#rest l \#key quux)} \\
\code{(a b c)} \quad \code{(d e f)} \\
\bottomrule
\end{tabular}
\end{center}

\noindent When \code{define-method} creates the generic function, the
generic function takes the same number of required arguments as the method,
and, if the method has \code{\#rest} or \code{\#key}, a rest argument.
\code{function-arguments} (Section~\ref{sec:functions-reflection}) shows this:

\begin{dylan}
(define-method foo (a #key x) x)
(function-arguments foo)       ; => 1, #t, #f
\end{dylan}

\noindent A generic function with no required parameters can hold only one
method, since all its methods have the same (empty) specializers; defining
another replaces the first.

\section{Declaring generic functions: define-generic-function}\label{sec:functions-define-generic}

\index{define-generic-function@\code{define-generic-function}}
\begin{quote}
\code{(define-generic-function \meta{name} \meta{parameter-list} \meta{keyword} \meta{option} \ldots)}
\end{quote}

\noindent \code{define-generic-function} creates a generic function with no
methods and stores it in the module variable \meta{name}, which is made
read-only. It is never required, since \code{define-method} creates generic
functions as needed, but it documents the protocol in one place and fixes
the parameter list that every method must be congruent with.

The parameter list has the form

\begin{quote}
\code{(\meta{required} \ldots{} [\#rest \meta{rest}] [\#key \meta{keyword} \ldots])}
\end{quote}

\noindent Required parameters are names only, without specializers. The
keys may be written as symbols or keywords, and \code{\#key} may appear with
no keys at all. \code{\#key} means that the arguments after the required
ones must be keyword/value pairs; \code{\#rest} means any number of further
arguments. Either way, every method must have \code{\#rest} or \code{\#key}.
The meaning of giving both is unspecified.

\begin{dylan}
(define-generic-function resize (shape #key scale))
(define-method resize ((n <integer>) #key scale)
  (* n (or scale 1)))
(resize 4 scale: 3)            ; => 12
(function-arguments resize)    ; => 1, #f, (scale:)
\end{dylan}

\noindent When keys are listed, every method is meant to accept them, by
naming them after \code{\#key} or by having \code{\#rest}. The Common Lisp
implementation checks this when the method is added; the compiler to C does
not.

The options are keyword/value pairs reserved by the manual for future use.
Both implementations accept them and ignore them. In the compiler to C a
generic function may be declared only once.

\section{Specializers and singletons}\label{sec:functions-singletons}

A specializer is an expression, evaluated once when the method is made, not
on each call. It must evaluate to a class or to a singleton.

\begin{itemize}
\item A \term{class} specializer matches any \emph{general instance} of the
  class: a direct instance, or an instance of any subclass. The specializer
  \code{<number>} matches integers, ratios and floats.
\item A \term{singleton} specializer matches one object, compared with
  \code{id?}. \code{(singleton \meta{object})} makes one; so does
  \code{(make <singleton> object: \meta{object})}.
\end{itemize}

\begin{dylan}
(define-method factorial ((n <integer>))
  (* n (factorial (- n 1))))
(define-method factorial ((n (singleton 0)))
  1)
(factorial 5)                  ; => 120
(define-method measure ((x (singleton 'cup))) 'pint)
(measure 'cup)                 ; => pint
(instance? 3 (singleton 3))    ; => #t
\end{dylan}

\noindent \index{singleton@\code{singleton}}\code{singleton} is a function;
\code{<singleton>}\index{<singleton>@\code{<singleton>}} is the class of
singletons, and a singleton prints as \code{\{singleton 3\}}. Because
singletons compare their objects with \code{id?}, they suit symbols,
keywords, characters, small integers, classes and other objects whose
identity is what matters; a singleton of a string literal matches only that
very string object.

Singletons can also give an individual object slots of its own
(Section~\ref{sec:classes-define-slot}).

\begin{warnbox}
In the compiler to C every specializer of a method added to a generic
function must be known when compiling: a class or a variable whose value is
a class, or \code{(singleton \meta{x})} or \code{(make <singleton> object:
\meta{x})} where \meta{x} is a literal, a class, or a variable whose
\code{define} gives it a literal value. \code{(singleton (+ 1 2))}, and a
singleton of an object made at run time, are errors there, though the
Common Lisp implementation accepts them.
\end{warnbox}

\section{Method dispatch}\label{sec:functions-dispatch}

When a generic function is called, it
\begin{enumerate}
\item finds the \emph{applicable} methods: those whose every specializer
  matches the corresponding argument;
\item sorts them from most to least specific;
\item calls the most specific, passing it the next one as its next method.
\end{enumerate}
\noindent If no method is applicable, an error is signaled.

\subsection{Specificity}\label{sec:functions-specificity}

For one argument, a singleton specializer is more specific than any class,
and a class is more specific than its superclasses. So a method on
\code{<integer>} overrides one on \code{<number>} for integer arguments, and
a method on \code{(singleton 5)} overrides both for the argument 5.

\begin{dylan}
(define-method describe ((x <number>)) 'number)
(define-method describe ((x <integer>)) 'integer)
(define-method describe ((x (singleton 5))) 'five)
(describe 2.5)                 ; => number
(describe 2)                   ; => integer
(describe 5)                   ; => five
\end{dylan}

\noindent Methods that specialize more than one parameter are
\term{multimethods}, and dispatch considers all the specialized arguments.
Following \dylanbook{chapter 10}, both implementations sort the applicable
methods by comparing their specializer lists \emph{lexicographically}: the
first parameter decides, and a later parameter is looked at only when the
earlier specializers are the same.

\begin{dylan}
(define-method meet ((a <number>) (b <integer>)) 'number-integer)
(define-method meet ((a <integer>) (b <number>)) 'integer-number)
(meet 1 2)                     ; => integer-number
\end{dylan}

\noindent Both methods apply to \code{(meet 1 2)}. The first parameter's
specializers differ, and \code{<integer>} is more specific than
\code{<number>}, so the second method wins, though the first is more
specific in the second parameter.

\subsection{The class precedence list}\label{sec:functions-cpl}

When a class has several direct superclasses, two specializers can both be
superclasses of an argument's class without either being a subclass of the
other. They are then ordered by the argument's class's \term{class
precedence list}: the class, then all its superclasses, in an order that
keeps each class before its superclasses, keeps the direct superclasses in
the order the class's definition lists them, and keeps each simple chain of
superclasses together. \code{all-superclasses} returns this list
(Section~\ref{sec:classes-reflection}). The algorithm is the one CLOS uses, given
in design note 31.

\begin{dylan}
(define-class <life-form> (<object>))
(define-class <sentient> (<life-form>))
(define-class <bipedal> (<life-form>))
(define-class <intelligent> (<sentient>))
(define-class <humanoid> (<bipedal>))
(define-class <vulcan> (<intelligent> <humanoid>))
(define-class <human> (<humanoid> <intelligent>))
(all-superclasses <vulcan>)
  ; => ({the class <vulcan>} {the class <intelligent>}
  ;     {the class <sentient>} {the class <humanoid>}
  ;     {the class <bipedal>} {the class <life-form>}
  ;     {the class <object>})
(define-method psychoanalyze ((being <intelligent>)) 'intelligent)
(define-method psychoanalyze ((being <humanoid>)) 'humanoid)
(psychoanalyze (make <vulcan>)) ; => intelligent
(psychoanalyze (make <human>))  ; => humanoid
\end{dylan}

\noindent A \code{<vulcan>} lists \code{<intelligent>} first among its
superclasses, so the \code{<intelligent>} method is more specific for it; a
\code{<human>} lists \code{<humanoid>} first. With multimethods the same
rule is applied lexicographically:

%! cont
\begin{dylan}
(define-method superior ((a <intelligent>) (b <intelligent>)) 'smarter)
(define-method superior ((a <humanoid>) (b <humanoid>)) 'better-looking)
(superior (make <vulcan>) (make <human>))   ; => smarter
(superior (make <human>) (make <vulcan>))   ; => better-looking
\end{dylan}

\begin{notebox}
Design note 31 later replaced the lexicographic rule with one that compares
every argument position and calls the methods \emph{ambiguous}, an error,
when the positions disagree, as they do in the last two calls above. Neither
implementation follows that note: both use the 1992 rule, under which a call
whose applicable methods all have class or singleton specializers is never
ambiguous.
\end{notebox}

\subsection{Calling the next method}\label{sec:functions-next-method}

\index{next-method@\code{next-method}}
A method often wants to add to the behavior it overrides rather than
replace it. Its next-method parameter, normally named \code{next-method},
holds a function that calls the next most specific applicable method. Called
with no arguments, it passes along the arguments the current method
received; called with arguments, it passes those instead.

\begin{dylan}
(define-method greet ((x <object>)) (list 'object x))
(define-method greet ((x <integer>)) (cons 'integer (next-method)))
(define-method greet ((x (singleton 5))) (cons 'five (next-method 7)))
(greet "hi")                   ; => (object "hi")
(greet 3)                      ; => (integer object 3)
(greet 5)                      ; => (five integer object 7)
\end{dylan}

\noindent When new arguments are given, they must lead to the same
applicable methods, in the same order, as the original ones; otherwise the
result is undefined. (In \code{greet} above, 7 is still an integer and an
object.) Keyword arguments passed to \code{next-method} with changed values
are not checked.

When there is no next method, the parameter is \code{\#f}, so a method can
test it before calling it; calling \code{\#f} is an error.

\begin{dylan}
(define-method last-one ((x <object>)) next-method)
(last-one 1)                   ; => #f
(define-method chain ((x <number>) #next more) (list 'number more))
(define-method chain ((x <integer>) #next more) (list 'integer (more)))
(chain 1)                      ; => (integer (number #f))
\end{dylan}

\noindent When methods take different keyword arguments, giving each a rest
parameter lets a method pass all of them on, changed or not, with
\code{apply} (Section~\ref{sec:functions-apply}):

\begin{dylan}
(define-method grow ((x <number>) #rest options #key (by 1))
  (+ x by))
(define-method grow ((x <integer>) #rest options #key by)
  (apply next-method x by: 10 options))
(grow 5)                       ; => 15
(grow 5 by: 2)                 ; => 15
\end{dylan}

\noindent The second call shows the leftmost-keyword rule at work: the
\code{by: 10} the integer method puts first wins over the caller's
\code{by: 2}.

\section{Calling functions: apply}\label{sec:functions-apply}

\index{apply@\code{apply}}
\begin{quote}
\code{(apply \meta{function} \meta{argument} \ldots{} \meta{sequence})}
\end{quote}

\noindent \code{apply} calls \meta{function} with the given arguments
followed by the elements of \meta{sequence}, the last argument, and returns
its values.

\begin{dylan}
(apply + 1 '(2 3))             ; => 6
(apply max (vector 3 9 2))     ; => 9
(define math-functions (list + * -))
(apply (first math-functions) 1 2 '(3 4))   ; => 10
(apply list '())               ; => ()
\end{dylan}

\section{Setters as generic functions}\label{sec:functions-setters}

A setter is an ordinary function, normally a generic function, stored in a
setter variable \code{(setter \meta{name})}. Its arguments are those of the
getter followed by the new value. \code{define-class} defines setter methods
for slots, and a program can define its own with \code{define-method}, so
that \code{set!} works on a computed property as it does on a slot:

\begin{dylan}
(define-class <temperature> (<object>)
  (kelvin init-value: 0 init-keyword: kelvin:))
(define-method celsius ((t <temperature>))
  (- (kelvin t) 273))
(define-method (setter celsius) ((t <temperature>) c)
  (set! (kelvin t) (+ c 273)))
(define t1 (make <temperature>))
(set! (celsius t1) 25)         ; => 25
(kelvin t1)                    ; => 298
(celsius t1)                   ; => 25
\end{dylan}

\noindent Section~\ref{sec:control-extended-set} describes how \code{set!}
calls a setter. Later Dylan changed both the name, to \code{celsius-setter},
and the argument order, putting the new value first; neither change is part
of this language.

\section{Reflective operations}\label{sec:functions-reflection}

These functions examine functions and generic functions
(\dylanbook{chapter 10}). Ordinary programs seldom need them, but they are
useful in tools and for understanding dispatch.

\begin{description}
\item[\code{(function-arguments \meta{function})}]
  \index{function-arguments@\code{function-arguments}}
  returns three values: the number of required arguments, whether the
  function takes a rest argument, and a sequence of the keywords it
  accepts, or \code{\#f} if it has no \code{\#key}. The sequence is empty
  for \code{\#key} with no keywords.
\item[\code{(applicable-method? \meta{function} \meta{argument} \ldots)}]
  \index{applicable-method?@\code{applicable-method?}}
  returns true if \meta{function} is a method applicable to the arguments,
  or a generic function with such a method.
\item[\code{(sorted-applicable-methods \meta{generic-function} \meta{argument} \ldots)}]
  \index{sorted-applicable-methods@\code{sorted-applicable-methods}}
  returns a sequence of the methods applicable to the arguments, most
  specific first.
\item[\code{(find-method \meta{generic-function} \meta{specializers})}]
  \index{find-method@\code{find-method}}
  returns the method whose specializers are exactly the classes and
  singletons in the sequence \meta{specializers}, or \code{\#f}.
\item[\code{(generic-function-methods \meta{generic-function})}]
  returns a sequence of all its methods, in no particular order.
\item[\code{(method-specializers \meta{method})}]
  returns a sequence of a method's specializers, one for each required
  parameter.
\end{description}

\noindent None of the sequences these functions return should be modified.

\begin{dylan}
(function-arguments (method (a b #rest r) a))     ; => 2, #t, #f
(function-arguments (method (a #key x (y: z 3)) a)) ; => 1, #f, (x: y:)
(define-method kind ((x <number>)) 'number)
(define-method kind ((x <integer>)) 'integer)
(define-method kind ((x (singleton 5))) 'five)
(applicable-method? kind 1/2)                     ; => #t
(applicable-method? kind "x")                     ; => #f
(map (method (m) (m 5))
     (sorted-applicable-methods kind 5))          ; => (five integer number)
((find-method kind (list <integer>)) 7)           ; => integer
(find-method kind (list <string>))                ; => #f
(size (generic-function-methods kind))            ; => 3
(method-specializers (find-method kind (list <integer>)))
                                                  ; => ({the class <integer>})
\end{dylan}

\subsection{Changing generic functions}\label{sec:functions-changing}

The 1992 manual also has functions that change a generic function while the
program runs. \code{add-method}, \code{remove-method} and
\code{freeze-methods} exist only in the Common Lisp implementation; in the
compiler to C, whose generic functions are fixed when the program is
compiled, using them is an error when compiling. \code{make-read-only}
exists in both, but in the compiler to C it changes nothing, since no
generic function can change there anyway; it returns its argument.

\begin{description}
\item[\code{(add-method \meta{generic-function} \meta{method})}] adds the
  method, replacing any with the same specializers, and returns two values:
  the method, and the method it replaced or \code{\#f}.
\item[\code{(remove-method \meta{generic-function} \meta{method})}] removes
  the method and returns it.
\item[\code{(make-read-only \meta{generic-function})}] forbids adding or
  removing methods from then on.
\item[\code{(freeze-methods \meta{generic-function} \meta{specializer} \ldots)}]
  forbids adding, replacing or removing methods whose specializers are the
  given ones or more specific.
\end{description}

%! lisp-only
\begin{dylan}
(define-method double ((x <number>)) (+ x x))
(define m (method ((s <string>)) (concatenate s s)))
(add-method double m)          ; => {a method on double}, #f
(double "ab")                  ; => "abab"
(remove-method double m)       ; => {an anonymous method}
(double "ab")                  ; => error: no applicable method
\end{dylan}

\noindent A generic function can also be made with \code{make}, with the
init-keywords \code{required:} (the number of required arguments, which must
be given), \code{rest?:}, \code{key?:} and \code{debug-name:}. It starts with
no methods. In the compiler to C such a generic function can never have any,
since methods cannot be added at run time, so it is of little use there;
that implementation also ignores \code{rest?:} and \code{key?:}.

%! lisp-only
\begin{dylan}
(define expand
  (make <generic-function> required: 1 debug-name: 'expand))
(add-method expand (method ((n <integer>)) (* n 10)))
(expand 5)                     ; => 50
\end{dylan}

\section{Proper tail calls}\label{sec:functions-tail-calls}

A call is in \term{tail position} when its value is the value of the
enclosing method: the last form of a method's body, either branch of an
\code{if} in tail position, the last form of a \code{begin} or \code{bind}
in tail position, and so on. The compiler to C guarantees \term{proper tail
calls}: every call in tail position, to a known method, a generic function,
a closure or a function in a variable, runs in constant stack space. A
recursive method whose recursive calls are all tail calls is therefore a
loop, and two methods can call each other in tail position indefinitely.
The \code{for}, \code{while} and other loops of
Chapter~\ref{ch:control} are themselves compiled this way.

\begin{dylan}
(define-method my-even? ((n <integer>))
  (if (= n 0) #t (my-odd? (- n 1))))
(define-method my-odd? ((n <integer>))
  (if (= n 0) #f (my-even? (- n 1))))
(my-even? 10000000)            ; => #t
\end{dylan}

\noindent The Common Lisp implementation compiles Dylan to Common Lisp, and
SBCL eliminates tail calls in that code too, so the example above runs there
as well; but the Common Lisp standard does not require it, and the Common
Lisp implementation makes no guarantee. Chapter~\ref{ch:dylanc} describes how the
compiler to C implements tail calls, with the C compiler's \code{musttail}
or with trampolines.

\chapter{Classes and Objects}\label{ch:classes}

Every Dylan object is a \emph{direct instance} of exactly one class, its
class, and an \emph{indirect instance} of all that class's superclasses; a
\term{general instance} of a class is a direct or an indirect instance.
Classes describe the structure of their instances, the \term{slots} that
hold an object's state, and together with generic functions they determine
the instances' behavior. Classes are themselves objects, instances of
\code{<class>}, and can be passed around and examined like any other.

This chapter describes defining classes and their slots, making and
initializing instances, the functions that examine classes, singletons and
sealing, and the classes built into the language. The standard is
\dylanbook{chapters 2, 6 and 11}. Dispatch on classes is described in
Chapter~\ref{ch:functions}.

\section{Defining classes: define-class}\label{sec:classes-define-class}

\index{define-class@\code{define-class}}
\begin{quote}
\code{(define-class \meta{name} (\meta{superclass} \ldots) \meta{slot-spec} \ldots)}
\end{quote}

\noindent \code{define-class} creates a class and stores it in the module
variable \meta{name}, which is made read-only. The superclass expressions
are evaluated and must give at least one class, with no class twice; the
new class inherits from all of them. Each \meta{slot-spec} adds a slot
(Section~\ref{sec:classes-slots}). At the interactive prompt
\code{define-class} returns the name.

\begin{dylan}
(define-class <point> (<object>)
  (x init-value: 0 init-keyword: x:)
  (y init-value: 0 init-keyword: y:))
(define p (make <point> x: 3))
(x p)                          ; => 3
(y p)                          ; => 0
(set! (y p) 4)                 ; => 4
(instance? p <point>)          ; => #t
\end{dylan}

\noindent By convention class names begin and end with angle brackets. A
class that needs no other superclass names \code{<object>}, the root of
the class hierarchy.

In the Common Lisp implementation, defining a class whose name already
holds a class modifies that class in place, as the 1992 manual says, so
that everything referring to it sees the new definition. In the compiler to
C a class may be defined only once.

\subsection{Multiple inheritance}\label{sec:classes-multiple-inheritance}

A class may have several direct superclasses. It inherits the slots of all
of them, and methods specialized on any of them apply to its instances.
The order of the superclasses matters: it determines the class's
\term{class precedence list}, which orders the class and all its
superclasses from most to least specific, and so decides which method runs
when methods on two unrelated superclasses both apply
(Section~\ref{sec:functions-cpl}).

\begin{dylan}
(define-class <named> (<object>)
  (name init-keyword: name: init-value: "anonymous"))
(define-class <located> (<object>)
  (place init-keyword: place: init-value: 'nowhere))
(define-class <landmark> (<named> <located>))
(define l (make <landmark> name: "Big Ben" place: 'london))
(list (name l) (place l))      ; => ("Big Ben" london)
(all-superclasses <landmark>)
  ; => ({the class <landmark>} {the class <named>}
  ;     {the class <located>} {the class <object>})
\end{dylan}

\noindent A class cannot be its own superclass, directly or indirectly. If
the superclasses' orderings cannot all be satisfied at once, no class
precedence list exists and defining the class is an error.

\section{Slots}\label{sec:classes-slots}

A slot is a place in an object that holds a value. A slot is not read by
name; it is read and written by \term{getter} and \term{setter} generic
functions, to which \code{define-class} adds methods. A slot named
\code{x} gets a getter method on the generic function \code{x} and a setter
method on \code{(setter x)}, and \code{(set! (x p) 10)} calls the setter.
Since accessors are generic functions, a slot looks to the rest of the
program exactly like a computed property, and a subclass can replace a
stored slot with a computed one without changing any caller.

\subsection{Slot specifications}

A slot specification takes one of three forms:

\begin{quote}
\meta{slot-name}\\
\code{(\meta{slot-name} \meta{keyword} \meta{value} \ldots)}\\
\code{(\meta{keyword} \meta{value} \ldots)}
\end{quote}

\noindent The slot name is a symbol and is not evaluated. It supplies the
default getter name, \meta{slot-name}, and the default setter name,
\code{(setter \meta{slot-name})}. In the third form there is no slot name,
and \code{getter:} must be given. Each option may be given at most once.

\begin{description}
\item[\code{getter:} \meta{name}] The variable to which the getter method
  is added. Not evaluated.
\item[\code{setter:} \meta{name}] The variable to which the setter method
  is added. Not evaluated. \code{setter: \#f} gives the slot no setter.
\item[\code{type:} \meta{type}] A class or singleton that the slot's values
  must be instances of; storing anything else is an error. Evaluated once.
  The default is \code{<object>}.
\item[\code{init-value:} \meta{value}] The slot's initial value, used for
  every new instance. Evaluated once, when the class is defined.
\item[\code{init-function:} \meta{function}] A function of no arguments,
  called to compute the initial value for each new instance. Evaluated once,
  to give the function. Cannot be combined with \code{init-value:}.
\item[\code{init-keyword:} \meta{keyword}] A keyword with which
  \code{make} may be given the initial value. Not evaluated.
\item[\code{required-init-keyword:} \meta{keyword}] Like
  \code{init-keyword:}, but \code{make} must be given it. Cannot be
  combined with \code{init-keyword:}, \code{init-value:} or
  \code{init-function:}.
\item[\code{allocation:} \meta{kind}] Where the slot's value is kept:
  \code{instance}, \code{class}, \code{each-subclass}, \code{constant} or
  \code{virtual} (Section~\ref{sec:classes-allocation}). Not evaluated. The
  default is \code{instance}.
\end{description}

\begin{dylan}
(define-class <rectangle> (<object>)
  (top type: <integer> init-value: 0 init-keyword: top:)
  (left type: <integer> init-value: 0 init-keyword: left:)
  (bottom type: <integer> init-value: 100 init-keyword: bottom:)
  (right type: <integer> init-value: 100 init-keyword: right:))
(define r (make <rectangle> top: 50 left: 50))
(top r)                        ; => 50
(bottom r)                     ; => 100
(set! (bottom r) 55)           ; => 55
(set! (bottom r) 'foo)         ; => error: foo is not an instance of <integer>
\end{dylan}

\noindent The getter and setter options let a slot have accessor names that
differ from its slot name, or no setter at all:

\begin{dylan}
(define-class <account> (<object>)
  (number required-init-keyword: number: setter: #f)
  (balance init-value: 0
           getter: current-balance setter: (setter current-balance))
  (getter: owner-name init-keyword: owner:))
(define a (make <account> number: 1001 owner: "Ann"))
(number a)                     ; => 1001
(current-balance a)            ; => 0
(owner-name a)                 ; => "Ann"
(make <account>)               ; => error: number: is required
\end{dylan}

\subsection{Initial values}\label{sec:classes-slot-init}

When \code{make} creates an instance, each slot is initialized from the
first of these that applies: the init-keyword's argument to \code{make}, if
given; the init-value; the result of calling the init-function. The
init-function is not called when the keyword is given. A slot with none of
these is left \term{uninitialized}, and reading it is an error until it has
been set. The initial value is stored directly into the slot; the setter
generic function is not called.

An init-value is evaluated once, when the class is defined, and the same
object becomes the initial value of every instance. To give each instance
its own fresh object, such as an empty collection, use an init-function:

\begin{dylan}
(define-class <bag> (<object>)
  (shared init-value: (make <stretchy-vector>))
  (own init-function: (method () (make <stretchy-vector>))))
(define b1 (make <bag>))
(define b2 (make <bag>))
(id? (shared b1) (shared b2))  ; => #t
(id? (own b1) (own b2))        ; => #f
\end{dylan}

\index{slot-initialized?@\code{slot-initialized?}}
\code{(slot-initialized? \meta{instance} \meta{getter})} tells whether the
slot that \meta{getter} reads in \meta{instance} has a value. It is an error
if \meta{getter} reads no slot of the instance. There is no way to make a
slot uninitialized again.

\begin{dylan}
(define-class <memo> (<object>) (note))
(define m (make <memo>))
(slot-initialized? m note)     ; => #f
(note m)                       ; => error: the slot note is not initialized
(set! (note m) "remember")     ; => "remember"
(slot-initialized? m note)     ; => #t
\end{dylan}

\subsection{Allocation}\label{sec:classes-allocation}

\begin{description}
\item[\code{instance}] Each instance has its own value. This is the
  default.
\item[\code{class}] There is one value, shared by all the general instances
  of the class, its subclasses' instances included. Setting it through one
  instance changes it for all.
\item[\code{each-subclass}] The class has one value, shared by its direct
  instances, and each subclass has a value of its own, shared by the
  subclass's direct instances.
\item[\code{constant}] The value is the init-value, which must be given,
  and it cannot be changed: the slot has no setter.
\item[\code{virtual}] No storage is allocated. The program defines methods
  on the getter, and on the setter if it wants one, that compute the value;
  \code{define-class} only makes sure the generic functions exist.
\end{description}

\noindent A \code{class} or \code{each-subclass} slot can be initialized
only with \code{init-value:}; without one it starts uninitialized, like
any other slot, until it is set; a \code{virtual} slot takes no initialization
options at all, and \code{slot-initialized?} is always true of it.

\begin{dylan}
(define-class <vehicle> (<object>)
  (count allocation: class init-value: 0)
  (kind allocation: each-subclass init-value: 'vehicle)
  (wheels allocation: constant init-value: 4)
  (description allocation: virtual))
(define-class <truck> (<vehicle>))
(define-method description ((v <vehicle>))
  (list (kind v) (wheels v)))
(define bike (make <vehicle>))
(define truck (make <truck>))
(set! (count bike) 2)          ; => 2
(count truck)                  ; => 2
(set! (kind truck) 'truck)     ; => truck
(kind bike)                    ; => vehicle
(kind (make <truck>))          ; => truck
(description truck)            ; => (truck 4)
(slot-initialized? bike description)   ; => #t
\end{dylan}

\noindent A virtual slot can also give an inherited stored slot a new
behavior in a subclass, a \term{filtered slot}: the subclass declares the
slot virtual for documentation and defines getter and setter methods that
call \code{next-method} to reach the stored value.

\begin{dylan}
(define-class <view> (<object>)
  (position init-value: 0 init-keyword: position:))
(define-class <displaced-view> (<view>)
  (offset init-value: 10))
(define-method position ((v <displaced-view>))
  (+ (next-method) (offset v)))
(define-method (setter position) ((v <displaced-view>) new)
  (next-method v (- new (offset v))))
(define dv (make <displaced-view> position: 5))
(position dv)                  ; => 15
(set! (position dv) 100)       ; => 100
(position dv)                  ; => 100
\end{dylan}

\subsection{Inherited slots}\label{sec:classes-inherited-slots}

Two slots are the same slot if they have the same getter. A subclass may
respecify a slot it inherits, to give it, for example, a different initial
value. The 1992 manual asks that the new specification repeat every feature
the inherited one specified (its type, an initial value, its init-keyword,
a non-instance allocation and a setter), with a type no wider than the
inherited one, and both implementations report a respecification that
leaves one out (``redefining an inherited slot must respecify its
init-keyword:''). The exception is a slot redeclared \code{virtual}, a
filtered slot (below), which respecifies nothing: the inherited slot keeps
its storage, its type and its initialization, so \code{make} still takes
its init-keyword. A class that inherits the same slot from two
superclasses with different specifications must respecify it.

\begin{dylan}
(define-class <base> (<object>)
  (color init-value: 'red init-keyword: color:))
(define-class <derived> (<base>)
  (color init-value: 'blue init-keyword: color:))
(color (make <base>))                 ; => red
(color (make <derived>))              ; => blue
(color (make <derived> color: 'green)) ; => green
\end{dylan}

\section{Making instances: make and initialize}\label{sec:classes-make}

\index{make@\code{make}}\index{initialize@\code{initialize}}
\begin{quote}
\code{(make \meta{class} \meta{keyword} \meta{value} \ldots)}
\end{quote}

\noindent \code{make} returns a general instance of \meta{class}. For a
class the program defines, it allocates a new instance, initializes its
slots (Section~\ref{sec:classes-slot-init}), then calls the generic function
\code{initialize} with the instance and all the keyword arguments, and
returns the instance; what \code{initialize} returns is ignored. Making an
instance of an abstract class, such as \code{<number>} or
\code{<collection>}, is an error. Built-in instantiable classes have
\code{make} methods of their own, described with them
(Chapter~\ref{ch:collections}).

The keywords \code{make} accepts are the init-keywords of the class's slots,
inherited slots included, and the keywords accepted by the applicable
\code{initialize} methods. An \code{initialize} method with \code{\#rest}
accepts any keyword. Any other keyword is an error, as the manual
specifies: the compiler to C reports \code{keyword argument not accepted:
y:}, the Common Lisp implementation \code{make -{}- illegal keyword
argument y:}.

Programs customize instance creation by adding methods to
\code{initialize}, not to \code{make}. An \code{initialize} method normally
calls \code{next-method} so that the superclasses' methods run too; the
method on \code{<object>} does nothing. When \code{initialize} runs, the
slots already have their initial values.

\begin{dylan}
(define-class <shape> (<object>)
  (name init-keyword: name: init-value: "shape")
  (area))
(define-class <square> (<shape>)
  (side init-keyword: side: init-value: 1))
(define-method initialize ((s <shape>) #rest options)
  (next-method)
  (print (list 'initializing (name s))))
(define-method initialize ((s <square>) #rest options)
  (next-method)
  (set! (area s) (* (side s) (side s))))
(define sq (make <square> side: 3 name: "sq"))
                               ; prints: (initializing "sq")
(area sq)                      ; => 9
\end{dylan}

\noindent The \code{<square>} method above takes \code{\#rest} rather than
\code{\#key side}: a parameter named \code{side} would shadow the getter
\code{side} within the method. To receive a keyword under another name,
write \code{\#key (side: the-side)} (Section~\ref{sec:functions-keywords}).

\section{Adding slots: define-slot}\label{sec:classes-define-slot}

\index{define-slot@\code{define-slot}}
\begin{quote}
\code{(define-slot \meta{owner} [\meta{slot-name}] \meta{keyword} \meta{value} \ldots)}
\end{quote}

\noindent \code{define-slot} adds a slot to an existing class or to a
single object. \meta{owner} is evaluated and must be a class or a
singleton; the slot name and options are those of a slot specification in
\code{define-class}. If the owner is a class, all its instances get the
slot. If it is a singleton, only that one object does, and its allocation
must be \code{instance}, \code{constant} or \code{virtual}, without
\code{init-keyword:} or \code{required-init-keyword:}.

\begin{dylan}
(define-class <pt> (<object>) (px init-value: 0))
(define p (make <pt>))
(define-slot (singleton p) note init-value: "only me")
(note p)                       ; => "only me"
(note (make <pt>))             ; => error: no applicable method
(define x (list 1 2))
(define-slot (singleton x) tag init-value: 'special)
(set! (tag x) 'changed)        ; => changed
(tag x)                        ; => changed
(define-slot (singleton 'cup) volume allocation: constant init-value: 8)
(volume 'cup)                  ; => 8
(define-slot <pt> weight init-value: 10)
(weight p)                     ; => 10
\end{dylan}

\noindent As the examples show, the object may be of any class, built-in or
not, and the slot's accessors are generic functions with a method
specialized on the singleton.

In the Common Lisp implementation \code{define-slot} takes effect when it is
evaluated, and instances that already exist gain the slot. In the compiler
to C the slot is part of the program from the start: a slot defined on a
class is in every instance ever made, and is listed by
\code{slot-descriptors} even before the \code{define-slot} form is reached.
The compiler to C also requires a class owner to be a class the program
itself defines, known when compiling, rather than a built-in class or a
value computed at run time; a singleton owner may be any object.

\section{Examining classes and objects}\label{sec:classes-reflection}

\begin{description}
\item[\code{(object-class \meta{object})}]
  \index{object-class@\code{object-class}}
  returns the class of which \meta{object} is a direct instance.
\item[\code{(instance? \meta{object} \meta{type})}]
  \index{instance?@\code{instance?}}
  returns true if \meta{object} is a general instance of the class
  \meta{type}, or is the object of the singleton \meta{type}.
\item[\code{(subclass? \meta{class1} \meta{class2})}]
  \index{subclass?@\code{subclass?}}
  returns true if \meta{class1} is \meta{class2} or a direct or indirect
  subclass of it.
\item[\code{(all-superclasses \meta{class})}]
  \index{all-superclasses@\code{all-superclasses}}
  returns the class precedence list: \meta{class} first, \code{<object>}
  last.
\item[\code{(direct-superclasses \meta{class})}]
  returns the superclasses named in the class's definition, in order.
\item[\code{(direct-subclasses \meta{class})}]
  returns the classes that name \meta{class} as a direct superclass.
\end{description}

\begin{dylan}
(define-class <animal> (<object>))
(define-class <dog> (<animal>))
(define rex (make <dog>))
(object-class rex)             ; => {the class <dog>}
(instance? rex <animal>)       ; => #t
(instance? 3 <real>)           ; => #t
(instance? 3 (singleton 3))    ; => #t
(subclass? <dog> <animal>)     ; => #t
(subclass? <dog> <dog>)        ; => #t
(direct-superclasses <dog>)    ; => ({the class <animal>})
(direct-subclasses <animal>)   ; => ({the class <dog>})
(all-superclasses <dog>)
  ; => ({the class <dog>} {the class <animal>} {the class <object>})
\end{dylan}

\noindent None of the sequences these functions return should be modified.
The two implementations give different answers for some built-in objects.
In the compiler to C, \code{\#t} and \code{\#f} are instances of an
internal class that prints as \code{\{the class <boolean>\}} but is not the
value of any variable a program can name, while in the Common Lisp
implementation their class is \code{<object>}; and the Common Lisp implementation seals the
built-in classes that cannot be subclassed, such as \code{<integer>} and
\code{<pair>}, so \code{all-superclasses} and \code{direct-superclasses} of
them return empty sequences (Section~\ref{sec:classes-sealing}).

\subsection{Slot descriptors}\label{sec:classes-slot-descriptors}

A \term{slot descriptor}, of class \code{<slot-descriptor>}, describes one
slot. Programs seldom need them, but debuggers and inspectors do. Both
implementations provide the slot-descriptor functions of
\dylanbook{chapter 11}. (Design note 9 later removed them from the
language.)

\begin{description}
\item[\code{(slot-descriptors \meta{class})}] all the slots of the class,
  inherited ones included.
\item[\code{(slot-descriptor \meta{instance} \meta{getter})}] the slot that
  \meta{getter} reads in \meta{instance}, or \code{\#f}.
\item[\code{(slot-getter \meta{sd})}, \code{(slot-setter \meta{sd})}] the
  getter and setter generic functions; \code{slot-setter} returns
  \code{\#f} for a slot without a setter.
\item[\code{(slot-type \meta{sd})}] the type, \code{<object>} if none was
  given.
\item[\code{(slot-allocation \meta{sd})}] the allocation, a symbol.
\item[\code{(init-value \meta{sd})}] two values: the init-value and
  \code{\#t}, or \code{\#f} and \code{\#f} if there is none.
\item[\code{(init-function \meta{sd})}, \code{(init-keyword \meta{sd})}]
  the init-function or init-keyword, or \code{\#f}.
\item[\code{(slot-value \meta{instance} \meta{sd})}] the value of the slot
  in the instance. Its setter, used as
  \code{(set! (slot-value \meta{instance} \meta{sd}) \meta{value})},
  sets it, checking the slot's type. Neither works on a virtual slot.
\end{description}

\begin{dylan}
(define-class <cell> (<object>)
  (contents type: <integer> init-value: 1 init-keyword: contents:))
(define c (make <cell> contents: 3))
(size (slot-descriptors <cell>))       ; => 1
(define sd (slot-descriptor c contents))
(id? (slot-getter sd) contents)        ; => #t
(id? (slot-setter sd) (setter contents)) ; => #t
(slot-type sd)                         ; => {the class <integer>}
(slot-allocation sd)                   ; => instance
(init-value sd)                        ; => 1, #t
(init-keyword sd)                      ; => contents:
(init-function sd)                     ; => #f
(slot-value c sd)                      ; => 3
(set! (slot-value c sd) 4)             ; => 4
(contents c)                           ; => 4
(slot-descriptor c object-class)       ; => #f
\end{dylan}

\subsection{Copying and converting}

\code{(shallow-copy \meta{object})} returns a new object with the same
contents, made with \code{make} on \code{(class-for-copy \meta{object})},
which by default is the object's class. The library has methods for
collections. For other objects the 1992 manual leaves the method to the
programmer, and that is so in the Common Lisp implementation; the compiler
to C has a default method that copies any instance's slots.

\code{(as \meta{class} \meta{object})} converts \meta{object} to
\meta{class}, returning it unchanged if it is already an instance. Its
methods for numbers, characters, symbols and collections are described in
Chapters~\ref{ch:numbers} and~\ref{ch:collections}.

\section{Singletons}\label{sec:classes-singletons}

A singleton is a specializer for one object (Section~\ref{sec:functions-singletons}).
\code{(singleton \meta{object})} and \code{(make <singleton> object:
\meta{object})} both make one. With methods on singletons a program can give
an individual object behavior of its own, and with \code{define-slot}
(Section~\ref{sec:classes-define-slot}) state of its own, without defining a
class for it. A method on a singleton is more specific than any method on
the object's classes.

%! lisp-only
\begin{dylan}
(define-class <window> (<object>))
(define main-window (make <window>))
(define-method title ((w <window>)) "untitled")
(define-method title ((w (singleton main-window))) "Main")
(title (make <window>))        ; => "untitled"
(title main-window)            ; => "Main"
\end{dylan}

\noindent The compiler to C does not accept that program: there, a
singleton used as a method's specializer must be of a literal, a class, or
a variable whose definition gives it a literal value, because the
specializer must be known when compiling
(Section~\ref{sec:functions-singletons}). A singleton of an instance made at
run time can still own slots, through \code{define-slot}, and singletons of
symbols, keywords, numbers and characters work everywhere:

\begin{dylan}
(define-method plural ((word <symbol>))
  (as <symbol> (concatenate (as <string> word) "s")))
(define-method plural ((word (singleton 'mouse))) 'mice)
(plural 'cat)                  ; => cats
(plural 'mouse)                ; => mice
(make <singleton> object: 4)   ; => {singleton 4}
\end{dylan}

\section{Sealing}\label{sec:classes-sealing}

\index{seal@\code{seal}}\index{make-read-only@\code{make-read-only}}
\code{(seal \meta{class})} seals a class: no subclass of it can be created
any more, and \code{all-superclasses}, \code{direct-superclasses} and
\code{slot-descriptors} return empty sequences for it. Sealing changes
nothing about existing instances or methods. \code{(make-read-only
\meta{class})} forbids changing the class's slots and superclasses without
forbidding subclasses. Both return the class.

\begin{dylan}
(define-class <a> (<object>) (a-slot init-value: 1))
(define-class <b> (<a>))
(seal <b>)                     ; => {the class <b>}
(all-superclasses <b>)         ; => ()
(direct-superclasses <b>)      ; => ()
(a-slot (make <b>))            ; => 1
(instance? (make <b>) <a>)     ; => #t
\end{dylan}

\noindent In the Common Lisp implementation, defining a subclass of a sealed
class is then an error. In the compiler to C all of a program's classes are
created when it is compiled, before \code{seal} runs, so \code{seal} affects
only what the reflective functions return afterwards; and there
\code{make-read-only} on a class does exactly what \code{seal} does. A
built-in class that cannot be subclassed, such as \code{<integer>} or
\code{<list>}, is reported when compiling, if a program tries.

\section{Changing classes at run time}\label{sec:classes-dynamic}

The 1992 manual lets a running program change its classes: \code{add-slot}
and \code{remove-slot} add and remove slots, returning slot descriptors;
\code{(set! (direct-superclasses \meta{class}) \meta{list})} gives a class
new superclasses; and \code{(make <class> superclasses: \meta{classes}
debug-name: \meta{name})} creates a new class. Existing instances are
updated to the new definition before they are next used.

All of this works in the Common Lisp implementation. A program for the
compiler to C is static: its classes and their slots are fixed when it is
compiled. There \code{add-slot} and \code{remove-slot} are errors when
compiling, \code{make} of \code{<class>} is an error when run, and
\code{define-slot} on a class, which the compiler handles as part of the
class's definition, is the way to add a slot (Section~\ref{sec:classes-define-slot}).

%! lisp-only
\begin{dylan}
(define-class <pt> (<object>) (px init-value: 0))
(define p (make <pt>))
(define-generic-function weight (x))
(define-generic-function (setter weight) (x v))
(define sd (add-slot <pt> getter: weight setter: (setter weight)
                          init-value: 3))
(weight p)                     ; => 3
(remove-slot <pt> sd)
(weight p)                     ; => error: no applicable method
(define c (make <class> superclasses: (list <object>) debug-name: '<dyn>))
(instance? (make c) c)         ; => #t
\end{dylan}

\section{The built-in classes}\label{sec:classes-built-in}

The classes below are those of the 1992 manual, as the compiler to C
arranges them. Each class is shown under its first direct superclass; a
class with more than one is marked with the others. Where the Common Lisp
implementation differs, a note says so.

%! skip
\begin{dylan}
<object>
  <boolean>                          (compiler to C only; internal)
  <number>
    <complex>
      <real>
        <rational>
          <integer>
          <ratio>
        <float>
          <single-float>
          <double-float>
          <extended-float>
      <rectangular-complex>          (Common Lisp only)
  <collection>
    <explicit-key-collection>
      <mutable-explicit-key-collection>   (also <mutable-collection>)
        <array>
          <vector>                   (also <mutable-sequence>)
            <simple-object-vector>
            <stretchy-vector>
            <byte-string>            (also <string>)
            <unicode-string>         (also <string>)
        <table>
    <mutable-collection>
      <mutable-sequence>             (also <sequence>)
        <string>
        <list>
          <pair>
          <empty-list>
        <deque>
    <sequence>
      <range>
  <function>
    <generic-function>
    <method>
  <class>
  <singleton>
  <slot-descriptor>
  <symbol>
  <keyword>
  <character>
  <condition>
    <serious-condition>
      <error>
        <simple-error>
        <type-error>
        <thread-error>               (compiler to C only)
    <warning>
      <simple-warning>
    <restart>
      <simple-restart>
      <abort>
\end{dylan}

\noindent Both implementations add the stream class \code{<stream>},
a direct subclass of \code{<object>} with the subclass
\code{<string-stream>} (Chapter~\ref{ch:printing}). The compiler to C
has further classes for its threads (Chapter~\ref{ch:extensions}), also
direct subclasses of \code{<object>}: \code{<thread>}, \code{<lock>},
\code{<condition-variable>} and \code{<atomic-integer>}. \code{<lock>}
has the subclass \code{<recursive-lock>}.

Numbers are described in Chapter~\ref{ch:numbers}, collections in
Chapter~\ref{ch:collections} and conditions in Chapter~\ref{ch:conditions}.

\subsection{Which classes can be subclassed}

A program may define subclasses of \code{<object>} and of the abstract
classes meant to be extended: \code{<number>}, the collection classes
\code{<collection>}, \code{<sequence>}, \code{<mutable-sequence>},
\code{<explicit-key-collection>} and the others above them, the
instantiable collection classes \code{<array>}, \code{<vector>},
\code{<string>}, \code{<table>}, \code{<deque>}, \code{<stretchy-vector>}
and \code{<range>}, and \code{<function>}, \code{<condition>},
\code{<error>} and the other condition classes. These classes are
\emph{sealed}, and subclassing them is an error in both implementations:
the number classes below \code{<number>}, from \code{<complex>} down;
the list classes \code{<list>}, \code{<pair>} and \code{<empty-list>}; and
the concrete vector and string classes \code{<simple-object-vector>},
\code{<byte-string>} and \code{<unicode-string>}.

The compiler to C also refuses, when compiling, to use as a superclass the
classes whose instances its runtime makes itself: \code{<method>},
\code{<generic-function>}, \code{<class>}, \code{<singleton>},
\code{<symbol>}, \code{<keyword>} and \code{<character>}. The Common Lisp
implementation accepts these.

\chapter{Numbers, Characters and Symbols}\label{ch:numbers}

This chapter describes Dylan's simplest objects: numbers, characters,
symbols and keywords. It follows chapters 14, 15 and 16 of the 1992
manual (\dylanbook{chapters 14--16}) and, for comparison, chapter~9
(\dylanbook{chapter 9}). The two implementations agree on almost
everything here, but numbers are where they differ most: the Common Lisp
implementation has the unbounded integers and the complex numbers the
manual describes, and the compiler to C has neither. Each difference is
pointed out where it arises; appendix~\ref{app:differences} collects them.

The examples show their results as comments, \code{; => value}. Where the
two implementations print a result differently, both are shown.

\section{The number classes}\label{sec:numbers-classes}

\index{number classes}
Numbers are arranged in a hierarchy of classes that follows the
mathematical kinds of number. The abstract classes have no direct
instances; they exist so that methods can specialize on them.

\begin{center}
\begin{tabular}{lll}
\toprule
Class & Direct superclass & Kind \\
\midrule
\code{<number>}\index{<number>} & \code{<object>} & abstract \\
\code{<complex>}\index{<complex>} & \code{<number>} & abstract \\
\code{<real>}\index{<real>} & \code{<complex>} & abstract \\
\code{<rational>}\index{<rational>} & \code{<real>} & abstract \\
\code{<integer>}\index{<integer>} & \code{<rational>} & concrete \\
\code{<ratio>}\index{<ratio>} & \code{<rational>} & concrete \\
\code{<float>}\index{<float>} & \code{<real>} & abstract \\
\code{<single-float>}\index{<single-float>} & \code{<float>} & concrete \\
\code{<double-float>}\index{<double-float>} & \code{<float>} & concrete \\
\code{<extended-float>}\index{<extended-float>} & \code{<float>} & concrete \\
\code{<rectangular-complex>}\index{<rectangular-complex>} & \code{<complex>} & concrete; Common Lisp only \\
\bottomrule
\end{tabular}
\end{center}

\noindent All of these classes but \code{<number>} are sealed: a program
may not define new subclasses of them, though it may define methods on
them. Every number
belongs to exactly one of the concrete classes, which
\code{object-class} reports:

\begin{dylan}
(object-class 3)        ; => <integer>
(object-class 1/2)      ; => <ratio>
(object-class (/ 4 2))  ; => <integer>
(object-class 1.5)      ; => <double-float>
(object-class 1.5s0)    ; => <single-float>
\end{dylan}

\noindent (A class prints as \code{\{the class <integer>\}}; this
chapter writes just the name.)

\subsection{Integers}

\index{integer!range}\index{overflow}
An \code{<integer>} is a whole number. The two implementations differ in
how large one may be:

\begin{itemize}
\item In the \textbf{Common Lisp implementation}, integers are unbounded,
  as the manual requires: \code{(expt 2 100)} is
  \code{1267650600228229401496703205376}, and an integer of any size is
  still an \code{<integer>}.
\item In the \textbf{compiler to C}, an integer is 63 bits, from
  $-2^{62}$ to $2^{62}-1$, that is, from \code{-4611686018427387904} to
  \code{4611686018427387903}. Every integer operation checks for
  overflow, and a result outside that range signals an error rather than
  giving a wrong answer:
\end{itemize}

\begin{dylan}
(+ 4611686018427387903 1)
; C:    Dylan error: integer overflow (integers are 63 bits)
; Lisp: => 4611686018427387904
\end{dylan}

\noindent The error is an ordinary \code{<simple-error>} (chapter~\ref{ch:conditions}),
so a program can handle it:

\begin{dylan}
(handler-case (* 4611686018427387903 2)
  ((<error>) 'overflow))
; C:    => overflow
; Lisp: => 9223372036854775806
\end{dylan}

\noindent The check applies to every operation that makes an integer:
\code{+}, \code{-}, \code{*}, \code{expt}, \code{ash}, \code{lcm}, and
the numerators and denominators of ratios.

\subsection{Ratios}

\index{ratio}
A \code{<ratio>} is an exact fraction, kept in lowest terms with a
positive denominator. Dividing integers that do not divide evenly makes
one, and a ratio whose denominator would be 1 is made an integer instead
(\term{rational canonicalization}):

\begin{dylan}
(/ 1 3)         ; => 1/3
(/ 6 3)         ; => 2
(/ -6 4)        ; => -3/2
(+ 1/3 1/6)     ; => 1/2
(+ 1/2 1/2)     ; => 1
(* 1/2 4)       ; => 2
\end{dylan}

\noindent In the compiler to C the numerator and the denominator are
each limited to 63 bits, with the same overflow error as integers;
intermediate results are computed exactly in 128 bits, so only a final
result that does not fit is an error. In the Common Lisp implementation
ratios are unbounded.

\subsection{Floats}

\index{float}
The three float classes are the IEEE formats the manual names. Their
representations are:

\begin{center}
\begin{tabular}{lll}
\toprule
Class & Compiler to C & Common Lisp implementation \\
\midrule
\code{<single-float>} & C \code{float} (32 bits) & SBCL \code{single-float} \\
\code{<double-float>} & C \code{double} (64 bits) & SBCL \code{double-float} \\
\code{<extended-float>} & C \code{double}, a class of its own & no instances \\
\bottomrule
\end{tabular}
\end{center}

\noindent \code{<double-float>} is the default: a float literal without
an exponent marker, or with \code{e} or \code{d}, is a double. In the
compiler to C an \code{<extended-float>} is a separate class with a
double's precision; the only way to make one is \code{as} (see
section~\ref{sec:numbers-conversion}). In the Common Lisp implementation,
SBCL's long float is its double float, so \code{<extended-float>} has no
instances at all, and \code{(as <extended-float> 1)} answers the
\code{<double-float>} \code{1.0}.

\subsection{Complex numbers}

\index{complex numbers}
The manual's \code{<complex>} numbers exist only in the Common Lisp
implementation, where they are Common Lisp's complex numbers and belong
to the class \code{<rectangular-complex>}:

\begin{dylan}
(make-rectangular 1 2)                   ; Lisp: => #c(1 2)
(make <complex> real: 1 imag: 2)         ; Lisp: => #c(1 2)
(* (make-rectangular 0 1)
   (make-rectangular 0 1))               ; Lisp: => -1
(abs (make-rectangular 3 4))             ; Lisp: => 5.0
(sqrt -4)                                ; Lisp: => #c(0.0 2.0)
\end{dylan}

\noindent The init-keywords of \code{<complex>} are \code{real:},
\code{imag:}, \code{magnitude:} and \code{angle:}, as in
\dylanbook{chapter 14}; \code{make-rectangular} and \code{make-polar}
make complex numbers from parts, and \code{real-part},
\code{imag-part} and \code{angle} take them apart. As the manual
requires, a complex number whose imaginary part is an exact zero is
canonicalized to a real: \code{(make-rectangular 3 0)} is \code{3}.

In the compiler to C, \code{<complex>} is kept, abstract, as the
superclass of \code{<real>}, but it has no instances of its own.
\code{make-rectangular} and \code{make-polar} accept only an exact zero
for the imaginary part or the angle, and answer the real; anything else
is an error, as is \code{make} of \code{<complex>} and the square root
of a negative number:

\begin{dylan}
(make-rectangular 3 0)   ; => 3
(make-rectangular 1 2)   ; C: error: complex numbers are not supported: 1 2
(sqrt -4)                ; C: error: sqrt: the result would be complex: -4
\end{dylan}

\section{Writing numbers}\label{sec:numbers-literals}

\index{number literals}
Numbers are written as in Common Lisp:

\begin{center}
\begin{tabular}{lll}
\toprule
Written & Value & Class \\
\midrule
\code{42}, \code{-7} & 42, $-7$ & \code{<integer>} \\
\code{1.} & 1 & \code{<integer>} (a trailing point means decimal) \\
\code{\#xFF}, \code{\#b101}, \code{\#o17} & 255, 5, 15 & \code{<integer>} \\
\code{6/4} & \code{3/2} & \code{<ratio>} \\
\code{1.5}, \code{.5}, \code{-0.0} & & \code{<double-float>} \\
\code{1e3}, \code{2d0} & \code{1000.0}, \code{2.0} & \code{<double-float>} \\
\code{1.5s2}, \code{1.5f2} & \code{150.0s0} & \code{<single-float>} \\
\code{1.5l0} & \code{1.5} & \code{<double-float>} \\
\bottomrule
\end{tabular}
\end{center}

\noindent The exponent markers \code{e} and \code{d} make a double,
\code{s} and \code{f} a single. The marker \code{l} also makes a double,
in both implementations; there is no way to write an
\code{<extended-float>} literal. A ratio is written with no spaces,
\code{numerator/denominator}, and is reduced as it is read.

\section{Arithmetic}\label{sec:numbers-arithmetic}

\index{arithmetic}\index{+}\index{-}\index{*}\index{/}
The four arithmetic functions take any number of arguments, as in
\dylanbook{chapter 14}:

\begin{dylan}
(+)             ; => 0
(*)             ; => 1
(+ 1 2 3)       ; => 6
(- 5)           ; => -5         the additive inverse
(- 10 1 2 3)    ; => 4          10 - 1 - 2 - 3
(/ 2)           ; => 1/2        the multiplicative inverse
(/ 0.5)         ; => 2.0
(/ 60 2 3)      ; => 10         60 / 2 / 3
(/ 7 2)         ; => 7/2
(/ 7.0 2)       ; => 3.5
\end{dylan}

\noindent \code{-} and \code{/} need at least one argument. With several
arguments, each function combines them from left to right by calling
\term{binary+}, \code{binary-}\index{binary-}, \code{binary*}\index{binary*}
or \code{binary/}\index{binary/}; with one, \code{-} calls
\code{unary-}\index{unary-} and \code{/} calls \code{unary/}\index{unary/}.
The variadic functions themselves are methods and cannot be extended,
but the binary and unary ones are generic functions: to make
\code{+} work on a class of your own, add a method to \code{binary+}.

\begin{dylan}
(define-class <money> (<object>)
  (cents init-keyword: cents:))

(define-method binary+ ((a <money>) (b <money>))
  (make <money> cents: (+ (cents a) (cents b))))

(cents (+ (make <money> cents: 150)
          (make <money> cents: 275)
          (make <money> cents: 5)))      ; => 430
\end{dylan}

\index{division by zero}
Dividing by an exact or a floating zero is an error in both
implementations, as design note~24 requires; so is \code{floor/},
\code{modulo} or any of the division family with a zero divisor:

\begin{dylan}
(handler-case (/ 1 0)   ((<error>) 'caught))   ; => caught
(handler-case (/ 1.0 0) ((<error>) 'caught))   ; => caught
\end{dylan}

\subsection{Contagion}

\index{contagion}
When the arguments of an arithmetic function are of different classes,
the rules are Common Lisp's (\dylanbook{chapter 14}): an operation
between a rational and a float answers a float of the float's format,
and an operation between two floats answers the wider format (single is
narrower than double, double than extended). Rational results are
canonicalized. Because a call with several arguments is a chain of
binary calls, the conversions happen from left to right.

\begin{dylan}
(+ 1 2.0)          ; => 3.0
(+ 1/2 0.5)        ; => 1.0
(* 2 0.5s0)        ; => 1.0s0 (C)   1.0f0 (Lisp)
(+ 1.0s0 1.0d0)    ; => 2.0         a <double-float>
(* 1.5s0 1.5d0)    ; => 2.25        a <double-float>
\end{dylan}

\begin{notebox}
The compiler to C converts a rational to the float's format by way of a
C \code{double}, so a single-float result can differ from the Common Lisp
implementation's in the last place: \code{(+ 1/3 0.5s0)} is
\code{0.8333333s0} in compiled code and \code{0.8333334f0} in the Common
Lisp implementation.
\end{notebox}

\section{Comparing numbers}\label{sec:numbers-comparison}

\index{comparison}\index{=}\index{<}
The comparison functions of \dylanbook{chapter 9} take two or more
arguments and compare them pairwise from left to right:

\begin{center}
\begin{tabular}{ll}
\toprule
Function & True when \\
\midrule
\code{(= a b c ...)} & every argument is \code{=} to the next \\
\code{(/= a b)} & the two arguments are not \code{=} (exactly two) \\
\code{(< a b c ...)} & each argument is less than the next \\
\code{(> a b c ...)} & each argument is greater than the next \\
\code{(<= a b c ...)} & each argument is less than or equal to the next \\
\code{(>= a b c ...)} & each argument is greater than or equal to the next \\
\bottomrule
\end{tabular}
\end{center}

\begin{dylan}
(< 1 2 3)       ; => #t
(< 1 3 2)       ; => #f
(<= 1 1 2)      ; => #t
(>= 3 3 1)      ; => #t
(= 3 3 3)       ; => #t
(/= 1 2)        ; => #t
\end{dylan}

\noindent Numbers are compared by mathematical value, whatever their
classes. A float is compared with a rational exactly (\term{rational
contagion for comparisons}): the float's exact value is used, not a
rounded one. Since the double nearest $0.1$ is not exactly $1/10$, they
are not \code{=}:

\begin{dylan}
(= 1 1.0)                     ; => #t
(= 1/2 0.5)                   ; => #t
(= 0.1 1/10)                  ; => #f
(< 1/3 0.3333333333333333)    ; => #f
\end{dylan}

\index{id?}
\code{=} is not \code{id?}. Two numbers of the same class and value are
\code{id?}; numbers of different classes never are, even when they are
\code{=}:

\begin{dylan}
(id? 1 1)       ; => #t
(id? 1 1.0)     ; => #f
\end{dylan}

\noindent The comparison functions are methods; to extend them to a
class of your own, add methods to the generic functions
\code{binary=}\index{binary=} and \code{binary<}\index{binary<}, on which
all six are built (section~\ref{sec:collections-equality} shows an
example).

\section{Predicates}\label{sec:numbers-predicates}

\index{predicates!numeric}
These generic functions test a number and answer \code{\#t} or
\code{\#f}:

\begin{center}
\begin{tabular}{lll}
\toprule
Function & Argument & True when \\
\midrule
\code{zero?}\index{zero?} & a number & it is zero (\code{0}, \code{0.0}, \code{-0.0}) \\
\code{positive?}\index{positive?} & a real & it is greater than zero \\
\code{negative?}\index{negative?} & a real & it is less than zero \\
\code{odd?}\index{odd?} & an integer & it is odd \\
\code{even?}\index{even?} & an integer & it is even \\
\code{integral?}\index{integral?} & a number & it has an integer value \\
\bottomrule
\end{tabular}
\end{center}

\begin{dylan}
(zero? 0.0)         ; => #t
(positive? -1)      ; => #f
(negative? -1/2)    ; => #t
(odd? 3)            ; => #t
(even? 3)           ; => #f
(integral? 3)       ; => #t
(integral? 2.0)     ; => #t
(integral? 2.5)     ; => #f
(integral? 1/2)     ; => #f
\end{dylan}

\noindent To ask which kind of number an object is, use \code{instance?}
with the classes of section~\ref{sec:numbers-classes}:

\begin{dylan}
(instance? 3 <rational>)     ; => #t
(instance? 1/2 <rational>)   ; => #t
(instance? 1.5 <real>)       ; => #t
(instance? 1.5 <integer>)    ; => #f
\end{dylan}

\section{The division family}\label{sec:numbers-division}

\index{division}
Dylan has the rounding and division functions of Common Lisp.
Each of the one-argument functions rounds a real to an integer and
answers two values: the integer and the remainder, which is the
difference between the argument and the integer. They differ in which
way they round:

\begin{center}
\begin{tabular}{ll}
\toprule
Function & Rounds toward \\
\midrule
\code{floor}\index{floor} & negative infinity \\
\code{ceiling}\index{ceiling} & positive infinity \\
\code{truncate}\index{truncate} & zero \\
\code{round}\index{round} & the nearest integer; a tie goes to the even one \\
\bottomrule
\end{tabular}
\end{center}

\begin{dylan}
(floor 7/2)      ; => 3 and 1/2
(floor -7/2)     ; => -4 and 1/2
(ceiling 2.5)    ; => 3 and -0.5
(truncate -2.7)  ; => -2 and -0.7000000000000002
(round 2.5)      ; => 2 and 0.5       the tie goes to 2, which is even
(round 3.5)      ; => 4 and -0.5
(round -2.5)     ; => -2 and -0.5
\end{dylan}

\noindent Use \code{bind} to receive both values (chapter~\ref{ch:control});
where one value is wanted, the first is used:

\begin{dylan}
(bind ((q r (floor 7/2)))
  (list q r))                ; => (3 1/2)
(+ (round 7/2) 1)            ; => 5
\end{dylan}

\index{floor/}\index{ceiling/}\index{round/}\index{truncate/}
The two-argument functions \code{floor/}, \code{ceiling/},
\code{round/} and \code{truncate/} divide their first argument by their
second, round the quotient as the one-argument function of the same name
does, and answer the rounded quotient and the remainder:

\begin{dylan}
(floor/ 7 2)        ; => 3 and 1
(floor/ -7 2)       ; => -4 and 1
(ceiling/ 7 2)      ; => 4 and -1
(truncate/ -7 2)    ; => -3 and -1
(round/ 7 2)        ; => 4 and -1
(round/ 5 2)        ; => 2 and 1
(floor/ 7.5 2)      ; => 3 and 1.5
\end{dylan}

\index{modulo}\index{remainder}
\code{modulo} answers the remainder of \code{floor/}, which has the sign
of the divisor; \code{remainder} answers the remainder of
\code{truncate/}, which has the sign of the dividend:

\begin{dylan}
(modulo 7 2)        ; => 1
(modulo -7 2)       ; => 1
(modulo 7 -2)       ; => -1
(remainder 7 2)     ; => 1
(remainder -7 2)    ; => -1
(remainder 7 -2)    ; => 1
\end{dylan}

\section{gcd, lcm, abs, min and max}\label{sec:numbers-misc}

\index{gcd}\index{lcm}
\code{gcd} answers the greatest common divisor of its integer arguments
and \code{lcm} their least common multiple. With no arguments they
answer \code{0} and \code{1}; with one, the argument. They are built on
the generic functions \code{binary-gcd}\index{binary-gcd} and
\code{binary-lcm}\index{binary-lcm}, which programs can extend.

\begin{dylan}
(gcd)            ; => 0
(gcd 12 18)      ; => 6
(gcd 12 18 8)    ; => 2
(gcd -12 18)     ; => 6
(lcm)            ; => 1
(lcm 4 6)        ; => 12
(lcm 4 6 10)     ; => 60
(lcm 0 5)        ; => 0
\end{dylan}

\index{abs}
\code{abs} answers the absolute value of a number (for a complex number,
in the Common Lisp implementation, its magnitude); \code{magnitude}\index{magnitude}
is the same for a real.

\begin{dylan}
(abs -5)       ; => 5
(abs -5/3)     ; => 5/3
(abs -2.5)     ; => 2.5
(magnitude -3) ; => 3
\end{dylan}

\index{min}\index{max}
\code{min} and \code{max} take one or more reals and answer the least or
the greatest of them, comparing with \code{binary<}. The answer is one
of the arguments, not converted:

\begin{dylan}
(min 3 1 2)       ; => 1
(max 3 1 2)       ; => 3
(max 1/2 0.4)     ; => 1/2
(max 1 2.0)       ; => 2.0
(min 3 1.5s0)     ; => 1.5s0 (C)   1.5f0 (Lisp)
(min 2)           ; => 2
\end{dylan}

\noindent When several arguments are equal the two implementations may
choose different ones: \code{(min 1 1.0)} is \code{1} in the compiler to
C, which keeps the first of equal arguments, and \code{1.0} in the
Common Lisp implementation, which uses Common Lisp's \code{min}.
Because \code{min} and \code{max} use \code{binary<}, they work on any
objects \code{<} can compare, characters among them:
\code{(max \#\textbackslash{}a \#\textbackslash{}z)} is \code{\#\textbackslash{}z}.

\section{Exponentials and trigonometry}\label{sec:numbers-transcendental}

\index{sqrt}\index{exp}\index{log}\index{expt}
\index{trigonometric functions}
The irrational functions of \dylanbook{chapter 14} are defined as in
Common Lisp:

\begin{center}
\begin{tabular}{ll}
\toprule
Function & Answers \\
\midrule
\code{(sqrt x)} & the square root of $x$ \\
\code{(exp x)} & $e^x$ \\
\code{(log x)} & the natural logarithm of $x$ \\
\code{(expt x y)} & $x^y$ \\
\code{(sin x)}, \code{(cos x)}, \code{(tan x)} & the trigonometric functions, $x$ in radians \\
\code{(asin x)}, \code{(acos x)}, \code{(atan x)} & their inverses \\
\code{(atan2 y x)} & the angle of the point $(x, y)$, from $-\pi$ to $\pi$ \\
\code{(sinh x)}, \code{(cosh x)}, \code{(tanh x)} & the hyperbolic functions \\
\code{(asinh x)}, \code{(acosh x)}, \code{(atanh x)} & their inverses \\
\bottomrule
\end{tabular}
\end{center}

\noindent For a rational argument they answer a \code{<double-float>}; for
a float argument, a float of the same format. So \code{(sqrt 25)} is
\code{5.0}, not \code{5}: the manual's examples show these functions
answering floats, and so do both implementations.

\begin{dylan}
(sqrt 25)        ; => 5.0
(sqrt 2)         ; => 1.4142135623730951
(sqrt 2.0s0)     ; => 1.4142135s0 (C)   1.4142135f0 (Lisp)
(exp 1)          ; => 2.718281828459045
(log 1)          ; => 0.0
(log 100)        ; => 4.605170185988092
(sin 0)          ; => 0.0
(cos 0)          ; => 1.0
(atan 1)         ; => 0.7853981633974483
(atan2 1 1)      ; => 0.7853981633974483
(atan2 -1 -1)    ; => -2.356194490192345
(asin 1)         ; => 1.5707963267948966
(acosh 1)        ; => 0.0
\end{dylan}

\noindent \code{expt} is exact when its base is rational and its power
is an integer, and answers a float otherwise:

\begin{dylan}
(expt 2 10)      ; => 1024
(expt 2 -2)      ; => 1/4
(expt 2/3 3)     ; => 8/27
(expt 10 0)      ; => 1
(expt 2.0 3)     ; => 8.0
(expt 4 1/2)     ; => 2.0
(expt 2.0 0.5)   ; => 1.4142135623730951
\end{dylan}

\noindent Where the true result is not a real number, or not finite, the
implementations behave differently:

\begin{center}
\begin{tabular}{lp{0.25\textwidth}p{0.43\textwidth}}
\toprule
Expression & Compiler to C & Common Lisp implementation \\
\midrule
\code{(sqrt -4)} & error: the result would be complex & \code{\#c(0.0 2.0)} \\
\code{(log 0)} & error: the argument must be positive & error (Lisp's division by zero) \\
\code{(exp 1000)} & \code{+infinity} & error (Lisp's floating-point overflow) \\
\code{(expt 2 100)} & error: integer overflow & \code{1267650600228229401496703205376} \\
\bottomrule
\end{tabular}
\end{center}

\noindent An error from Common Lisp's arithmetic reaches a Dylan
handler as a \code{<simple-error>} whose message begins
\code{Lisp condition:} (chapter~\ref{ch:conditions}).

\section{Integers as bits}\label{sec:numbers-logical}

\index{logical operations}
These functions treat an integer as an infinite two's-complement string
of bits, as Common Lisp does. \code{logior}, \code{logxor},
\code{logand} and \code{logeqv} take any number of integers; with none,
they answer their identities, \code{0}, \code{0}, \code{-1} and
\code{-1}. The rest take exactly the arguments shown.

\begin{center}
\begin{tabular}{ll}
\toprule
Function & Answers \\
\midrule
\code{(logior i ...)}\index{logior} & inclusive or \\
\code{(logxor i ...)}\index{logxor} & exclusive or \\
\code{(logand i ...)}\index{logand} & and \\
\code{(logeqv i ...)}\index{logeqv} & equivalence (exclusive nor) \\
\code{(lognot i)}\index{lognot} & the complement \\
\code{(lognand i j)}\index{lognand} & \code{(lognot (logand i j))} \\
\code{(lognor i j)}\index{lognor} & \code{(lognot (logior i j))} \\
\code{(logandc1 i j)}\index{logandc1} & \code{(logand (lognot i) j)} \\
\code{(logandc2 i j)}\index{logandc2} & \code{(logand i (lognot j))} \\
\code{(logorc1 i j)}\index{logorc1} & \code{(logior (lognot i) j)} \\
\code{(logorc2 i j)}\index{logorc2} & \code{(logior i (lognot j))} \\
\code{(logbit? index i)}\index{logbit?} & whether bit \meta{index} of \meta{i} is 1 \\
\code{(ash i count)}\index{ash} & \meta{i} shifted left \meta{count} bits (right if negative) \\
\bottomrule
\end{tabular}
\end{center}

\begin{dylan}
(logior 1 2 4)      ; => 7
(logand 12 10)      ; => 8
(logxor 12 10)      ; => 6
(logeqv 12 10)      ; => -7
(lognot 5)          ; => -6
(lognand 12 10)     ; => -9
(lognor 12 10)      ; => -15
(logandc1 12 10)    ; => 2
(logandc2 12 10)    ; => 4
(logorc1 12 10)     ; => -5
(logorc2 12 10)     ; => -3
(logbit? 0 5)       ; => #t
(logbit? 1 5)       ; => #f
(ash 1 10)          ; => 1024
(ash 1024 -3)       ; => 128
(ash -8 -1)         ; => -4
\end{dylan}

\noindent In the compiler to C a shift whose result does not fit in 63
bits is an overflow error: \code{(ash 1 70)} is an error there and
\code{1180591620717411303424} in the Common Lisp implementation.

\section{Parts of numbers}\label{sec:numbers-parts}

\index{numerator}\index{denominator}\index{rationalize}
\code{numerator} and \code{denominator} answer the parts of a rational
in lowest terms; an integer's denominator is 1. For a float, both
implementations answer the parts of the fraction the float holds exactly,
as floats:

\begin{dylan}
(numerator 6/4)       ; => 3
(denominator 6/4)     ; => 2
(denominator 5)       ; => 1
(numerator 1.5)       ; => 3.0
(denominator 1.5)     ; => 2.0
(denominator 0.1)     ; => 3.602879701896397e16
\end{dylan}

\noindent The last line shows that the double nearest $0.1$ is a fraction
with a large power of two below it. \code{rationalize} answers the
simplest rational that is within the float's precision of it, which is
usually the fraction you meant:

\begin{dylan}
(rationalize 0.5)     ; => 1/2
(rationalize 0.1)     ; => 1/10
(rationalize 0.333)   ; => 333/1000
(rationalize 1/3)     ; => 1/3
\end{dylan}

\index{real-part}\index{imag-part}\index{angle}
\code{real-part}, \code{imag-part} and \code{angle} work on reals in
both implementations (and on complex numbers in the Common Lisp
implementation): a real is its own real part, its imaginary part is
\code{0}, and its angle is zero for a non-negative number and $\pi$ for a
negative one.

\begin{dylan}
(real-part 5)    ; => 5
(imag-part 5)    ; => 0
(angle 5)        ; => 0 (C)   0.0 (Lisp)
(angle -5)       ; => 3.141592653589793
\end{dylan}

\section{Converting numbers}\label{sec:numbers-conversion}

\index{as!numbers}
\code{as} converts a real to a float of a given format:

\begin{dylan}
(as <float> 1/2)           ; => 0.5        a <double-float>
(as <double-float> 1)      ; => 1.0
(as <single-float> 1/3)    ; => 0.33333334s0 (C)   0.33333334f0 (Lisp)
(as <extended-float> 1)    ; => 1.0l0 (C)   1.0, a <double-float> (Lisp)
\end{dylan}

\noindent To go from a float to an integer, round it with \code{floor},
\code{ceiling}, \code{round} or \code{truncate}; to go to a rational,
use \code{rationalize}, or \code{numerator} and \code{denominator}. The
compiler to C accepts \code{as} only for these float conversions (and
for a number already of the class asked for): \code{(as <integer> 3.0)}
and \code{(as <rational> 0.5)} are errors there. The Common Lisp
implementation also converts a float with an integer value to an
integer, \code{(as <integer> 3.0)} answering \code{3}, and any float to
its exact rational, \code{(as <rational> 0.5)} answering \code{1/2}; it
signals an error for \code{(as <integer> 2.5)}.

\section{How floats print}\label{sec:numbers-printing}

\index{float!printing}
Both implementations print a float with the fewest digits that read back
as the same float, so a float printed and read again is the same float.
A float prints in fixed notation when its magnitude is at least
$10^{-5}$ and less than $10^{16}$, or it is zero, and otherwise with an
exponent. A double-float, the default format, is marked only by its
exponent, written with \code{e}; a single-float ends in \code{s0} (or
\code{s}\meta{n} with an exponent), an extended-float in \code{l0}. An
infinity prints as \code{+infinity} or \code{-infinity} and a NaN as
\code{+nan}; neither can be read back.

\begin{center}
\begin{tabular}{ll}
\toprule
Value & Prints as \\
\midrule
\code{0.1} & \code{0.1} \\
\code{(/ 1.0 3)} & \code{0.3333333333333333} \\
\code{100.0} & \code{100.0} \\
\code{123456789.0} & \code{123456789.0} \\
\code{1e15} & \code{1000000000000000.0} \\
\code{1e16} & \code{1.0e16} \\
\code{1e-4} & \code{0.0001} \\
\code{1e-7} & \code{1.0e-7} \\
\code{(/ 2.0s0 3)} & \code{0.6666667s0} \\
\code{(* 1.5s0 2)} & \code{3.0s0} \\
\bottomrule
\end{tabular}
\end{center}

\noindent Integers print in decimal and ratios as
\code{numerator/denominator}. Printing in general is
the subject of chapter~\ref{ch:printing}.

\section{Characters}\label{sec:characters}

\index{<character>}\index{character}
A \code{<character>} is written \code{\#\textbackslash{}} followed by the
character, or by one of the names \code{space}, \code{newline},
\code{tab}, \code{page}, \code{rubout} and \code{backspace}:

\begin{dylan}
(list #\a #\A #\( #\space #\newline #\tab)
; => (#\a #\A #\( #\space #\newline #\tab)
\end{dylan}

\index{as!characters}
A character's integer is its Unicode code point, in both
implementations. \code{as} converts in either direction
(\dylanbook{chapter 16}):

\begin{dylan}
(as <integer> #\A)        ; => 65
(as <integer> #\newline)  ; => 10
(as <character> 97)       ; => #\a
\end{dylan}

\index{as-uppercase}\index{as-lowercase}
\code{as-uppercase} and \code{as-lowercase} answer the other case of a
character, or the character itself if it has no other case. Both
implementations follow Unicode, not only ASCII: the uppercase of Greek
small lambda, code point 955, is capital lambda, 923.

\begin{dylan}
(as-uppercase #\a)    ; => #\A
(as-lowercase #\Z)    ; => #\z
(as-uppercase #\1)    ; => #\1
(as <integer> (as-uppercase (as <character> 955)))   ; => 923
\end{dylan}

\index{character!comparison}
Characters compare by code point, so case is significant and every
capital ASCII letter is less than every small one:

\begin{dylan}
(< #\a #\b)       ; => #t
(< #\A #\a)       ; => #t
(> #\a #\B)       ; => #t
(= #\a #\A)       ; => #f
(id? #\a #\a)     ; => #t
\end{dylan}

\noindent \code{as-uppercase} and \code{as-lowercase} also work on whole
strings (section~\ref{sec:collections-strings}), which is how to compare
strings without regard to case. A string of characters that do not all
fit in eight bits is a \code{<unicode-string>} in the compiler to C; in
the Common Lisp implementation every string is a \code{<byte-string>},
which holds any character.

\section{Symbols and keywords}\label{sec:symbols}

\index{<symbol>}\index{<keyword>}\index{symbol}\index{keyword}
Symbols and keywords are objects identified by their names
(\dylanbook{chapter 15}). A symbol is written as a name, quoted when it
is to be a value rather than a variable: \code{'apple}. A keyword is a
name followed by a colon, \code{apple:}, and evaluates to itself; the
colon is not part of its name. Symbols name variables; keywords name
keyword arguments and slot options.

\begin{dylan}
(object-class 'apple)     ; => <symbol>
(object-class apple:)     ; => <keyword>
\end{dylan}

\index{symbol!case}
Names are not case sensitive: \code{Foo}, \code{foo} and \code{FOO} are
one symbol, and two symbols with the same name are always \code{id?}.
A symbol does remember a spelling, the one it was first written with,
and prints that way. Converting between symbols, keywords and strings is
done with \code{as}:

\begin{dylan}
(list 'Hello 'hello 'HELLO)           ; => (Hello Hello Hello)
(id? 'Hello 'HELLO)                   ; => #t
(as <symbol> "hello")                 ; => Hello
(as <string> 'hello)                  ; => "Hello"
(as <keyword> "colour")               ; => colour:
(as <string> 'colour:)                ; => "colour"
(as <symbol> "two words")             ; => two words
\end{dylan}

\noindent The spelling remembered is the first anywhere, and the
implementation's own names come first. In the compiler to C that
includes every name the library uses, so \code{'Size} prints as
\code{size} and \code{Test:} as \code{test:} there, while the Common
Lisp implementation, which predefines fewer names, prints them
\code{Size} and \code{Test:}. Either way \code{(id? 'size 'Size)} is
\code{\#t}; only the printing differs.

\index{symbol!comparison}
Symbols are compared only by identity: \code{(= 'abc 'ABC)} and
\code{(id? 'abc 'ABC)} are both \code{\#t}, and there is no ordering, so
\code{(< 'a 'b)} is an error (no applicable method for
\code{binary<}). To sort symbols, sort their names:

\begin{dylan}
(sort '(pear apple fig)
      test: (method (a b) (< (as <string> a) (as <string> b))))
; => (apple fig pear)
\end{dylan}

\chapter{Collections}\label{ch:collections}

A \term{collection} is an object that holds other objects, its
\term{elements}, each under a \term{key}: lists, vectors, strings,
tables, ranges and the rest. Dylan defines collections by a protocol of
generic functions rather than by their representations, so that one
function, \code{map} or \code{sort} or \code{reduce}, works on all of
them, and a class you define joins them by providing a handful of
methods. This chapter follows chapter~12 of the 1992 manual
(\dylanbook{chapter 12}) and the parts of chapter~9 that concern
collections.

The chapter first describes the classes and the protocol every
collection shares, then the functions on all collections and on
sequences, and then each built-in class in turn. Examples show their
results as comments. A list, a vector and a string print the same in
both implementations; the other collections print differently, and
where that matters the examples convert them to lists first. A class in
a result is written by its name alone, as in chapter~\ref{ch:numbers}.

\section{The collection classes}\label{sec:collections-classes}

\index{collection classes}
Collections are classified along two independent lines. A collection
either finds its elements by integer position, as a \term{sequence}
does, or by keys of any kind, as an \term{explicit-key collection} does;
and it either can be changed after it is made, as a \term{mutable
collection} can, or cannot. The abstract classes that express this, and
the concrete classes below them, are:

\begin{itemize}[nosep]
\item \code{<collection>}\index{<collection>}
  \begin{itemize}[nosep]
  \item \code{<explicit-key-collection>}\index{<explicit-key-collection>}
  \item \code{<mutable-collection>}\index{<mutable-collection>}
  \item \code{<sequence>}\index{<sequence>}
    \begin{itemize}[nosep]
    \item \code{<range>}
    \end{itemize}
  \item \code{<mutable-explicit-key-collection>}\index{<mutable-explicit-key-collection>}
    (an explicit-key and a mutable collection)
    \begin{itemize}[nosep]
    \item \code{<table>}
    \item \code{<array>}
    \end{itemize}
  \item \code{<mutable-sequence>}\index{<mutable-sequence>}
    (a mutable collection and a sequence)
    \begin{itemize}[nosep]
    \item \code{<list>}, with subclasses \code{<pair>} and \code{<empty-list>}
    \item \code{<deque>}
    \item \code{<string>}
    \item \code{<vector>} (also an \code{<array>})
      \begin{itemize}[nosep]
      \item \code{<simple-object-vector>}
      \item \code{<stretchy-vector>}
      \item \code{<byte-string>} (also a \code{<string>})
      \item \code{<unicode-string>} (also a \code{<string>})
      \end{itemize}
    \end{itemize}
  \end{itemize}
\end{itemize}

\noindent The concrete classes, the ones a program makes instances of,
are these:

\begin{center}
\begin{tabular}{p{0.32\textwidth}p{0.6\textwidth}}
\toprule
Class & What it is \\
\midrule
\code{<list>}, \code{<pair>}, \code{<empty-list>} & linked lists of pairs (sealed) \\
\code{<simple-object-vector>} & a fixed-size vector of any objects (sealed) \\
\code{<byte-string>} & a fixed-size string of 8-bit characters (sealed) \\
\code{<unicode-string>} & a fixed-size string of any characters (sealed; compiler to C only) \\
\code{<stretchy-vector>} & a vector that grows and shrinks \\
\code{<deque>} & a double-ended queue \\
\code{<range>} & an arithmetic progression, possibly unbounded; not mutable \\
\code{<table>} & a hash table from any keys to any elements \\
\code{<array>} & a multi-dimensional array \\
\bottomrule
\end{tabular}
\end{center}

\noindent \code{<vector>}, \code{<string>} and \code{<array>} are
instantiable too, but \code{make} of them answers an instance of a
concrete subclass: \code{(make <vector> size: 3)} is a
\code{<simple-object-vector>}, \code{(make <string>)} a
\code{<byte-string>}, and an \code{<array>} of one dimension is a
\code{<simple-object-vector>}.

\section{Keys, elements and size}\label{sec:collections-keys}

\index{element}\index{key}
Every collection is a mapping from keys to elements. The keys of a
sequence are the integers from 0 up to its size; the keys of a table are
whatever objects were stored in it; the keys of an array are sequences
of indices. \code{element} answers the element under a key, and
\code{(setter element)} replaces it:

\begin{dylan}
(element '(a b c) 1)                       ; => b
(element "abc" 0)                          ; => #\a
(element #(a b c) 5 default: 'none)        ; => none
(define v (vector 1 2 3))
(set! (element v 0) 'first)
v                                          ; => #(first 2 3)
\end{dylan}

\noindent Asking for a key the collection does not have is an error,
unless a \code{default:} is given, which is then the answer. Setting the
element of a key a collection does not have is an error too, except for
a table, which adds the key.

\index{size}\index{empty?}\index{key-sequence}
\code{size} answers the number of elements, or \code{\#f} for a
collection of unbounded size; \code{empty?} answers whether there are
none; and \code{key-sequence} answers a sequence of the keys:

\begin{dylan}
(size '(1 2 3))           ; => 3
(size "abc")              ; => 3
(size (range from: 1))    ; => #f
(empty? '())              ; => #t
(empty? #(1))             ; => #f
(key-sequence "ab")       ; => a sequence of 0 and 1
\end{dylan}

\noindent The key sequence of a sequence holds the integers from 0 up to
its size (a range, except that the Common Lisp implementation answers a
list for a vector or a string); that of a table is a list of its keys, in
no particular order.

\index{class-for-copy}
\code{class-for-copy} answers the class used when a function such as
\code{map} or \code{reverse} makes a new collection like an existing
one. For most collections it is the collection's own class; a range,
which is not mutable, answers \code{<list>}, and a string answers
\code{<byte-string>}. A collection class that is not mutable must define
a method for it (section~\ref{sec:collections-own}).

\section{The iteration protocol}\label{sec:collections-iteration}

\index{iteration protocol}
All iteration over collections, by \code{for-each}, \code{map},
\code{reduce} and every other function in this chapter, is done with
four generic functions. A collection chooses its own representation for
an \term{iteration state}, an object recording where an iteration has
got to; for every collection, the state \code{\#f} means that the
iteration is finished.

\begin{description}
\item[\code{(initial-state \meta{collection})}]\index{initial-state}
  answers the state for the first element, or \code{\#f} if the
  collection is empty.
\item[\code{(next-state \meta{collection} \meta{state})}]\index{next-state}
  answers the state for the next element, or \code{\#f} if there is
  none. It may change the state it is given, which must not be used
  again.
\item[\code{(current-element \meta{collection} \meta{state})}]\index{current-element}
  answers the element at a state.
\item[\code{(copy-state \meta{collection} \meta{state})}]\index{copy-state}
  answers a state that \code{next-state} will not change, for a program
  that wants to return to a place in an iteration.
\end{description}

\noindent Explicit-key collections add one more, and mutable collections
another:

\begin{description}
\item[\code{(current-key \meta{collection} \meta{state})}]\index{current-key}
  answers the key of the element at a state.
\item[\code{((setter current-element) \meta{collection} \meta{state} \meta{value})}]
  \index{current-element!setter}
  replaces the element at a state; it is written
  \code{(set! (current-element \meta{c} \meta{s}) \meta{value})}.
\end{description}

\noindent And collections that can be walked backward, the vectors, the
strings, the deques and the bounded ranges, also have:

\begin{description}
\item[\code{(final-state \meta{collection})}]\index{final-state}
  answers the state for the last element, or \code{\#f} if the
  collection is empty.
\item[\code{(previous-state \meta{collection} \meta{state})}]\index{previous-state}
  answers the state for the element before, or \code{\#f}.
\end{description}

Here is the protocol used directly. A sequence's keys are counted
alongside its states, since a sequence is walked in key order:

\begin{dylan}
(define v (vector 10 20 30))
(for ((state (initial-state v) (next-state v state))
      (key 0 (+ key 1)))
     ((not state))
  (print (list key (current-element v state))))
; prints (0 10), (1 20), (2 30)

(for ((state (final-state v) (previous-state v state)))
     ((not state))
  (print (current-element v state)))
; prints 30, 20, 10
\end{dylan}

\noindent A table supplies its keys through \code{current-key}, and any
mutable collection can be changed in place as it is walked:

\begin{dylan}
(define t (make <table>))
(set! (element t 'a) 1)
(for ((state (initial-state t) (next-state t state)))
     ((not state))
  (print (list (current-key t state) (current-element t state))))
; prints (a 1)

(define l (list 1 2 3))
(for ((state (initial-state l) (next-state l state)))
     ((not state))
  (set! (current-element l state) (* 10 (current-element l state))))
l    ; => (10 20 30)
\end{dylan}

\noindent \code{copy-state} keeps a place while the iteration goes on:

\begin{dylan}
(bind ((s (initial-state l)))
  (bind ((saved (copy-state l s)))
    (set! s (next-state l s))
    (list (current-element l s) (current-element l saved))))
; => (20 10)
\end{dylan}

\noindent The protocol assumes that a collection does not change while it
is iterated over, except through \code{(setter current-element)}. A
\term{sequence} must also be \term{stable}: every iteration visits its
elements in the same order, the order of their keys, which is its
\term{natural order}. A table need not be: two iterations over the same
table may visit its elements in different orders.

\begin{notebox}
A 1994 design note, number~26, later replaced this protocol with two
functions, \code{forward-iteration-protocol} and
\code{backward-iteration-protocol}. Both implementations follow the 1992 manual: the protocol is the one
described here, and the newer functions do not exist.
\end{notebox}

Most programs never call these functions; they use \code{for-each} and
the other loops of chapter~\ref{ch:control}, or the functions in the
rest of this chapter, all of which are built on them.

\section{Defining your own collection}\label{sec:collections-own}

\index{collection!defining}
To make a new kind of collection, define a subclass of one of the
abstract classes of section~\ref{sec:collections-classes}: a sequence
class for a collection keyed by position, an explicit-key class for one
with keys of its own. Then give it methods for the iteration protocol:

\begin{itemize}
\item every collection: \code{initial-state}, \code{next-state},
  \code{current-element} and \code{copy-state};
\item an explicit-key collection: also \code{current-key};
\item a mutable collection: also \code{(setter current-element)};
\item a collection that is not mutable: also \code{class-for-copy},
  naming a mutable class for copies.
\end{itemize}

\noindent Nothing else is required: \code{element}, \code{size},
\code{map}, \code{reduce}, \code{sort} and the rest all have default
methods that work through the protocol. They may be slow (the default
\code{element} on a sequence walks from the start), so a class that can
do better should also define methods for \code{element},
\code{(setter element)}, \code{size} or \code{member?}.

\subsection{A sequence}

This class is an immutable sequence of the first \meta{count} squares,
computed when asked for. Its iteration state is simply the index:

\begin{dylan}
(define-class <squares> (<sequence>)
  (count init-keyword: count: init-value: 0))

(define-method initial-state ((s <squares>))
  (if (> (count s) 0) 0 #f))

(define-method next-state ((s <squares>) (state <integer>))
  (bind ((next (+ state 1)))
    (if (< next (count s)) next #f)))

(define-method current-element ((s <squares>) (state <integer>))
  (* state state))

(define-method copy-state ((s <squares>) (state <integer>))
  state)

(define-method class-for-copy ((s <squares>)) <list>)
\end{dylan}

\noindent With only these, it is a full collection:

\begin{dylan}
(define sq (make <squares> count: 6))
(size sq)                          ; => 6
(element sq 4)                     ; => 16
(map (method (x) (+ x 1)) sq)      ; => (1 2 5 10 17 26)
(reduce + 0 sq)                    ; => 55
(member? 9 sq)                     ; => #t
(choose odd? sq)                   ; => (1 9 25)
(reverse sq)                       ; => (25 16 9 4 1 0)
(as <vector> sq)                   ; => #(0 1 4 9 16 25)
(= sq '(0 1 4 9 16 25))            ; => #t
(for-each ((x sq)) ()
  (display x) (display " "))       ; prints 0 1 4 9 16 25
\end{dylan}

\noindent \code{map}, \code{choose} and \code{reverse} answer lists,
because \code{class-for-copy} says so.

\subsection{A mutable explicit-key collection}

This class keeps its keys and elements in an association list, a list
of pairs. Its iteration state is the rest of that list, from the current
pair on; it supplies \code{current-key}, and, being mutable,
\code{(setter current-element)}:

\begin{dylan}
(define-class <alist> (<mutable-explicit-key-collection>)
  (entries init-value: '()))

(define-method initial-state ((a <alist>))
  (if (empty? (entries a)) #f (entries a)))
(define-method next-state ((a <alist>) (state <pair>))
  (if (empty? (cdr state)) #f (cdr state)))
(define-method current-element ((a <alist>) (state <pair>))
  (cdr (car state)))
(define-method current-key ((a <alist>) (state <pair>))
  (car (car state)))
(define-method (setter current-element) ((a <alist>) (state <pair>) value)
  (set! (cdr (car state)) value))
\end{dylan}

\noindent The default \code{element} and \code{(setter element)} find a
key by comparing with \code{=}; \code{key-sequence}, \code{size},
\code{do} and the others work as for any collection:

\begin{dylan}
(define a (make <alist>))
(set! (entries a) (list (cons 'red 1) (cons 'green 2) (cons 'blue 3)))
(element a 'green)                ; => 2
(element a 'pink default: 0)      ; => 0
(set! (element a 'blue) 30)
(element a 'blue)                 ; => 30
(size a)                          ; => 3
(reduce + 0 a)                    ; => 33
(any? (method (v) (> v 10)) a)    ; => #t
(do print a)                      ; prints 1, 2 and 30
\end{dylan}

\noindent (The class need not define \code{copy-state}: the default
method answers the state itself, which is right for a state that
\code{next-state} never changes, as here.)

\section{Functions on all collections}\label{sec:collections-functions}

\index{collection functions}
The functions in this section work on every collection. Those that take
several collections process corresponding elements together: elements
with the same key. For sequences that means elements at the same
position, and processing stops at the end of the shortest. Other
collections are \term{aligned}: only the keys all of them have are
processed (\dylanbook{chapter 12}, ``Collection Alignment'').

\subsection{do, map, map-as and map-into}

\index{do}
\code{(do \meta{function} \meta{collection} ...)} calls the function on
each group of corresponding elements, for effect, and answers
\code{\#f}:

\begin{dylan}
(do (method (a b) (print (+ a b)))
    '(100 100 200 200)
    '(1 2 3 4))
; prints 101, 102, 203 and 204, then answers #f
\end{dylan}

\index{map}\index{map-as}
\code{map} answers a new collection of the results, of the
\code{class-for-copy} of the first collection; \code{map-as} takes the
class of the result as its first argument:

\begin{dylan}
(map + '(100 100 200 200) '(1 2 3 4))              ; => (101 102 203 204)
(map + '(1 2 3) #(10 20))                          ; => (11 22)
(map + (range from: 1) '(10 20 30))                ; => (11 22 33)
(map as-uppercase "abc")                           ; => "ABC"
(map-as <vector> + '(100 100 200 200) '(1 2 3 4))  ; => #(101 102 203 204)
\end{dylan}

\noindent Mapping over a table makes a table with the same keys:

\begin{dylan}
(define prices (make <table>))
(set! (element prices 'tea) 3)
(set! (element prices 'cake) 5)
(define doubled (map (curry * 2) prices))
(element doubled 'cake)                            ; => 10
\end{dylan}

\noindent The class of the result decides what it can hold. \code{map}
over a string makes a string, so a function that answers something other
than a character is an error there; ask for another class with
\code{map-as}:

\begin{dylan}
(map-as <list> (method (c) 1) "ab")                ; => (1 1)
(map (method (c) 'x) "ab")                         ; => error
\end{dylan}

\noindent The compiler to C reports the error as \code{a <byte-string>
holds only 8-bit characters: x}; the Common Lisp implementation as
\code{as -{}- not specialized for this class type}. (The Common Lisp
implementation turns an integer result into the character with that code
instead, so \code{(map (method (c) 1) "ab")} answers a string of two
characters there and is an error in the compiler to C.)

\index{map-into}
\code{(map-into \meta{target} \meta{function} \meta{collection} ...)}
stores the results into the mutable collection \meta{target}, replacing
its elements, and answers it. The target may be one of the sources.

\begin{dylan}
(define x (list 100 100 200 200))
(map-into x + '(1 2 3 4) '(10 20 30 40))   ; => (11 22 33 44)
x                                          ; => (11 22 33 44)
(map-into x - x)                           ; => (-11 -22 -33 -44)
\end{dylan}

\noindent The function receives only the elements of the source
collections. The 1992 manual's own example,
\code{(map-into x + '(1 2 3 4))} answering \code{(101 102 203 204)},
also passes the target's elements, contrary to its prose; in both
implementations it answers \code{(1 2 3 4)} (appendix~\ref{app:differences}).
Mapping into a table stores each result under its key, adding keys the
table lacks.

\subsection{any?, every?, reduce and reduce1}

\index{any?}\index{every?}
\code{any?} answers the first true value the function returns for a
group of elements, or \code{\#f}; \code{every?} answers \code{\#t} if the
function is true for every group and \code{\#f} otherwise. Both stop as
soon as the answer is known.

\begin{dylan}
(any? > '(1 2 3 4) '(5 4 3 2))       ; => #t
(any? even? '(1 3 5 7))              ; => #f
(any? (method (x) (and (even? x) (* x 10))) '(1 2 3))   ; => 20
(every? > '(1 2 3 4) '(5 4 3 2))     ; => #f
(every? odd? '(1 3 5 7))             ; => #t
\end{dylan}

\index{reduce}\index{reduce1}
\code{(reduce \meta{function} \meta{initial} \meta{collection})}
combines the elements from the first to the last: it calls the function
on the initial value and the first element, then on that result and the
second element, and so on, and answers the last result, or the initial
value if the collection is empty. \code{reduce1} uses the first element
as the initial value; it is an error on an empty collection.

\begin{dylan}
(reduce max 10 '(3 1 4 1 5 9))    ; => 10
(reduce max 10 '(3 12 9 8 8 6))   ; => 12
(reduce - 0 '(1 2 3))             ; => -6      ((0 - 1) - 2) - 3
(reduce + 0 '())                  ; => 0
(reduce1 + '(1 2 3 4 5))          ; => 15
\end{dylan}

\subsection{member? and find-key}

\index{member?}
\code{(member? \meta{value} \meta{collection} test: \meta{test})}
answers whether the collection has an element that matches the value.
The test defaults to \code{id?} and is called with the value first and
an element second.

\begin{dylan}
(define flavors #(chocolate pistachio pumpkin))
(member? 'chocolate flavors)         ; => #t
(member? 'banana flavors)            ; => #f
(member? 2.0 #(1 2 3))               ; => #f      not id? to 2
(member? 2.0 #(1 2 3) test: =)       ; => #t
\end{dylan}

\noindent On a list, as the 1992 manual says, the true value is the rest
of the list from the matching element on: \code{(member? 2 '(1 2 3))}
answers \code{(2 3)}. (A 1993 design note made it a boolean; neither
implementation follows that.)

\index{find-key}
\code{(find-key \meta{collection} \meta{predicate})} answers the key of
the first element for which the predicate is true, or the
\code{failure:} value (default \code{\#f}). \code{skip:} \meta{n} skips
the first \meta{n} matches.

\begin{dylan}
(find-key flavors (method (f) (id? f 'pistachio)))   ; => 1
(find-key '(1 2 3 4) even?)                          ; => 1
(find-key '(1 2 3 4) even? skip: 1)                  ; => 3
(find-key '(1 2 3 4) even? skip: 2)                  ; => #f
(find-key '(1 3) even? failure: 'none)               ; => none
\end{dylan}

\subsection{replace-elements! and fill!}

\index{replace-elements"!}
\code{(replace-elements! \meta{collection} \meta{predicate}
\meta{new-value-function})} replaces each element for which the
predicate is true by the result of calling the function on it, at most
\code{count:} of them if that is given, and answers the collection:

\begin{dylan}
(replace-elements! (list 10 13 16 19) odd? (method (x) (* 2 x)))
; => (10 26 16 38)
(replace-elements! (list 1 3 5 7) odd? (method (x) 0) count: 2)
; => (0 0 5 7)
\end{dylan}

\index{fill"!}
\code{(fill! \meta{collection} \meta{value})} stores the value under
every key. For a sequence, \code{start:} (inclusive, default 0) and
\code{end:} (exclusive, default the size) limit it to a part:

\begin{dylan}
(fill! (list 'a 'b 'c 'd 'e 'f) 3 start: 2)      ; => (a b 3 3 3 3)
(fill! (vector 1 2 3 4 5) 0 start: 1 end: 3)     ; => #(1 0 0 4 5)
\end{dylan}

\section{Functions on sequences}\label{sec:collections-sequences}

\index{sequence functions}
These functions work on any sequence. Unless a description says
otherwise, a function that answers a new sequence answers one of the
\code{class-for-copy} of its (first) argument and leaves the argument
unchanged; a function whose name ends in \code{!} may change its argument
and reuse its storage, and its result may or may not be the argument
itself, so always use the result. Functions that take a \code{test:}
default it to \code{id?}.

\subsection{Adding and removing}

\index{add}\index{add"!}\index{add-new}\index{add-new"!}
\code{(add \meta{sequence} \meta{element})} answers a new sequence with
the element added; \code{add!} may add it to the sequence itself. Where
the element goes depends on the class: at the front of a list, a vector,
a string and a deque. \code{add!} on a stretchy vector adds at the end;
\code{add} on one adds at the end in the compiler to C and at the front
in the Common Lisp implementation.
\code{add-new} and \code{add-new!} add the element only if the sequence
does not already have it (by \code{test:}).

\begin{dylan}
(define numbers '(3 4 5))
(add numbers 1)             ; => (1 3 4 5)
numbers                     ; => (3 4 5)
(add #(3 4 5) 1)            ; => #(1 3 4 5)
(add "bc" #\a)              ; => "abc"
(add! (list 3 4 5) 1)       ; => (1 3 4 5)
(add-new '(3 4 5) 1)        ; => (1 3 4 5)
(add-new '(3 4 5) 4)        ; => (3 4 5)
\end{dylan}

\index{remove}\index{remove"!}\index{choose}\index{choose-by}
\code{(remove \meta{sequence} \meta{value})} answers the sequence
without the elements that match the value, or without the first
\code{count:} of them; \code{remove!} may change the sequence.
\code{choose} keeps the elements for which a predicate is true, and
\code{choose-by} keeps the elements of one sequence whose partners in
another satisfy a predicate.

\begin{dylan}
(remove '(3 1 4 1 5 9) 1)                 ; => (3 4 5 9)
(remove '(3 1 4 1 5 9) 1 count: 1)        ; => (3 4 1 5 9)
(remove "banana" #\a)                     ; => "bnn"
(remove! (list 3 1 4 1 5 9) 1)            ; => (3 4 5 9)
(choose even? '(3 1 4 1 5 9))             ; => (4)
(choose even? #(1 2 3 4))                 ; => #(2 4)
(choose-by even? (range from: 1) '(a b c d e f g h i))
; => (b d f h)
\end{dylan}

\subsection{Sets}

\index{intersection}\index{union}\index{remove-duplicates}
\index{remove-duplicates"!}
\code{intersection} answers the elements of the first sequence that are
also in the second; \code{union} answers the elements of both, those of
the second only if the first lacks them; \code{remove-duplicates}
answers a sequence without repeated elements, keeping the first of each.
The test of \code{intersection} and \code{union} is called with an
element of the first sequence and one of the second.

\begin{dylan}
(intersection '(john paul george ringo)
              '(richard george edward charles john))
; => (john george)
(union '(butter flour sugar salt eggs)
       '(eggs butter mushrooms onions salt))
; => (butter flour sugar salt eggs mushrooms onions)
(remove-duplicates '(spam eggs spam sausage spam spam spam))
; => (spam eggs sausage)
(remove-duplicates '(1 1.0 2) test: =)   ; => (1 2)
\end{dylan}

\subsection{Copying and joining}

\index{copy-sequence}\index{concatenate}\index{concatenate-as}
\index{replace-subsequence"!}
\code{copy-sequence} answers a new sequence of the elements from
\code{start:} (default 0) up to \code{end:} (default the size).
\code{concatenate} joins sequences into a new one of the first one's
\code{class-for-copy}; \code{concatenate-as} takes the class.
\code{replace-subsequence!} overwrites the elements of a sequence from
\code{start:} on with those of another, which must fit.

\begin{dylan}
(define hamlet '(to be or not to be))
(id? hamlet (copy-sequence hamlet))          ; => #f
(copy-sequence hamlet start: 2 end: 4)       ; => (or not)
(copy-sequence "hello" start: 1)             ; => "ello"
(concatenate "low-" "calorie")               ; => "low-calorie"
(concatenate '(0 1 2) '(3 4 5) '(6 7 8))     ; => (0 1 2 3 4 5 6 7 8)
(concatenate #(1) '(2))                      ; => #(1 2)
(concatenate-as <string> '(#\n #\o #\n) '(#\f #\a #\t))   ; => "nonfat"
(concatenate-as <vector> '(0 1 2) '(3 4 5))  ; => #(0 1 2 3 4 5)
(define phrase (copy-sequence "I hate oatmeal."))
(replace-subsequence! phrase "like" start: 2)  ; => "I like oatmeal."
\end{dylan}

\noindent (\code{phrase} is a copy because a string literal is a
constant; see section~\ref{sec:collections-strings}.)

\subsection{Reversing and sorting}

\index{reverse}\index{reverse"!}\index{sort}\index{sort"!}
\code{reverse} answers the elements in the opposite order and
\code{reverse!} may reuse the sequence's storage. \code{sort} answers the
elements in ascending order by \code{test:}, which defaults to \code{<};
\code{sort!} may reuse the storage. Both sorts are stable, keeping equal
elements in their original order, whether or not \code{stable:} is given.

\begin{dylan}
(define x '(bim bam boom))
(reverse x)                  ; => (boom bam bim)
x                            ; => (bim bam boom)
(reverse "abc")              ; => "cba"
(reverse! (list 1 2 3))      ; => (3 2 1)
(define nums '(3 1 4 1 5 9))
(sort nums)                  ; => (1 1 3 4 5 9)
nums                         ; => (3 1 4 1 5 9)
(sort nums test: >)          ; => (9 5 4 3 1 1)
(sort #(3 1 2))              ; => #(1 2 3)
(sort "hello")               ; => "ehllo"
(sort '("pear" "apple" "fig"))           ; => ("apple" "fig" "pear")
(sort '((b . 2) (a . 1) (c . 2) (d . 1))
      test: (method (x y) (< (cdr x) (cdr y))))
; => ((a . 1) (d . 1) (b . 2) (c . 2))
\end{dylan}

\subsection{Particular elements and searching}

\index{first}\index{second}\index{third}\index{last}
\code{first}, \code{second}, \code{third} and \code{last} answer those
elements, and with \code{set!} replace them. Asking for an element that
is not there is an error.

\begin{dylan}
(first '(a b c))              ; => a
(third '(a b c))              ; => c
(last '(emperor of china))    ; => china
(last #(1 2 3))               ; => 3
(define l (list 1 2 3))
(set! (first l) 'one)
(set! (last l) 'three)
l                             ; => (one 2 three)
\end{dylan}

\index{subsequence-position}
\code{(subsequence-position \meta{big} \meta{pattern})} answers the index
at which \meta{pattern} first occurs in \meta{big}, or the
\code{count:}th occurrence, or \code{\#f}:

\begin{dylan}
(subsequence-position "Ralph Waldo Emerson" "Waldo")   ; => 6
(subsequence-position "banana" "an" count: 2)          ; => 3
(subsequence-position "banana" "x")                    ; => #f
(subsequence-position '(a b c b c) '(b c))             ; => 1
\end{dylan}

\section{Lists and pairs}\label{sec:collections-lists}

\index{<list>}\index{<pair>}\index{<empty-list>}\index{list}
A list is either the empty list, \code{'()}, the only instance of
\code{<empty-list>}, or a \code{<pair>}, whose \term{car} is an element
and whose \term{cdr} is the rest of the list. The classes are sealed. The
empty list is not \code{\#f}: it is a true value.

\index{cons}\index{list*}\index{car}\index{cdr}
\begin{dylan}
(cons 1 2)                  ; => (1 . 2)
(cons 1 '(2 3))             ; => (1 2 3)
(list 1 2 3)                ; => (1 2 3)
(list* 1 2 3 '(4 5 6))      ; => (1 2 3 4 5 6)      the last argument is the tail
(list* 1 2)                 ; => (1 . 2)
(car '(4 5 6))              ; => 4
(cdr '(4 5 6))              ; => (5 6)
(car '())                   ; => ()
(cdr '())                   ; => ()
(object-class '())          ; => <empty-list>
(object-class '(1))         ; => <pair>
\end{dylan}

\noindent \code{(setter car)} and \code{(setter cdr)} change a pair. The
twelve functions \code{caar} through \code{cdddr} compose \code{car} and
\code{cdr} as their names spell (\code{(caddr x)} is
\code{(car (cdr (cdr x)))}), and each has a setter:

\begin{dylan}
(define x (list 4 5 6))
(set! (car x) 9)
(set! (cdr (cdr x)) '(a b))
x                          ; => (9 5 a b)
(cadr '(1 2 3))            ; => 2
(caddr '(1 2 3))           ; => 3
(cddr '(1 2 3))            ; => (3)
(define y (list 1 2 3))
(set! (cadr y) 'two)
y                          ; => (1 two 3)
\end{dylan}

\index{append}\index{find-pair}
\code{append} joins lists; its result shares the last list.
\code{(find-pair \meta{item} \meta{list})} finds the first element of a
list of pairs whose car is the item (by \code{test:}):

\begin{dylan}
(append '(1 2) '(3) '(4 5))               ; => (1 2 3 4 5)
(find-pair 'b '((a . 1) (b . 2)))         ; => (b . 2)
(find-pair 'z '((a . 1)))                 ; => #f
\end{dylan}

\noindent \code{(make <list> size: \meta{n} fill: \meta{x})} makes a list
of \meta{n} copies of \meta{x} (default \code{\#f}). \code{add!} on a list
is \code{cons} and \code{remove!} may splice the list. The \code{size}
of a circular list is \code{\#f}:

\begin{dylan}
(define circle (list 1 2 3))
(set! (cdr (cddr circle)) circle)
(size circle)    ; => #f
\end{dylan}

\begin{warnbox}
Do not let the Common Lisp implementation's read-eval-print loop print a
circular list, as it would by showing the value of the \code{set!} above
typed on its own: it prints forever. Wrap such a form in
\code{(begin ... \#f)}.
\end{warnbox}

\noindent The 1993 design note that renamed \code{cons}, \code{car} and
\code{cdr} to \code{pair}, \code{head} and \code{tail} and removed
\code{list*}, \code{append} and the \code{c...r} functions is not
followed: both implementations have the 1992 names only.

\section{Vectors and arrays}\label{sec:collections-vectors}

\index{<vector>}\index{<simple-object-vector>}\index{vector}
A vector is a one-dimensional array with constant-time access by index.
\code{\#(...)} writes a constant vector, \code{vector} makes one of its
arguments, and \code{make} makes one of a given size:

\begin{dylan}
(vector 1 2 3)                                   ; => #(1 2 3)
(make <vector> size: 3)                          ; => #(#f #f #f)
(make <simple-object-vector> size: 2 fill: 'x)   ; => #(x x)
(as <vector> '(1 2))                             ; => #(1 2)
(as <list> #(1 2))                               ; => (1 2)
\end{dylan}

\noindent A \code{<simple-object-vector>} holds any objects and cannot
change size. Its elements are read and set by \code{element}, or by the
functions of arrays below.

\index{<array>}\index{aref}\index{dimensions}
An \code{<array>} has any number of dimensions; its keys are sequences
of indices. \code{make} takes \code{dimensions:} (required) and
\code{fill:} (default \code{\#f}). \code{aref} reads the element at
indices given as separate arguments, \code{(setter aref)} sets it, and
\code{dimensions} answers the dimensions as a list. A vector is an array
of one dimension.

\begin{dylan}
(define a (make <array> dimensions: '(2 3) fill: 0))
(set! (aref a 1 2) 'x)
(aref a 1 2)            ; => x
(element a '(1 2))      ; => x
(dimensions a)          ; => (2 3)
(size a)                ; => 6
(aref #(7 8 9) 1)       ; => 8
(dimensions #(1 2 3))   ; => (3)
(define vv (vector 7 8 9))
(set! (aref vv 1) 5)
vv                      ; => #(7 5 9)
\end{dylan}

\noindent An array is iterated over in row-major order. A wrong number of
indices, or an index out of bounds, is an error. \code{map} over an
array answers an array with the same dimensions:

\begin{dylan}
(define b (map list a))
(dimensions b)          ; => (2 3)
(aref b 1 2)            ; => (x)
\end{dylan}

\section{Strings}\label{sec:collections-strings}

\index{<string>}\index{<byte-string>}\index{<unicode-string>}\index{string}
A string is a mutable sequence of characters. A \code{<byte-string>}
holds characters with codes below 256, a \code{<unicode-string>} any
characters. \code{"..."} writes a constant string; \code{make} of
\code{<string>} or \code{<byte-string>} takes \code{size:} and
\code{fill:} (default \code{\#\textbackslash{}space}).

\begin{dylan}
(make <string> size: 3 fill: #\z)   ; => "zzz"
(object-class "abc")                ; => <byte-string>
(as <list> "abc")                   ; => (#\a #\b #\c)
(as <string> '(#\h #\i))            ; => "hi"
(size "abc")                        ; => 3
\end{dylan}

\noindent The implementations differ on Unicode. In the compiler to C a
string literal with a character above code 255 is a
\code{<unicode-string>}, and a byte string cannot hold such a character
(storing one is an error). In the Common Lisp implementation every string
is a \code{<byte-string>} that holds any character, and \code{make} of
\code{<unicode-string>} is an error.

\begin{warnbox}
A string literal is a constant, and the compiler to C may share one
literal among all its occurrences, so \code{(id? "abc" "abc")} can be
\code{\#t} there (it is \code{\#f} in the Common Lisp implementation).
Do not change a literal: copy it first with \code{copy-sequence}.
\end{warnbox}

\index{as-uppercase}\index{as-lowercase}\index{as-uppercase"!}
\index{as-lowercase"!}
Strings compare lexicographically by their characters, a prefix being
less than the longer string; \code{=} compares their characters. The
case functions of characters extend to strings: \code{as-uppercase} and
\code{as-lowercase} answer new strings, and \code{as-uppercase!} and
\code{as-lowercase!} change the string they are given.

\begin{dylan}
(< "apple" "banana")       ; => #t
(< "app" "apple")          ; => #t
(< "Zebra" "apple")        ; => #t      Z (90) is less than a (97)
(= "abc" "abc")            ; => #t
(as-uppercase "Van Gogh")  ; => "VAN GOGH"
(as-lowercase "Van Gogh")  ; => "van gogh"
(define s (copy-sequence "Van Gogh"))
(as-uppercase! s)
s                          ; => "VAN GOGH"
(sort '("pear" "Apple" "fig")
      test: (method (a b) (< (as-lowercase a) (as-lowercase b))))
; => ("Apple" "fig" "pear")
\end{dylan}

\section{Stretchy vectors}\label{sec:collections-stretchy}

\index{<stretchy-vector>}
A \code{<stretchy-vector>} is a vector that grows and shrinks.
\code{add!} adds an element at its end and answers the same vector;
\code{remove!} removes elements from the vector itself. Its size can be
set: a larger size adds \code{\#f} elements and a smaller one drops
elements from the end.

\begin{dylan}
(define sv (make <stretchy-vector>))
(add! sv 'a)
(add! sv 'b)
(add! sv 'c)
(size sv)           ; => 3
(element sv 1)      ; => b
(set! (size sv) 5)
(as <list> sv)      ; => (a b c #f #f)
(set! (size sv) 2)
(as <list> sv)      ; => (a b)
(remove! sv 'a)
(as <list> sv)      ; => (b)
\end{dylan}

\noindent \code{make} takes \code{size:} and \code{fill:}. Setting an
element beyond the end is an error in both implementations; use
\code{add!}, or set the size first.

\section{Deques}\label{sec:collections-deques}

\index{<deque>}\index{push}\index{pop}\index{push-last}\index{pop-last}
A \code{<deque>}, a double-ended queue, adds and removes elements at
either end in constant time. \code{push} adds at the front and
\code{push-last} at the back, each answering the deque; \code{pop} and
\code{pop-last} remove and answer the first or the last element, and it
is an error to pop an empty deque. \code{add!} is \code{push}, and
\code{remove!} removes from the deque itself.

\begin{dylan}
(define d (make <deque>))
(push d 2)
(push d 1)
(push-last d 3)
(add! d 0)
(as <list> d)       ; => (0 1 2 3)
(pop d)             ; => 0
(pop-last d)        ; => 3
(as <list> d)       ; => (1 2)
(remove! d 2)
(as <list> d)       ; => (1)
\end{dylan}

\noindent A deque prints as \code{\{<deque> 0 1 2 3\}}.
\code{make} takes \code{size:} and \code{fill:}. The Common Lisp
implementation can also set a deque's size; the compiler to C cannot.

\section{Ranges}\label{sec:collections-ranges}

\index{<range>}\index{range}
A \code{<range>} is an arithmetic progression, computed rather than
stored, and possibly unbounded. \code{range} makes one from keywords:
\code{from:} (default 0) is the first value and \code{by:} (default 1)
the step; \code{up-to:} is an exclusive bound, \code{through:} an
inclusive one, and \code{size:} a number of elements. With none of the
last three the range is unbounded. A negative \code{by:} makes a range
that runs downward, with the bounds below it. \code{make} of
\code{<range>} takes the same keywords.

\begin{dylan}
(as <list> (range from: 1 up-to: 4))              ; => (1 2 3)
(as <list> (range from: 1 through: 4))            ; => (1 2 3 4)
(as <list> (range from: 10 through: 1 by: -3))    ; => (10 7 4 1)
(as <list> (range from: 0 by: 1/2 up-to: 2))      ; => (0 1/2 1 3/2)
(as <list> (make <range> from: 2 size: 3))        ; => (2 3 4)
(size (range from: 1))                            ; => #f
(element (range from: 10 by: 5) 100)              ; => 510
\end{dylan}

\noindent Several functions treat ranges specially so that they finish
even on an unbounded one. \code{member?} computes its answer;
\code{copy-sequence} and \code{reverse} answer ranges (an unbounded range
cannot be reversed); \code{intersection} of two ranges is a range; and
\code{=} compares two ranges by their values:

\begin{dylan}
(member? 1000000 (range))                     ; => #t
(member? 6 (range from: 1 by: 2))             ; => #f
(as <list> (copy-sequence (range from: 5) start: 2 end: 5))  ; => (7 8 9)
(as <list> (intersection (range from: 0 by: 2)
                         (range from: 0 by: 3 up-to: 20)))   ; => (0 6 12 18)
(= (range from: 1 up-to: 4) (range from: 1 through: 3))      ; => #t
(= (range from: 1 up-to: 4) '(1 2 3))                        ; => #t
\end{dylan}

\noindent A range is not mutable; its \code{class-for-copy} is
\code{<list>}, so \code{map} and \code{choose} over a range answer lists:

\begin{dylan}
(map (curry * 2) (range from: 1 through: 3))   ; => (2 4 6)
(choose odd? (range from: 1 through: 6))       ; => (1 3 5)
\end{dylan}

\noindent A range prints as \code{\{range 1 by 1 size 3\}}.

\section{Tables}\label{sec:collections-tables}

\index{<table>}\index{table}\index{remove-key"!}
A \code{<table>} maps keys of any kind to elements. Keys are compared
with \code{=} and hashed with \code{=hash}, so two strings with the same
characters are the same key, and so are \code{1} and \code{1.0}. Setting
the element of a new key adds it; \code{remove-key!} removes one (and
does nothing if the key is absent).

\begin{dylan}
(define t (make <table>))
(set! (element t 'one) 1)
(set! (element t "two") 2)
(element t 'one)                     ; => 1
(element t (copy-sequence "two"))    ; => 2
(element t 'three default: 0)        ; => 0
(size t)                             ; => 2
(remove-key! t 'one)
(key-sequence t)                     ; => ("two")
\end{dylan}

\noindent Asking for a missing key without a \code{default:} is an
error. A table is unstable under iteration: its elements come in an
order of the implementation's choosing, which may change as it grows.
A table prints as \code{\{table "two" => 2\}}.

\section{Equality: = and id?}\label{sec:collections-equality}

\index{=}\index{id?}\index{equality}
Dylan has two equalities. \code{id?} answers whether two objects are the
same object, indistinguishable by any program; \code{=} answers whether
they are equal in value, by the generic function
\code{binary=}\index{binary=}:

\begin{itemize}
\item numbers are \code{=} when they have the same mathematical value
  (chapter~\ref{ch:numbers});
\item collections are \code{=} when they have the same keys and the
  elements under each key are \code{=}; for sequences, the same elements
  in the same positions, whatever the classes of the sequences;
\item pairs, including those of dotted lists, are \code{=} when their
  cars and their cdrs are;
\item anything else is \code{=} only when it is \code{id?}.
\end{itemize}

\begin{dylan}
(= '(1 2) (list 1 2))                 ; => #t
(id? '(1 2) (list 1 2))               ; => #f
(= #(1 2) '(1 2))                     ; => #t
(= "abc" #(#\a #\b #\c))              ; => #t
(= (list 1 (vector 2)) (list 1.0 (vector 2.0)))   ; => #t
(id? 'a 'a)                           ; => #t
(id? #\a #\a)                         ; => #t
\end{dylan}

\noindent The functions that search take a \code{test:} defaulting to
\code{id?}, so pass \code{test: =} to find a value equal to another but
not the same object.

\index{=hash}
To give a class of your own a notion of equality, add a method to
\code{binary=}. If its instances are to be table keys, also add one to
\code{=hash}, which must answer the same integer for any two objects
that are \code{=}:

\begin{dylan}
(define-class <point> (<object>)
  (x init-keyword: x:)
  (y init-keyword: y:))

(define-method binary= ((a <point>) (b <point>))
  (and (= (x a) (x b)) (= (y a) (y b))))

(define-method =hash ((p <point>))
  (+ (* 31 (=hash (x p))) (=hash (y p))))

(define p1 (make <point> x: 1 y: 2))
(define p2 (make <point> x: 1 y: 2))
(id? p1 p2)                           ; => #f
(= p1 p2)                             ; => #t
(member? p2 (list p1))                ; => #f
(member? p2 (list p1) test: =)        ; => a list holding p1, a true value
(define names (make <table>))
(set! (element names p1) "home")
(element names p2)                    ; => "home"
\end{dylan}

\chapter{Functional Operations}\label{ch:functional}

Functions in Dylan are objects: they can be passed as arguments, stored
in variables and collections, and returned as results
(chapter~\ref{ch:functions}). The functions in this chapter make new
functions out of existing ones. They come from chapter~17 of the 1992
manual (\dylanbook{chapter 17}), with \code{identity} from chapter~18
and \code{apply} from chapter~8. Combined with the collection functions
of chapter~\ref{ch:collections}, they let many loops be written as a
single expression.

Each function here answers a new function that accepts any number of
arguments (as many as the functions it was made from need). In both
implementations they are ordinary closures; the examples show their
results as comments, and all of them give the same results in the
compiler to C and the Common Lisp implementation.

\section{compose}\label{sec:functional-compose}

\index{compose}
\code{(compose \meta{f} \meta{g} ...)} answers a function that calls the
last function on its arguments, then each function before it on the
previous result, from right to left. With one argument, \code{compose}
answers that argument. Thus \code{cadr} could be written
\code{(compose car cdr)}:

\begin{dylan}
((compose car cdr) '(1 2 3))              ; => 2
((compose (curry * 2) (curry + 1)) 5)     ; => 12     (5 + 1) * 2
(map (compose - abs) '(-1 2 -3))          ; => (-1 -2 -3)
\end{dylan}

\noindent Only the first value of each inner call is passed on:
\code{((compose list floor) 7/2)} is \code{(3)}, though \code{floor}
answers two values.

This is the manual's example: the distance of a point from the origin,
built from three small functions.

\begin{dylan}
(define-method sum ((numbers <sequence>))
  (reduce1 + numbers))
(define-method square ((x <number>)) (* x x))
(define-method square-all ((coords <sequence>))
  (map square coords))

(define distance (compose sqrt sum square-all))
(distance '(3 4 5))      ; => 7.0710678118654755
\end{dylan}

\section{complement, disjoin and conjoin}\label{sec:functional-logic}

\index{complement}
\code{(complement \meta{predicate})} answers a function that answers
\code{\#t} where the predicate answers \code{\#f}, and \code{\#f}
otherwise:

\begin{dylan}
(map (complement even?) '(1 2 3))          ; => (#t #f #t)
((complement id?) 'a 'b)                   ; => #t
((complement member?) 5 '(1 2))            ; => #t
(define not-empty? (complement empty?))
(choose not-empty? '("a" "" "bc"))         ; => ("a" "bc")
\end{dylan}

\index{disjoin}
\code{(disjoin \meta{p} \meta{q} ...)} answers the \term{disjunction}
of its predicates: a function that calls them in order on its arguments
and answers the first true value any of them returns, without calling
the rest, or \code{\#f} if none does. It is like an \code{or} of calls.

\begin{dylan}
(define small-or-even? (disjoin (curry > 3) even?))
(choose small-or-even? '(1 2 3 4 5 6 7))           ; => (1 2 4 6)
((disjoin (method (x) #f) (method (x) (* x 10))) 4) ; => 40
(choose (disjoin zero? (curry < 10)) '(0 5 11 3 20)) ; => (0 11 20)
\end{dylan}

\index{conjoin}
\code{(conjoin \meta{p} \meta{q} ...)} answers the \term{conjunction}:
a function that calls the predicates in order and answers \code{\#f} as
soon as one does, or else the value of the last. It is like an
\code{and} of calls.

\begin{dylan}
(define positive-integer? (conjoin (rcurry instance? <integer>) positive?))
(map positive-integer? '(5 -5 5.0))                    ; => (#t #f #f)
((conjoin (rcurry instance? <integer>) (method (x) (* x 2))) 4)   ; => 8
(find-key '(3 8 12 5) (conjoin even? (curry < 10)))    ; => 2
\end{dylan}

\section{curry and rcurry}\label{sec:functional-curry}

\index{curry}\index{rcurry}
\code{(curry \meta{function} \meta{arg} ...)} answers a function that
calls \meta{function} with the given arguments followed by its own.
\code{rcurry}, ``right curry'', puts the given arguments after its own
instead:

\begin{dylan}
((curry list 1 2) 3 4)             ; => (1 2 3 4)
((rcurry list 1 2) 3 4)            ; => (3 4 1 2)
(map (curry + 1) '(3 4 5))         ; => (4 5 6)
(map (rcurry - 1) '(5 6 7))        ; => (4 5 6)
\end{dylan}

\noindent So \code{(curry id? 'x)} tests for identity with the symbol
\code{x}; \code{(curry > 6)} is true of numbers less than 6, and
\code{(rcurry > 6)} of numbers greater than 6:

\begin{dylan}
((curry id? 'x) 'x)                ; => #t
((curry > 6) 3)                    ; => #t
((rcurry > 6) 3)                   ; => #f
((curry concatenate "set-") "size" "!")    ; => "set-size!"
(define yuppify (rcurry concatenate ", ayup"))
(yuppify "I'm from New Hampsha")   ; => "I'm from New Hampsha, ayup"
\end{dylan}

\noindent \code{rcurry} is the natural way to fix the second argument of
a two-argument function, as in \code{(rcurry instance? <integer>)}, a
predicate for integers, or \code{(rcurry element 0)}, which takes the
first element of a sequence:

\begin{dylan}
(map (rcurry element 0) '("abc" "xyz"))    ; => (#\a #\x)
\end{dylan}

\section{always and identity}\label{sec:functional-always}

\index{always}
\code{(always \meta{object})} answers a function that ignores any
arguments and answers the object:

\begin{dylan}
((always 1) 'x 'y 'z)       ; => 1
((always #t) #f #f)         ; => #t
((always 'k))               ; => k
(map (always 0) '(a b c))   ; => (0 0 0)
\end{dylan}

\index{identity}
\code{(identity \meta{object})} answers its argument. It is useful
where a function is required and nothing should be done, such as the
start of a chain of compositions:

\begin{dylan}
(identity 5)                ; => 5
(map identity '(1 2))       ; => (1 2)
(define add1 (curry + 1))
((reduce (method (f g) (compose f g)) identity (list add1 add1 add1)) 0)
; => 3
\end{dylan}

\section{apply}\label{sec:functional-apply}

\index{apply}
\code{(apply \meta{function} \meta{arg} ... \meta{sequence})} calls the
function with the given arguments followed by the elements of the
sequence, which may be any sequence (\dylanbook{chapter 8}):

\begin{dylan}
(apply + 1 2 '(3 4))                    ; => 10
(apply max '(3 1 4))                    ; => 4
(apply list '())                        ; => ()
(apply + #(1 2 3))                      ; => 6
(apply + 1 (range from: 1 through: 3))  ; => 7
(apply (curry + 100) '(1 2))            ; => 103
(apply concatenate '("a" "b" "c"))      ; => "abc"
\end{dylan}

\noindent Since \code{map} takes any number of collections, applying it
to a list of lists transposes them:

\begin{dylan}
(apply map list '((1 2 3) (a b c)))     ; => ((1 a) (2 b) (3 c))
\end{dylan}

\section{Idioms}\label{sec:functional-idioms}

\index{idioms}
These operations are at their best as the arguments of \code{map},
\code{choose}, \code{reduce}, \code{any?}, \code{every?} and
\code{find-key}, where they replace small anonymous methods.

\paragraph{Filtering and transforming.} Choose the odd numbers from 1 to
10, double them, and add them up:

\begin{dylan}
(reduce + 0 (map (curry * 2) (choose odd? (range from: 1 through: 10))))
; => 50
\end{dylan}

\paragraph{Negating a test.} Keep the lists that do not contain 2.
Because \code{member?} takes the value first, \code{curry} fixes it:

\begin{dylan}
(choose (complement (curry member? 2)) '((1 2) (3 4) (2 5)))
; => ((3 4))
\end{dylan}

\paragraph{Testing membership in a fixed set.} \code{rcurry} fixes the
collection instead. On a list \code{member?} answers the tail where the
element was found, which \code{any?} passes on as its true value:

\begin{dylan}
(any? (rcurry member? '(a b)) '(x y b))     ; => (b)
\end{dylan}

\paragraph{Counting.} Count the odd numbers greater than 3:

\begin{dylan}
(define count-if (method (pred seq) (size (choose pred seq))))
(count-if (conjoin odd? (curry < 3)) '(1 2 3 4 5 6 7))   ; => 2
\end{dylan}

\paragraph{Comparing by a key.} Sort strings without regard to case,
and find the length of the longest:

\begin{dylan}
(sort '("pear" "Apple" "fig")
      test: (method (a b) (< (as-lowercase a) (as-lowercase b))))
; => ("Apple" "fig" "pear")
(reduce1 max (map size '("pear" "banana" "fig")))    ; => 6
\end{dylan}

\paragraph{Checking several classes.} \code{curry} fixes the object and
leaves the class to vary:

\begin{dylan}
(every? (curry instance? 3) (list <integer> <rational> <number>))   ; => #t
\end{dylan}

\chapter{Conditions}\label{ch:conditions}

A \term{condition} is an object that represents an exceptional
situation: an error, a warning, or a request to recover in some
particular way. A program \term{signals} a condition when the situation
arises, and the condition system finds a \term{handler} for it, a
function that some caller established in advance, and calls it. This
chapter follows chapter~13 of the 1992 manual (\dylanbook{chapter 13}),
which both implementations provide complete; it also describes what each
implementation does with a condition no handler takes, and the errors
the implementations themselves signal. As in chapter~\ref{ch:numbers}, a
class in a result is written by its name alone.

\section{How the condition system works}\label{sec:conditions-overview}

\index{signal}\index{handler}
Dylan's conditions are objects, instances of subclasses of
\code{<condition>}, and a handler is chosen by the condition's class: a
handler for \code{<error>} applies to every error, one for
\code{<type-error>} only to type errors. Handlers are established
dynamically, for the duration of a form's evaluation, and the most
recently established applicable handler is tried first.

Signalling is a function call. When a condition is signalled, the
signaller is still active and the handler runs ``on top of'' it. The
handler then does one of three things:

\begin{itemize}
\item It \textbf{exits} non-locally, abandoning the computation that
  signalled, to a point established outside it. This is the common case,
  and \code{handler-case} does it for you.
\item It \textbf{returns} values, which become the values of
  \code{signal} in the signaller. This is allowed only for conditions
  whose \term{recovery protocol} says what the values mean, such as
  warnings.
\item It \textbf{declines}, by calling the next handler, letting an
  older handler decide.
\end{itemize}

\noindent If every handler declines, or there are none, the condition's
\term{default handler} runs. For a warning it prints the message and
returns; for an error it stops the program, or in the Common Lisp
implementation enters the debugger.

Recovery can also be formal. A \term{restart} is a condition that a
handler signals to ask for a particular recovery, and a \term{restart
handler}, established by the code that knows how to recover, handles it.
Section~\ref{sec:conditions-restarts} describes restarts.

\section{The condition classes}\label{sec:conditions-classes}

\index{condition classes}
The predefined condition classes are these:

\begin{itemize}[nosep]
\item \code{<condition>}\index{<condition>}
  \begin{itemize}[nosep]
  \item \code{<serious-condition>}\index{<serious-condition>}
    \begin{itemize}[nosep]
    \item \code{<error>}\index{<error>}
      \begin{itemize}[nosep]
      \item \code{<simple-error>}\index{<simple-error>}
      \item \code{<type-error>}\index{<type-error>}
      \end{itemize}
    \end{itemize}
  \item \code{<warning>}\index{<warning>}
    \begin{itemize}[nosep]
    \item \code{<simple-warning>}\index{<simple-warning>}
    \end{itemize}
  \item \code{<restart>}\index{<restart>}
    \begin{itemize}[nosep]
    \item \code{<simple-restart>}\index{<simple-restart>}
    \item \code{<abort>}\index{<abort>}
    \end{itemize}
  \end{itemize}
\end{itemize}

\begin{description}
\item[\code{<condition>}] The class of all conditions. Its default
  handler returns \code{\#f}, so a plain condition nobody handles is
  ignored.
\item[\code{<serious-condition>}] Conditions that cannot safely be
  ignored. Its default handler is the debugger of last resort
  (section~\ref{sec:conditions-unhandled}).
\item[\code{<error>}] Serious conditions that represent something wrong
  with the program, as distinct from its environment; a handler for
  \code{<error>} does not catch other serious conditions.
\item[\code{<simple-error>}] An error that is just a message: a format
  string and arguments. They are given by the init-keywords
  \code{format-string:} and \code{format-arguments:}, and the functions
  \code{condition-format-string}\index{condition-format-string} and
  \code{condition-format-arguments}\index{condition-format-arguments}
  read them. \code{error} makes one from a string.
\item[\code{<type-error>}] An error saying a value was not of the
  expected type. Its init-keywords are \code{value:} and \code{type:},
  read by \code{type-error-value}\index{type-error-value} and
  \code{type-error-expected-type}\index{type-error-expected-type}.
  \code{check-type} signals one, and so do the implementations when a
  typed slot or variable is given a value of the wrong type.
\item[\code{<warning>}] Conditions worth reporting but not stopping
  for. The default handler prints the message and returns \code{\#f}.
\item[\code{<simple-warning>}] A warning that is a format string and
  arguments, like \code{<simple-error>}; \code{signal} makes one from a
  string.
\item[\code{<restart>}] Requests to recover. The default handler signals
  an error saying no handler was established for the restart. The
  init-keyword \code{condition:} is accepted, so that a restart can name
  the condition it recovers from.
\item[\code{<simple-restart>}] The restart \code{cerror} and
  \code{break} establish handlers for.
\item[\code{<abort>}] The restart that asks to abandon the current
  command and return to something like a command loop.
\end{description}

\noindent Conditions are made with \code{make} like any object. A
condition prints as its message:

\begin{dylan}
(print (make <simple-error> format-string: "bad value: ~s"
                            format-arguments: '(42)))
; prints bad value: 42
\end{dylan}

\section{Signalling}\label{sec:conditions-signalling}

\index{error}
\code{(error \meta{condition})} signals a condition and never returns.
\code{(error \meta{string} \meta{argument} ...)} makes a
\code{<simple-error>} from a format string and arguments and signals it.
This is the usual way to report an error:

\begin{dylan}
(define-method withdraw (balance amount)
  (if (> amount balance)
      (error "cannot withdraw ~a from a balance of ~a" amount balance)
      (- balance amount)))
\end{dylan}

\noindent If a handler for the error returns instead of exiting,
\code{error} does not return either: it reports the error as unhandled
(section~\ref{sec:conditions-unhandled}).

\index{signal}
\code{(signal \meta{condition})} signals a condition and answers whatever
values the handler returns. \code{(signal \meta{string} \meta{argument}
...)} makes a \code{<simple-warning>} and signals it. With no handler, a
warning is printed and \code{signal} answers \code{\#f}:

\begin{dylan}
(signal "disk is ~a% full" 91)
; prints a warning:  disk is 91% full
; => #f
\end{dylan}

\noindent The compiler to C prints the warning on standard error as
\code{Dylan warning: disk is 91\% full}; the Common Lisp implementation
uses Common Lisp's \code{warn}, which prints
\code{WARNING: disk is 91\% full}.

\index{format directives}
The format strings of \code{error}, \code{signal}, \code{cerror} and
\code{break} understand the directives \code{\textasciitilde{}a} and
\code{\textasciitilde{}d} (the argument as \code{display} prints it),
\code{\textasciitilde{}s} (as \code{print} prints it),
\code{\textasciitilde{}\%} (a new line) and \code{\textasciitilde{}\textasciitilde{}}
(a tilde). Other directives are copied as they are, and arguments no
directive uses are printed after the message, separated by spaces.

\index{cerror}
\code{(cerror \meta{restart-description} \meta{condition-or-string}
\meta{argument} ...)} is a \term{continuable error}. It signals the error
as \code{error} does, but first establishes a handler for
\code{<simple-restart>} described by \meta{restart-description}. If a
handler (or, in the Common Lisp implementation, a person at the
debugger) signals a \code{<simple-restart>}, \code{cerror} returns
\code{\#f} and the program carries on with the correction the
description promised:

\begin{dylan}
(define-method careful-sqrt (x)
  (if (negative? x)
      (begin (cerror "Use the absolute value" "negative argument ~a" x)
             (sqrt (abs x)))
      (sqrt x)))

(handler-bind (<error> (method (c next) (signal (make <simple-restart>))))
  (careful-sqrt -16))                            ; => 4.0
\end{dylan}

\index{check-type}
\code{(check-type \meta{value} \meta{type})} answers the value if it is
an instance of the type, and signals a \code{<type-error>} if not:

\begin{dylan}
(check-type 5 <integer>)        ; => 5
(handler-case (check-type "abc" <integer>)
  ((<type-error> condition: c)
   (list (type-error-value c) (type-error-expected-type c))))
; => ("abc" <integer>)
\end{dylan}

\noindent The implementations differ when a handler returns from the
type error: the compiler to C signals it with \code{error}, so a
returning handler makes it unhandled; the Common Lisp implementation
signals it with \code{signal}, so \code{check-type} answers whatever the
handler returned.

\index{abort}
\code{(abort)} signals an \code{<abort>} with \code{error}; it never
returns. A handler for \code{<abort>} decides what abandoning means:

\begin{dylan}
(handler-case (begin (abort) 'not-reached)
  ((<abort>) 'aborted))         ; => aborted
\end{dylan}

\index{break}
\code{(break)} or \code{(break \meta{string} \meta{argument} ...)} is a
breakpoint. In the Common Lisp implementation it enters the debugger
with the message, offering a restart that continues; in the compiler to
C, which has no debugger, it prints \code{Dylan break:} and the message
on standard error and continues. Either way it then answers \code{\#f}.

\section{Handling with handler-case}\label{sec:conditions-handler-case}

\index{handler-case}
\code{handler-case} is the usual way to handle conditions: it evaluates a
form, and if a condition of one of the listed classes is signalled during
it, abandons the form and evaluates the matching clause instead.

\begin{dylan}
(handler-case protected-form
  ((type option ...) form ...)
  ...)
\end{dylan}

\noindent If the protected form finishes, \code{handler-case} answers its
values. If a condition is signalled that is an instance of a clause's
\meta{type} (a class, or a singleton), the protected form is exited and
that clause's forms are evaluated, their last value being the answer, or
\code{\#f} if there are none. The clauses are tried in order, so put the
more specific classes first. The options are:

\begin{description}
\item[\code{condition: \meta{name}}] binds the variable \meta{name} to the
  condition in the clause's forms;
\item[\code{test: \meta{function}}] a predicate the condition must
  also satisfy for the clause to apply;
\item[\code{description: \meta{string}}] a description of the handler,
  which introspection (\code{do-handlers}) can see.
\end{description}

\begin{dylan}
(handler-case (error "bad value: ~s" 42)
  ((<simple-error> condition: c)
   (list (condition-format-string c) (condition-format-arguments c))))
; => ("bad value: ~s" (42))

(handler-case (error "code ~a" 7)
  ((<simple-error> test: (method (c) (= (first (condition-format-arguments c)) 8)))
   'eight)
  ((<error>) 'other))
; => other

(handler-case (error "x")
  ((<error>) 'first)
  ((<simple-error>) 'second))
; => first

(handler-case (values 1 2)
  ((<error>) 'oops))
; => 1 and 2
\end{dylan}

\noindent The clause's forms run after the exit, in the dynamic
environment of the \code{handler-case} itself, so its handlers are no
longer in force there. A singleton type handles one particular
condition object:

\begin{dylan}
(define c1 (make <simple-error> format-string: "one"))
(define c2 (make <simple-error> format-string: "two"))
(define-method try (c)
  (handler-case (error c)
    (((singleton c1)) 'got-c1)
    ((<simple-error> condition: e) (list 'other (condition-format-string e)))))
(list (try c1) (try c2))     ; => (got-c1 (other "two"))
\end{dylan}

\section{Handling with handler-bind}\label{sec:conditions-handler-bind}

\index{handler-bind}\index{next-handler}
\code{handler-bind} is the primitive form, and both implementations
have it. It establishes one handler, a function, while its body is
evaluated, and answers the body's values:

\begin{dylan}
(handler-bind (type function option ...)
  form ...)
\end{dylan}

\noindent The type, function and options are evaluated before the body;
the options are \code{test:} and \code{description:}, as for
\code{handler-case}. When an applicable condition is signalled, the
function is called with two arguments, the condition and a
\term{next-handler} function, without exiting anything. It may exit
non-locally (with \code{bind-exit}, chapter~\ref{ch:control}); it may
return values, which \code{signal} then answers; or it may decline by
calling the next-handler with no arguments, which tries the next older
applicable handler.

\begin{dylan}
(handler-bind (<warning> (method (c next) (print 'outer) 'outer-value))
  (handler-bind (<warning> (method (c next) (print 'inner-declines) (next)))
    (signal "low battery")))
; prints inner-declines, then outer
; => outer-value
\end{dylan}

\noindent There is no ``condition wall'': while a handler runs, every
handler established before it is still in force, including those
established inside the signalling code. That is what lets a handler
signal a restart that the signaller handles
(section~\ref{sec:conditions-restarts}). It also means that a handler
which signals a condition of its own class, or causes an error it also
handles, can recurse until the stack is exhausted.

\section{Default handlers}\label{sec:conditions-default-handler}

\index{default-handler}
When no dynamic handler handles a condition, \code{signal} calls
\code{default-handler}, a generic function, on it. Its predefined methods
are:

\begin{center}
\begin{tabular}{ll}
\toprule
Method on & Does \\
\midrule
\code{<condition>} & answers \code{\#f} \\
\code{<serious-condition>} & reports the condition as unhandled \\
\code{<warning>} & prints the message, answers \code{\#f} \\
\code{<restart>} & signals an error: no handler was established for the restart \\
\bottomrule
\end{tabular}
\end{center}

\noindent A condition class of your own inherits the method of its
superclass. A condition that is neither serious nor a warning is simply
ignored when unhandled:

\begin{dylan}
(define-class <note> (<condition>)
  (text init-keyword: text:))

(signal (make <note> text: "hello"))          ; => #f
(handler-case (signal (make <note> text: "hi"))
  ((<note> condition: n) (text n)))           ; => "hi"
\end{dylan}

\noindent A program may call \code{default-handler} itself, and may add
methods to it for its own classes.

\section{Restarts}\label{sec:conditions-restarts}

\index{restart}
A restart is a condition that asks for a particular way of recovering.
The code that knows how to recover establishes a handler for a restart
class, usually with \code{handler-case}, around the code that may fail;
a handler for the original condition, which knows which recovery it
wants, signals the restart. Values a recovery needs travel in the
restart's slots.

This program parses digits. When a character is not a digit it signals
an error, but offers a \code{<use-value>} restart that supplies a value
instead; the caller's handler chooses that recovery:

\begin{dylan}
(define-class <use-value> (<restart>)
  (value init-keyword: value:))

(define-method parse-digit ((c <character>))
  (if (and (>= c #\0) (<= c #\9))
      (- (as <integer> c) (as <integer> #\0))
      (handler-case (error "not a digit: ~s" c)
        ((<use-value> condition: r) (value r)))))

(handler-bind (<simple-error>
               (method (c next) (signal (make <use-value> value: -1))))
  (map parse-digit '(#\4 #\x #\2)))
; => (4 -1 2)
\end{dylan}

\noindent The handler for \code{<simple-error>} runs while
\code{parse-digit}'s \code{handler-case} is still in force, so its
\code{signal} reaches the restart handler, which exits from the
\code{handler-case} with the value.

\index{<simple-restart>}
\code{cerror} establishes a \code{<simple-restart>} handler, as shown in
section~\ref{sec:conditions-signalling}; signalling a
\code{<simple-restart>} where none is established is an error. The
\code{condition:} init-keyword lets a restart be tied to one condition,
so that a \code{test:} on its handler can accept only restarts meant
for that condition (\dylanbook{chapter 13} has an example).

\section{Introspection}\label{sec:conditions-introspection}

\index{do-handlers}
\code{(do-handlers \meta{function})} calls the function on each handler
in force, the most recent first, with four arguments: the handler's type,
test, function and description. An absent test is given as a function
that answers \code{\#t}, and an absent description as \code{\#f}.

\begin{dylan}
(handler-bind (<warning> (method (c n) #f) description: "ignore warnings")
  (handler-bind (<error> (method (c n) #f))
    (do-handlers (method (type test function description)
                   (print (list type description))))))
; prints (<error> #f), then (<warning> "ignore warnings")
\end{dylan}

\noindent At the Common Lisp implementation's prompt the list ends with
the loop's own handler, an \code{<abort>} described as
\code{"Return to the read-eval-print loop."}

\index{return-allowed?}\index{return-description}
\index{restart-query}\index{return-query}
Four generic functions let an interactive debugger ask whether a
condition may be handled by returning (\code{return-allowed?}), what
returning would mean (\code{return-description}), and the user for the
contents of a restart (\code{restart-query}) or for values to return
(\code{return-query}). The compiler to C, having no debugger, defines
only methods that do nothing: all four answer \code{\#f}. In the Common
Lisp implementation \code{return-allowed?} answers \code{\#t} for a
warning, whose recovery protocol allows returning, and
\code{return-description} then answers the string
\code{"Return; the value is ignored."}; its \code{return-query} reads
values typed by the user.

\section{Defining your own conditions}\label{sec:conditions-defining}

\index{condition!defining}
Define a condition class as a subclass of the predefined class closest
in meaning: \code{<error>} for a program error, \code{<warning>} for
something worth reporting, \code{<restart>} for a recovery. Slots carry
whatever a handler needs to know.

The simplest way to give a condition a message that both
implementations report is to subclass \code{<simple-error>} (or
\code{<simple-warning>}) and pass \code{format-string:} and
\code{format-arguments:} when making it:

\begin{dylan}
(define-class <bad-reading> (<simple-error>)
  (sensor init-keyword: sensor:))

(define-method report-bad (sensor value)
  (error (make <bad-reading> sensor: sensor
               format-string: "sensor ~a read ~a"
               format-arguments: (list sensor value))))

(handler-case (report-bad 'north 9999)
  ((<bad-reading> condition: c) (list (sensor c) c)))
; => (north sensor north read 9999)
\end{dylan}

\noindent Unhandled, it is reported as \code{sensor south read -1}. A
direct subclass of \code{<error>} with no message is reported by its
class only: in the compiler to C as
\code{Dylan error: \{instance of <oops>\}}, even if the class has a
\code{print-object} method (chapter~\ref{ch:printing}); that method is
used when the program prints the condition itself.

\section{Unhandled conditions}\label{sec:conditions-unhandled}

\index{unhandled condition}\index{debugger}
What happens to a serious condition that no handler takes, the
``debugger of last resort'' of the manual, is where the implementations
differ most.

\subsection{In the Common Lisp implementation}

An unhandled error enters the Common Lisp debugger, which shows the
message and offers restarts: Dylan's restart handlers in force, then
restarts that return to the Dylan prompt, then Common Lisp's own. Type
the number of a restart to choose it. Here a \code{cerror} offers its
restart:

\begin{repl}
? (define-method checked (n)
    (when (> n 10)
      (cerror "Use 10 instead" "~a is too big" n)
      (set! n 10))
    n)

;Value: checked

? (checked 99)

debugger invoked on a DYLAN::DYLAN-CONDITION in thread
#<THREAD tid=712328 "main thread" RUNNING {1200030003}>:
  Error: 99 is too big

Type HELP for debugger help, or (SB-EXT:EXIT) to exit from SBCL.

restarts (invokable by number or by possibly-abbreviated name):
  0: [DYLAN-RESTART] Use 10 instead
  1:                 Return to the read-eval-print loop.
  2: [DYLAN-REP    ] Return to the Dylan read-eval-print loop.
  3: [ABORT        ] Exit debugger, returning to top level.

((LAMBDA NIL :IN DYLAN::IMPLEMENTATION-SPECIFIC/INDUCE-ERROR))
0] 0

;Value: 10
\end{repl}

\noindent Choosing \code{Return to the read-eval-print loop} abandons the
form and prints \code{;Aborted}; \code{DYLAN-REP} abandons it
without the message. \code{break} offers a \code{CONTINUE}
restart, which makes \code{break} answer \code{\#f}.

An error signalled by Common Lisp itself (a Lisp type error inside a
built-in function, a floating-point overflow) is offered first to Dylan's
handlers as a \code{<simple-error>} whose message begins
\code{Lisp condition:}. If none takes it, it reaches the debugger as the
Lisp error it is, and the debugger offers \code{DYLAN-REP}, which
returns to the Dylan prompt, and Common Lisp's own restarts. In a program compiled to Lisp with
\code{dylan:compile-file} and loaded into the runtime, an unhandled error
likewise goes to the Common Lisp debugger (or, when SBCL runs with its
debugger disabled, ends SBCL with a backtrace). Chapter~\ref{ch:lisp} describes the debugger further.

\subsection{In the compiler to C}

A compiled program has no debugger. An unhandled serious condition in
the main thread prints \code{Dylan error:} and the condition's message
on standard error and ends the program with exit status 1. Anything the
program had written to standard output is flushed first.

\begin{shell}
$ cat range.dyl
(print "before")
(error "the value ~s is out of range" 42)
(print "after")
$ bin/dylanc range.dyl
$ ./range
"before"
Dylan error: the value 42 is out of range
$ echo $?
1
\end{shell}

\noindent The cleanup forms of \code{unwind-protect} that are in force are
\emph{not} run when a program stops this way. Some other cases:

\begin{center}
\begin{tabular}{p{0.38\textwidth}p{0.55\textwidth}}
\toprule
Program & Reported \\
\midrule
\code{(check-type "x" <integer>)} & \code{Dylan error: "x" is not an instance of \{the class <integer>\}} \\
a handler returns from \code{(error "nope")} & \code{Dylan error: a handler returned from error: nope} \\
\code{(abort)} with no handler & \code{Dylan error: no handler was established for the restart \{instance of <abort>\}} \\
\bottomrule
\end{tabular}
\end{center}

\noindent A condition a thread other than the main one does not handle
ends that thread instead, as chapter~\ref{ch:extensions} describes.

\section{Errors the implementations signal}\label{sec:conditions-runtime-errors}

\index{runtime errors}
Errors the implementations detect themselves are signalled as Dylan
conditions, so programs can handle them like any other. They are
\code{<type-error>}s when a value of the wrong type is stored in a typed
slot or variable, and \code{<simple-error>}s otherwise. In the compiler
to C:

\begin{center}
\begin{tabular}{p{0.36\textwidth}p{0.57\textwidth}}
\toprule
Cause & Message \\
\midrule
no method applies to the arguments & \code{no applicable method for f} \\
wrong number of arguments & \code{wrong number of arguments (2)} \\
a keyword the function does not take & \code{keyword argument not accepted: b:} \\
calling something not a function & \code{not a function: 5} \\
integer overflow & \code{integer overflow (integers are 63 bits)} \\
division by zero & \code{division by zero} \\
a typed slot or variable given a wrong value & \code{"s" is not an instance of \{the class <integer>\}} (a \code{<type-error>}) \\
reading an uninitialized slot & \code{the slot s is not initialized} \\
a variable with no value & \code{the variable x@dylan-user has no value} \\
\code{car} of a non-list & \code{not a list: 5} \\
an element that does not exist & \code{element: there is no element at the index 5} \\
\code{make} of an abstract class & \code{cannot make an instance of this class: \{the class <integer>\}} \\
\bottomrule
\end{tabular}
\end{center}

\noindent The compiler also warns, when compiling, about a variable that
is not defined anywhere (\code{warning: x is not defined}); the program
is still built, and signals the error above if the reference is
evaluated.

The Common Lisp implementation detects the same errors with messages of
its own, for example \code{generic-dispatch -- no applicable methods},
\code{invalid number of arguments: 3}, \code{illegal keyword argument b:
(a:)} and \code{unbound variable x}. An error that arises inside Common
Lisp, such as \code{(car 5)}, has a message beginning
\code{Lisp condition:} followed by Common Lisp's own description.

\begin{dylan}
(define-method f ((x <integer>)) x)

(handler-case (f "s")
  ((<error> condition: c) c))
; C:    no applicable method for f
; Lisp: generic-dispatch -- no applicable methods {the generic function f} ("s")

(define-class <account> (<object>)
  (balance type: <integer> init-value: 0))

(handler-case (set! (balance (make <account>)) "lots")
  ((<type-error> condition: c) (type-error-value c)))
; => "lots"
\end{dylan}

\chapter{Modules}\label{ch:modules}

\index{module}
A small program can live comfortably in one namespace: every class,
generic function and variable it defines sits side by side, and every
name means the same thing everywhere. A larger program soon outgrows
that. Two parts of it want to use the same name for different things, or
one part wants to keep its helpers private so that nobody else comes to
depend on them. Dylan's answer is the \term{module}.

This chapter explains what a module is, how the \code{define-module} form
declares one, how a source file says which module its code belongs to,
and how each implementation handles programs spread over several modules
and several files. It ends with a complete multi-file program.

\section{What a module is}\label{sec:modules-what}

\index{module variable}
A module is a mapping from names to \term{module variables}. It is not a
container for classes or functions as such: Dylan modules deal only in
variables. Since classes and functions are ordinary values stored in
variables, controlling who can see a variable controls who can use the
class or function in it. If you export the variable \code{<point>}, other
modules can make points; if you do not, the class is private
\dylanbook{chapter 5}.

A module gets its variables in two ways.

\begin{itemize}
\item It may \term{own} a variable. A definition compiled in the
  module creates the variable there if the name does not already refer
  to one. The defining forms are \code{define}, \code{define-class},
  \code{define-method} and \code{define-generic-function}.
\item It may \term{import} a variable that another module exports, by
  \term{using} that module. An imported variable is the very same
  variable as in the module that owns it, not a copy: an assignment made
  through one module is seen through every other.
\end{itemize}

Each variable is owned by exactly one module, but it may be visible in
many. Within one module a name refers to at most one variable; it is an
error to import two different variables under the same name, or to import
a name and also define it.

Only \term{exported} variables can be imported. A module exports a
variable it owns by naming it in an \code{export} or \code{create} clause,
and it can pass on variables it imported with the \code{export:} option
of a \code{use} clause.

\index{lexical variable}
Modules affect only module variables. A local variable made by
\code{bind}, by a method's parameters or by any other binding form is
looked up lexically, as always, and hides a module variable of the same
name.

\section{The predefined modules}\label{sec:modules-predefined}

\index{dylan module}\index{dylan-user module}
Two modules exist before any program runs.

\begin{description}
\item[\code{dylan}] owns every predefined variable of the language:
  \code{+}, \code{make}, \code{<object>}, \code{map} and the rest
  (Appendix~\ref{app:library} lists them all). It exports them all. It
  cannot be redefined. In the compiler to C the library itself is
  written in Dylan, in the module \code{dylan}, and it also owns a few
  internal variables, whose names begin with \code{\%}, which it does not
  export.
\item[\code{dylan-user}] uses \code{dylan} and owns nothing else to begin
  with. Code goes here unless something says otherwise, which is why the
  examples in the rest of this book never mention modules at all.
\end{description}

\section{Declaring a module}\label{sec:modules-define}

\index{define-module}
The 1992 manual says that modules are ``manipulated by the programming
environment'' and gives no syntax for declaring them. Apple's later
design note~32 supplied one, in the infix syntax. Both implementations
accept it in prefix form:

\begin{dylan}
(define-module module-name clause ...)
\end{dylan}

\noindent
Each \meta{clause} is a \code{use} clause, an \code{export} clause or a
\code{create} clause. The module name is a symbol. Module names live in a
namespace of their own, separate from variables, so \code{define-module}
creates no variable, and a module may share its name with a variable.
\code{define-module} is a defining form: it belongs at top level or
inside a top-level \code{begin}. The compiler to C reports one inside a
method (``define-module must be at top level''); the interactive
implementation does not check.

\subsection{\code{use} clauses}\label{sec:modules-use}

\index{use clause}
A \code{use} clause imports variables from another module:

\begin{dylan}
(use used-module
     import:  all | (name-or-rename ...)
     exclude: (name ...)
     prefix:  "string"
     rename:  ((old-name => new-name) ...)
     export:  all | (name ...))
\end{dylan}

\noindent
All the options are optional. With none, every variable the used module
exports is imported under its own name, and nothing is exported again.
The options work as follows.

\begin{description}
\item[\code{import:}]\index{import: (module option)} Which variables to
  import. The default, \code{all}, imports everything the used module
  exports. A list imports only the names listed. An entry of the form
  \code{(old => new)} imports \code{old} under the name \code{new}; it
  is shorthand for listing \code{old} and adding the same pair to
  \code{rename:}.
\item[\code{exclude:}]\index{exclude: (module option)} Names not to
  import. It is allowed only when \code{import:} is \code{all}.
\item[\code{prefix:}]\index{prefix: (module option)} A string put in
  front of each imported name. With \code{prefix: "gfx\$"} the used
  module's \code{draw-line} is \code{gfx\$draw-line} here. Names given in
  \code{rename:} do not get the prefix.
\item[\code{rename:}]\index{rename: (module option)} Pairs
  \code{(old => new)}. Each \code{old} is imported, whatever
  \code{import:} says, under the name \code{new}.
\item[\code{export:}]\index{export: (module option)} Variables imported
  by this clause that the module being defined exports again, named by
  their names in this module; \code{all} re-exports everything the
  clause imports. The default is to re-export nothing.
\end{description}

\noindent
Renaming and excluding work on symbols, not on single variables. A name
\code{foo} and the setter name \code{(setter foo)} travel together: if
you import a slot getter \code{x} under the name \code{px}, its setter
comes with it, and \code{(set! (px p) 3)} works in the importing module.

A module may use the same module in more than one clause; it gets the
union of what the clauses import. Each clause's options apply to that
clause alone.

\subsection{\code{export} and \code{create} clauses}\label{sec:modules-export}

\index{export clause}\index{create clause}
\begin{dylan}
(export name ...)
(create name ...)
\end{dylan}

\noindent
An \code{export} clause names variables that the module owns, which a
definition in the module will create, and makes them visible to modules
that use it. A \code{create} clause names variables that the module owns
and exports but does \emph{not} define: they are to be defined by a module
that uses this one. That is how one module can publish an interface whose
implementation lives in another. It is an error to name in either clause
a variable the module imports. A name may appear more than once.

\subsection{Rules}\label{sec:modules-rules}

\begin{itemize}
\item The uses must not form a cycle: if \code{a} uses \code{b}, \code{b}
  may not use \code{a}, directly or through other modules.
\item A module cannot use itself, and \code{dylan} cannot be redefined.
\item A module may not define a variable it imports, with one exception:
  it may define a variable that the owning module named in a
  \code{create} clause. Adding a \emph{method} to an imported generic
  function with \code{define-method} is not defining the variable; it is
  the normal way one module extends another's generic functions, and it
  is always allowed.
\end{itemize}

\section{An example of each option}\label{sec:modules-example-options}

Design note~32 gives a set of four modules that between them use every
option. Written in prefix form and cut down a little, they are:

\begin{dylan}
(define-module graphics
  (use dylan)
  (create draw-line erase-line skew-line))

(define-module lines
  (use dylan)
  (use graphics import: (draw-line erase-line skew-line)))

(define-module rects
  (use dylan)
  (use graphics prefix: "g$" exclude: (skew-line)))

(define-module gx
  (use dylan export: all)
  (use graphics rename: ((skew-line => warp-line)) export: all))
\end{dylan}

\begin{itemize}
\item \code{graphics} owns and exports three variables and defines none
  of them.
\item \code{lines} imports all three and supplies their definitions,
  which \code{create} allows.
\item \code{rects} sees \code{g\$draw-line} and \code{g\$erase-line},
  and not \code{skew-line} under any name.
\item \code{gx} sees \code{draw-line}, \code{erase-line} and
  \code{warp-line}, and re-exports them together with everything in
  \code{dylan}. A module that says only \code{(use gx)} therefore gets the
  whole language plus the three graphics operations.
\end{itemize}

\noindent
Given these declarations, the file for \code{lines} might read

\begin{dylan}
module: lines

(define-method draw-line (from to)
  (display "draw ") (display from) (display "-") (display to) (newline))
(define-method erase-line (from to)
  (display "erase ") (display from) (display "-") (display to) (newline))
(define-method skew-line (line)
  (display "skew ") (display line) (newline))
\end{dylan}

\noindent
a file in \code{rects} calls \code{(g\$draw-line 1 2)}, and a file in a
module that uses \code{gx} calls \code{(warp-line 'l1)}, which runs the
\code{skew-line} method above and prints \code{skew l1}. Both
implementations run this program as shown.

\section{Which module code belongs to}\label{sec:modules-header}

\index{module: (file header)}\index{file header}
Every expression is compiled in some module, and that module decides
what each free name in it means. Design note~33 says how a source file
names its module: by a \term{file header}, a set of lines of the form
\code{keyword: value} at the very top of the file, ended by a blank line.

\begin{dylan}
module: shapes
author: A. Programmer
        Example Software, Inc.
version: 1.2

(define-class <shape> (<object>))
...
\end{dylan}

\noindent
The rules for the header are these.

\begin{itemize}
\item Each keyword starts at the beginning of a line. Keywords are not
  case-sensitive, so \code{Module:} is the same as \code{module:}.
\item A line that begins with white space continues the value of the
  line before it.
\item The first blank line ends the header. It is required: without it
  the header runs on into the code, and both implementations silently
  skip every line up to the first blank one.
\item \code{module:} names the module. The other keywords
  (\code{author:}, \code{copyright:}, \code{version:} and any others) are
  read and ignored by both implementations.
\end{itemize}

\noindent
A file without a header, one whose first line is not of the form
\code{keyword: value}, is compiled in \code{dylan-user}. Design note~33
makes \code{module:} required; both implementations instead treat its
absence as \code{module: dylan-user}, so that a one-file program needs no
header at all.

A \code{define-module} form is itself compiled in some module, normally
\code{dylan-user}, but the names in its clauses are not looked up there:
they name variables of the module being declared or of the modules it
uses. It is usual to put all the \code{define-module} forms of a program
in one file with no header.

\section{Modules in the compiler to C}\label{sec:modules-dylanc}

\index{dylanc!modules}
The compiler to C compiles a whole program at once, so it sees every
module before it compiles any code. Give \cmd{dylanc} all the files of the
program:

\begin{shell}
$ bin/dylanc -o app modules.dyl shapes.dyl app.dyl
\end{shell}

\noindent
It first finds every top-level \code{define-module} form in every file,
then defines each module after the modules it uses. The order of the
\code{define-module} forms, and of the files that contain them, therefore
does not matter. Then it compiles each file's forms in the module its
header names. The top-level forms of all the files run in the order the
files were given on the command line, and within each file from top to
bottom.

Mistakes in modules are reported when compiling, with the file and line:

\begin{shell}
e4.dyl:2:1: error: define-module: modules m1 use each other
e5.dyl:1:1: error: unknown module nosuch
e7.dyl:1:1: error: define-module: cannot export car, which is imported
e8b.dyl:5:5: error: define of f: the variable is imported from module a
\end{shell}

\noindent
A name that is used in a module but defined nowhere it can see draws a
warning, not an error, and an error only if the program actually reaches
it:

\begin{shell}
e1b.dyl:1:8: warning: g is not defined (module b)
$ ./e1
Dylan error: the variable g@b has no value
\end{shell}

\noindent
As the message shows, the compiler writes a variable as
\meta{name}\code{@}\meta{module} when the module matters. Here \code{g}
was defined in module \code{a} but not exported, so the \code{g} in module
\code{b} is a different, undefined variable.

\index{NAME-MODULE.c}
The generated C is organized by module. With \cmd{-k} (keep the C), a
program \code{app} with code in the modules \code{shapes} and \code{app}
leaves \file{app.c} (data, dispatch and \code{main}), \file{app.h},
\file{app-shapes.c} and \file{app-app.c}, one file of functions for each
module that has any, and \file{app-dylan.c}, the library's own functions.
Chapter~\ref{ch:dylanc} describes these files and the compiler's other
options.

\section{Modules in the interactive implementation}\label{sec:modules-lisp}

\index{interactive implementation!modules}
The Common Lisp implementation compiles one form at a time, so a module
must be declared before any code that uses it is compiled, and a module
must be declared after the modules it uses:

\begin{repl}
? (define-module m1 (use dylan) (use m2))

debugger invoked on a DYLAN::DYLAN-CONDITION ...
  Error: unknown module m2
\end{repl}

\noindent
Declaring the modules in order is enough. Redefining a module with
\code{define-module} replaces its clauses, and every module that uses it
is brought up to date at once, as \dylanbook{chapter 5} asks: variables
newly exported become visible immediately in the modules that import
them.

\subsection{At the read-eval-print loop}

\index{in-module@\code{dylan:in-module}}
The loop starts in \code{dylan-user}. Typing
\code{(dylan:in-module \meta{name})} at the \code{?} prompt switches the
module that later input is compiled in; from Lisp, the function
\code{(dylan:in-module "\meta{name}")}, given the name as a string, does
the same. The loop answers with
\code{;Module:} and the new module's name.

\begin{repl}
? (define-module counter (use dylan) (export next-count))
;Value: counter

? (dylan:in-module counter)
;Module: counter

? (define count 0)
;Value: count

? (define-method next-count () (set! count (+ count 1)) count)
;Value: next-count

? (dylan:in-module dylan-user)
;Module: dylan-user

? (define-module client (use dylan) (use counter))
;Value: client

? (dylan:in-module client)
;Module: client

? (next-count)
;Value: 1

? (next-count)
;Value: 2
\end{repl}

\noindent
Module \code{client} can call \code{next-count} but cannot see
\code{count}, which \code{counter} keeps to itself.
(Chapter~\ref{ch:lisp} describes the loop in full.)

\subsection{Files}

\code{dylan:compile-file} reads a file's header and compiles the file in
the module it names. Names are resolved to variables as the file is
compiled, so the compiled Lisp needs no module information to run, and a
compiled \code{define-module} form defines its module again when the
file is loaded. To
compile a program of several files, compile them in one Lisp session in
the order a reader would need them, the file of \code{define-module}
forms first, and then load the results in the same order:

\begin{shell}
$ sbcl --non-interactive --eval '(load "src/load-compiler.lisp")' \
    --eval '(dylan:compile-file "modules.dyl" "modules.lisp")' \
    --eval '(dylan:compile-file "shapes.dyl" "shapes.lisp")' \
    --eval '(dylan:compile-file "app.dyl" "app.lisp")'
$ sbcl --non-interactive --eval '(load "src/load-runtime.lisp")' \
    --eval '(load "modules.lisp")' --eval '(load "shapes.lisp")' \
    --eval '(load "app.lisp")'
\end{shell}

\noindent
Compiling \file{shapes.dyl} in a session that has not seen
\file{modules.dyl} fails with ``unknown module shapes''. When the Lisp
files are loaded, SBCL may print \code{WARNING: redefining ...} lines for
the small accessor functions that each file declares for the module
variables it shares with the others; they are harmless.

\section{Read-only variables}\label{sec:modules-readonly}

\index{read-only variable}
\dylanbook{Chapter 5} makes some module variables \term{read-only}: the
variable always holds the same object, though that object may itself be
changed. \code{define-class} and \code{define-generic-function} create
read-only variables, so a class or generic function cannot be replaced by
assignment. Variables made by \code{define} (or by \code{define-method}
alone) are ordinary and can be assigned with \code{set!}.

The two implementations check this at different times. The compiler to C
refuses the program:

\begin{shell}
e3.dyl:2:1: error: <p> holds a class or generic function and cannot be assigned
\end{shell}

\noindent
The interactive implementation signals an error when the assignment is
run:

\begin{repl}
? (define-class <p> (<object>))
;Value: <p>

? (set! <p> 3)

debugger invoked on a DYLAN::DYLAN-CONDITION ...
  Error: variable is read-only <p>
\end{repl}

\noindent
Methods can still be added to a read-only generic function: the variable
does not change, only the function in it.

\section{A complete multi-file program}\label{sec:modules-example}

This program has three files. The first declares the modules; it has no
header, so it is compiled in \code{dylan-user}.

\begin{dylan}
;;; modules.dyl
(define-module shapes
  (use dylan)
  (export <shape> <circle> <square> area describe-shape))

(define-module app
  (use dylan)
  (use shapes import: (<circle> <square> (describe-shape => show))))
\end{dylan}

\noindent
Module \code{shapes} exports five variables. Module \code{app} imports
only two classes and one function, and calls the function \code{show}.
It does not import \code{area} or \code{<shape>}.

\begin{dylan}
;;; shapes.dyl
module: shapes
author: A. Programmer

(define-class <shape> (<object>))
(define-class <circle> (<shape>)
  (radius init-keyword: radius:))
(define-class <square> (<shape>)
  (side init-keyword: side:))

(define-generic-function area (shape))
(define-method area ((c <circle>)) (* 3 (radius c) (radius c)))
(define-method area ((s <square>)) (* (side s) (side s)))

(define-method describe-shape ((s <shape>))
  (display "a shape of area ")
  (display (area s))
  (newline))
\end{dylan}

\noindent
The slot getters \code{radius} and \code{side} are owned by \code{shapes}
but not exported, so no other module can read the slots directly.

\begin{dylan}
;;; app.dyl
module: app

(show (make <circle> radius: 2))
(show (make <square> side: 3))
\end{dylan}

\noindent
With the compiler to C:

\begin{shell}
$ bin/dylanc -o app modules.dyl shapes.dyl app.dyl
$ ./app
a shape of area 12
a shape of area 9
\end{shell}

\noindent
With the interactive implementation, compile the three files and load
them as shown in Section~\ref{sec:modules-lisp}; the output is the same.

If \file{app.dyl} also called \code{(area (make <square> side: 1))}, the
compiler to C would warn that \code{area} is not defined in module
\code{app}, and the program would stop with an error on reaching the
call, because \code{app} did not import it.

A program meant to run in both implementations must keep to the names
both provide; Appendix~\ref{app:library} marks the ones that only one of
them has.

\chapter{Printing and Streams}\label{ch:printing}

\index{printing}
The 1992 manual says remarkably little about output. Its examples print
values, its chapter on conditions passes streams to functions and calls
\code{format}, but it defines no printing functions and no streams at
all. Both implementations fill the gap, and they fill it differently.
They share a small common core, \code{print}, \code{display},
\code{write-line}, \code{newline} and \code{pp}, which Thomas provided in
1992 and which are enough for most programs. The compiler to C adds
streams, \code{format}, \code{write}, and a generic function,
\code{print-object}, through which a program decides how its own objects
print.

This chapter describes the common core first, then how each kind of
object prints, then the compiler's streams and \code{format}, and
finally what the interactive implementation offers instead.

\section{Two ways to print an object}\label{sec:printing-escape}

\index{escape (printing)}
Every printing function prints an object in one of two styles.

\begin{itemize}
\item \term{Escaped} output, the style of \code{print} and \code{write},
  is meant to be read by a programmer, and where possible by the Dylan
  reader. Strings appear in double quotes with \code{"} and \code{\textbackslash}
  escaped by a backslash, and characters appear as \code{\#\textbackslash a}
  or \code{\#\textbackslash space}.
\item \term{Plain} output, the style of \code{display}, is meant for the
  person running the program. Strings and characters appear as
  themselves, with no quotes or \code{\#\textbackslash}.
\end{itemize}

\noindent
The two styles differ only for strings and characters, and for
collections and objects that contain them. A number, a symbol or a list
of numbers looks the same either way.

\section{The common printing functions}\label{sec:printing-functions}

These five functions exist in both implementations. Each answers its
argument (\code{newline} answers \code{\#f}), so a \code{print} can be
wrapped around an expression without changing its value.

\begin{description}
\item[\code{(print \meta{object})}]\index{print} prints the object
  escaped, on a line of its own.
\item[\code{(write-line \meta{object})}]\index{write-line} does the
  same as \code{print}.
\item[\code{(display \meta{object})}]\index{display} prints the object
  plain, and nothing else: no newline, no space.
\item[\code{(newline)}]\index{newline} ends the line.
\item[\code{(pp \meta{object})}]\index{pp} prints the object escaped on
  a line of its own; in the interactive implementation it is
  \term{pretty-printed}, broken over lines and indented as Dylan code is.
\end{description}

\begin{dylan}
(print "a \"quoted\" string")     ; prints "a \"quoted\" string"
(display "a \"quoted\" string")   ; prints a "quoted" string
(print #\a)                       ; prints #\a
(display #\a)                     ; prints a
(display 42) (newline)            ; prints 42 and ends the line
\end{dylan}

\begin{warnbox}
The two implementations put the newline of \code{print} and
\code{write-line} in different places. The compiler to C prints the
object and \emph{then} a newline, as most languages do. The interactive
implementation, following Thomas, prints a newline \emph{first} and then
the object, leaving the line unfinished. A sequence of \code{print}s
therefore gives the same lines in both, but output that mixes
\code{print} with \code{display} does not:
\begin{dylan}
(print 1) (display "x") (newline)
;; compiler to C:          1
;;                         x
;; interactive:            (an empty line)
;;                         1x
\end{dylan}
(At the read-eval-print loop the \code{;Value:} lines come in between;
this is the output of a compiled file.)
A program meant to run in both should build its lines with
\code{display} and \code{newline}, which behave identically.
\end{warnbox}

In the interactive implementation each of the five takes exactly one
argument (\code{newline} none). In the compiler to C each also takes a
\code{stream:} keyword argument (Section~\ref{sec:printing-streams}), and
\code{print} takes a \code{verbose:} keyword, which it accepts and
ignores. The compiler's \code{pp} does no pretty-printing; it is
\code{print} under another name.

\section{How objects print}\label{sec:printing-classes}

The table below shows how each kind of object prints in each
implementation; ``the same'' means that the interactive implementation
prints exactly what the compiler to C does. Where they differ, the
difference is only in the printed form; the objects are the same.

{\frenchspacing
\begin{longtable}{@{}>{\raggedright\arraybackslash}p{0.27\textwidth}>{\raggedright\arraybackslash}p{0.33\textwidth}>{\raggedright\arraybackslash}p{0.33\textwidth}@{}}
\toprule
Object & Compiler to C & Interactive \\
\midrule
\endhead
\code{42}, \code{-7} & \code{42}, \code{-7} & the same \\
\code{3/4} & \code{3/4} & the same \\
\code{1.5}, \code{9.0}, \code{0.1} & \code{1.5}, \code{9.0}, \code{0.1} & the same \\
\code{1.0e20}, \code{(sqrt 2)} & \code{1.0e20}, \code{1.4142135623730951} & the same \\
\code{\#t}, \code{\#f}, \code{'()} & \code{\#t}, \code{\#f}, \code{()} & the same \\
\code{'foo}, \code{'Hello} & \code{foo}, \code{Hello} & the same \\
\code{foo:} & \code{foo:} & the same \\
\code{'(1 (2 3) . 4)} & \code{(1 (2 3) . 4)} & the same \\
\code{'{}'x} & \code{(quote x)} & \code{'x} \\
\code{(vector 1 "two" \#\textbackslash 3)} & \code{\#(1 "two" \#\textbackslash 3)} & the same \\
a stretchy vector of 1 and 2 & \code{\{<stretchy-vector> 1 2\}} & \code{\#[instance-of <stretchy-vector>]} \\
a deque of 1 and 2 & \code{\{<deque> 1 2\}} & \code{\#[instance-of <deque>]} \\
\code{(range from: 0 up-to: 5)} & \code{\{range 0 by 1 size 5\}} & \code{\#[instance-of <range>]} \\
\code{(range from: 0)} & \code{\{range 0 by 1\}} & \code{\#[instance-of <range>]} \\
a table mapping \code{a} to 1 & \code{\{table a => 1\}} & \code{\#[instance-of <table>]} \\
an array & \code{\{instance of <\%general-array>\}} & \code{\#[instance-of <array>]} \\
the class \code{<point>} & \code{\{the class <point>\}} & \code{\#[dylan-class <point>]} \\
an instance of \code{<point>} & \code{\{instance of <point>\}} & \code{\#[instance-of <point>]} \\
a generic function & \code{\{the generic function print\}} & \code{\#[generic function area]} \\
a method & \code{\{method\}} & \code{\#[method]} \\
\code{(singleton 3)} & \code{\{singleton 3\}} & \code{\#[dylan-singleton 3]} \\
a \code{<simple-error>} & its message, \code{bad thing} & the same \\
\bottomrule
\end{longtable}}

\noindent
Some details:

\begin{itemize}
\item \textbf{Floats} print in the shortest form that reads back as the
  same number, with at least one digit after the point (\code{9.0},
  never \code{9.}). Very large and very small values take an exponent,
  but the implementations switch at different points. The compiler to C
  writes a float without an exponent when its magnitude is at least
  $10^{-5}$ and below $10^{16}$; the interactive implementation, which
  uses Common Lisp's printer, only when it is at least $10^{-3}$ and below
  $10^{7}$. So \code{123456789.0} prints as \code{123456789.0} in the
  first and \code{1.23456789e8} in the second, and \code{0.0001} as
  \code{0.0001} and \code{1.0e-4}.
\item \textbf{Ratios} print as \meta{numerator}\code{/}\meta{denominator}
  in lowest terms; \code{(/ 6 4)} prints as \code{3/2}.
\item \textbf{Symbols} print as they were first spelled. Dylan symbols
  are not case-sensitive, but the reader remembers the spelling a symbol
  was first entered in, so \code{'Hello} prints as \code{Hello}
  \dylanbook{chapter 15}. \textbf{Keywords} print with their colon.
\item \textbf{Characters} with names print escaped by name:
  \code{\#\textbackslash space}, \code{\#\textbackslash newline},
  \code{\#\textbackslash tab}, \code{\#\textbackslash return},
  \code{\#\textbackslash page}, \code{\#\textbackslash null},
  \code{\#\textbackslash rubout}, \code{\#\textbackslash backspace}.
\item \textbf{Lists and vectors} print their elements in the same style
  as the whole: \code{(display (list "a" \#\textbackslash b))} prints
  \code{(a b)}.
\item \textbf{Conditions} print as their message, the format string
  filled in with the format arguments, as \dylanbook{chapter 13}
  suggests. A \code{<type-error>} prints as
  ``\meta{value} is not an instance of \meta{type}'', where the type is
  printed as a class prints: \code{3 is not an instance of \{the class
  <string>\}} in the compiler to C, \code{3 is not an instance of
  <string>} in the interactive implementation.
\item \textbf{Functions} that the language predefines as plain methods,
  such as \code{car}, print as methods. Which predefined functions are
  generic differs between the implementations: \code{print} is a generic
  function in the compiler to C and a method in the interactive one.
\end{itemize}

\section{Customizing printing with \code{print-object}}\label{sec:printing-print-object}

\index{print-object}
In the compiler to C every printing function, \code{format} included,
works by calling one generic function:

\begin{dylan}
(print-object object stream escape?)
\end{dylan}

\noindent
\code{print-object} writes \meta{object} to \meta{stream}; \meta{escape?}
is \code{\#t} for escaped output (\code{print}, \code{write}, \code{\textasciitilde s})
and \code{\#f} for plain output (\code{display}, \code{\textasciitilde a}). The library
defines a method for each of its own classes, and a default method on
\code{<object>} that prints \code{\{instance of \meta{class}\}}. To make
the instances of your own class print differently, add a method:

\begin{dylan}
(define-class <point> (<object>)
  (x init-keyword: x: init-value: 0)
  (y init-keyword: y: init-value: 0))

(define-method print-object ((p <point>) stream escape?)
  (format stream "{point ~s, ~s}" (x p) (y p)))

(define p (make <point> x: 3 y: 4))
(print p)                     ; prints {point 3, 4}
(print (list p p))            ; prints ({point 3, 4} {point 3, 4})
(format #t "p is ~a~%" p)     ; prints p is {point 3, 4}
\end{dylan}

\noindent
The method is used wherever a point is printed, including inside lists,
vectors and other collections, and by \code{format}. Write the output
to the \meta{stream} argument, never to \code{*standard-output*}
directly, or printing into a string will not work.

A method can honor the \meta{escape?} argument, so that \code{print}
shows something a programmer can recognize and \code{display} shows what
a user should see:

\begin{dylan}
(define-class <name> (<object>) (text init-keyword: text:))

(define-method print-object ((n <name>) stream escape?)
  (if escape?
      (format stream "#<name ~s>" (text n))
      (display (text n) stream: stream)))

(define n (make <name> text: "Ada"))
(print n)                       ; prints #<name "Ada">
(display n)                     ; prints Ada
(format #t "~a / ~s~%" n n)     ; prints Ada / #<name "Ada">
\end{dylan}

\noindent
Methods are inherited and can call \code{next-method} like any other:

\begin{dylan}
(define-class <temp> (<object>) (deg init-keyword: deg:))
(define-method print-object ((t <temp>) stream escape?)
  (display (deg t) stream: stream)
  (display " degrees" stream: stream))

(define-class <hot> (<temp>))
(define-method print-object ((t <hot>) stream escape?)
  (display "hot! " stream: stream)
  (next-method))

(print (make <hot> deg: 99))    ; prints hot! 99 degrees
\end{dylan}

\noindent
\code{print-object} is an extension: the 1992 manual has no such
function, and the interactive implementation does not provide one
(Section~\ref{sec:printing-lisp}).

\section{Streams in the compiler to C}\label{sec:printing-streams}

\index{stream}
The compiler to C provides a small set of streams, enough for the uses
the 1992 manual makes of them. They are an extension, described with the
compiler's other extensions in Chapter~\ref{ch:extensions}.

\begin{description}
\item[\code{<stream>}]\index{<stream>@\code{<stream>}} The class of all
  streams.
\item[\code{*standard-output*}]\index{*standard-output*@\code{*standard-output*}}
  A stream that writes to the program's standard output.
\item[\code{*standard-error*}]\index{*standard-error*@\code{*standard-error*}}
  A stream that writes to the program's standard error.
\item[\code{<string-stream>}]\index{<string-stream>@\code{<string-stream>}}
  A stream that collects what is written to it, in memory. You make
  one by calling \code{make}: \code{(make <string-stream>)}.
\item[\code{(stream-contents \meta{string-stream})}]\index{stream-contents}
  Answers a new string of everything written to the stream so far. The
  stream keeps its contents; calling \code{stream-contents} again answers
  them again, with anything written since.
\end{description}

\noindent
Every printing function takes a \code{stream:} keyword argument naming
the stream to write to. It defaults to \code{*standard-output*}, and
\code{\#t} (or \code{\#f}) also means \code{*standard-output*}.

\begin{dylan}
(define s (make <string-stream>))
(print 42 stream: s)
(display "hello" stream: s)
(write #\x stream: s)
(display (stream-contents s))
;; prints  42
;;         hello#\x

(display "warning: disk nearly full" stream: *standard-error*)
(newline stream: *standard-error*)
\end{dylan}

\noindent
\code{write}\index{write} prints its argument escaped, like
\code{print} but with no newline:

\begin{description}
\item[\code{(write \meta{object} \#key stream)}] prints the object
  escaped and answers it.
\end{description}

\noindent
\code{<stream>} cannot usefully be subclassed by a program. The
operation that puts characters on a stream is internal to the library,
so printing to an instance of your own subclass of \code{<stream>} fails
with ``no applicable method''.

\section{\code{format}}\label{sec:printing-format}

\index{format}
\begin{dylan}
(format stream control-string argument ...)
\end{dylan}

\noindent
\code{format} writes \meta{control-string} to \meta{stream}, replacing
each \term{directive}, a tilde and a letter, by output for the next
argument. Its directives are a small subset of Common Lisp's, the ones
the 1992 manual's examples use:

\begin{center}
\begin{tabular}{@{}ll@{}}
\toprule
Directive & Output \\
\midrule
\code{\textasciitilde a} & the next argument, plain, as \code{display} prints it \\
\code{\textasciitilde s} & the next argument, escaped, as \code{print} and \code{write} print it \\
\code{\textasciitilde d} & the next argument, plain (the same as \code{\textasciitilde a}) \\
\code{\textasciitilde\%} & a newline \\
\code{\textasciitilde\&} & a newline (the same as \code{\textasciitilde\%}) \\
\code{\textasciitilde\textasciitilde} & a tilde \\
\bottomrule
\end{tabular}
\end{center}

\noindent
The letters may be upper or lower case. Any other tilde sequence is
copied to the output unchanged. There are no field widths, no
padding and no number bases: \code{\textasciitilde d} prints an integer
exactly as \code{display} would, and \code{\textasciitilde\&} always
starts a new line, even at the beginning of one.

The \meta{stream} argument is a stream, or one of two booleans:

\begin{itemize}
\item \code{\#t} writes to \code{*standard-output*}, and \code{format}
  answers \code{\#f}.
\item \code{\#f} writes to a new string, and \code{format} answers the
  string.
\end{itemize}

\begin{dylan}
(format #t "~a + ~a = ~a~%" 1 2 (+ 1 2))    ; prints 1 + 2 = 3
(format #f "~a + ~a = ~a" 1 2 (+ 1 2))      ; => "1 + 2 = 3"
(format #f "~A ~S ~D ~z" "x" "x" 5)         ; => "x \"x\" 5 ~z"
(format *standard-error* "to stderr~%")
\end{dylan}

\noindent
Too few arguments for the directives is an error
(``\code{format: too few arguments for "\textasciitilde a and \textasciitilde a"}''); arguments left
over after the last directive are ignored.

\section{Printing in the interactive implementation}\label{sec:printing-lisp}

\index{interactive implementation!printing}
The interactive implementation has exactly the five functions of
Section~\ref{sec:printing-functions}, each with one argument and no
\code{stream:} keyword. It has no streams, no \code{format}, no
\code{write} and no \code{print-object}, so the way an object prints
cannot be changed from Dylan: an instance of a class of your own always
prints as \code{\#[instance-of \meta{class}]}. To show its contents,
print them yourself:

\begin{dylan}
(define-method show ((p <point>))
  (display "{point ") (display (x p))
  (display ", ") (display (y p)) (display "}"))
\end{dylan}

\noindent
The read-eval-print loop prints each result escaped after
\code{;Value:}, so at the \code{?} prompt you seldom need \code{print}
at all (Chapter~\ref{ch:lisp}).

\subsection{Pretty-printing}

\index{pp}
\code{pp} uses the Common Lisp pretty printer, set up to show Dylan's
notation, to lay out a long list as code:

\begin{repl}
? (pp '(define-method foo ((x <integer>))
         (bind ((y (* x x)) (z (+ x x x x x x x x x x x x x x x x x x x x)))
           (list y z y z y z))))

(define-method foo ((x <integer>))
 (bind ((y (* x x)) (z (+ x x x x x x x x x x x x x x x x x x x x)))
  (list y z y z y z)))
\end{repl}

\subsection{Formatted output through Common Lisp}

\index{lisp-procedure}
When you need formatted output in the interactive implementation, Common
Lisp's own \code{format} can be reached with \code{lisp-procedure}
(Section~\ref{sec:lisp-interop}). It has every Common Lisp directive,
field widths included. Pass \code{\#t} for standard output, or the empty
list, which is Lisp's \code{NIL}, to get a string back:

\begin{repl}
? (define cl-format (lisp-procedure 'format))
;Value: cl-format

? (cl-format #t "~5,2f|~8d|~r~%" 3.14159 42 12)
 3.14|      42|twelve
;Value: ()

? (cl-format '() "~a-~a" 1 2)
;Value: "1-2"
\end{repl}

\begin{warnbox}
Common Lisp's \code{format} prints Dylan values as Common Lisp sees
them, not as Dylan's printer does: \code{\#t} comes out as \code{T}, a
keyword \code{foo:} as \code{FOO}, and \code{\#f}, which is not a Lisp
value at all, as \code{\#f} only because the implementation teaches Lisp
to print it so. Use it for numbers and strings. Do not pass \code{\#f}
as the destination: to Common Lisp it is an ordinary true object, not
\code{NIL}.
\end{warnbox}

\section{Writing portable output}\label{sec:printing-portable}

A program that must run in both implementations should

\begin{itemize}
\item build its output from \code{display} and \code{newline}, which
  behave the same in both;
\item use \code{print} only for whole lines, remembering where each
  implementation puts the newline;
\item avoid \code{format}, \code{write}, streams and \code{print-object},
  which only the compiler to C has; and
\item not depend on the printed form of collections other than lists,
  vectors and strings, or of classes, functions and instances, which
  differ (Section~\ref{sec:printing-classes}).
\end{itemize}


\part{The Implementations}
\chapter{The Interactive Implementation}\label{ch:lisp}

\index{interactive implementation}
The interactive implementation runs Dylan inside a Common Lisp system,
SBCL. You type a Dylan expression, it is compiled to Common Lisp at once,
the Lisp compiles that to machine code and runs it, and the value is
printed. Nothing about the program is fixed in advance: classes can gain
slots, generic functions can gain and lose methods, and definitions can
be replaced while the program runs, all as the 1992 manual describes.
That makes it the natural place to learn the language, to try out an
idea, and to follow the DIRM's examples.

This chapter describes how the implementation is built, how to load and
use it, how Dylan values look to Lisp and how to call Lisp from Dylan,
how it is made fast, and how it differs from the compiler to C.

\section{How it works}\label{sec:lisp-architecture}

\index{interactive implementation!architecture}
The implementation began in 1992 as Thomas, a Dylan system written in
Scheme at Digital Equipment Corporation's Cambridge Research Laboratory.
This project translated it into Common Lisp and has since corrected it
against the manual and rebuilt much of it for speed. Every source file in
\file{src/} still names the Thomas file it grew from.

An expression goes through four stages.

\begin{enumerate}
\item \textbf{Reading.} Common Lisp's reader cannot read Dylan: a colon
  is a package marker to it, and it knows nothing of \code{\#t} and
  \code{\#f}. A small reader of Dylan's own (in \file{compat.lisp}) turns
  text into lists, symbols, numbers and strings. A Dylan symbol becomes a
  Lisp symbol in the \code{DYLAN} package; a keyword \code{foo:} becomes
  the Lisp keyword \code{:FOO}.
\item \textbf{Compiling to Lisp.} The compiler (\file{compiler.lisp}
  and the \file{comp-*.lisp} files) turns each Dylan form into a Lisp
  form: \code{begin} into \code{progn}, \code{set!} into \code{setq}, a
  method into a \code{lambda}, a call into a call through the runtime's
  dispatcher. Each module variable becomes a Lisp special variable
  with a name of its own, \code{!MODULE/X/} for a variable \code{x} of
  \code{dylan-user} and \code{!MODULE/X@M/} for one owned by module
  \code{m}.
\item \textbf{Compiling to machine code.} The Lisp form is handed to
  SBCL's own compiler, which turns it into native code.
\item \textbf{Running.} The code runs against the runtime library, the
  \file{runtime-*.lisp} files, which implement classes, generic
  dispatch, collections, numbers and conditions.
\end{enumerate}

\noindent
The files of \file{src/}, in the order they are loaded, do these jobs:

\begin{longtable}{@{}>{\raggedright\arraybackslash}p{0.34\textwidth}>{\raggedright\arraybackslash}p{0.6\textwidth}@{}}
\toprule
File & Role \\
\midrule
\endhead
\file{package.lisp} & the \code{DYLAN} package and its exported entry points \\
\file{globals.lisp} & declares the runtime's global variables, which nearly every file refers to before the file that defines them is loaded \\
\file{compat.lisp} & the Dylan reader and printer, weak tables, escapes, sorting and the interface to Common Lisp's condition system \\
\file{support.lisp} & small utilities shared by the compiler and runtime \\
\file{comp-util.lisp} & the compiler's utilities: names, variables, parameter lists \\
\file{comp-module.lisp} & modules, \code{define-module} and the \code{module:} file header (Chapter~\ref{ch:modules}) \\
\file{compiler.lisp} & the heart of the compiler, the table of special forms, the predefined names, \code{compile-file} and \code{compile-to-file} \\
\file{comp-method.lisp} & \code{method}, \code{define-method}, \code{define-generic-function} \\
\file{comp-class.lisp} & \code{define-class} and \code{define-slot} \\
\file{comp-sf.lisp} & the other special forms: \code{bind}, \code{for}, \code{select} and the rest \\
\file{comp-exc.lisp} & \code{handler-bind}, \code{handler-case}, \code{bind-exit}, \code{unwind-protect} \\
\file{class.lisp} & classes, slots, instances and \code{make} \\
\file{generic.lisp} & generic functions and method dispatch, with its caches \\
\file{class-structure.lisp} & the predefined class heterarchy \\
\file{runtime-top.lisp}, \file{runtime-internal.lisp} & common parameter lists and the procedures compiled code calls \\
\file{runtime-methods.lisp}, \file{runtime-functions.lisp} & the predefined methods and functions \\
\file{runtime-extensions.lisp} & \code{lisp-variable}, \code{lisp-procedure} \\
\file{runtime-bitstrings.lisp} & \code{logand} and the other bit operations \\
\file{runtime-collections*.lisp} & collections: the iteration protocol, the generic operations, and one file for each kind (lists, vectors, strings, deques, ranges, tables, arrays) \\
\file{runtime-exceptions.lisp} & conditions, handlers and restarts \\
\file{runtime-printing.lisp} & streams, \code{print-object}, \code{format} and the other printing functions (Chapter~\ref{ch:printing}) \\
\file{runtime-extended-float.lisp} & \code{<extended-float>} arithmetic (Chapter~\ref{ch:numbers}) \\
\file{runtime-fast-paths.lisp} & turns on the guarded fast paths (Section~\ref{sec:lisp-performance}) \\
\file{rep.lisp} & the read-eval-print loop \\
\file{load-dylan.lisp}, \file{load-compiler.lisp}, \file{load-runtime.lisp} & loaders \\
\file{build-dylan.lisp} & saves the executable \file{bin/dylan} \\
\file{full-test.lisp} & the regression test driver \\
\bottomrule
\end{longtable}

\noindent
\file{src/README} describes the translation from Scheme in more detail.

\section{Loading it}\label{sec:lisp-loading}

\index{dylan@\cmd{dylan}}\index{build-dylan.lisp@\file{build-dylan.lisp}}
The usual way in is the executable \file{bin/dylan}. Running \cmd{make}
in the repository's top directory builds it: \file{src/build-dylan.lisp}
loads the whole system into SBCL and saves it, with SBCL, as one
executable that starts in the read-eval-print loop.

\begin{shell}
$ make
$ bin/dylan

Entering Dylan read-eval-print-loop.
Leave by typing "(exit)"

?
\end{shell}

\noindent
The executable needs nothing else from the repository, so it may be
copied or linked anywhere. It does not see later changes to \file{src/};
run \cmd{make} again to rebuild it. On Windows, without \cmd{make}, run
\cmd{sbcl -{}-script src/build-dylan.lisp}, which writes
\file{bin\textbackslash dylan.exe}.

\index{load-dylan.lisp@\file{load-dylan.lisp}}
The other way is to start a plain SBCL in the repository's top directory
and load the system into it, which is how \cmd{make} itself begins:

\begin{shell}
$ sbcl
* (load "src/load-dylan.lisp")
;Dylan loaded.  (dylan:rep) starts a read-eval-print loop.
\end{shell}

\index{fasl cache}
The first time, this compiles every source file into a \term{fasl}
(Lisp's compiled-file format), which takes a little while. The fasls are
kept in \file{src/fasl/\meta{lisp}-\meta{version}/}, for example
\file{src/fasl/SBCL-2-6-3/}, one directory for each Lisp and version
since their fasls are incompatible. Later loads use them directly. When a
source file changes, it is recompiled together with every file loaded
after it, since those may depend on its macros and inline functions.

Three loaders exist, for the three ways of using the system that Thomas
offered:

\begin{description}
\item[\file{load-dylan.lisp}] the interactive environment: the
  read-eval-print loop.
\item[\file{load-compiler.lisp}] the compiler to Lisp:
  \code{dylan:compile-file} and \code{dylan:compile-to-file}.
\item[\file{load-runtime.lisp}] what compiled programs need to run.
\end{description}

\noindent
In this version the three load the same thing, the whole system;
\file{load-compiler.lisp} and \file{load-runtime.lisp} simply load
\file{load-dylan.lisp}. Use whichever names your purpose. A loader can
be given on the command line as well:

\begin{shell}
$ sbcl --eval '(load "src/load-dylan.lisp")' --eval '(dylan:rep)'
\end{shell}

\section{The read-eval-print loop}\label{sec:lisp-rep}

\index{read-eval-print loop}\index{rep@\code{dylan:rep}}
\file{bin/dylan} starts the loop at once; in a Lisp with the system
loaded, \code{(dylan:rep)} starts it. It prompts with \code{?}, reads one
Dylan form, compiles and runs it, and prints the value after
\code{;Value:}. A form may span several lines; the loop reads until the
form is complete.

\begin{repl}
* (dylan:rep)

Entering Dylan read-eval-print-loop.
Leave by typing "(exit)"

? (define-method double ((n <number>)) (+ n n))

;Value: double

? (double 21)

;Value: 42

? (values 1 2 3)

;Value[1]: 1
;Value[2]: 2
;Value[3]: 3

? (values)


? dylan:done
DYLAN::|DYLAN:DONE|
*
\end{repl}

\noindent
A form that returns several values shows each on its own
\code{;Value[\meta{n}]:} line; one that returns none shows nothing.
Values are printed as \code{print} prints them (Chapter~\ref{ch:printing}),
and output the form itself writes appears before the \code{;Value:} line.
A definition answers the name it defined.

Typing \code{dylan:done}, or ending the input (Ctrl-D), leaves the
loop. In \file{bin/dylan} that ends the program. In a Lisp you loaded
the system into, it returns to Lisp's \code{*} prompt; the Lisp prints
the symbol the loop returned. Everything the loop defined remains: call
\code{(dylan:rep)} again and carry on where you left off. \code{dylan:done}
may also be typed in parentheses, as \code{(dylan:done)}.

\index{exit@\code{(exit)}}
Typing \code{(exit)} quits altogether: it ends \file{bin/dylan}, and in a
Lisp you loaded the system into, it ends the Lisp as well. It is a command
of the loop, like \code{dylan:done}, not a Dylan function; a variable
named \code{exit} is still evaluated when typed bare. Typing \code{(dylan:in-module \meta{name})} at the
\code{?} prompt changes the module later input is compiled in
(Section~\ref{sec:modules-lisp}).

Redefining things is normal at the loop. A second \code{define} of a
variable replaces its value with a warning, a second \code{define-method}
with the same specializers replaces the method, and a second
\code{define-class} redefines the class in place:

\begin{repl}
? (define v 1)
;Value: v

? (define v 2)
WARNING: redefining variable v
;Value: v

? (define-method describe ((x <integer>)) "an integer")
;Value: describe

? (define-method describe ((x <integer>)) "a whole number")
;Value: describe

? (describe 3)
;Value: "a whole number"
\end{repl}

\index{empty-environment"!@\code{dylan:empty-environment"!}}
From Lisp, \code{(dylan:empty-environment!)} makes the loop forget the
module variables it has defined, and their values: a variable defined
earlier is unbound afterwards, and using it is an error. It is rarely
needed.

\subsection{Errors and the Lisp debugger}\label{sec:lisp-errors}

\index{debugger}\index{restart}
An error that the Dylan program does not handle itself reaches SBCL's
debugger. The debugger lists the \term{restarts} available, and the Dylan
restarts in effect appear among them. For an error Dylan signals, the
loop provides ``Return to the read-eval-print loop'', which abandons the
form, prints \code{;Aborted} and gives you a new \code{?} prompt; for
every error, whatever signalled it, it also provides \code{DYLAN-REP},
``Return to the Dylan read-eval-print loop'', which does the same
without the \code{;Aborted}:

\begin{repl}
? (define x 10)
;Value: x

? (error "Oops: ~a" x)

debugger invoked on a DYLAN::DYLAN-CONDITION in thread
#<THREAD tid=690943 "main thread" RUNNING {1200030003}>:
  Error: Oops: 10

Type HELP for debugger help, or (SB-EXT:EXIT) to exit from SBCL.

restarts (invokable by number or by possibly-abbreviated name):
  0: [DYLAN-RESTART] Return to the read-eval-print loop.
  1: [DYLAN-REP    ] Return to the Dylan read-eval-print loop.
  2: [ABORT        ] Exit debugger, returning to top level.

((LAMBDA NIL :IN DYLAN::IMPLEMENTATION-SPECIFIC/INDUCE-ERROR))
0] 0
;Aborted

? x
;Value: 10
\end{repl}

\noindent
While in the debugger you can use all of SBCL's debugger commands, such
as \code{backtrace}; the frames show the Lisp the Dylan was compiled to.

A restart the program itself establishes is offered too. \code{cerror}
makes one whose description is its first argument; choosing it makes
\code{cerror} return \code{\#f} and the program carry on:

\begin{repl}
? (cerror "Use 0 instead" "bad value ~a" 7)

debugger invoked on a DYLAN::DYLAN-CONDITION ...
  Error: bad value 7
...
restarts (invokable by number or by possibly-abbreviated name):
  0: [DYLAN-RESTART] Use 0 instead
  1:                 Return to the read-eval-print loop.
  2: [DYLAN-REP    ] Return to the Dylan read-eval-print loop.
  3: [ABORT        ] Exit debugger, returning to top level.
\end{repl}

\noindent
Calling \code{(abort)} from Dylan, or signaling any \code{<abort>},
returns straight to the prompt with \code{;Aborted}, as
\dylanbook{chapter 13} describes. \code{(break "look at \textasciitilde a" 42)} enters the
debugger with the message \code{Dylan breakpoint: look at 42}; its
\code{CONTINUE} restart makes \code{break} return \code{\#f}.

\begin{notebox}
Some errors are detected by Common Lisp itself rather than by the
Dylan runtime, for example \code{(car 5)}, which reaches Lisp's
\code{car} with a non-list. Such an error arrives in the debugger as a
Lisp \code{TYPE-ERROR}, and the restarts offered are \code{DYLAN-REP}
and Lisp's own:
\begin{repl}
restarts (invokable by number or by possibly-abbreviated name):
  0: [DYLAN-REP] Return to the Dylan read-eval-print loop.
  1: [ABORT    ] Exit debugger, returning to top level.
\end{repl}
\noindent Choose \code{DYLAN-REP} to carry on at the \code{?} prompt.
(In \file{bin/dylan}, \code{ABORT} also returns to the loop. In a Lisp you
loaded the system into, it lands at the Lisp \code{*} prompt; nothing is
lost: type \code{(dylan:rep)} to go back in.) A Dylan
handler still catches these errors as instances of \code{<error>}:
\code{(handler-case (car 5) ((<error>) 'caught))} answers \code{caught}.
\end{notebox}

\subsection{Running a file at the loop}\label{sec:lisp-test-file}

\index{test@\code{dylan:test}}
The loop reads from standard input, so \cmd{bin/dylan < fact.dyl} runs
a file, printing each value as if it had been typed; at the end of the
file the program exits. The loop has no command to load a file of Dylan.
The regression test driver can do it: \code{dylan:test} reads a
file form by form, echoes each form, runs it, and prints its value, or
the message of an error, and goes on to the next form.

\begin{shell}
$ cd src
$ sbcl --eval '(load "load-dylan.lisp")' --eval '(load "full-test.lisp")' \
       --eval '(dylan:test "/path/to/fact.dyl")'
\end{shell}

\begin{repl}
(define-method fact ((n <integer>))
 (if (= n 0)
     1
     (* n (fact (- n 1)))))
;Value: fact

(print (fact 30))
265252859812191058636308480000000
;Value: 265252859812191058636308480000000

(car 3)
;Error: Lisp condition:  The value
                   3
                 is not of type
                   LIST
\end{repl}

\noindent
The definitions remain afterwards, so \code{(dylan:rep)} can follow in
the same session. Note that \file{full-test.lisp} must be loaded after
\file{load-dylan.lisp}.

\section{Compiling Dylan to Lisp files}\label{sec:lisp-compile-file}

\index{compile-file@\code{dylan:compile-file}}\index{compile-to-file@\code{dylan:compile-to-file}}
The compiler can also write its output to a file of Lisp instead of
running it, so that a program can be compiled once and loaded later
without being compiled again.

\begin{description}
\item[\code{(dylan:compile-file \meta{input} \meta{output})}] compiles
  the file of Dylan \meta{input} into the file of Lisp \meta{output}, in
  the module the input's \code{module:} header names, or
  \code{dylan-user}.
\item[\code{(dylan:compile-to-file \meta{output} \meta{form} ...)}]
  compiles Dylan forms, already read, into \meta{output}.
\end{description}

\noindent
\code{dylan:compile-file} shadows Common Lisp's function of the same
name: inside the \code{DYLAN} package, \code{compile-file} compiles
Dylan. The forms given to \code{dylan:compile-to-file} must be Dylan
forms whose symbols are in the \code{DYLAN} package, which is what the
Dylan reader \code{dylan::dylan-read} produces; Lisp forms typed at the
\code{*} prompt would have their symbols in the wrong package.

Here is a complete round trip. The program:

\begin{dylan}
;;; hello.dyl
(define-method greet ((name <string>))
  (display "Hello, ")
  (display name)
  (display "!")
  (newline))
(greet "world")
\end{dylan}

\noindent
is compiled in one Lisp session and run in another:

\begin{shell}
$ sbcl --non-interactive --eval '(load "src/load-compiler.lisp")' \
       --eval '(dylan:compile-file "hello.dyl" "hello.lisp")'
$ sbcl --non-interactive --eval '(load "src/load-runtime.lisp")' \
       --eval '(load "hello.lisp")'
Hello, world!
\end{shell}

\noindent
The output file is ordinary Common Lisp. It begins with the Dylan it
came from, as comments, and then the compiled forms:

\begin{lisp}
(in-package #:dylan)

(defvar !module/greet/ '<<the-unassigned-value>>)
...
(dylan%catch-all-conditions
 (lambda ()
   (progn
    (progn
     (cond
      ((dylan%eq? !module/greet/ '<<the-unassigned-value>>)
       (setq !module/greet/ (dylan%create-generic-function 'greet 1 nil nil)))
      ...)
     (dylan%add-method !module/greet/ ...)
     'greet)
    (dylan%call nil (dylan%free-variable-ref !module/greet/ 'greet) "world"))))
\end{lisp}

\noindent
\code{compile-to-file} is the way to compile forms you build in Lisp:

\begin{lisp}
(dylan:compile-to-file "two.lisp"
  (dylan::dylan-read (make-string-input-stream "(define x 41)"))
  (dylan::dylan-read (make-string-input-stream "(print (+ x 1))")))
\end{lisp}

\noindent
Loading \file{two.lisp} prints \code{42}. A program of several files
is compiled file by file in one session and loaded in the same order
(Section~\ref{sec:modules-lisp}).

\section{Dylan values as Lisp objects}\label{sec:lisp-representation}

\index{interactive implementation!representation}
Most Dylan values are the corresponding Lisp values, which is what makes
calling between the two languages easy.

\begin{longtable}{@{}>{\raggedright\arraybackslash}p{0.34\textwidth}>{\raggedright\arraybackslash}p{0.6\textwidth}@{}}
\toprule
Dylan & Common Lisp \\
\midrule
\endhead
integers & integers, fixnums or bignums: there is no limit on size \\
ratios & ratios \\
floats & \code{double-float}s \\
complex numbers & complex numbers; \code{(sqrt -4)} is \code{\#c(0.0 2.0)} \\
characters & characters \\
strings & Lisp strings \\
\code{<simple-object-vector>} & simple vectors \\
lists & lists; the empty list is \code{NIL} \\
\code{\#t} & \code{T} \\
\code{\#f} & a distinct object of its own, the value of \code{DYLAN::DYLAN\%FALSE} \\
symbols & symbols in the \code{DYLAN} package \\
keywords & Lisp keywords: \code{foo:} is \code{:FOO} \\
functions & Lisp functions, which take two hidden leading arguments \\
classes, instances & structures of the runtime's own \\
tables, stretchy vectors, ranges, deques & Dylan instances \\
\bottomrule
\end{longtable}

\paragraph{\code{\#f} and \code{NIL}.}\index{false@\code{\#f}}
Thomas ran on a Scheme in which false and the empty list were the same
object. The 1992 manual makes them different: \code{\#f} is the only
false value, and the empty list is true \dylanbook{chapters 3 and 8}. This
implementation follows the manual. The empty list is Lisp's \code{NIL},
because Dylan lists are Lisp lists, so \code{\#f} has to be something
else, an object created for the purpose. \code{\#t} is Lisp's \code{T}.
The consequence for Lisp code is that a Lisp ``false'' (\code{NIL}) is a
Dylan empty list and therefore \emph{true} in Dylan, and Dylan's
\code{\#f} is a non-\code{NIL} object and therefore \emph{true} in Lisp.

\paragraph{Symbols.} The Dylan reader interns each name in the
\code{DYLAN} package, folded to upper case as Lisp's reader would, so
Dylan's \code{foo} is the Lisp symbol \code{DYLAN::FOO}. (The names
\code{t} and \code{nil} are kept in lower case, so as not to be Lisp's
\code{T} and \code{NIL}.) Because the \code{DYLAN} package uses
\code{COMMON-LISP}, a Dylan symbol whose name is that of a Common Lisp
function or variable \emph{is} the Common Lisp symbol: Dylan's
\code{'expt} is \code{CL:EXPT}. The exceptions are the few names the
package shadows: \code{add-method}, \code{adjoin}, \code{compile-file},
\code{find-method}, \code{last}, \code{reduce}, \code{set-difference} and
\code{union}. The spelling a symbol was first typed in is remembered for
printing.

\section{Calling Common Lisp from Dylan}\label{sec:lisp-interop}

\index{lisp-variable}\index{lisp-procedure}
Two functions, extensions carried over from Thomas, reach into the
host Lisp.

\begin{description}
\item[\code{(lisp-variable \meta{symbol})}] answers the value of the
  Lisp global variable named by \meta{symbol}.
\item[\code{(lisp-procedure \meta{symbol})}] answers a Dylan function
  that calls the Lisp function named by \meta{symbol} (or, if the symbol
  has no function, the function that is the symbol's value). The Dylan
  function accepts any number of arguments; the Lisp function reports
  the wrong number.
\end{description}

\begin{repl}
? (lisp-variable 'most-positive-fixnum)
;Value: 4611686018427387903

? (lisp-variable 'pi)
;Value: 3.141592653589793

? ((lisp-procedure 'expt) 2 100)
;Value: 1267650600228229401496703205376

? ((lisp-procedure 'string-upcase) "hello")
;Value: "HELLO"

? (define cl-sort (lisp-procedure 'sort))
;Value: cl-sort

? (cl-sort (list 3 1 2) '<)
;Value: (1 2 3)
\end{repl}

\noindent
Functions of your own written in Lisp can be called the same way. Define
them in the \code{DYLAN} package, so that the Dylan symbol names them:

\begin{shell}
$ sbcl --eval '(load "src/load-dylan.lisp")' \
       --eval '(defun dylan::square (x) (* x x))' --eval '(dylan:rep)'
\end{shell}

\begin{repl}
? (define square (lisp-procedure 'square))
;Value: square

? (square 12)
;Value: 144

? (square 1 2)

debugger invoked on a DYLAN::DYLAN-CONDITION ...
  Error: invalid number of arguments: 2
\end{repl}

\noindent
Keep these rules in mind:

\begin{itemize}
\item \textbf{Lisp's false is Dylan's true.} A Lisp predicate answers
  \code{NIL} for false, which Dylan sees as the empty list, a true
  value:
\begin{dylan}
(if ((lisp-procedure 'evenp) 3) 'yes 'no)    ; => yes
\end{dylan}
  Test the result against the empty list instead:
\begin{dylan}
(define lisp-true? (method (x) (not (id? x '()))))
(lisp-true? ((lisp-procedure 'evenp) 3))      ; => #f
\end{dylan}
  Likewise, do not pass \code{\#f} to a Lisp function that expects
  \code{NIL}; pass \code{'()}.
\item \textbf{Dylan functions cannot be passed to Lisp.} A Dylan
  function, including one made by \code{lisp-procedure}, expects two
  hidden arguments that Lisp does not supply, so handing one to a Lisp
  function such as \code{mapcar} fails with ``invalid number of
  arguments''. Pass a symbol naming a Lisp function, as with \code{'<}
  above.
\item \textbf{Lisp values print as Dylan sees them.} A Lisp value that
  has no Dylan meaning, such as a package or a hash table, prints in
  Lisp's notation.
\end{itemize}

\section{How it is made fast}\label{sec:lisp-performance}

\index{interactive implementation!performance}
Thomas's authors did not aim at speed, and a direct translation of it
was slow. The Common Lisp implementation has since been rebuilt around a
few techniques, all of which keep the language's full dynamism.

\begin{itemize}
\item \textbf{Dispatch caches.} Each generic function caches, for each
  combination of argument classes it has seen, the sorted list of
  applicable methods, along with the keyword checks those methods need.
  Adding or removing a method, or changing a class's superclasses,
  invalidates the cache. A generic function of one, two or three
  required arguments calls the cached method directly without building
  an argument list, and calls with several applicable methods
  (\code{next-method} chains) and with keyword arguments go through the
  cache too. Singleton specializers are found through a hash table.
\item \textbf{Direct calls.} A call compiles to a direct call of the
  function, without wrapping each argument in a closure as Thomas did. A
  method whose specializers are all \code{<object>} is called as its
  compiled body, without an argument-checking wrapper.
\item \textbf{Guarded fast paths.} Calls of \code{+}, \code{-},
  \code{*}, \code{<}, \code{>}, \code{<=}, \code{>=}, \code{=},
  \code{modulo}, \code{remainder}, \code{element} and
  \code{(setter element)} compile inline when the arguments are Lisp
  numbers or characters, or a simple vector or string and an index.
  \code{for-each}, \code{map}, \code{reduce} and their kin walk lists,
  simple vectors and ranges directly. Each fast path is guarded by a
  flag that is true only while no method a program has added could apply
  to those arguments; add a method to \code{binary+} that could apply to
  integers, and the inline path for \code{+} turns itself off, so the
  language's meaning never changes.
\item \textbf{Direct slot access.} Slots the runtime's own classes keep
  privately are read straight from the instance, and each class caches
  its slots' required init keywords for \code{make}.
\end{itemize}

\noindent
\file{generic.lisp} holds the dispatcher and its caches;
\file{runtime-fast-paths.lisp} turns the fast paths on once every
predefined method exists. \file{bench/run.lisp} measures the
implementation against the same programs written in plain Common Lisp;
Chapter~\ref{ch:performance} gives the results.

\section{The book's examples and the regression test}\label{sec:lisp-examples}

\index{regression test}\index{dylan-examples.dyl@\file{dylan-examples.dyl}}
The file \file{dylan-examples.dyl} at the top of the repository holds the
examples of the 1992 manual, extracted and edited into runnable form.
\code{dylan:test-examples} runs it with \code{dylan:test}
(Section~\ref{sec:lisp-test-file}):

\begin{shell}
$ cd src
$ sbcl --eval '(load "load-dylan.lisp")' \
       --eval '(load "full-test.lisp")' \
       --eval '(dylan:test-examples)'
\end{shell}

\noindent
The transcript it prints is compared with \file{src/reference/MIT.output},
the transcript Thomas produced from the same file on MIT Scheme in 1992
(the transcripts from Gambit and Scheme->C are kept beside it). The two
agree on every value except a handful of places where the manual and
Thomas disagree and the manual is followed: \code{\#f} and \code{()} are
different objects, floats print as \code{9.0} rather than \code{9.},
\code{(sqrt 25)} is \code{5.0}, \code{make} given a keyword that
nothing accepts is an error (the book's \code{<displaced-view>} example
does this, so it and the forms after it fail), and \code{map-into}
passes only the source elements. \file{src/README} and
Appendix~\ref{app:testing} list them.

\section{Differences from the compiler to C}\label{sec:lisp-vs-dylanc}

\index{interactive implementation!versus compiler to C}
The two implementations accept the same language, but each goes beyond
the other in places.

\paragraph{What only the interactive implementation does.} It is fully
dynamic, as the 1992 manual allows:

\begin{itemize}
\item \code{add-method}, \code{remove-method}, \code{add-slot},
  \code{remove-slot} and \code{freeze-methods} work. A program can add a
  slot to a class that already has instances:
\begin{repl}
? (define-class <sq> (<object>) (side init-keyword: side:))
;Value: <sq>

? (define s (make <sq> side: 2))
;Value: s

? (define-generic-function color (object))
;Value: color

? (add-slot <sq> getter: color init-value: 'red)
;Value: {slot color}

? (color s)
;Value: red
\end{repl}
\item New code can be compiled and run at any time, and any definition
  can be replaced (Section~\ref{sec:lisp-rep}).
\item Integers have no limit on size, and complex numbers exist:
  \code{(* 99999999999 99999999999)} is \code{9999999999800000000001},
  and \code{(sqrt -4)} is \code{\#c(0.0 2.0)}. The compiler to C has
  63-bit integers, where overflow is an error, and no complex numbers.
\item \code{lisp-variable} and \code{lisp-procedure} reach the host
  Lisp.
\end{itemize}

\paragraph{What only the compiler to C does.} Its extensions are
absent here:

\begin{itemize}
\item no threads, locks, condition variables, atomic integers, thread
  variables or \code{dynamic-bind};
\item no macros: \code{define-macro} and \code{define-for-macros} are
  not special forms here;
\item no \code{current-milliseconds}.
\end{itemize}

\noindent
Appendix~\ref{app:library} marks every name that only one
implementation provides.

\paragraph{What behaves differently.}

\begin{itemize}
\item Modules must be declared before they are used, in order; the
  compiler to C sees the whole program first (Chapter~\ref{ch:modules}).
\item Errors are reported when the code runs; the compiler to C reports
  many of them, such as an assignment to a read-only variable, when it
  compiles.
\end{itemize}

\section{Limitations}\label{sec:lisp-limitations}

\begin{itemize}
\item There is no way to load a file of Dylan source at the
  read-eval-print loop itself; use \code{dylan:test} or compile the file
  with \code{dylan:compile-file} and load the result.
\item Thomas's MIT Scheme restart example, written against MIT Scheme's
  condition system, has no counterpart; it would have to be rewritten
  against Common Lisp's.
\end{itemize}

\chapter{The Compiler to C}\label{ch:dylanc}
\sloppy

The second implementation in this system is a compiler. It reads a whole
Dylan program, translates it into C, and runs a C compiler to make a
native executable. The executable needs nothing from SBCL at run time: it
is an ordinary program linked with a small runtime library written in C,
with a garbage collector --- the Memory Pool System, or by choice the
Boehm-Demers-Weiser collector (Section~\ref{sec:dylanc-gc}) --- and
with the system's thread library.

This chapter is about using the compiler: the \cmd{bin/dylanc} command
and every one of its options, the files it writes, the cache it keeps,
the kinds of build it can make, how to call it from Common Lisp, what a
compiled program does when something goes wrong, and what a program must
not do because it is compiled whole and in advance. How the compiler
works inside is the subject of Chapter~\ref{ch:internals}, and how to
make compiled programs fast is the subject of
Chapter~\ref{ch:performance}.

\section{A first build}\label{sec:dylanc-first}

\index{dylanc@\cmd{dylanc}}
A program is one or more files of Dylan source, conventionally with the
extension \file{.dyl}. Here is the smallest useful one:

\begin{dylan}
(define-method greet ((name <string>))
  (format #t "Hello, ~a!~%" name))
(greet "world")
\end{dylan}

Compile it and run the result:

\begin{shell}
$ bin/dylanc hello.dyl
$ ./hello
Hello, world!
\end{shell}

\cmd{bin/dylanc} is a short shell script that runs
\file{compiler/dylanc.lisp} under \cmd{sbcl -{}-script}, so SBCL must be
on your \code{PATH}; the script can be run from any directory, and can be
linked into a directory on your \code{PATH}. The executable is named
after the first file, less its \file{.dyl}, and is written beside that
file: \cmd{bin/dylanc src/hello.dyl} makes \file{src/hello}.

The top-level forms of the program run in order, from the first file to
the last, when the executable starts. There is no \code{main} function to
write: the program \emph{is} its top-level forms, as it would be if you
typed them into the interactive implementation's read-eval-print loop. A
top-level form's value is not printed; use \code{print}, \code{format}
and the other printing functions (Chapter~\ref{ch:printing}) to produce
output. When the last top-level form has run, the program exits with
status~0.

\paragraph{Several files.}
Give several files to make one program from them all:

\begin{shell}
$ bin/dylanc -o app utilities.dyl shapes.dyl main.dyl
\end{shell}

The files are read in the order given, and their top-level forms run in
that order. A file whose first line is a \code{module:} header (design
note~33; see Chapter~\ref{ch:modules}) is compiled in the module it
names; any other file is compiled in \code{dylan-user}. Every
\code{define-module} form in every file is known before any other form is
compiled, so a module may be defined after the files that use it.

\paragraph{What a first build costs.}
The compiler is SBCL code, compiled into \file{compiler/fasl/} the first
time it is used, and the runtime library is compiled by \cmd{make} the
first time a build needs it. After that, a small program builds in about a
second and a half on a current machine. The first build with a new
combination of C compiler and options takes several seconds longer,
because the library's C must be compiled for it
(Section~\ref{sec:dylanc-cache}).

\section{The options}\label{sec:dylanc-options}

\index{dylanc@\cmd{dylanc}!options}
\cmd{bin/dylanc -h} lists the options:

\begin{shell}
Usage: dylanc [options] file.dyl ...

Compile the Dylan files, in order, into one executable.

Options:
  -o FILE          the executable (default: the first file's name, less .dyl)
  -O0 -O1 -O2 -O3  how much the C compiler optimizes (default -O2)
  -Os              a small executable: optimized for size, unused code dropped
                   (the library compiled with the program, as --whole)
  --whole          compile the library with the program, not apart (the
                   library's C, compiled apart, is compiled once and kept)
  --lto            optimize the program and the runtime library together when
                   linking (link-time optimization)
  -g, --debug      for debugging: no optimization, debugging information, and
                   every inline operation checking its arguments when run
  -k, --keep-c     keep the C: FILE.c, FILE-MODULE.c for each module, FILE.h
  -C, --c-only     write the C only (FILE.c and the rest); do not compile it
  -E, --expand     write each use of the program's macros and its expansion;
                   no C
  --trampoline     tail calls by trampolines, for C compilers without musttail
  --cc CC          the C compiler (default: $CC, or cc)
  --no-lines       no #line directives pointing to the Dylan source
  --sanitize KIND  for testing: a debug build with the sanitizers, KIND
                   address (AddressSanitizer and UndefinedBehaviorSanitizer)
                   or thread (ThreadSanitizer)
  --gc NAME        the garbage collector: mps (the default), the Memory Pool
                   System, generational and moving; or boehm, the installed
                   Boehm-Demers-Weiser collector (libgc)
  --gc-stress      for testing: a collection every DYLAN_GC_STRESS (default
                   100) allocations
  -h, --help       this
\end{shell}

Options and file names may be mixed in any order. Each option is
described below, with the occasions for using it.

\begin{description}
\item[\cmd{-o} \meta{file}] \index{dylanc@\cmd{dylanc}!-o@\cmd{-o}}
  Names the executable. Without it the executable is named after the
  first file. The C files, when they are kept, are named after the
  executable: \cmd{-o build/app} writes \file{build/app.c} and the rest.

\item[\cmd{-O0}, \cmd{-O1}, \cmd{-O2}, \cmd{-O3}]
  How hard the C compiler optimizes. The default, \cmd{-O2}, is right for
  almost every use. \cmd{-O0} compiles fastest, which matters only for
  large programs while you are developing them; measure before assuming
  that \cmd{-O3} helps, since it makes larger code. An
  ordinary build is also compiled with \cmd{-g}, so a debugger can show
  the Dylan source even in an optimized program (see
  Section~\ref{sec:dylanc-debug}). The Dylan compiler's own optimizations
  (Chapter~\ref{ch:internals}) are always on; these options affect only
  the C compiler.

\item[\cmd{-Os}] \index{size builds}
  Makes a small executable. The library is compiled with the program,
  and only the parts of it the program can reach are kept (\term{tree
  shaking}); the C is compiled with \cmd{-Os}, each function and datum in
  its own section, the linker drops the sections nothing uses, and the
  executable is stripped of its symbols. Use it for programs you ship or
  copy elsewhere, when size matters more than build time. Section~\ref{sec:dylanc-size}
  shows what it gains.

\item[\cmd{-{}-whole}]
  Compiles the library with the program, as \cmd{-Os} does, and shakes it
  the same way, but optimizes for speed as usual. Normally the library is
  compiled apart from the program, once, and kept
  (Section~\ref{sec:dylanc-cache}); with \cmd{-{}-whole} its C is compiled
  afresh with every build. The result is smaller than an ordinary build,
  and a build takes about twice as long. It is mostly a check that the
  separately compiled library and the shaken one agree, which the test
  suites make (Appendix~\ref{app:testing}).

\item[\cmd{-{}-lto}] \index{link-time optimization}
  Link-time optimization: the C compiler optimizes the program and the C
  runtime library together when it links them, so that small runtime
  functions can be inlined into the program. The runtime library is built
  for this the first time it is asked for, with \cmd{-O2 -g -flto}, in
  \file{runtime/build/lto-\meta{cc}/}. The effect on speed is usually
  small; try it on your own program and measure. It is ignored by
  \cmd{-g} and by the testing builds below.

\item[\cmd{-g}, \cmd{-{}-debug}]
  A build for debugging: the C is compiled with \cmd{-O0 -g
  -DDY\_DEBUG}, and the compiler trusts none of the types it infers, so
  every operation it would have done inline tests its arguments when the
  program runs. Section~\ref{sec:dylanc-debug} has more.

\item[\cmd{-k}, \cmd{-{}-keep-c}]
  Keeps the C files, which are otherwise deleted once the executable is
  made. Use it to read what the compiler wrote
  (Section~\ref{sec:internals-reading}), or to see what a debugger is
  showing you.

\item[\cmd{-C}, \cmd{-{}-c-only}]
  Writes the C files and stops: no C compiler is run and no executable is
  made. Use it to compile the C elsewhere, with a build system of your
  own, or on a machine without SBCL. Section~\ref{sec:dylanc-files}
  shows the command that compiles the result by hand.

\item[\cmd{-E}, \cmd{-{}-expand}]
  Writes each use of the program's own macros, with the place it occurs
  and what it expands into, and stops; no C is written. Use it to see
  what your macros do (Chapter~\ref{ch:extensions}).

\item[\cmd{-{}-trampoline}] \index{trampoline}
  Makes tail calls with trampolines instead of the \code{musttail}
  attribute. Every Dylan call in tail position must run in constant
  stack space, and C does not promise that; GCC~15 and later and Clang~13
  and later do, through \code{musttail}. With an older or different C
  compiler, \cmd{-{}-trampoline} makes each tail call return to a loop
  that makes the call, which any C11 compiler can do. Programs built this
  way are correct but slower: a call-heavy benchmark (\code{tak}) took
  about 1.7 times as long. Without this option, a C compiler that lacks
  \code{musttail} stops with the error ``This C compiler cannot guarantee
  tail calls (musttail): define DY\_TRAMPOLINE.''

\item[\cmd{-{}-cc} \meta{cc}]
  The C compiler to run, such as \cmd{gcc} or \cmd{clang}, or a full path.
  Without it, the compiler is the value of the \code{CC} environment
  variable, or \cmd{cc} if that is not set. It is worth trying both GCC
  and Clang on a program whose speed matters.

\item[\cmd{-{}-no-lines}]
  Leaves out the \code{\#line} directives that point the C back to the
  Dylan source (Section~\ref{sec:dylanc-files}). C compilers' messages,
  and debuggers, then refer to the C files themselves. Use it when you
  want to step through the generated C rather than the Dylan.

\item[\cmd{-{}-gc} \meta{name}] \index{garbage collection!choosing}
  The garbage collector the program is built with: \cmd{mps}, the Memory
  Pool System, the default, which collects by generations and moves
  objects; or \cmd{boehm}, the Boehm-Demers-Weiser collector, which must
  be installed. See Section~\ref{sec:dylanc-gc}. The environment variable
  \code{DYLAN\_GC=boehm} makes \cmd{boehm} the default.

\item[\cmd{-{}-sanitize} \meta{kind}]
  For testing the compiler and runtime: a debug build with
  \meta{kind} \cmd{address} (AddressSanitizer and
  UndefinedBehaviorSanitizer) or \cmd{thread} (ThreadSanitizer). See
  Section~\ref{sec:dylanc-testing-builds}.

\item[\cmd{-{}-gc-stress}]
  For testing: a garbage collection every few allocations. See
  Section~\ref{sec:dylanc-testing-builds}.

\item[\cmd{-h}, \cmd{-{}-help}]
  Prints the summary above and exits with status 0.
\end{description}

When options conflict, the more specific build wins: a testing build
(\cmd{-{}-sanitize}, \cmd{-{}-gc-stress}) uses its own C flags whatever
else is given; otherwise \cmd{-g} wins over \cmd{-Os}, and \cmd{-Os} over
\cmd{-O}\meta{n}.

\section{The files the compiler writes}\label{sec:dylanc-files}

\index{C files}
For a program whose executable is \meta{name}, the compiler writes these
C files:

\begin{description}
\item[\file{\meta{name}.c}] The program's static data --- every class,
  symbol, keyword and constant, laid out as the runtime expects --- with
  the object system's functions (the dispatch function of each generic
  function, slot accessors, and the makers \code{make} calls), and
  \code{main}.
\item[\file{\meta{name}-\meta{module}.c}] The functions written in each
  module, one file for each module. A program written in
  \code{dylan-user} has \file{\meta{name}-dylan-user.c}; a program with
  modules \code{counting} and \code{app} has
  \file{\meta{name}-counting.c} and \file{\meta{name}-app.c} too.
\item[\file{\meta{name}-dylan.c}] The library's functions: the module
  \code{dylan}, written in Dylan in \file{runtime/dylan/} and compiled
  with every program. This is the largest file by far, about a megabyte.
\item[\file{\meta{name}.h}] What the files share: a declaration of every
  object \file{\meta{name}.c} defines and a prototype of every function.
\end{description}

The C is linked with \file{runtime/build/libdylan.a}, the runtime library
compiled from \file{runtime/c/} with the collector in it, and with
\cmd{-lm -lpthread} (with \cmd{-{}-gc boehm}, the library is
\file{runtime/build/boehm-\meta{cc}/libdylan.a} and \cmd{-lgc} is
linked too). With \cmd{-k} the files stay after the build:

\begin{shell}
$ bin/dylanc -k hello.dyl
$ ls hello*
hello  hello.c  hello.dyl  hello-dylan.c  hello-dylan-user.c  hello.h
\end{shell}

\paragraph{\code{\#line} directives.} \index{\#line directives}
The generated C carries \code{\#line} directives that name the Dylan file
and line each statement came from, and after each function a directive
back to the C file itself. A C compiler's warnings, a debugger's
breakpoints and stack traces, and a profiler's reports therefore name
places in your Dylan source. Here is the start of the function
\file{hello-dylan-user.c} holds for \code{greet}:

\begin{ccode}
D fn1190_greet(D self, D a0, D a1, D a2, D a3, intptr_t argc) {
    (void)self; (void)a0; (void)a1; (void)a2; (void)a3; (void)argc;
    D v105609, v105610, v105611;
#line 1 "hello.dyl"
    if (DY_UNLIKELY(argc != 1)) dy_wrong_argument_count(self, argc);
    v105609 = a0;
#line 2 "hello.dyl"
    ...
\end{ccode}

\paragraph{Compiling the C yourself.}
With \cmd{-C}, the compiler writes the C and stops. The files can then be
compiled with the C compiler's options and the runtime library's
directory:

\begin{shell}
$ bin/dylanc -C hello.dyl
$ cc -O2 -std=c11 -pthread -DDY_MPS \
     -I$DYLAN/runtime/c -I$DYLAN/runtime/mps/code \
     hello.c hello-dylan.c hello-dylan-user.c \
     $DYLAN/runtime/build/libdylan.a -lm -lpthread -o hello
$ ./hello
Hello, world!
\end{shell}

Here \code{\$DYLAN} stands for the root of the source tree. The runtime
library \file{libdylan.a} must exist; \cmd{make -C runtime} makes it, and
any build with \cmd{dylanc} makes it too. Add \cmd{-DDY\_TRAMPOLINE} to
the C compiler's options, and link \file{libdylan-trampoline.a} instead,
if your C compiler lacks \code{musttail}. For the Boehm collector, leave
out \cmd{-DDY\_MPS} and the MPS include directory, link
\file{runtime/build/boehm-cc/libdylan.a} (made by any
\cmd{dylanc -{}-gc boehm} build), and add \cmd{-lgc}. The C is C11 with POSIX
threads; \file{runtime/c/dylan.h} checks what it assumes (64-bit
pointers, arithmetic right shifts) with \code{\_Static\_assert}, and
checks that the runtime it is compiled against is the version the
compiler wrote for.

\section{The library's cache}\label{sec:dylanc-cache}

\index{library cache}
Every program carries the whole Dylan library, about a megabyte of C. To
avoid compiling that with every build, the compiler writes the library's
C so that it does not depend on the program: \file{\meta{name}-dylan.c}
is the same text for almost every program. The driver compiles it once
for each C compiler and set of options and keeps the object in
\file{runtime/build/library/}, named for a digest (MD5) of the C text,
the C compiler's name and \cmd{-{}-version} output, the options, and the
runtime headers \file{dylan.h} and \file{primitive.h}. A later build that
would compile the same text the same way uses the object instead. This
halves the time of a small build.

The newest 40 objects are kept, with the C each was compiled from; older
ones are deleted. The cache is safe to delete at any time
(\cmd{rm -r runtime/build/library}); it is rebuilt as needed. Builds
running at the same time share it safely.

A program can change the library's C --- for example by defining a method
that replaces one of the library's own. Then the digest differs and the
library is simply compiled for that build. \cmd{-Os} and
\cmd{-{}-whole} do not use the cache at all: they compile the library
with the program, keeping only what the program can reach.

\section{Size}\label{sec:dylanc-size}

These are the sizes of the \code{greet} program above, built with GCC~16
on x86-64 Linux, as reported by \cmd{size} (the code, the initialized
data and the uninitialized data, excluding the shared C and collector
libraries):

\begin{center}
\begin{tabular}{lrr}
\toprule
Build & text + data (KB) & file size (KB) \\
\midrule
default (\cmd{-O2}) & 463 & 2765 \\
\cmd{-{}-whole} & 307 & 1798 \\
\cmd{-{}-lto} & 334 & 2374 \\
\cmd{-Os} & 221 & 231 \\
\bottomrule
\end{tabular}
\end{center}

The file sizes of all but the \cmd{-Os} build include debugging
information. Hello world built \cmd{-Os} is about 220~KB of code and data;
most of it is library code that a call of a generic function such as
\code{print-object}, \code{make} or \code{as} could reach.
Chapter~\ref{ch:performance} says more about size.

\section{Debug builds}\label{sec:dylanc-debug}

\index{debugging}
A program built with \cmd{-g} differs from an ordinary one in three ways:

\begin{itemize}
\item The C is compiled with \cmd{-O0 -g}, so a debugger can show every
  variable and step through every line.
\item The compiler trusts none of the classes it has inferred. Where an
  ordinary build adds two integers with no test because it has proved
  them integers, a debug build tests them anyway. A debug build therefore
  catches a mistake in the compiler's reasoning as an error, rather than
  as a crash or a wrong answer.
\item \code{DY\_DEBUG} is defined, which makes the runtime's accessors
  check the representation of each value they take apart: that what is
  read as a pair \emph{is} a pair, and so on. A violation stops the
  program with a message naming what was expected.
\end{itemize}

None of these change what a correct program does: a debug build reports
the same Dylan errors, in the same way, as an ordinary one. Debug builds
are for finding bugs in the compiler or runtime, and for using a
debugger. Because of the \code{\#line} directives, a debugger works in
Dylan terms --- breakpoints can be set at lines of the Dylan file. Given
this \file{sum.dyl},

\begin{dylan}
(define-method sum-to ((n <integer>))
  (for ((i 0 (+ i 1))
        (s 0 (+ s i)))
       ((> i n) s)))
(print (sum-to 100))
\end{dylan}

\noindent a debug build stops where it is asked to:

\begin{shell}
$ bin/dylanc -g -o sum sum.dyl
$ gdb -batch -ex "break sum.dyl:3" -ex run sum
...
Thread 1 "sum" hit Breakpoint 1, fn1190_sum_to (self=4692864, a0=201,
    a1=0, a2=0, a3=0, argc=1) at /home/me/sum.dyl:3
3               (s 0 (+ s i)))
\end{shell}

The names of the C functions are those the compiler gave the Dylan
functions (\code{fn1190\_sum\_to} is the method \code{sum-to}), and the
values are Dylan values in their machine representation: \code{a0=201}
is the integer 100, stored shifted left one bit with the low bit set.
Chapter~\ref{ch:internals} explains the representation.

\section{Testing builds}\label{sec:dylanc-testing-builds}

Two options build programs for testing the compiler and runtime rather
than for use. Each builds its own variant of the runtime library, the
first time it is needed, in \file{runtime/build/\meta{variant}-\meta{cc}/},
with the same C flags as the program.

\begin{description}
\item[\cmd{-{}-sanitize address}] Compiles with \cmd{-O1 -g
  -fno-omit-frame-pointer -fsanitize=address,undefined -DDY\_DEBUG}, as a
  debug build: AddressSanitizer reports any read or write outside an
  object, and UndefinedBehaviorSanitizer any undefined behavior in C,
  such as a signed overflow or a misaligned access.
\item[\cmd{-{}-sanitize thread}] Compiles with \cmd{-O1 -g
  -fsanitize=thread -DDY\_DEBUG}: ThreadSanitizer reports data races
  between threads.
\item[\cmd{-{}-gc-stress}] Compiles with \cmd{-O2 -g -DDY\_GC\_STRESS}:
  the program forces a garbage collection every \meta{n} allocations,
  where \meta{n} is the value of the environment variable
  \code{DYLAN\_GC\_STRESS}, or 100. A pointer the collector cannot see
  --- a bug in the compiler or runtime --- shows up quickly as a crash or
  a wrong answer. Such a program runs very slowly.
\end{description}

\begin{shell}
$ bin/dylanc --gc-stress -o hgc hello.dyl
$ DYLAN_GC_STRESS=10 ./hgc
Hello, world!
\end{shell}

One of AddressSanitizer's checks is turned off by the runtime itself: its
detection of a frame used after its function returned, which keeps the
frames of functions apart from the stack. The collector finds the
references a program holds by scanning its stack, and would not see those
frames; an object referred to only from one would be freed, or moved,
while in use. The runtime supplies the option
\code{detect\_stack\_use\_after\_return=0} as the sanitizer's default,
so if you set \code{ASAN\_OPTIONS} yourself, leave that one off.

The test suites use these builds (Appendix~\ref{app:testing}). An
ordinary Dylan program has no need of them, but a program that misbehaves
in a way you cannot explain is worth running under them: what they find
is a bug to report.

\section{The garbage collector}\label{sec:dylanc-gc}

\index{garbage collection!choosing}\index{Memory Pool System}
\index{Boehm collector}
A compiled program's memory is managed by one of two collectors, chosen
when it is built.

\begin{description}
\item[\cmd{-{}-gc mps}] (the default) The Memory Pool System (MPS),
  Ravenbrook Limited's collector, written in the first place for
  Harlequin's Dylan. Its source is in the tree (\file{runtime/mps/},
  under its own BSD license) and is compiled into the runtime library;
  nothing need be installed. It is \emph{generational} and
  \emph{moving}: new objects go in a nursery, a collection of the nursery
  copies the few that survive out of it and leaves the rest behind
  without visiting them, and the older generations are collected rarely.
  The stacks and the program's static data are scanned conservatively,
  so an object they point to stays in place; the heap itself is scanned
  exactly. Each thread allocates by bumping a pointer in its own region,
  without locking.
\item[\cmd{-{}-gc boehm}] The Boehm-Demers-Weiser collector,
  \file{libgc}, which must be installed on the system
  (Section~\ref{sec:start-requirements}). It is \emph{conservative}: it
  treats any word that looks like a pointer as one, so it needs no
  knowledge of the program's objects, and objects never move. It
  collects by marking and sweeping the whole heap, in parallel on two
  threads, and allocates from per-thread free lists. It runs on more
  platforms than the MPS, and is the choice where the MPS has no port.
\end{description}

The MPS is the faster. A program that makes many short-lived objects
--- lists built and discarded, boxed floats, strings, conditions ---
gains most, two to three times on such programs, since a collection
costs in proportion to what survives, not to what was allocated; a
program that keeps a large heap alive gains from the old generations
being collected rarely; and a program that allocates little sees no
difference. The price is memory: a program that allocates heavily uses
about twice as much under the MPS, whose nursery is 128~MB and whose
heap may double between collections of everything, and
Section~\ref{sec:perf-gc} tells how to trade that back. The two
collectors are the same to the Dylan program, with one visible
exception: under the MPS an object's identity hash is a number given it
the first time it is hashed, not its address, since the address changes.

Every object of a program must be built for the same collector, and the
runtime library for each is kept apart: the MPS one is
\file{runtime/build/libdylan.a}, and the Boehm one is in
\file{runtime/build/boehm-\meta{cc}/}, with the variant and
\cmd{-{}-trampoline} in the name as for testing builds
(Section~\ref{sec:dylanc-testing-builds}). The testing builds work with
either collector. Under the MPS, \cmd{-{}-gc-stress} forces a
collection of every generation, and a ThreadSanitizer build uses the MPS
without its memory barrier, whose signal handling and the sanitizer's
would deadlock, so its collections there are not incremental. Setting
the environment variable \code{DYLAN\_GC=boehm} makes the Boehm
collector the default, which is how the test suites
(Appendix~\ref{app:testing}) are run with it. The MPS's behavior can be
tuned by environment variables, in Section~\ref{sec:perf-gc}. It has
been built here on Linux on x86-64 and macOS on Apple silicon, and has
ports to FreeBSD, to Windows with Visual C, and to Linux on ARM64.

\section{Calling the compiler from Lisp}\label{sec:dylanc-lisp}

\index{dylanc@\code{dylanc} package}
The compiler is a Common Lisp program, and it can be called from a Lisp
session. Load \file{compiler/load.lisp}, which loads the interactive
implementation's reader and the compiler, compiling any source newer
than its fasl; the package is \code{DYLAN-COMPILER}, nicknamed
\code{DYLANC}. Two functions do the work:

\begin{description}
\item[\code{(dylanc:build \meta{files} \meta{executable} \&key ...)}]
  Compiles the list of Dylan \meta{files} to C and on to the
  \meta{executable}. The keyword arguments correspond to the options of
  \cmd{dylanc}: \code{:optimize} (a string such as \code{"-O3"},
  default \code{"-O2"}), \code{:size}~\code{t} (\cmd{-Os}),
  \code{:whole}~\code{t}, \code{:lto}~\code{t}, \code{:debug}~\code{t},
  \code{:trampoline}~\code{t}, \code{:cc} (default \code{"cc"}; the
  \code{CC} environment variable is not consulted here),
  \code{:variant} (\code{:address}, \code{:thread} or
  \code{:gc-stress}), \code{:gc} (\code{:mps} or \code{:boehm}), and
  \code{:c-file}, the name of the main C file (default
  \meta{executable}\code{.c}). It answers \code{T} if the
  executable was made and \code{NIL} if not. Diagnostics are written to
  \code{*error-output*}. Unlike \cmd{dylanc}, it leaves the C files in
  place.
\item[\code{(dylanc:compile-files \meta{files} \meta{c-file})}]
  Writes the C only, as \cmd{-C} does. It answers the list of C files
  written, or \code{NIL} after reporting errors.
\end{description}

\begin{lisp}
(load "compiler/load.lisp")
(dylanc:build '("hello.dyl") "hello")              ; => T
(dylanc:build '("hello.dyl") "hello-small" :size t) ; => T
(dylanc:compile-files '("hello.dyl") "cf.c")
;; => (#P"cf.c" #P"cf-dylan.c" #P"cf-dylan-user.c")
\end{lisp}

Relative file names are taken relative to the Lisp's current directory
(\code{*default-pathname-defaults*}).

\section{Errors when compiling}\label{sec:dylanc-diagnostics}

\index{diagnostics}
The compiler reports problems as \term{diagnostics}, one to a line, on
the standard error stream, in the form \meta{file}\code{:}\meta{line}%
\code{:}\meta{column}\code{: }\meta{kind}\code{: }\meta{message}, where
\meta{kind} is \code{error} or \code{warning}:

\begin{shell}
$ cat r1.dyl
(print (undefined-thing 3))
$ bin/dylanc r1.dyl
r1.dyl:1:8: warning: undefined-thing is not defined (module dylan-user)
\end{shell}

A \textbf{warning} does not stop the build. The commonest is a variable
that is referred to but defined nowhere in the program; the program is
built, and fails only if the reference is reached (see
Section~\ref{sec:dylanc-runtime-errors}). Treat such a warning as an
error unless you know why it is there: it is usually a misspelling.

An \textbf{error} stops the build: no executable is made and
\cmd{dylanc} exits with status~1. Errors found while reading and
converting the source --- an unbalanced parenthesis, a malformed special
form, an integer too large --- are all reported in one run; the compiler
abandons the top-level form that has the error and goes on to the next.
Errors found later, when the program is analysed as a whole, stop at the
first. For example:

\begin{shell}
$ cat err2.dyl
(define f (method (x) (+ x 1))
(print (f 1))
$ bin/dylanc err2.dyl
err2.dyl:2:14: error: End of file inside a Dylan list.
\end{shell}

The position of a reading error is where the reader stopped, which for a
missing closing parenthesis is the end of the file; the unclosed list
began earlier.

\paragraph{Exit status.} \index{exit status!of dylanc@of \cmd{dylanc}}
\cmd{dylanc} exits with status 0 when it has made what was asked; 1 when
the program had errors or the C compiler failed (``\code{dylanc: the C
compiler failed.}''); and 2 for a bad command line --- an unknown option,
an option missing its value, no files, or a file that does not exist ---
after printing the usage summary.

\section{Errors when running}\label{sec:dylanc-runtime-errors}

\index{unhandled conditions}
A compiled program signals conditions as Chapter~\ref{ch:conditions}
describes, and handlers established with \code{handler-case} or
\code{handler-bind} see them as they would in the interactive
implementation. The runtime's own errors --- a wrong number of
arguments, no applicable method, a type error, integer overflow,
division by zero, an unbound variable --- are signalled as Dylan
conditions too, \code{<simple-error>} or \code{<type-error>}, and can be
handled.

There is no debugger to fall into when a condition is not handled.
Instead, the default handler (\code{default-handler}) acts as follows:

\begin{itemize}
\item An unhandled \textbf{serious condition} --- any \code{<error>},
  among others --- writes \code{Dylan error:} and the condition's message
  to the standard error stream, and the program exits with status~1.
  The standard output is flushed first, so everything printed before the
  error appears.
\item An unhandled \textbf{warning} writes \code{Dylan warning:} and its
  message to the standard error stream, and \code{signal} returns
  \code{\#f}; the program continues.
\item \code{break} writes \code{Dylan break:} and its message, and
  returns \code{\#f}.
\item \code{restart-query} and \code{return-query} do nothing.
\end{itemize}

\begin{dylan}
(define f (method (x) (+ x 1)))
(print (f 1))
(print (f "one"))
(print 'never)
\end{dylan}

\begin{shell}
$ bin/dylanc err1.dyl && ./err1; echo "status $?"
2
Dylan error: no applicable method for binary+
status 1
\end{shell}

The message names the generic function \code{binary+}: a call of
\code{+} with two arguments calls the library's binary addition directly
(Chapter~\ref{ch:internals}). In a thread other than the main one, an
unhandled serious condition ends only that thread; see
Chapter~\ref{ch:extensions}.

\begin{warnbox}
A recursion that is \emph{not} in tail position uses C stack, and a
program that recurses too deeply --- millions of frames --- overflows the
stack and is killed by the operating system with a segmentation fault,
not a Dylan error. Output it had printed but not yet flushed is lost.
Calls in tail position never do this: they run in constant space however
deep they go. Write deep recursions as tail calls, or raise the stack
limit with \cmd{ulimit -s}.
\end{warnbox}

\section{What a compiled program cannot do}\label{sec:dylanc-static}

\index{static program}
The compiler compiles a \term{static program}: everything the program
defines is known when it is compiled, and its classes and generic
functions do not change while it runs. That is what lets the compiler
lay out every class, compile every generic function's dispatch in
advance, and call methods directly. The 1992 manual describes a dynamic
language in which a running program may add methods and slots and make
new classes, and the interactive implementation allows all of that. A
compiled program does not. The restrictions, and how each shows itself,
are these.

\paragraph{Changing classes and generic functions.}
\code{add-method}, \code{remove-method}, \code{add-slot},
\code{remove-slot} and \code{freeze-methods} cannot be used. Any use of
them --- even one that would never run --- is a compile error:

\begin{shell}
$ cat s.dyl
(define-generic-function g (x))
(add-method g (method ((x <integer>)) x))
$ bin/dylanc s.dyl
s.dyl:2:1: error: add-method cannot be used: a program's classes and
generic functions are fixed when it is compiled
\end{shell}

(The message is on one line; it is broken here to fit the page.)
Define methods with \code{define-method} instead, which adds them when
the program is compiled. A class cannot be defined twice
(``\code{<a> is defined twice (classes cannot be redefined)}'').

\paragraph{Specializers and superclasses known when compiling.}
Each specializer of a method made by \code{define-method}, and each
superclass of a class, must be something the compiler can find: a class
named directly, a variable defined as a constant class, \code{(singleton
\meta{x})}, or \code{(make <singleton> object: \meta{x})} of a constant.
A variable whose value is computed when the program runs is a compile
error:

\begin{dylan}
(define c (if (> 2 1) <integer> <string>))
(define-method m ((x c)) x)
; error: a specializer must be a class or a singleton known at compile time
\end{dylan}

\noindent while this is accepted, because \code{c} is defined as a class:

\begin{dylan}
(define c <integer>)
(define-method m ((x c)) 'yes)
(print (m 3))        ; prints yes
\end{dylan}

The owner of a slot added with \code{define-slot} may be any singleton,
but a class owner must be known when compiling.

\paragraph{Making classes and generic functions when running.}
\code{(make <class> ...)} is an error when it runs (``cannot make an
instance of this class''). \code{(make <generic-function> ...)}
answers a generic function with no methods; since no method can be
added, every call of it signals ``no applicable method''. \code{seal}
works as the manual says, and \code{make-read-only} answers its
argument.

\paragraph{Numbers.}
Integers are 63 bits: from $-2^{62}$ to $2^{62}-1$, that is from
\code{-4611686018427387904} to \code{4611686018427387903}. There are no
bignums. An integer literal outside that range is a compile error; an
operation whose result falls outside it signals an error when it runs:

\begin{dylan}
(print (* 4611686018427387903 2))
; Dylan error: integer overflow (integers are 63 bits)
\end{dylan}

The same overflow rule applies to the numerator and denominator of a
ratio. \code{<complex>} exists, as the abstract superclass of
\code{<real>}, but has no instances, and a complex literal such as
\code{\#c(1 2)} is a compile error. The interactive implementation has
both bignums and complex numbers (Chapter~\ref{ch:numbers}).

\paragraph{No evaluation at run time.}
A compiled program cannot evaluate or load code while it runs: there is
no \code{eval} and no \code{load}. Any code it runs must be in its
source when it is compiled. (Macros, which run at compile time, can
compute code; see Chapter~\ref{ch:extensions}.)

\paragraph{Undefined variables.}
A reference to a variable that nothing defines is only a warning when
compiling, as shown in Section~\ref{sec:dylanc-diagnostics}. If the
reference is reached, it is an error when the program runs:

\begin{shell}
$ ./r1
Dylan error: the variable undefined-thing@dylan-user has no value
\end{shell}

Appendix~\ref{app:differences} collects these restrictions with every
other difference from the 1992 manual.

\fussy

\chapter{Macros and Threads}\label{ch:extensions}
\sloppy

The 1992 manual defines a language without macros and without threads.
The compiler to C adds both. They are \emph{extensions}: a program that uses them is
not a program of the 1992 language, and the interactive Common Lisp
implementation does not have them (it has its own extensions, described
in Chapter~\ref{ch:lisp}). Everything in this chapter applies to
programs compiled with \cmd{dylanc}.

\section{Macros}\label{sec:ext-macros}

\index{macros}
A \term{macro} is a function that runs while the program is being
compiled. It receives the forms of a call, unevaluated, and answers a new
form, its \term{expansion}, which the compiler compiles in place of the
call. Macros let you add new syntax to the language: a new looping form,
a definition form that writes several definitions, a form that evaluates
its argument in a special way.

The macros here are Common Lisp's kind. If you know \code{defmacro}, you
know them already: the parameters are bound to the argument forms, the
body builds the expansion, usually with backquote, and nothing protects
the expansion's names from the names around the call except the care of
the macro's writer and \code{gensym}.

The manual calls \code{and}, \code{or}, \code{cond}, \code{case},
\code{for} and the other syntax of Chapter~\ref{ch:control} macros. In
this system those are built into the compiler, not written with
\code{define-macro}, and the program cannot redefine them: the compiler
accepts a \code{define-macro} of one of their names without complaint,
but a call of the name still means the built-in form.

\subsection{\code{define-macro}}\label{sec:ext-define-macro}

\index{define-macro@\code{define-macro}}
\begin{quote}
\code{(define-macro \meta{name} \meta{parameter-list} \meta{form} ...)}
\end{quote}

\noindent at top level makes \meta{name}, a variable of the current
module, a macro. From then on, a call \code{(\meta{name}
\meta{argument} ...)} is replaced, while compiling, by the value of the
\meta{form}s, run with the parameters bound to the argument forms. The
parameter list is a method's: required parameters, \code{\#rest} and
\code{\#key} are all allowed.

\begin{dylan}
(define-macro my-unless (test #rest body)
  `(if ,test #f (begin ,@body)))

(my-unless (= 1 2)
  (print 'one)
  (print 'two))
(print (my-unless #t 'never))
\end{dylan}

\noindent prints

\begin{dylan}
one
two
#f
\end{dylan}

Here \code{test} is bound to the form \code{(= 1 2)} and \code{body} to
the list of the two \code{print} forms; the backquote template builds an
\code{if} form from them. Backquote follows the usual Lisp rules
(\code{,} inserts a value, \code{,@} splices a list, and backquotes may
nest); see Appendix~\ref{app:differences}.

Because the argument forms arrive unevaluated, a macro can use a form
both as data and as code:

\begin{dylan}
(define-macro show (form)
  `(format #t "~s => ~s~%" ',form ,form))

(show (+ 1 2))
(show (reverse '(a b c)))
\end{dylan}

\noindent prints

\begin{dylan}
(+ 1 2) => 3
(reverse (quote (a b c))) => (c b a)
\end{dylan}

A macro's body is ordinary Dylan, and it may compute its expansion any
way it likes; it need not use backquote, and its result need not even be
a list. Keyword parameters work as in a method:

\begin{dylan}
(define-macro scaled (#key (scale 1) value)
  (* scale value))
(print (list (scaled value: 5) (scaled scale: 3 value: 5)))
; prints (5 15)
\end{dylan}

Here the multiplication happens when the program is compiled; the
expansions are the constants \code{5} and \code{15}.

Macros may expand into uses of other macros, or of themselves; the
compiler expands again until no macro call is left. An expansion may also
be a definition, or a \code{begin} of several, at top level. A macro that
never stops expanding is reported (``macros expanded 500 deep --- one
never stops?'').

\subsection{Rules}

\begin{itemize}
\item \code{define-macro} must be at top level. Inside a method it is a
  compile error (``define-macro must be at top level'').
\item A macro must be defined before its first use in the program. A
  call that comes earlier is compiled as a call of a variable that is not
  defined: the compiler warns, and the call fails when it runs.
\item A name cannot be both a macro and a variable. \code{(define m 2)}
  after \code{(define-macro m () 1)}, or a use of \code{m} other than as
  the operator of a call, is a compile error (``m is a macro'').
\item A call of a macro must be a proper list, and must supply arguments
  its parameter list accepts (``macro m: wrong number of arguments'').
\item An error signalled by a macro's body is a compile error, reported
  at the call with the macro's name: \code{(define-macro m (a) (error "no
  good: \textasciitilde{}s" a))} and a call \code{(m 1)} report
  ``\code{macro m: no good: 1}''.
\end{itemize}

\subsection{Names in expansions: \code{gensym} and capture}
\label{sec:ext-hygiene}

\index{gensym@\code{gensym}}\index{macros!hygiene}
The names in an expansion mean what they mean \emph{where the macro is
used}. This is simple and predictable, but it means that the expansion
and the code around it can interfere. A macro that binds a variable of
its own risks capturing a variable of the same name at the call:

\begin{dylan}
(define-macro bad-swap! (a b)
  `(bind ((tmp ,a)) (set! ,a ,b) (set! ,b tmp)))

(bind ((tmp 1) (other 2))
  (bad-swap! tmp other)
  (print (list tmp other)))
; prints (1 2): nothing was swapped
\end{dylan}

The expansion binds a new \code{tmp}, and \code{(set! tmp other)} then
assigns that one, not the caller's. The cure is \code{(gensym)}, which
answers a fresh symbol that no program can write, for the expansion to
bind:

\begin{dylan}
(define-macro swap! (a b)
  (bind ((tmp (gensym "tmp")))
    `(bind ((,tmp ,a))
       (set! ,a ,b)
       (set! ,b ,tmp))))

(define tmp 1)
(define other 2)
(swap! tmp other)
(print (list tmp other))      ; prints (2 1)
\end{dylan}

\code{gensym} takes an optional string, the start of the new name, and
appends a number. The name only \emph{looks} like one a program could
write: \cmd{dylanc -E} shows it as \code{tmp1}, but it is not the symbol
\code{tmp1}, and nothing in the program can refer to it.

Capture can happen the other way too: an expansion that calls a function
by name gets whatever that name means at the call. If the caller has a
local variable of that name, the expansion uses it:

\begin{dylan}
(define-macro pair-of (x) `(list ,x ,x))
(print (pair-of 1))                     ; prints (1 1)
(bind ((list vector))
  (print (pair-of 2)))                  ; prints #(2 2)
\end{dylan}

\noindent There is no protection against this short of not binding, near
a macro's use, the names its expansions call. The built-in syntax ---
\code{for}, \code{case}, \code{with-lock} and the rest --- is not exposed
to either kind of capture: its expansions refer to the library's
variables directly and make their temporaries with names no program can
write.

A local variable or parameter with the name of a macro hides the macro,
as it would hide a module variable:

\begin{dylan}
(define-macro double (x) `(* 2 ,x))
(print (double 4))                              ; prints 8
(define apply-it (method (double) (double 4)))
(print (apply-it (method (n) (+ n 100))))       ; prints 104
\end{dylan}

\subsection{The compile-time environment and \code{define-for-macros}}
\label{sec:ext-compile-time}

\index{define-for-macros@\code{define-for-macros}}
Macro bodies run inside the compiler, in the interactive Common Lisp
implementation of Dylan (Chapter~\ref{ch:lisp}), which is loaded with
the compiler. They run in a \emph{compile-time environment} of their own,
separate from the program being compiled. It holds the whole library of
the interactive implementation, \code{gensym}, and whatever
\code{define-for-macros} has defined --- but nothing the program itself
defines. A macro body that calls one of the program's functions fails:

\begin{dylan}
(define helper (method (x) (list x x)))
(define-macro use-helper (x) (helper x))
(print (use-helper 3))
; error: macro use-helper: unbound variable helper
\end{dylan}

Functions and variables for macro bodies to use are defined with

\begin{quote}
\code{(define-for-macros \meta{form} ...)}
\end{quote}

\noindent at top level, which runs the \meta{form}s in the compile-time
environment while compiling, and leaves nothing in the compiled program.
A variable defined there keeps its value from one expansion to the next.
Here a helper builds the names of accessors for a macro that defines a
class:

\begin{dylan}
(define-for-macros
  (define accessor-name
    (method (class-name slot)
      (as <symbol>
          (concatenate (as <string> class-name) "-" (as <string> slot))))))

(define-macro define-record (name #rest slots)
  `(define-class ,name (<object>)
     ,@(map (method (s)
              `(,(accessor-name name s)
                init-keyword: ,(as <keyword> (as <string> s))))
            slots)))

(define-record point x y)
(bind ((p (make point x: 3 y: 4)))
  (print (list (point-x p) (point-y p))))   ; prints (3 4)
\end{dylan}

Two consequences of running macro bodies in the interactive
implementation are worth knowing. First, \code{gensym} and anything
\code{define-for-macros} defines exist only at compile time; a compiled
program cannot call \code{gensym}. Second, where the two implementations
differ, a macro body sees the interactive one: it has bignums, for
instance (Appendix~\ref{app:differences}).

\subsection{Macros and modules}\label{sec:ext-macro-modules}

A macro is a variable of the module it is defined in, and it is
exported, imported, renamed and excluded exactly as a variable is
(Chapter~\ref{ch:modules}). Here the module \code{macros} defines and
exports \code{repeat}, and \code{user} uses it:

\begin{dylan}
; macmods.dyl
(define-module macros (use dylan) (export repeat))
(define-module user (use dylan) (use macros))
\end{dylan}

\begin{dylan}
; macdef.dyl
module: macros

(define-macro repeat (n #rest body)
  (bind ((i (gensym)))
    `(dotimes (,i ,n) ,@body)))
\end{dylan}

\begin{dylan}
; macuse.dyl
module: user

(repeat 3 (print 'hi))
\end{dylan}

\begin{shell}
$ bin/dylanc -o macs macmods.dyl macdef.dyl macuse.dyl
$ ./macs
hi
hi
hi
\end{shell}

(The first lines of \file{macdef.dyl} and \file{macuse.dyl} are the
\code{module:} headers; the comments above them are not part of the
files.) The names in an expansion are looked up in the module where the
macro is \emph{used}, so \code{dotimes} must mean the same there as in
the macro's module --- here both use \code{dylan}.

\subsection{Seeing expansions: \cmd{dylanc -E}}\label{sec:ext-expand}

\index{dylanc@\cmd{dylanc}!-E@\cmd{-E}}
\cmd{dylanc -E} writes each use of the program's own macros, where it
occurs, and its expansion, and makes no C. It is the way to debug a
macro:

\begin{shell}
$ bin/dylanc -E m1.dyl
m1.dyl:4:1: (my-unless (= 1 2) (print 'one) (print 'two))
  => (if (= 1 2) #f (begin (print 'one) (print 'two)))
m1.dyl:7:8: (my-unless #t 'never)
  => (if #t #f (begin 'never))
m1.dyl:12:1: (show (+ 1 2))
  => (format #t "~s => ~s~%" '(+ 1 2) (+ 1 2))
m1.dyl:13:1: (show (reverse '(a b c)))
  => (format #t "~s => ~s~%" '(reverse '(a b c)) (reverse '(a b c)))
\end{shell}

Only the program's macros are shown, not the built-in syntax. Each step
of a macro that expands into another macro is shown separately.

\section{Threads}\label{sec:ext-threads}

\index{threads}
A compiled program can run several \term{threads} at once. Each Dylan
thread is a POSIX thread, so threads really do run in parallel on a
machine with several processors. The garbage collector is built for
threads and collects while they all allocate. Threads share the
program's module variables and every object they can reach; each has
its own stack, its own handlers and exits, and its own values of thread
variables (Section~\ref{sec:ext-thread-variables}).

The threads API is small. Table~\ref{tab:ext-threads} lists it; the
sections after it explain each part.

\begin{table}[htbp]
\centering
\small
\begin{tabular}{>{\raggedright\arraybackslash}p{0.47\textwidth}%
                >{\raggedright\arraybackslash}p{0.45\textwidth}}
\toprule
Form & What it does \\
\midrule
\code{(make-thread \meta{function})} & start a thread calling
  \meta{function}, which takes no arguments; answer the \code{<thread>} \\
\code{(join-thread \meta{thread})} & wait for \meta{thread} to end;
  answer its function's values \\
\code{(current-thread)} & the \code{<thread>} running \\
\code{(thread-yield)} & let other threads run; answer \code{\#f} \\
\code{(make <lock>)}, \code{(make <recursive-lock>)} & a new lock \\
\code{(acquire-lock \meta{lock})}, \code{(release-lock \meta{lock})} &
  hold and let go of a lock \\
\code{(with-lock (\meta{lock}) \meta{form} ...)} & the forms, holding
  the lock \\
\code{(make <condition-variable>)} & a new condition variable \\
\code{(wait-condition-variable \meta{cv} \meta{lock})} & release the
  lock, wait to be notified, acquire it again \\
\code{(notify-condition-variable \meta{cv})} & wake one waiting thread \\
\code{(notify-all-condition-variable \meta{cv})} & wake every waiting
  thread \\
\code{(make <atomic-integer> value: \meta{n})} & an integer cell for
  atomic updates (\meta{n} defaults to 0) \\
\code{(atomic-value \meta{a})} & its value \\
\code{(atomic-add! \meta{a} \meta{n})} & add \meta{n}; answer the new
  value \\
\code{(atomic-compare-and-swap! \meta{a} \meta{old} \meta{new})} & if
  it holds \meta{old}, store \meta{new} and answer \code{\#t}; else
  \code{\#f} \\
\code{(define-thread-variable \meta{name} \meta{value})} & a module
  variable with a value for each thread \\
\code{(dynamic-bind ((\meta{name} \meta{value}) ...) \meta{form} ...)} &
  the forms, with thread variables given new values in this thread \\
\code{<thread-error>}, \code{thread-error-thread},
  \code{thread-error-condition} & what \code{join-thread} signals when a
  thread ended with an unhandled condition \\
\bottomrule
\end{tabular}
\caption{The threads API.}\label{tab:ext-threads}
\end{table}

\subsection{Making and joining threads}\label{sec:ext-make-thread}

\index{make-thread@\code{make-thread}}\index{join-thread@\code{join-thread}}
\code{(make-thread \meta{function})} starts a new thread that calls
\meta{function} with no arguments, and answers at once with the
\code{<thread>} object. \code{(join-thread \meta{thread})} waits until
that thread's function returns, and answers the values it returned. A
thread may be joined more than once, and by any thread but itself; each
join answers the same values.

This program sums the integers below a million in four threads, each
summing a quarter, and adds the four results:

\begin{dylan}
(define sum-range
  (method (from below)
    (for ((i from (+ i 1))
          (s 0 (+ s i)))
         ((= i below) s))))

(define workers
  (map (method (k)
         (make-thread
           (method () (sum-range (* k 250000) (* (+ k 1) 250000)))))
       '(0 1 2 3)))

(print (reduce + 0 (map join-thread workers)))
; prints 499999500000
\end{dylan}

The threads run in whatever order the operating system chooses, but the
program's output does not depend on it: each thread computes alone, and
the main thread waits for all four before it adds. Write threaded
programs this way where you can --- threads that share nothing while they
work and are joined to collect their results.

A thread's function may return several values, and \code{join-thread}
answers them all:

\begin{dylan}
(define t (make-thread (method () (values 1 2 3))))
(print (bind ((a b c (join-thread t))) (list a b c)))   ; prints (1 2 3)
\end{dylan}

\code{(current-thread)} answers the thread that calls it. In the main
thread it answers a \code{<thread>} too, but the main thread cannot be
joined. \code{(thread-yield)} asks the operating system to run other
threads before this one continues; it answers \code{\#f}. A thread
prints as \code{\{thread\}}.

\subsection{When the program ends}\label{sec:ext-program-exit}

\index{threads!program exit}
The program ends when the main thread has run the last top-level form,
\emph{whatever its other threads are doing}. It does not wait for them;
they simply stop. This program prints \code{main-is-done} and exits with
status 0, although its other thread would loop forever:

\begin{dylan}
(define spinner
  (make-thread
    (method ()
      (bind-methods ((spin () (thread-yield) (spin)))
        (spin)))))
(print 'main-is-done)
\end{dylan}

Join every thread whose work must finish. Equally, an unhandled error in
the main thread ends the whole program at once, with the other threads.

\subsection{Unhandled conditions in threads}\label{sec:ext-thread-errors}

\index{thread-error@\code{<thread-error>}}
A serious condition that a thread does not handle ends \emph{that thread
only}. The program goes on. The condition is kept with the thread, and
\code{join-thread} on it signals a \code{<thread-error>} --- a subclass
of \code{<error>} --- in the joining thread. \code{thread-error-thread}
answers the thread that ended, and \code{thread-error-condition} the
condition that ended it:

\begin{dylan}
(define worker (make-thread (method () (error "disk ~a is full" 3))))

(handler-case (join-thread worker)
  ((<thread-error> condition: e)
   (bind ((original (thread-error-condition e)))
     (format #t "worker failed: ~a ~s~%"
             (condition-format-string original)
             (condition-format-arguments original))
     (print (id? (thread-error-thread e) worker)))))
\end{dylan}

\noindent prints

\begin{dylan}
worker failed: disk ~a is full (3)
#t
\end{dylan}

If the \code{<thread-error>} is not handled either, it ends the joining
thread in turn --- in the main thread, the program --- with a message
that carries the original one:

\begin{shell}
Dylan error: a thread ended with an unhandled condition: disk 3 is full
\end{shell}

A thread that is never joined ends silently when a condition is not
handled in it. Warnings are not serious conditions: an unhandled warning
in a thread is reported on the standard error stream, as in the main
thread, and the thread goes on.

\subsection{Locks}\label{sec:ext-locks}

\index{lock@\code{<lock>}}\index{recursive-lock@\code{<recursive-lock>}}
\index{with-lock@\code{with-lock}}
Threads that change shared data must take turns, or their changes will
be lost. A \term{lock} lets one thread at a time hold it. \code{(make
<lock>)} makes one; \code{(acquire-lock \meta{lock})} waits until no other
thread holds it and then holds it; \code{(release-lock \meta{lock})} lets
it go. Usually you will use

\begin{quote}
\code{(with-lock (\meta{lock}) \meta{form} ...)}
\end{quote}

\noindent which holds the lock while the forms run and releases it
however they are left --- by returning, by a non-local exit, or by an
error --- and answers the forms' values. Three threads each adding
100\,000 to a shared counter:

\begin{dylan}
(define counter 0)
(define counter-lock (make <lock>))
(define bump
  (method ()
    (dotimes (i 100000)
      (with-lock (counter-lock)
        (set! counter (+ counter 1))))))
(define threads (list (make-thread bump) (make-thread bump)
                      (make-thread bump)))
(for-each ((t threads)) () (join-thread t))
(print counter)                                  ; prints 300000
\end{dylan}

Without the lock, two threads can read the same old value of
\code{counter} and both store one more than it, and the total comes out
short.

A \code{<lock>} checks its use. Acquiring a \code{<lock>} the thread
already holds would wait forever, so it is an error instead
(``acquire-lock: the thread holds the lock already''), and so is
releasing a lock the thread does not hold (``release-lock: the thread
does not hold the lock''). A \code{<recursive-lock>} --- a subclass of
\code{<lock>} --- may be acquired again by the thread holding it, and is
released when it has been released as many times:

\begin{dylan}
(define r (make <recursive-lock>))
(print (with-lock (r) (with-lock (r) 'nested)))   ; prints nested
\end{dylan}

Neither kind of lock takes init-keywords.

\subsection{Condition variables}\label{sec:ext-condition-variables}

\index{condition-variable@\code{<condition-variable>}}
A thread that must wait until another thread has done something --- put
an item in a queue, say --- waits on a \term{condition variable}.
\code{(wait-condition-variable \meta{cv} \meta{lock})}, called while
holding \meta{lock}, which must be a \code{<lock>} (not a
\code{<recursive-lock>}), releases the lock and waits; when another
thread calls \code{(notify-condition-variable \meta{cv})}, one waiting
thread wakes, acquires the lock again, and returns.
\code{notify-all-condition-variable} wakes every waiting thread.

A wait can also end without a notification, and another thread may get
to the data first. So always wait in a loop that tests what is awaited,
with the lock held. Here a consumer thread takes ten numbers from a
queue that the main thread fills:

\begin{dylan}
(define lock (make <lock>))
(define not-empty (make <condition-variable>))
(define queue (make <deque>))

(define put!
  (method (item)
    (with-lock (lock)
      (push-last queue item)
      (notify-condition-variable not-empty))))

(define take!
  (method ()
    (with-lock (lock)
      (while (empty? queue)
        (wait-condition-variable not-empty lock))
      (pop queue))))

(define consumer
  (make-thread
    (method ()
      (for ((n 0 (+ n 1))
            (total 0 (+ total (take!))))
           ((= n 10) total)))))

(dotimes (i 10) (put! (* i i)))
(print (join-thread consumer))                  ; prints 285
\end{dylan}

However the two threads interleave, the consumer takes the ten squares
and answers their sum.

\subsection{Atomic integers}\label{sec:ext-atomic}

\index{atomic-integer@\code{<atomic-integer>}}
For a counter or a flag, a lock is more than is needed. An
\code{<atomic-integer>} holds an integer that threads can update without
one. \code{(make <atomic-integer> value: \meta{n})} makes one holding
\meta{n} (0 if \code{value:} is not given); \code{atomic-value} reads it;
\code{(atomic-add! \meta{a} \meta{n})} adds \meta{n} and answers the new
value, as one indivisible step; and \code{(atomic-compare-and-swap!
\meta{a} \meta{old} \meta{new})} stores \meta{new} only if \meta{a}
still holds \meta{old}, answering \code{\#t} if it did and \code{\#f} if
not.

\begin{dylan}
(define hits (make <atomic-integer>))
(define bump-atomically
  (method () (dotimes (i 100000) (atomic-add! hits 1))))
(do join-thread (list (make-thread bump-atomically)
                      (make-thread bump-atomically)))
(print (atomic-value hits))                          ; prints 200000
(print (atomic-compare-and-swap! hits 200000 0))     ; prints #t
(print (atomic-compare-and-swap! hits 200000 0))     ; prints #f
(print (atomic-add! hits 5))                         ; prints 5
\end{dylan}

The values must be integers, and an addition that overflows the 63-bit
range signals an error, as ordinary arithmetic does.

\subsection{Thread variables and \code{dynamic-bind}}
\label{sec:ext-thread-variables}

\index{define-thread-variable@\code{define-thread-variable}}
\index{dynamic-bind@\code{dynamic-bind}}
\begin{quote}
\code{(define-thread-variable \meta{name} \meta{value})}
\end{quote}

\noindent at top level defines a module variable that has a separate
value in each thread. Reading it, and assigning it with \code{set!},
affect only the current thread's value. The main thread starts with
\meta{value}, and so does every thread made after the definition has run
--- a new thread does \emph{not} inherit its creator's current value.

\begin{quote}
\code{(dynamic-bind ((\meta{name} \meta{value}) ...) \meta{form} ...)}
\end{quote}

\noindent gives each thread variable \meta{name} the new \meta{value} in
the current thread while the forms run, and restores the old value
however they are left. This is how such variables are normally used: as
settings that a computation and everything it calls can see, without
passing them as arguments, and without one thread's settings disturbing
another's.

\begin{dylan}
(define-thread-variable *indent* 0)

(define show
  (method (text)
    (format #t "~a~a~%" (make <string> size: *indent* fill: #\space) text)))

(define nested
  (method (depth)
    (show depth)
    (when (< depth 2)
      (dynamic-bind ((*indent* (+ *indent* 2)))
        (nested (+ depth 1))))))

(nested 0)
(print *indent*)
(print (join-thread (make-thread (method () (set! *indent* 99) *indent*))))
(print *indent*)
\end{dylan}

\noindent prints

\begin{dylan}
0
  1
    2
0
99
0
\end{dylan}

The other thread's \code{set!} changed only its own value.

\subsection{What threads may share}\label{sec:ext-thread-sharing}

Assigning a module variable or a slot is a single store of one machine
word, so a thread never sees a value half written. Beyond that, data
that threads share and change must be protected with locks. In
particular the library's collections --- tables, stretchy vectors,
deques and the rest --- do no locking of their own: two threads adding to
one table at once can leave it in a wrong state. The runtime's own shared state, such as the table of
symbols, is locked internally, and since the program's classes and
generic functions never change, method dispatch needs no locks at all.

\fussy

\chapter{Inside the Compiler}\label{ch:internals}
\sloppy

This chapter describes how the compiler to C works: the stages a program
passes through, how Dylan values look in memory, how functions are
called and how tail calls are kept proper, how generic functions
dispatch, how exits, cleanups and conditions are built, and how the
garbage collector and threads fit in. You do not need any of it to use
the compiler. It helps when you want to know why one way of writing a
program is faster than another (Chapter~\ref{ch:performance}), when you
read the C the compiler writes, or when you want to work on the compiler
itself. The design, phase by phase, is also recorded in
\file{CompilerPlan.md} at the top of the source tree.

\section{The stages}\label{sec:internals-pipeline}

\index{compiler!stages}
The compiler is about 9\,000 lines of Common Lisp in \file{compiler/}.
It is a \term{whole-program compiler}: it sees the entire program, with
the library, before it writes any C, and since nothing changes once the
program runs, everything it learns about the program is true for good.
A build passes through these stages:

\begin{enumerate}
\item \textbf{Reading} (\file{front/source.lisp}). Each file is read
  with the reader of the interactive implementation
  (\file{src/compat.lisp}), the only code the two implementations share
  besides the macros' run time. The reader records the file, line and
  column where every list begins, so that every later diagnostic, and
  every \code{\#line} directive, can name a place in the source. A file's
  \code{module:} header (design note~33) says which module its forms
  belong to.

\item \textbf{Modules} (\file{front/modules.lisp}). Every
  \code{define-module} form in the program is found first, and the
  modules are defined, each after the modules it uses. Then every name
  in the program is resolved to a \emph{module variable}: one object per
  variable, owned by one module, however many modules import it and under
  whatever names. From here on names are gone, and nothing can be
  confused by renaming or shadowing.

\item \textbf{Macros and core forms} (\file{front/convert.lisp},
  \file{front/macros.lisp}). Each top-level form is converted to
  \emph{core Dylan}, about a dozen kinds of node: constants, variable
  references and assignments, \code{if}, sequences, \code{bind} (with
  multiple values), mutually recursive local methods, methods, calls,
  calls of runtime primitives, \code{bind-exit} and
  \code{unwind-protect}, and the definitions. Every other special form
  and built-in macro --- \code{cond}, \code{case}, \code{for},
  \code{for-each}, \code{handler-case}, backquote, \code{with-lock} and
  the rest --- is rewritten into these. A loop becomes a local method
  that calls itself in tail position. The program's own macros are
  expanded here too, by running their bodies in the interactive
  implementation (Chapter~\ref{ch:extensions}).

\item \textbf{The IR} (\file{ir/anf.lisp}, \file{ir/ir.lisp}). Core
  Dylan is converted to an intermediate representation in
  \term{A-normal form}: every intermediate value is given a name, and
  every operand of a call or primitive is a constant or a variable.
  Every variable is bound exactly once (the IR is \emph{alpha-converted}),
  so a variable is a unique object that analyses can annotate. Operands
  are evaluated left to right, as in the source.

\item \textbf{Analysis} (\file{ir/analyze.lisp}). For each variable:
  whether it is assigned, whether a function other than the one binding
  it uses it (it is \emph{captured}), and how often it is referred to.
  For each local function: where it is called, and whether it
  \emph{escapes} --- is used other than by being called. For each exit
  procedure of a \code{bind-exit}: whether it escapes its function. And
  for each call: whether it is in \emph{tail position}, that is, whether
  its value is the value of its function.

\item \textbf{Control lowering} (\file{ir/control.lisp}). A
  \code{bind-exit} whose exit is only called directly within its own
  function stays, and will become a jump. Every other \code{bind-exit},
  and every \code{unwind-protect}, is rewritten as a call of a runtime
  function, with the body moved into a new closure
  (Section~\ref{sec:internals-exits}).

\item \textbf{Contification} (\file{ir/contify.lisp}). A local function
  that never escapes and whose every call is a jump --- a call in tail
  position from the function that defines it, or from another such
  function --- is not compiled as a function at all, but as a labelled
  block of the function that contains it, its calls becoming
  assignments to its parameters and a \code{goto}. This is how loops
  become C loops.

\item \textbf{Type inference} (\file{backend/types.lisp},
  Section~\ref{sec:internals-types}) finds a class for as many
  variables as it can, and decides which \code{<double-float>} variables
  can be plain C \code{double}s.

\item \textbf{Calls} (\file{backend/calls.lisp},
  Section~\ref{sec:internals-calls}). Each call whose function is known
  is resolved: to a method, to a generic function's dispatch function,
  to a runtime primitive, or to inline C.

\item \textbf{Escape analysis} (in \file{backend/codegen.lisp}) finds
  \code{\#rest} lists, closures and instances that cannot outlive the
  function that makes them, and puts them on the C stack instead of the
  heap.

\item \textbf{Code generation} (\file{backend/codegen.lisp},
  \file{backend/objects.lisp}, \file{backend/cdata.lisp}). Every
  remaining function becomes a C function; every class, symbol and
  constant becomes static C data; every generic function gets a dispatch
  function.

\item \textbf{Output} (\file{backend/output.lisp}). The C is split into
  \file{\meta{name}.c}, one file per module, and \file{\meta{name}.h}
  (Section~\ref{sec:dylanc-files}), and \file{driver.lisp} runs the C
  compiler.
\end{enumerate}

When the library is compiled with the program (\cmd{-Os},
\cmd{-{}-whole}), tree shaking (\file{backend/shake.lisp},
Section~\ref{sec:internals-shaking}) runs between control lowering and
code generation.

The library itself is Dylan, in module \code{dylan}, in the files of
\file{runtime/dylan/} read in name order: conditions, numbers,
characters, collections, sequences, lists, vectors, strings, stretchy
vectors, deques, ranges, tables, arrays, printing and threads. It is
compiled with every program, by the same stages, on top of the
\code{\%} primitives written in C in \file{runtime/c/library.c} and the
other files of \file{runtime/c/}.

\section{How values are represented}\label{sec:internals-representation}

\index{representation of values}\index{tagged words}
Every Dylan value is one 64-bit machine word, the C type \code{D}
(\file{runtime/c/dylan.h}). The low bits of the word say what it is:

\begin{center}
\begin{tabular}{ll}
\toprule
Low bits & Value \\
\midrule
\code{...xx1} & an \code{<integer>}: the value shifted left one bit, so
  63 bits signed \\
\code{...010} & a \code{<character>}: the Unicode code point shifted left
  three bits \\
\code{...000} & a pointer to an object, 8-byte aligned \\
\bottomrule
\end{tabular}
\end{center}

Integers and characters are \term{immediate}: they need no memory of
their own, and allocating one costs nothing. The integer $n$ is the word
$2n+1$; that is why a debugger shows the integer 100 as 201. Adding two
tagged integers $2a+1$ and $2b+1$ is one machine addition and a
subtraction of 1, with the overflow checked by the C compiler's
\code{\_\_builtin\_add\_overflow}; an overflow signals an error rather
than wrapping.

Every object a pointer points to begins with a \term{header} word: a
pointer to its class's \term{wrapper}. The wrapper holds what the
runtime needs to know about the class --- the class object, the class's
numeric id, the instance size, how instances are represented (an ordinary
instance, a pair, a byte string, a double float, and so on), and the set
of the class's ancestors. Pairs are a header and two words; strings a
header, a length and the bytes; a \code{<double-float>} a header and a C
\code{double}; an instance of a class the program defines a header and
one word per slot.

\code{\#t}, \code{\#f} and \code{()} are three static objects, compared
by address. \code{\#t} and \code{\#f} are instances of an internal class
\code{<boolean>}, which the manual does not name and the \code{dylan}
module does not export.

\paragraph{Class ids.} \index{class ids}
Since the program is static, the compiler knows every class. It numbers
them in a preorder walk of the tree that each class makes with its
\emph{first} superclass, so that a class's descendants in that tree have
the ids from its own up to a limit. Where every subclass of a class lies
in that range, \code{(instance? x C)} is a range check on \code{x}'s
class id: one subtraction and one comparison. Where multiple inheritance
puts a subclass outside it, the test reads a bit in the ancestor set
every wrapper carries. The classes, their wrappers and the symbols are
all static C data emitted by the compiler, so a program does no work at
startup to build its class hierarchy.

\paragraph{Memory.}
Objects are allocated by the garbage collector
(Section~\ref{sec:internals-gc}), through the functions
\code{dy\_alloc} and its kin in \file{dylan.h}, which hide which
collector it is: those that can hold pointers apart from those that
cannot --- strings, float boxes, ratios --- so that the collector does
not scan the latter.

\section{The calling convention}\label{sec:internals-calls-convention}

\index{calling convention}\index{uniform entry}
Every Dylan function has a \term{uniform entry}, a C function of one
prototype:

\begin{ccode}
D f(D self, D a0, D a1, D a2, D a3, intptr_t argc);
\end{ccode}

\code{self} is the function object being called --- a closure, a
generic function, or a method --- through which a closure reaches its
free variables. The first four arguments travel in \code{a0} to
\code{a3}, in machine registers; \code{argc} says how many there are.
Arguments after the fourth go in a per-thread \term{argument area}, which
the callee reads before it calls anything else. Every function checks
\code{argc} on entry, so calling a function with the wrong number of
arguments is an error, not a crash.

One prototype for every function is what makes proper tail calls
possible (Section~\ref{sec:internals-tail}). It costs a little: a
function of two arguments is still passed six. So a function that is
called from places the compiler can see, other than in tail position,
also gets a \term{specialized entry}, named with \code{\_spec}, with
exactly its own parameters: \code{double}s where the arguments are
unboxed floats, and one parameter per keyword, filled in when compiling
from the keywords the call passes (\code{DY\_UNSUPPLIED} for those it
does not). A call through a specialized entry passes no count, checks
no arity and parses no keywords when it runs.

\paragraph{Multiple values.}
\index{multiple values}
A function returns its first value as its C return value. When it
returns some other number of values, it stores the count and all the
values in a per-thread values buffer. A caller that wants all the values
sets the count to 1 before the call; a function that returns one value
leaves the buffer alone unless it has itself made a call that may have
changed the count. A caller that wants only the first value --- the
common case --- reads the C return value and ignores the buffer.

\paragraph{\code{next-method}.}
A method that calls \code{next-method} receives its next method through
the thread: the dispatch function (or the method before it) stores it
there, and the method reads it on entry. The chain of methods for each
combination of argument classes is known when compiling, so the
next-method objects are static data, one for each tail of each chain.
Methods that do not use \code{next-method} ignore it, and a method
whose \code{next-method} is only called, with no arguments, reuses its
own arguments rather than capturing them in a closure.

\section{Proper tail calls}\label{sec:internals-tail}

\index{tail calls}\index{musttail}
Dylan requires that a call in tail position not use stack: a loop
written as a recursion must run in constant space, however many times it
goes round. C promises nothing of the kind. The compiler meets the
requirement in three ways.

\begin{enumerate}
\item A function's tail call of itself, and calls among local methods
  that have been contified, become \code{goto}s within one C function.
  This covers almost every loop: \code{for}, \code{while},
  \code{dotimes} and \code{for-each} all expand into local methods
  calling themselves in tail position.
\item Every other tail call --- through a module variable, a closure
  passed as an argument, a generic function, \code{next-method},
  \code{apply} --- is a C tail call marked
  \code{\_\_attribute\_\_((musttail))}, which GCC~15 and later and
  Clang~13 and later \emph{guarantee} to compile as a jump. The
  attribute requires caller and callee to have the same prototype, which
  the uniform entry provides. The generated C writes it with the macro
  \code{DY\_TAIL\_CALL}.
\item With \cmd{-{}-trampoline}, a tail call instead stores the function
  and its arguments in the thread and returns a marker; the nearest
  caller that is not making a tail call runs a loop that makes the
  pending call until a real value comes back. Any C11 compiler can
  compile this.
\end{enumerate}

A specialized entry cannot make a \code{musttail} call, because its
prototype differs from its callee's. When a function makes tail calls,
either its uniform entry holds the only body, or its specialized entry
gets a second copy of the body whose tail calls are ordinary calls. That
is safe: a specialized entry is only ever entered by a call that is not
a tail call, so each holds at most one extra C frame.

Calls inside an \code{unwind-protect}, inside a \code{bind-exit} whose
exit escapes, and inside a handler's extent are not in tail position:
their frame must remain until they return, as the language's semantics
already imply.

\section{Closures}\label{sec:internals-closures}

\index{closures}
Closures are \emph{flat}: a closure is a function object holding its C
entry and a copy of the value of each free variable it uses. A variable
that is both captured by a closure and assigned lives instead in a
one-word heap \term{cell}, which every closure that uses it shares, so
all of them see each assignment. A function with no free variables ---
every method defined at top level, for instance --- is a static object
and costs nothing to make. A local method that is only called, and never
escapes, is not made into a closure at all: it is contified or called
directly.

\section{Generic dispatch}\label{sec:internals-dispatch}

\index{generic functions!dispatch}\index{dispatch}
Because the program is static, the compiler knows every method of every
generic function. For each generic function it computes, at compile
time, which methods apply to each combination of argument classes and
singletons, sorts each such set by specificity, and writes a C
\term{dispatch function}: a decision tree that tests the arguments ---
singletons first, by comparing pointers, then a \code{switch} on the
class id --- and ends at a leaf for each distinct sorted list of
methods. A leaf checks keyword arguments, if the methods take any, sets
the thread's next method, and tail-calls the first method. A call that
no method fits reaches \code{dy\_no\_applicable\_method}, which signals
the error.

Here is the dispatch function for a generic function \code{area} with
one method on \code{<square>} and one on \code{<circle>}:

\begin{ccode}
D dy_gf_area_51a90f52_dispatch(D self, D a0, D a1, D a2, D a3,
                               intptr_t argc) {
    (void)a0; (void)a1; (void)a2; (void)a3;
    if (DY_UNLIKELY(argc != 1)) dy_wrong_argument_count(self, argc);
    D x0 = a0; (void)x0;
    switch (dy_wrapper_of(x0)->id) {
    case 61:
        goto leaf_0;
    case 62:
        goto leaf_1;
    default: goto no_method;
    }
leaf_0:
    dy_thread->next_method = DY_FALSE;
    DY_TAIL_CALL_DIRECT(fn1197_area, DY_OBJ(&dy_method_326),
                        a0, a1, a2, a3, argc);
leaf_1:
    dy_thread->next_method = DY_FALSE;
    DY_TAIL_CALL_DIRECT(fn1199_area, DY_OBJ(&dy_method_327),
                        a0, a1, a2, a3, argc);
no_method:
    dy_no_applicable_method(self);
}
\end{ccode}

\noindent (Three long lines have been broken to fit the page.)

There are no caches and no locks: the tree is fixed, and threads share
it freely.

Methods are ordered by the rule of the manual's chapter~10: arguments
compared left to right, a singleton before any class, a class before its
superclasses. The manual leaves open the order of two classes neither of
which is a subclass of the other; this compiler, like CLOS and like the
interactive implementation, puts first the one that comes earlier in the
\emph{argument's} class precedence list. Two methods that remain
unordered are ambiguous, and calling them is an error.

Better still is not to dispatch at all. Where type inference shows that
only one method can apply, the call goes straight to that method, or the
method is done inline (Section~\ref{sec:internals-calls}). A generic
function with a single method, specialized on \code{<object>} for every
argument, is never dispatched.

\paragraph{Classes and \code{make}.}
For each class that is instantiated, the compiler writes a
\term{maker}, which allocates the instance, fills each slot from its
init-keyword or its default, checks the init-keywords, and calls
\code{initialize} only if the program defines methods on it. A
\code{make} of a class the program defines, where only the library's
\code{make} method on \code{<class>} applies, calls the maker directly.
Slot offsets are fixed where every possible class of the receiver agrees
on them; otherwise accessors look the offset up in a small table in the
wrapper.

\section{Exits, cleanups and conditions}\label{sec:internals-exits}

\index{bind-exit@\code{bind-exit}!implementation}
\index{unwind-protect@\code{unwind-protect}!implementation}
A \code{bind-exit} whose exit procedure is only called directly, within
the function that binds it (or within local methods contified into it),
compiles to a jump: calling the exit assigns the values and does a
\code{goto} past the body. This is the common case --- an early return
from a loop --- and it costs nothing.

Every other \code{bind-exit}, and every \code{unwind-protect}, needs a
\term{frame} in the thread's \term{dynamic chain}, the list of active
exits, cleanups and handlers, innermost first. The frames live on the C
stack of three runtime functions in \file{runtime/c/control.c}:
\code{dy\_call\_with\_exit}, \code{dy\_unwind\_protect} and
\code{dy\_handler\_bind}. The compiler moves the body into a closure and
passes it to one of these:

\begin{dylan}
(unwind-protect p c)  =>  (%primitive unwind-protect (method () p)
                                                     (method () c))
(bind-exit (x) b)     =>  (%primitive call-with-exit (method (x) b))
\end{dylan}

\code{dy\_call\_with\_exit} pushes an exit frame, calls \code{\_setjmp},
and calls the body with an exit procedure. Calling that exit procedure
stores its arguments as the values to return, runs the cleanups of every
frame between it and the target, innermost first, and \code{longjmp}s to
the target frame. An exit called after its \code{bind-exit} has returned
is detected --- each frame carries a serial number --- and signalled as
an error. Because only these runtime functions call \code{setjmp},
generated code never does, and can make tail calls everywhere.

\paragraph{Conditions.}
\index{conditions!implementation}
Handlers are frames in the same chain. \code{signal} walks the chain
for a handler whose type matches and whose test, if any, is true, and
calls it with the condition and a \code{next-handler} that continues the
search further out; if none applies it calls \code{default-handler}. The
condition classes and \code{signal}, \code{error}, \code{cerror},
\code{check-type}, \code{abort}, \code{break} and
\code{default-handler} are written in Dylan, in
\file{runtime/dylan/01-conditions.dyl}; only the frame operations and
the transfer of control are C. When the runtime detects an error of its
own --- a wrong argument count, a type error, an overflow --- it calls
back into that Dylan code, so the error is an ordinary
\code{<simple-error>} or \code{<type-error>} that handlers can catch.

\section{The garbage collector}\label{sec:internals-gc}

\index{garbage collection}\index{Memory Pool System}
Memory is managed by the Memory Pool System in \file{runtime/mps/},
compiled as one C file into the runtime library, which is built with
\code{DY\_MPS}. It is a moving collector, so it must know every object's
layout: \file{runtime/c/gc-mps.c} gives it an \emph{object format} ---
for each kind of object, its size, the words of it that are references,
and how to leave a forwarding pointer or padding where an object was ---
and the \emph{roots}: the program's static data (found from the linker's
symbols on Linux, and from the Mach-O segments on macOS) and each
thread's stack and registers, scanned conservatively, and each thread's
state and thread variables, allocated outside the heap and registered
as a root. A tagged integer or character on a stack that happens to look
like a pointer into the heap can only keep an object alive, and in
place, longer than necessary; it can do no harm.

Objects go in five \emph{pools}: AMC, generational and moving, for
objects that may hold references; AMCZ for those that do not (strings,
floats, ratios), which the collector never scans; AMS, which does not
move, for \code{<thread>} objects, since a thread's state points to its
own; LO, which is neither moved nor protected, for objects holding a
mutex or condition variable the kernel uses; and AWL, weak, for the
table of identity hashes. Each thread has its own \emph{allocation
points} in AMC and AMCZ, from which it allocates by bumping a pointer;
the slow path, when the point's buffer is empty, enters the collector.
An object's fields are filled in before it is committed to the
collector, which may move it or ask for it to be allocated again.
Because the collector moves objects, a Dylan object's identity hash
cannot be its address: \code{\%object-hash} gives each object a number
the first time it is asked, kept in a table with weak keys, rebuilt when
the collector reports that keys have moved. The collector's behavior can
be tuned with environment variables (Chapter~\ref{ch:performance}).

\paragraph{The Boehm-Demers-Weiser collector.} \index{Boehm collector}
With \cmd{dylanc -{}-gc boehm} (Section~\ref{sec:dylanc-gc}) the
runtime is built without \code{DY\_MPS} and memory is managed instead by
\file{libgc}, built for threads. It is a \emph{conservative} collector:
it treats any word on a stack, in a register or in a scanned object that
looks like a pointer into the heap as one, so it needs no object format.
Strings and float boxes are allocated atomic so that their contents are
not mistaken for pointers. Each thread's state --- the values buffer, the
argument area, the dynamic chain --- is allocated uncollectable, so the
collector scans it and sees pointers held there. The collector does
mark-and-sweep collection, in parallel on two threads, and allocates
small objects from thread-local free lists without locking.

The functions \code{dy\_alloc} and \code{dy\_alloc\_init} in
\file{dylan.h} allocate with either collector.

\section{Threads}\label{sec:internals-threads}

\index{threads!implementation}
Each thread registers itself and its stack with the collector as it
starts, and makes its allocation points (under the Boehm collector,
\file{dylan.h} defines \code{GC\_THREADS} and includes \file{gc.h}
before \file{pthread.h}, so \code{pthread\_create} goes through the
collector, which registers each thread's stack). Each thread has one
\code{\_Thread\_local} pointer, \code{dy\_thread}, to its state: the
values buffer, the argument area, the dynamic chain, its next method,
its pending tail call in trampoline mode, its \code{<thread>} object,
and its values of the program's thread variables.

A thread starts by setting up that state and an exit frame at the top of
its dynamic chain, then calls its function. An unhandled serious
condition reaches \code{default-handler}, which in a thread other than
the main one records the condition in the \code{<thread>} and unwinds to
that top frame, ending the thread; \code{join-thread}, written in Dylan
in \file{runtime/dylan/20-threads.dyl}, finds it there and signals a
\code{<thread-error>}. Threads are detached when made; \code{join-thread}
waits on a condition variable in the \code{<thread>}. A \code{<lock>}
is an error-checking POSIX mutex, a \code{<recursive-lock>} a recursive
one, and an \code{<atomic-integer>} a C11 \code{\_Atomic} integer.

\section{Separate compilation of the library}\label{sec:internals-separate}

\index{separate compilation}\index{library cache}
Compiling the library's megabyte of C with every program would double
the time of a small build. So, by default, the compiler compiles the
library \emph{as if it did not know the program}. Every function the
\code{dylan} module exports, or that the compiler's expansions name, may
be called from anywhere with anything: no argument types and no
specialized entries are derived from the program's calls, and nothing of
the library is dropped. The C names of the library's constants,
variables and generic functions are derived from what they are, not
from the order in which the compiler met them, and its C temporaries are
numbered within each function. \file{\meta{name}-dylan.c} declares only
what it refers to. The result is that the library's C is the same text
for almost every program, and the driver compiles it once per C compiler
and option set and caches the object (Section~\ref{sec:dylanc-cache}).

What is lost is specialization of library code to the program, and the
dropping of unused library code. Calls \emph{from} the program
\emph{into} the library lose little, since the compiler still inlines
the library's arithmetic, \code{element}, \code{size} and iteration on
known collections at the call site.

\section{Tree shaking}\label{sec:internals-shaking}

\index{tree shaking}
With \cmd{-Os} or \cmd{-{}-whole} the library is compiled with the
program, and \file{backend/shake.lisp} keeps only what the program can
reach. The program's own top-level forms are all kept, as are the
library's top-level forms that are not definitions and every
\code{define-class} (classes are cheap static data). A library
definition of a variable, generic function or method is kept when
something kept refers to its variable. A method is kept, besides, only
if each class its specializers name could have instances --- a class
the runtime itself makes instances of, or one that something kept uses
as a value (makes, say, or answers). Naming a class as a specializer or
a superclass does not count.

Shaking is by name and class, not by data flow, and so it keeps more
than a program needs: a call of a generic function such as
\code{print-object}, \code{make} or \code{as} keeps every method that
could apply to a class with instances. That is why hello world is still
about 220~KB when built \cmd{-Os}.

\section{Type inference and inlining}\label{sec:internals-types}
\label{sec:internals-calls}

\index{type inference}
The type of a variable, for the compiler, is a class all its values are
instances of, or nothing known. Types come from:

\begin{itemize}
\item constants;
\item the result classes of the operations done inline --- an integer
  addition answers an integer, a float multiplication a
  \code{<double-float>};
\item \code{make} of a class whose maker the call goes to;
\item the getters of slots declared with a \code{type:};
\item the type a \code{bind} declares, as in
  \code{(bind (((x <integer>) \meta{init})) ...)}, which is checked when
  the binding is made;
\item the specializers of a method defined with \code{define-method},
  which dispatch has already checked;
\item for local functions, global functions and sole methods whose every
  call is known --- a loop, most often --- the classes of the arguments
  at those calls, and the classes of what the functions return.
\end{itemize}

These are found together, starting from nothing, until none changes.
The types are also \emph{flow-sensitive} where a test proves more: in
the branch an \code{(instance? x C)} test chooses, \code{x} is known to
be a \code{C}.

Calls whose function is known are then resolved (\file{calls.lisp}). A
call of \code{+}, \code{-}, \code{*}, \code{/}, \code{<}, \code{>} or
\code{=} with two arguments calls the binary generic function
(\code{binary+} and so on) directly. The library's methods for integer
and float arithmetic and comparison, for \code{size}, \code{element}
and the iteration protocol on lists, simple object vectors and byte
strings, and \code{id?}, \code{not}, \code{car}, \code{cdr},
\code{instance?} and slot getters and setters are done \emph{inline} ---
when the static program shows that no other method could apply to such
arguments. The manual's generic functions that the runtime writes in C
--- \code{instance?}, \code{object-class}, \code{subclass?} and the
other reflective functions of \dylanbook{chapters 10 and 11} --- each
have one library method, which calls a \code{\%} primitive; while the
program adds no method of its own to one, its calls are calls of the
primitive, and \code{instance?} of a known class is done inline. Each inline case is guarded by a test of the arguments' classes,
and the test is left out where their types are known. A float
operation with one argument known to be a \code{<double-float>} answers a
\code{<double-float>} (float contagion), so its result's type is known
even if the other argument's is not.

\paragraph{Unboxing.} \index{unboxing}
A \code{<double-float>} variable that no closure holds, and every use of
which wants a C \code{double}, is \emph{unboxed}: it is a C
\code{double} local, and computing it allocates nothing. A float
variable that is returned, stored in a slot or passed to an unknown
function needs its box, and is allocated on the heap each time it
changes. Integers need no unboxing; they are immediate.

\paragraph{Escape analysis.} \index{escape analysis}
A \code{\#rest} list that is only walked, through variables no closure
holds, is built on the C stack (up to eight pairs); keyword arguments
are parsed where they were passed, with no list at all. A closure, or an
instance whose class's maker \code{make} calls (and which has no
\code{initialize} methods), is made on the stack of the function that
makes it when nothing can reach it after that function returns: it is
only tested, called, given to a slot accessor sure to apply, or passed to
a known function whose parameter is used the same way. Constants
propagate through \code{bind}s, integer arithmetic on constants is done
when compiling, an \code{if} on a constant keeps one branch, and unused
pure bindings are dropped.

In a debug build (\cmd{-g}) none of the inferred types is used, so every
inline operation tests its arguments.

\section{Reading the generated C}\label{sec:internals-reading}

\index{C files!reading}
\cmd{dylanc -k} keeps the C. Here is a small program, \file{sum.dyl}:

\begin{dylan}
(define-method sum-to ((n <integer>))
  (for ((i 0 (+ i 1))
        (s 0 (+ s i)))
       ((> i n) s)))
(print (sum-to 100))
\end{dylan}

\noindent and the C function the compiler writes for \code{sum-to}, in
\file{sum-dylan-user.c}:

\begin{ccode}
D fn1190_sum_to(D self, D a0, D a1, D a2, D a3, intptr_t argc) {
    (void)self; (void)a0; (void)a1; (void)a2; (void)a3; (void)argc;
    D v105609, v105610, v105612, v105613, v105614, v105615, v105616,
      v105617, v105618, v105619;
    D t1;
    D t2;
    D t3;
    D t4;
#line 1 "sum.dyl"
    if (DY_UNLIKELY(argc != 1)) dy_wrong_argument_count(self, argc);
    v105609 = a0;
#line 2 "sum.dyl"
    t1 = DY_INT(0);
    t2 = DY_INT(0);
    v105612 = t1;
    v105613 = t2;
    goto for_1;
    for_1:;
#line 4 "sum.dyl"
    v105615 = DY_BOOL((intptr_t)(v105609) < (intptr_t)(v105612));
#line 2 "sum.dyl"
    if (v105615 != DY_FALSE) {
        dy_thread->mv_count = 1;
        return v105613;
    } else {
        v105617 = dy_int_add(v105612, DY_INT(1));
#line 3 "sum.dyl"
        v105619 = dy_int_add(v105613, v105612);
#line 2 "sum.dyl"
        t3 = v105617;
        t4 = v105619;
        v105612 = t3;
        v105613 = t4;
        goto for_1;
    }
}
\end{ccode}

\noindent (one line of declarations has been broken to fit the page).
Reading it from the top:

\begin{itemize}
\item The function has the uniform entry. Its name, \code{fn1190\_sum\_to},
  is a number and the Dylan name; every variable \code{v}\meta{n} is
  one IR variable, and every \code{t}\meta{n} a temporary.
\item It checks that it was given one argument; the method is the only
  method of \code{sum-to}, and its dispatch function has already checked
  that the argument is an \code{<integer>}.
\item The \code{for} loop has become a label, \code{for\_1}, and a
  \code{goto}. The step forms are evaluated into temporaries before any
  loop variable is assigned, because \code{for} steps its variables in
  parallel.
\item \code{(> i n)} is a single C comparison of the tagged words, with
  no test of their classes: \code{n} is an integer by its specializer,
  and \code{i} by type inference, since it starts at 0 and is only ever
  incremented. \code{(> i n)} became \code{(< n i)}, and on tagged
  integers $2a+1 < 2b+1$ exactly when $a<b$.
\item \code{(+ i 1)} and \code{(+ s i)} are \code{dy\_int\_add},
  an inline function in \file{dylan.h} that adds the tagged words and
  signals an error on overflow.
\item Returning one value sets the values count to 1 first, because a
  value count might be left over from a call made earlier in the
  function.
\item The \code{\#line} directives name the Dylan line each statement came
  from.
\end{itemize}

If \code{sum-to} were instead written as \code{(define sum-to (method
((n <integer>)) ...))}, \code{n} would not be known to be an integer
inside it: a method made by \code{method} and stored in a variable is
called through the variable, and must check its own argument. The
comparison would then be preceded by a test, \code{if
(DY\_IS\_INT(v105609))}, with a call of the generic function
\code{binary<} for the other case. Small differences like this are the
subject of the next chapter.

\fussy

\chapter{Performance}\label{ch:performance}
\sloppy

Dylan was designed to be both dynamic and efficient, and the compiler to
C aims at its original targets: small and fast. This chapter reports how
fast compiled programs are on a set of benchmarks, explains what makes a
Dylan program fast or slow under each implementation, and shows how to
measure a program and tune the garbage collector. Most of the advice
follows from how the compiler works (Chapter~\ref{ch:internals}); the
measurements here show how much each point matters.

\begin{notebox}
Every time in this chapter was measured on one machine: an Intel Core
i9-14900KF running Fedora Linux, with SBCL~2.6.3 and GCC~16.2.1. Times
depend on the processor, the C compiler, the collector's version and the
load on the machine, and vary by several percent from run to run.
Treat them as an indication of proportions, not as absolute figures, and
run the benchmarks on your own machine for numbers that apply to you.
\end{notebox}

\section{The benchmarks}\label{sec:perf-benchmarks}

\index{benchmarks}
The directory \file{bench/} holds eleven small programs, each written
three times: in Dylan (\file{benchmarks.dyl}), in plain Common Lisp
(\file{reference.lisp}) and in plain C (\file{reference.c}). Each answers
a value, and the three versions must agree.

\begin{center}
\small
\begin{tabular}{>{\raggedright\arraybackslash}p{0.17\textwidth}%
                >{\raggedright\arraybackslash}p{0.72\textwidth}}
\toprule
Name & What it does \\
\midrule
\code{table} & insert into and look up in a \code{<table>} \\
\code{sort} & sort a list of pseudo-random integers \\
\code{string} & build and compare strings \\
\code{deque} & push and pop at both ends of a \code{<deque>} \\
\code{stretchy} & grow a \code{<stretchy-vector>} and walk it \\
\code{range} & iterate over ranges \\
\code{exit} & leave a computation early with \code{bind-exit} \\
\code{condition} & signal errors and handle them with \code{handler-case} \\
\code{next-method} & call down a chain of methods with \code{next-method} \\
\code{tak} & the Takeuchi function: integer arithmetic and calls \\
\code{float} & points of a grid inside the Mandelbrot set: floating
  point \\
\bottomrule
\end{tabular}
\end{center}

Two scripts run them. Each prints, for every benchmark, the best of
three runs.

\begin{shell}
$ sbcl --script bench/run-c.lisp          # compiled Dylan vs C and Lisp
$ sbcl --script bench/run.lisp            # interactive Dylan vs Lisp
$ sbcl --script bench/run.lisp table tak  # just these
\end{shell}

\noindent Their builds go in \file{bench/build/}. A line marked
\code{MISMATCH} means the versions answered differently, which would be
a bug.

\index{BENCH\_SCALE@\code{BENCH\_SCALE}}
At their usual sizes several benchmarks finish in a few milliseconds,
too quickly to compare reliably. The environment variable
\code{BENCH\_SCALE}, if set to an integer, multiplies every size, in
both scripts and in \file{reference.c}:

\begin{shell}
$ BENCH_SCALE=10 sbcl --script bench/run-c.lisp
\end{shell}

\noindent The file \file{SpeedComparison.md} at the top of the
repository puts the three ways of running the benchmarks --- plain SBCL,
the Common Lisp implementation and the compiler to C --- side by side,
with hand-written C, at \code{BENCH\_SCALE=10}, and discusses the
results. The tables below were run at the usual sizes.

\subsection{The compiler to C}

\cmd{bench/run-c.lisp} compiles the Dylan with \cmd{dylanc} (at
\cmd{-O2}), compiles \file{reference.c} with \cmd{cc -O2}, and loads
\file{reference.lisp} into SBCL. A run gave:

\begin{shell}
benchmark     Dylan->C ms     C ms      D/C    Lisp ms   D/Lisp
table                  13        6     2.17        8.0     1.63
sort                   29       13     2.23       26.0     1.12
string                 11        2     5.50        7.0     1.57
deque                   6        2     3.00        2.0     3.00
stretchy                3        0        -        2.0     1.50
range                   5        0        -        2.0     2.50
exit                    1        3     0.33        1.0     1.00
condition               4        0        -        1.0     4.00
next-method             5        2     2.50        8.0     0.62
tak                    22       17     1.29       35.0     0.63
float                   3        2     1.50       55.0     0.05
\end{shell}

The columns are the compiled Dylan's time, the C's, their ratio, plain
SBCL's time and the ratio to it. A ratio is shown as \code{-} when the
other time is under a millisecond and the ratio would mean nothing.

Read it this way. On calls and integer arithmetic (\code{tak}) and on
floating point (\code{float}), compiled Dylan takes 1.3 to 1.5 times as
long as hand-written C. On the collections it is two to three times the C's
time, and on strings about five times. Non-local exits (\code{exit}) are
\emph{faster} than the C, because the Dylan exit becomes a jump where
the C uses \code{setjmp}. Against plain SBCL, compiled Dylan is faster
on calls, \code{next-method} and floats, and slower on the collection
benchmarks by up to three times. The very short benchmarks
(\code{stretchy}, \code{range}, \code{condition}) take a few
milliseconds where the C takes under one, so their ratios vary most from
run to run; in other runs the \code{string} ratio was 6.5 and the
\code{table} ratio 1.6.

\subsection{The interactive implementation}

\cmd{bench/run.lisp} compiles \file{benchmarks.dyl} with the interactive
implementation's \code{dylan:compile-file}, compiles the resulting Lisp
with SBCL, and compares it with \file{reference.lisp}:

\begin{shell}
benchmark         size      dylan       lisp    ratio
table           200000     0.076s     0.007s    10.9x
sort            100000     0.083s     0.025s     3.3x
string          200000     0.050s     0.005s    10.0x
deque           300000     0.029s     0.002s    14.5x
stretchy        300000     0.026s     0.002s    13.0x
range          1000000     0.048s     0.003s    16.0x
exit            300000     0.027s     0.001s    27.0x
condition        30000     0.010s     0.000s    10.0x
next-method     300000     0.039s     0.008s     4.9x
tak                300     0.469s     0.036s    13.0x
float              300     0.089s     0.054s     1.6x
\end{shell}

The interactive implementation is typically ten to fifteen times slower
than plain Lisp, and so ten to twenty times slower than the compiler to
C. That is the price of a language whose classes and methods can change
while it runs: every call of a generic function must allow for methods
added since the last call. The interactive implementation is for
developing and exploring; the compiler is for running.

\section{Writing fast Dylan for the compiler}\label{sec:perf-c}

The compiler is fastest where it knows the classes of the values it
works with. Then it calls methods directly, does arithmetic and
collection access inline without tests, and keeps floats unboxed. Each
point below says how to give it that knowledge, with a measurement. The
programs are small, and each was timed as a whole, run as a separate
executable built at \cmd{-O2}.

\subsection{Prefer \code{define-method} to \code{define} of a \code{method}}

How a function is defined matters more than you might expect. These
four versions of \code{tak}, computing \code{(tak 30 22 12)}, differ
only in their first lines:

\begin{center}
\small
\begin{tabular}{lr}
\toprule
Definition & Time \\
\midrule
\code{(define tak (method (x y z) ...))} & 0.04 s \\
\code{(define-method tak (x y z) ...)} & 0.04 s \\
\code{(define-method tak ((x <integer>) (y <integer>) (z <integer>)) ...)}
  & 0.08 s \\
\code{(define tak (method ((x <integer>) (y <integer>) (z <integer>)) ...))}
  & 0.30 s \\
\bottomrule
\end{tabular}
\end{center}

\begin{itemize}
\item An unspecialized \code{define-method} with only one method is not
  dispatched at all, and a \code{define} of a \code{method} with no
  specializers is called directly; both are as fast as the program can
  be.
\item Specializing the only method on \code{<integer>} makes every call
  go through the dispatch function, which tests the argument classes.
  Inside the method the arguments are then known to be integers, which
  helps a method that does much work with them, but here the method is
  tiny and the dispatch doubles its cost.
\item A \code{method} with specializers stored with \code{define} is the
  slowest: it is a closure that checks its own arguments on every call,
  is called through the variable, and does not pass on to the compiler
  the knowledge its specializers give.
\end{itemize}

So: specialize methods when you need dispatch, or when the method does
enough with its arguments to profit from knowing their classes; do not
specialize a lone method merely to document its argument types; and
define functions with \code{define-method}, or with \code{define} of an
unspecialized \code{method}.

\subsection{Let the compiler know classes: declarations in \code{bind}}

\index{bind@\code{bind}!type declarations}
A value fetched from a collection, or returned by a function the
compiler cannot see into, has no known class. Every operation on it is
then preceded by tests of its class, and a float result must be boxed
because the compiler cannot be sure it is a float. A type declaration in
\code{bind} (\dylanbook{chapter 4}) tells the compiler the class; it is
checked once, when the binding is made. Summing squares of the floats in
a vector of 1000 elements, 100\,000 times:

\begin{dylan}
(define-method sum-squares (v)
  (bind ((s 0.0))
    (for-each ((x v)) ()
      (set! s (+ s (* x x))))
    (> s 0.0)))                                   ; 1.06 s
\end{dylan}

\begin{dylan}
(define-method sum-squares (v)
  (bind ((s 0.0))
    (for-each ((x v)) ()
      (bind (((y <double-float>) x))
        (set! s (+ s (* y y)))))
    (> s 0.0)))                                   ; 0.20 s
\end{dylan}

In the first version \code{(* x x)} tests whether \code{x} is an integer
or a float and allocates a box for the product; in the second it is one
machine multiplication, and nothing is allocated. Specializing a method
on the class of its argument has the same effect inside the method; so
does a \code{type:} on a slot, for the values its getter answers.

Collections are no different. A \code{for-each} over a parameter of
unknown class tests, on each step, whether it is walking a list or a
vector; over one known to be a \code{<simple-object-vector>} it is a
plain indexed loop. Summing a vector of 1000 integers 100\,000 times took
0.16~s with an unspecialized parameter and 0.10~s with the parameter
specialized on \code{<simple-object-vector>}.

\subsection{Keep floats unboxed}

\index{unboxing}
A \code{<double-float>} is a heap object, so a float computation that
creates a new one on every step spends most of its time allocating. The
compiler avoids this by keeping a float variable as a C \code{double}
when \emph{every} use of it wants a \code{double}
(Section~\ref{sec:internals-types}). Returning the variable, storing it,
or passing it to a function the compiler cannot see into each need a
box, and force one on every step. Compare (100 million steps each):

\begin{dylan}
(define-method norm-sum (scale n)
  (for ((i 0 (+ i 1))
        (s 0.0 (+ s (* scale i))))
       ((= i n) s)))                    ; returns s itself: 0.87 s

(define-method norm-sum (scale n)
  (for ((i 0 (+ i 1))
        (s 0.0 (+ s (* scale i))))
       ((= i n) (> s 1.0)))             ; only compares s: 0.04 s
\end{dylan}

The difference is not the work but the boxes. A hot loop should keep its
float accumulators in variables used only arithmetically, and box the
result once, after the loop, by passing it to a function or returning it
from an enclosing one.

\paragraph{Use float contagion, not \code{as}.}
\index{float contagion}
Mixing an integer into float arithmetic converts it automatically
(\dylanbook{chapter 14}), and the compiler does the conversion inline.
An explicit \code{(as <double-float> i)} calls the generic function
\code{as} instead. In the loop above, writing \code{(* scale (as
<double-float> i))} instead of \code{(* scale i)} made it two and a half
times slower (2.2~s against 0.87~s).

\subsection{Typed slots}

A slot declared with \code{type: <double-float>} lets the compiler know
the class of what its getter answers, so arithmetic on it needs no test.
The value stored in the slot is still a boxed float, so the gain is
modest when the program writes the slot often: stepping a particle's
position 100 million times took 1.18~s with untyped slots and 1.05~s
with typed ones.

\subsection{Local exits are free; escaping ones are not}

\index{bind-exit@\code{bind-exit}!cost}
A \code{bind-exit} whose exit procedure is called only directly within
the function that binds it --- an early return from a loop --- compiles
to a jump. One whose exit is passed to another function needs a frame in
the dynamic chain, \code{setjmp}, \code{longjmp}, and a closure
(Section~\ref{sec:internals-exits}). Searching a ten-element vector 30
million times, returning early:

\begin{dylan}
(define-method find-first (v (x <integer>))
  (bind-exit (return)
    (for-each ((e v)) ()
      (when (= e x) (return e)))
    #f))                                              ; 0.08 s

(define-method check (e x return) (when (= e x) (return e)))
(define-method find-first (v (x <integer>))
  (bind-exit (return)
    (for-each ((e v)) ()
      (check e x return))
    #f))                                              ; 1.48 s
\end{dylan}

The second is nearly twenty times slower. Keep the calls of an exit
procedure inside the function that binds it where you can. The same
applies to \code{unwind-protect} (and so to \code{with-lock} and
\code{dynamic-bind}, which use it), whose cleanup always needs a frame,
and to \code{handler-case}: a handled error costs a few hundred
nanoseconds.

\subsection{First-class calls}

A call of a function held in a variable or passed as an argument goes
through the function's uniform entry, with no inlining and no
specialized entry. In a hot loop it costs more than a direct call or
inline code, but often less than you might fear: summing squares of a
vector with \code{(reduce (method (acc x) (+ acc (* x x))) 0 v)} took
0.19~s where the equivalent \code{for-each} loop took 0.16~s. It matters
most when the function called is small and could otherwise have been
inlined. \code{apply} with a long list, and \code{\#rest} lists that
escape, allocate.

\subsection{\code{\#rest} and keywords}

A \code{\#rest} list that is only walked --- not stored, returned, or
captured by a closure --- is built on the C stack, for up to eight
elements, and costs no allocation. Keyword arguments to a function the
compiler can see are matched when compiling, with no list built at all.
So these are cheap in the common cases; a \code{\#rest} list that is
stored or returned is allocated on the heap.

\subsection{What the static program gives}

Everything in a compiled program is effectively sealed: the compiler
knows every class and every method. That is what allows it to do
\code{+} on integers inline, or \code{element} on a vector: it can see
that no method but the library's could apply to such arguments. If your
program defines methods on the library's generic functions for the
library's own classes --- a method on \code{binary+} for two integers,
say, or on \code{element} for \code{<simple-object-vector>} --- the
inline code for those classes is dropped, and every such call in the
program goes through dispatch. (With a method on \code{binary+} for two
integers, \code{+} of two integers calls it, as it must; the inline
cases for floats remain.) Methods for your own classes, or for
combinations the library does not handle inline, do not have this
effect.

\section{Writing fast Dylan for the interactive implementation}
\label{sec:perf-lisp}

The interactive implementation compiles Dylan to Common Lisp, which SBCL
compiles to machine code. Generic functions dispatch through caches;
calls of known functions are made directly; arithmetic on numbers,
\code{element} on simple vectors and strings, and iteration over simple
vectors, lists and ranges are done inline, behind a test of the
arguments' classes; and slots are read directly. Two things defeat the
fast paths:

\begin{itemize}
\item Each inline operation is guarded by a flag that is true only while
  the generic functions behind it have no methods but the predefined
  ones that could apply to those arguments. Adding a method to
  \code{+}, \code{element}, \code{<} and the like for a built-in class
  turns the fast path off for the rest of the session.
\item Redefining things as the program runs (adding methods, adding
  slots) invalidates dispatch caches, which are then rebuilt on the next
  calls.
\end{itemize}

Beyond that, the advice for the compiler applies here too, in smaller
measure: local exits are cheaper than escaping ones, and code that works
on known collection classes is faster. For a program whose speed matters,
develop it interactively and compile it with \cmd{dylanc}.

\section{Size}\label{sec:perf-size}

\index{size builds}
By default the library is compiled apart from the program and linked
whole (Section~\ref{sec:dylanc-cache}), which makes builds fast but
executables large: hello world is about 460~KB of code and data. With
\cmd{-Os}, the library is compiled with the program and shaken
(Section~\ref{sec:internals-shaking}), and hello world is about 220~KB,
stripped; Section~\ref{sec:dylanc-size} compares the builds. The
original goal was under 100~KB. The shaking is by name and class, so a
call of a generic function such as \code{print-object}, \code{make} or
\code{as} keeps every method that could apply to a class with instances,
and those reach much of the library; getting further needs an analysis
of which classes flow to which calls, which the compiler does not yet
do.

A program built \cmd{-Os} may also run slower than one built
\cmd{-O2}, as C optimized for size often does: the vector-summing loop
of Section~\ref{sec:perf-c} took 0.27~s built \cmd{-Os} against 0.16~s,
though \code{tak} and the float loop ran at nearly the same speed.

\section{Measuring}\label{sec:perf-measuring}

\index{perf@\cmd{perf}}\index{profiling}
An ordinary build is compiled with \cmd{-g}, so the usual tools for C
programs work on compiled Dylan, and its \code{\#line} directives make
them name Dylan source lines. On Linux, \cmd{perf} is the simplest:

\begin{shell}
$ bin/dylanc takp.dyl
$ perf stat -e task-clock ./takp
15
             78.05 msec task-clock:u
       0.078615932 seconds time elapsed
$ perf record ./takp
$ perf report --stdio
    50.04%  takp     takp                  [.] fn1190_tak
    49.68%  takp     takp                  [.] dy_gf_tak_e10c5970_dispatch
...
\end{shell}

The functions are named as the compiler names them
(Section~\ref{sec:internals-reading}): \code{fn1190\_tak} is a method
\code{tak}; \code{dy\_gf\_tak\_e10c5970\_dispatch} is the dispatch
function of the generic function \code{tak}. Here, with the method
specialized on \code{<integer>}, half the time goes to dispatch --- the
case of the second table above. Other names to recognize:
\code{\_spec} marks a specialized entry, \code{dy\_maker\_}\meta{class}
a class's maker, and \code{GC\_}\meta{name} the collector; time in the
collector's functions is time spent allocating and collecting.
\cmd{perf annotate} shows the hot instructions with the Dylan lines they
came from.

For timing inside a program, the \code{dylan} module exports
\code{(current-milliseconds)}, an extension in both implementations that
answers a millisecond count from a fixed but arbitrary moment, as the
benchmark driver uses it.

\section{Tuning the garbage collector}\label{sec:perf-gc}

\index{garbage collection!tuning}\index{Memory Pool System!tuning}
The collector is configured by environment variables, which it reads
when the program starts; nothing in Dylan needs to change. The Memory
Pool System, the default collector (Section~\ref{sec:dylanc-gc}), reads
these:

\begin{description}
\item[\code{DYLAN\_MPS\_NURSERY}] The nursery's size in kilobytes,
  131\,072 (128~MB) by default. New objects go here, and it is collected
  when full: a larger nursery collects less often, and more of what it
  holds has died by then. A much smaller one collects too often, each
  time protecting again the older objects the program is writing into
  and taking a fault on each page it then writes; with 32~MB, programs
  that build and keep writing into a large table or vector ran slower
  than under the Boehm collector, and with 128~MB faster. The nursery is
  also the floor of the memory a program that allocates heavily uses: a
  program that must run in little memory can set, say, 32\,768.
\item[\code{DYLAN\_MPS\_EXTEND}] How much memory, in bytes, a pool
  takes from the system at a time: a megabyte by default. The
  collector's own default, which is small, costs a system call for each.
\item[\code{DYLAN\_MPS\_SPARE}] The fraction of the memory the
  collector holds that it may keep unused rather than return to the
  system, between 0 and 1; the default is 1, keeping everything. Memory
  returned must be faulted in afresh when the nursery next fills, which
  for an allocation-heavy program took longer than the collecting: with
  0.75, the MPS's own default, the benchmarks of
  Section~\ref{sec:perf-benchmarks} took 15 to 20 percent longer and used
  a quarter to a half less memory. A program's memory under the MPS is otherwise
  about twice what it is under the Boehm collector, since the collector
  lets the heap double between collections of everything.
\item[\code{DYLAN\_MPS\_PAUSE}] The longest the collector should
  work at once, in seconds, before giving the program a turn; the
  collector's default is 0.1. The MPS collects incrementally, protecting
  the memory it has yet to scan; a shorter pause means more interruptions
  and more protection faults.
\end{description}

\paragraph{The Boehm-Demers-Weiser collector.} \index{Boehm collector!tuning}
A program built with \cmd{-{}-gc boehm} reads the Boehm collector's
variables instead. These were checked with that collector (version
8.2.6), with a program that builds and discards 200 lists of 100\,000
elements. \code{GC\_PRINT\_STATS=1} makes the collector report each
collection on the standard error, which is how the counts below were
made.

\begin{description}
\item[\code{GC\_INITIAL\_HEAP\_SIZE}] The heap's starting size, in bytes
  or with a suffix such as \code{64M} or \code{1G}. The heap starts small
  and grows; a larger start means fewer collections. The test program
  made 202 collections by default and 10 with \code{64M}. Fewer
  collections are not automatically faster: here the larger heap made
  the run slower (0.24~s against 0.18~s), because touching fresh memory
  costs more than the small collections it saved. It helps programs that
  keep much live data.
\item[\code{GC\_FREE\_SPACE\_DIVISOR}] How eagerly the heap grows rather
  than collecting; the default is 3, and smaller values grow the heap
  more and collect less. With 1, the test program made 101 collections
  instead of 202 and used about 13~MB instead of 7~MB.
\item[\code{GC\_MARKERS}] How many threads mark in parallel during a
  collection. The runtime sets two, since more spend longer waiting on
  each other than marking; \code{GC\_MARKERS=1} turns parallel marking
  off, and a larger number is allowed.
\item[\code{GC\_PRINT\_STATS}] Set to 1 to report each collection.
\end{description}

\begin{shell}
$ GC_PRINT_STATS=1 ./alloc 2>&1 | grep -c 'Marking for collection'
202
$ GC_INITIAL_HEAP_SIZE=64M GC_PRINT_STATS=1 ./alloc 2>&1 |
      grep -c 'Marking for collection'
10
\end{shell}

The collectors' own documentation lists further variables. As with the
compiler's options, measure your program before and after: the right
settings depend on how much it allocates and how much it keeps.

\fussy


\appendix
\chapter{Differences from the 1992 Manual}\label{app:differences}
\sloppy

The standard for both implementations is the 1992 manual,
\emph{Dylan: an object-oriented dynamic language} (\textsl{DIRM}). Where
an implementation departs from it --- by necessity, by a decision where
the manual is silent or contradicts itself, or by adding something ---
this appendix says so. It collects in one place what the other chapters
mention as they go. Each entry says which implementation it concerns:
\emph{both}, the \emph{interactive} Common Lisp implementation, or the
\emph{compiler} to C.

\section{Different behavior, in both implementations}
\label{sec:diff-both}

\begin{description}
\item[\code{!key}, \code{!rest} and \code{!next}.]
  \index{"!key@\code{"!key}}
  These are read as \code{\#key}, \code{\#rest} and \code{\#next}, as
  Thomas read them, because Thomas's programs --- the book's examples
  among them --- use them. So they cannot be used as variable names.
  Every other name the manual allows is allowed.
  \code{((method (a !key (b 2)) (list a b)) 1 b: 5)} \yields{}
  \code{(1 5)}.

\item[Modules.] \index{modules}
  The 1992 manual describes modules but leaves the way they are declared
  to the programming environment. Both implementations declare them with
  \code{define-module}, as design note~32 describes, written in prefix
  form, and name a file's module with a \code{module:} header, as design
  note~33 describes. Chapter~\ref{ch:modules} has the details.

\item[Backquote.] \index{backquote}
  Chapter 3 of the manual defers backquote to ``a separate section of
  this manual'' that the manual does not contain. Both implementations
  follow the usual Lisp rules for \code{`}, \code{,} and \code{,@},
  nesting included: \code{`(1 ,(+ 1 1) ,@(list 3 4))} \yields{}
  \code{(1 2 3 4)}.

\item[\code{map-into}.] \index{map-into@\code{map-into}}
  The function is passed only the elements of the source collections, as
  chapter 12's prose and design note~27 say. The chapter's one example,
  \code{(map-into x + '(1 2 3 4))} with \code{x} holding
  \code{(100 100 200 200)}, shows \code{(101 102 203 204)}, which would
  require passing the target's elements too; here it answers
  \code{(1 2 3 4)}.

\item[Irrational functions.]
  \code{sqrt}, \code{exp}, \code{log}, the trigonometric functions and
  their kin answer floats for exact arguments, as in Common Lisp, which
  the manual cites for them, and as the manual's examples show:
  \code{(sqrt 25)} \yields{} \code{5.0}, \code{(exp 0)} \yields{}
  \code{1.0}. Rationals are otherwise exact: \code{(/ 6 4)} \yields{}
  \code{3/2}.

\item[Method specificity among unrelated classes.]
  \index{specificity}
  The manual orders applicable methods by comparing their specializers,
  but leaves open the order of two classes neither of which is a
  subclass of the other. Both implementations follow CLOS: the class
  that comes first in the \emph{argument's} class precedence list is
  more specific.
\begin{dylan}
(define-class <a> (<object>))
(define-class <b> (<object>))
(define-class <c> (<a> <b>))
(define-method which ((x <a>)) 'a)
(define-method which ((x <b>)) 'b)
(which (make <c>))                 ; => a
\end{dylan}
\end{description}

\section{Additions in both implementations}\label{sec:diff-additions}

\begin{description}
\item[Printing.] \code{display}, \code{newline}, \code{print},
  \code{write-line}, \code{write} and \code{pp}, and the generic
  function \code{print-object}, to which programs add methods to control
  how their objects print (Chapter~\ref{ch:printing}). The manual uses
  \code{print} in its examples without defining it.
\item[Streams.] \code{<stream>}, \code{*standard-output*},
  \code{*standard-error*}, \code{<string-stream>} with
  \code{stream-contents}, and \code{format}; every printing function
  takes \code{stream:}. The manual uses streams and \code{format} but
  does not define them. There is no input.
\end{description}

\section{The interactive implementation}\label{sec:diff-interactive}

\begin{description}
\item[Lisp access.] \code{lisp-variable} and \code{lisp-procedure}
  reach Common Lisp values and functions from Dylan. Lisp's \code{NIL}
  arrives as the empty list, not \code{\#f}
  (Chapter~\ref{ch:lisp}).
\item[\code{\#t} and \code{\#f}.] Their class, as \code{object-class}
  answers it, is \code{<object>}; the manual names no class for them.
  (The compiler differs; see below.)
\item[No extensions of the compiler.] The interactive implementation has
  none of the compiler's macros and threads.
\end{description}

\section{The compiler to C}\label{sec:diff-compiler}

\subsection{What a static program cannot do}

The compiler compiles whole, static programs, whose classes and generic
functions are fixed when they are compiled
(Section~\ref{sec:dylanc-static}). The manual allows a running program
to change them; a compiled program cannot:

\begin{itemize}
\item \code{add-method}, \code{remove-method}, \code{add-slot},
  \code{remove-slot} and \code{freeze-methods} are compile errors,
  wherever they appear.
\item Each specializer of a method made with \code{define-method}, and
  each superclass of a class, must be known when compiling: a class, a
  constant variable holding one, \code{(singleton \meta{x})}, or
  \code{(make <singleton> object: \meta{x})} of a constant. A
  \code{define-slot} owner may be any singleton, but a class owner must
  be known when compiling.
\item A class cannot be defined twice.
\item \code{(make <class> ...)} is an error when it runs. A
  \code{<generic-function>} made by \code{make} has no methods, and none
  can be added.
\item There is no evaluation or loading of code at run time.
\end{itemize}

\subsection{Numbers}

\begin{itemize}
\item \textbf{Integers are 63 bits}, from $-2^{62}$ to $2^{62}-1$. There
  are no bignums. A literal out of range is a compile error; an
  arithmetic result out of range signals ``integer overflow (integers are
  63 bits)'' rather than giving a wrong answer. The numerator and
  denominator of a ratio have the same limit. The interactive
  implementation has bignums: \code{(* 4611686018427387904 2)}
  \yields{} \code{9223372036854775808} there, while in a compiled program
  even the literal is an error.
\item \textbf{No complex numbers.} \code{<complex>} exists, abstract, as
  the superclass of \code{<real>} (so \code{(instance? 3 <complex>)}
  \yields{} \code{\#t}), but it has no instances. A complex literal is a
  compile error, and a function whose result would be complex signals an
  error: \code{(sqrt -4.0)} signals ``sqrt: the result would be complex''
  where the interactive implementation answers \code{\#c(0.0 2.0)}.
\item \code{<single-float>} is a C \code{float}; \code{<double-float>}
  and \code{<extended-float>} are both a C \code{double}. Floats print
  as the shortest decimal that reads back: \code{21.0}, and \code{1.5s0}
  for a single float.
\end{itemize}

\subsection{Conditions without a debugger}

There is no interactive debugger. An unhandled serious condition prints
\code{Dylan error:} and its message on the standard error and ends the
program with status~1 --- or, in a thread other than the main one, ends
that thread (Chapter~\ref{ch:extensions}). An unhandled warning prints
\code{Dylan warning:} and its message and continues. \code{break}
prints its message and returns \code{\#f}. \code{restart-query} and
\code{return-query} do nothing.

\subsection{Representation}

\code{\#t} and \code{\#f} are instances of an internal class,
\code{<boolean>}, which \code{object-class} answers but which the
\code{dylan} module does not export (the manual names no class for
them). \code{seal} works as the manual says, and \code{make-read-only}
answers its argument, a class, unchanged, since a compiled program's
classes are already fixed.

\subsection{Additions}

All of these are extensions; the 1992 manual has none of them.

\begin{itemize}
\item \textbf{Macros}: \code{define-macro}, \code{define-for-macros},
  and \code{gensym} at compile time; \cmd{dylanc -E} to see expansions
  (Chapter~\ref{ch:extensions}).
\item \textbf{Threads}: \code{make-thread}, \code{join-thread},
  \code{current-thread}, \code{thread-yield}; \code{<lock>},
  \code{<recursive-lock>}, \code{acquire-lock}, \code{release-lock} and
  \code{with-lock}; \code{<condition-variable>} with
  \code{wait-condition-variable}, \code{notify-condition-variable} and
  \code{notify-all-condition-variable}; \code{<atomic-integer>} with
  \code{atomic-value}, \code{atomic-add!} and
  \code{atomic-compare-and-swap!}; \code{define-thread-variable} and
  \code{dynamic-bind}; \code{<thread-error>} with
  \code{thread-error-thread} and \code{thread-error-condition}
  (Chapter~\ref{ch:extensions}).
\item \textbf{Timing}: \code{current-milliseconds}.
\end{itemize}

\section{Where the two implementations disagree}\label{sec:diff-between}

Besides the differences above, the implementations disagree where the
manual leaves a choice open or where the host shows through. A portable
program depends on none of these; each is described where it arises.

\begin{itemize}
\item The operator of a call is evaluated first by the compiler to C
  and last by the interactive implementation; arguments are evaluated
  left to right in both (Section~\ref{sec:syntax-calls}).
\item A method whose parameter list has \code{\#key} with no keywords
  accepts no keywords in the compiler to C and any keyword in the
  interactive implementation (Section~\ref{sec:functions-keyword-checking}).
\item A handler that returns from the error \code{check-type} signals
  leaves it unhandled in the compiler to C; in the interactive
  implementation \code{check-type} answers what the handler returned
  (Section~\ref{sec:conditions-signalling}).
\item Equal string literals may be one object in the compiler to C;
  in the interactive implementation each is its own
  (Section~\ref{sec:collections-strings}).
\item \code{min} and \code{max} of equal numbers of different classes,
  and functions whose true result is complex or infinite, answer
  differently (Sections~\ref{sec:numbers-misc}
  and~\ref{sec:numbers-transcendental}).
\item Assigning a read-only module variable is a compile error in the
  compiler to C and an error when it runs in the interactive
  implementation, which also needs modules declared before use
  (Sections~\ref{sec:modules-readonly} and~\ref{sec:modules-lisp}).
\item Error messages are worded differently
  (Section~\ref{sec:conditions-runtime-errors}).
\end{itemize}

\fussy

\chapter{Library Reference}\label{app:library}

\index{library reference}
This appendix lists every name the module \code{dylan} exports, in
either implementation, and the special forms of the language. Each entry
gives the name with its parameters, as a call would be written, and a
one-line description. The full definitions are in the DIRM, cited by
chapter at the head of each section, and in the chapters of this book.

In a signature, \code{\#rest} and \code{\#key} introduce rest and keyword
parameters as in a method's parameter list (Chapter~\ref{ch:functions}).
A name in angle brackets is a class. Where a function answers several
values, the description says so.

The narrow last column says which implementations provide the name.
\begin{itemize}
\item An empty column means both do.
\item \textsc{c} means only the compiler to C provides it.
\item \textsc{l} means only the interactive (Common Lisp) implementation
  provides it.
\end{itemize}

\noindent
Every entry was checked by referring to it in a small program compiled
with \cmd{bin/dylanc} and by evaluating it at the interactive
implementation's read-eval-print loop. A name marked \textsc{l} is either
unknown to the compiler to C or rejected by it when compiling: the
compiler reports \code{add-method}, \code{remove-method},
\code{add-slot}, \code{remove-slot} and \code{freeze-methods} as
impossible in a static program, and \code{<rectangular-complex>} and
\code{(setter direct-superclasses)} as not supported.

\section{Objects, equality and comparison}\label{sec:lib-objects}

{\sloppy The root class and the operations every object supports.\par}

{\small\frenchspacing
\begin{longtable}{@{}>{\raggedright\arraybackslash}p{0.43\textwidth}>{\raggedright\arraybackslash}p{0.47\textwidth}>{\centering\arraybackslash}p{0.04\textwidth}@{}}
\toprule
Name & Description & \\
\midrule
\endhead
\code{<object>}\index{<object>} & The class of all objects; every class has it as a superclass. &  \\
\code{(id? object1 object2 \#rest objects)}\index{id?} & True if the arguments are all the same object. &  \\
\code{(= object1 object2 \#rest objects)}\index{=} & True if the arguments are all equal, as \code{binary=} decides. &  \\
\code{(/= object1 object2)}\index{/=} & True if the arguments are not equal. &  \\
\code{(binary= object1 object2)}\index{binary=} & The generic equality test on which \code{=} is built; add methods to it to define equality for a class. &  \\
\code{(=hash object)}\index{=hash} & A hash code for \code{object} consistent with \code{=}: equal objects answer equal codes. &  \\
\code{(< object1 object2 \#rest objects)}\index{<} & True if each argument is less than the next, as \code{binary<} decides. &  \\
\code{(> object1 object2 \#rest objects)}\index{>} & True if each argument is greater than the next. &  \\
\code{(<= object1 object2 \#rest objects)}\index{<=} & True if each argument is less than or equal to the next. &  \\
\code{(>= object1 object2 \#rest objects)}\index{>=} & True if each argument is greater than or equal to the next. &  \\
\code{(binary< object1 object2)}\index{binary<} & The generic ordering on which \code{<}, \code{>}, \code{<=} and \code{>=} are built; methods exist for reals, characters and strings. &  \\
\code{(not object)}\index{not} & \code{\#t} if \code{object} is \code{\#f}, otherwise \code{\#f}. &  \\
\code{(identity object)}\index{identity} & Answers \code{object}. &  \\
\bottomrule
\end{longtable}}

\section{Classes, instances and slots}\label{sec:lib-classes}

{\sloppy Making objects and examining classes (\dylanbook{chapter 11}). The compiler to C fixes a program's classes when it compiles it, so it rejects \code{add-slot} and \code{remove-slot} (Chapter~\ref{ch:dylanc}).\par}

{\small\frenchspacing
\begin{longtable}{@{}>{\raggedright\arraybackslash}p{0.43\textwidth}>{\raggedright\arraybackslash}p{0.47\textwidth}>{\centering\arraybackslash}p{0.04\textwidth}@{}}
\toprule
Name & Description & \\
\midrule
\endhead
\code{<class>}\index{<class>} & The class of classes. &  \\
\code{<singleton>}\index{<singleton>} & The class of singleton specializers, each matching exactly one object. &  \\
\code{(singleton object)}\index{singleton} & A singleton specializer for \code{object}. &  \\
\code{<slot-descriptor>}\index{<slot-descriptor>} & The class of the objects that describe slots. &  \\
\code{(make class \#rest key-value-pairs)}\index{make} & Makes and initializes a new instance of \code{class}, passing it the init arguments. &  \\
\code{(initialize instance \#rest init-arguments)}\index{initialize} & Called by \code{make} on each new instance; add methods to it to finish initializing instances of a class. &  \\
\code{(instance? object type)}\index{instance?} & True if \code{object} is a general instance of \code{type} (a class or singleton). &  \\
\code{(subclass? class1 class2)}\index{subclass?} & True if \code{class1} is \code{class2} or one of its subclasses. &  \\
\code{(object-class object)}\index{object-class} & The class of which \code{object} is a direct instance. &  \\
\code{(all-superclasses class)}\index{all-superclasses} & The class and all its superclasses, in precedence order. &  \\
\code{(direct-superclasses class)}\index{direct-superclasses} & The class's direct superclasses. &  \\
\code{((setter direct-superclasses) class superclasses)} & Changes the class's direct superclasses. & \textsc{l} \\
\code{(direct-subclasses class)}\index{direct-subclasses} & The class's direct subclasses. &  \\
\code{(as class object)}\index{as} & Converts \code{object} to an instance of \code{class}, for example a list to a vector or a character to an integer. &  \\
\code{(shallow-copy object)}\index{shallow-copy} & A new object like \code{object}, sharing its contents. &  \\
\code{(class-for-copy object)}\index{class-for-copy} & The class that \code{map}, \code{copy-sequence} and the like use for a copy of \code{object}. &  \\
\code{(seal class)}\index{seal} & Makes the class read-only and forbids new subclasses of it; answers it. &  \\
\code{(make-read-only object)}\index{make-read-only} & Makes a class or generic function read-only, so that its slots or methods cannot change; answers it. &  \\
\code{(slot-descriptors class)}\index{slot-descriptors} & The descriptors of the class's slots. &  \\
\code{(add-slot owner \#key getter setter type init-value init-function init-keyword allocation)}\index{add-slot} & Adds a slot to a class (or singleton) that may already have instances. & \textsc{l} \\
\code{(remove-slot owner slot-descriptor)}\index{remove-slot} & Removes a slot. & \textsc{l} \\
\code{(slot-descriptor instance getter)}\index{slot-descriptor} & The descriptor of the slot of \code{instance} that \code{getter} reads. &  \\
\code{(slot-getter slot-descriptor)}\index{slot-getter} & The slot's getter, or \code{\#f}. &  \\
\code{(slot-setter slot-descriptor)}\index{slot-setter} & The slot's setter, or \code{\#f}. &  \\
\code{(slot-type slot-descriptor)}\index{slot-type} & The slot's type. &  \\
\code{(slot-allocation slot-descriptor)}\index{slot-allocation} & The slot's allocation, a symbol such as \code{instance} or \code{class}. &  \\
\code{(init-value slot-descriptor)}\index{init-value} & The slot's init value, and a second value saying whether it has one. &  \\
\code{(init-function slot-descriptor)}\index{init-function} & The slot's init function, or \code{\#f}. &  \\
\code{(init-keyword slot-descriptor)}\index{init-keyword} & The slot's init keyword, or \code{\#f}. &  \\
\code{(slot-value instance slot-descriptor)}\index{slot-value} & The value of the slot in \code{instance}. &  \\
\code{((setter slot-value) instance slot-descriptor new-value)} & Sets the slot's value. &  \\
\code{(slot-initialized? instance getter)}\index{slot-initialized?} & True if the slot \code{getter} reads has a value in \code{instance}. &  \\
\bottomrule
\end{longtable}}

\section{Functions and generic functions}\label{sec:lib-functions}

{\sloppy Calling functions and examining and changing generic functions (\dylanbook{chapter 10}). The compiler to C rejects \code{add-method}, \code{remove-method} and \code{freeze-methods}, and a \code{<generic-function>} it makes with \code{make} has no methods.\par}

{\small\frenchspacing
\begin{longtable}{@{}>{\raggedright\arraybackslash}p{0.43\textwidth}>{\raggedright\arraybackslash}p{0.47\textwidth}>{\centering\arraybackslash}p{0.04\textwidth}@{}}
\toprule
Name & Description & \\
\midrule
\endhead
\code{<function>}\index{<function>} & The class of all functions. &  \\
\code{<generic-function>}\index{<generic-function>} & The class of generic functions. &  \\
\code{<method>}\index{<method>} & The class of methods. &  \\
\code{(apply function \#rest arguments)}\index{apply} & Calls \code{function}; the last argument is a sequence of further arguments. &  \\
\code{(values \#rest objects)}\index{values} & Answers its arguments as multiple values. &  \\
\code{(add-method generic-function method)}\index{add-method} & Adds a method to a generic function, replacing one with the same specializers; answers the method and the one replaced, or \code{\#f}. & \textsc{l} \\
\code{(remove-method generic-function method)}\index{remove-method} & Removes a method from a generic function. & \textsc{l} \\
\code{(freeze-methods generic-function \#rest specializers)}\index{freeze-methods} & Forbids changes to the methods applicable to arguments of the given classes. & \textsc{l} \\
\code{(generic-function-methods generic-function)}\index{generic-function-methods} & The generic function's methods. &  \\
\code{(method-specializers method)}\index{method-specializers} & The specializers of the method's required parameters. &  \\
\code{(function-arguments function)}\index{function-arguments} & Three values: the number of required arguments, whether there is a \code{\#rest} argument, and the keywords accepted. &  \\
\code{(applicable-method? function \#rest sample-arguments)}\index{applicable-method?} & True if \code{function} could be called with such arguments. &  \\
\code{(sorted-applicable-methods generic-function \#rest sample-arguments)}\index{sorted-applicable-methods} & The methods applicable to such arguments, most specific first, and those whose order is ambiguous. &  \\
\code{(find-method generic-function specializers)}\index{find-method} & The method with exactly these specializers, or \code{\#f}. &  \\
\bottomrule
\end{longtable}}

\section{Functional operations}\label{sec:lib-functional}

{\sloppy Functions that make functions (\dylanbook{chapter 17}).\par}

{\small\frenchspacing
\begin{longtable}{@{}>{\raggedright\arraybackslash}p{0.43\textwidth}>{\raggedright\arraybackslash}p{0.47\textwidth}>{\centering\arraybackslash}p{0.04\textwidth}@{}}
\toprule
Name & Description & \\
\midrule
\endhead
\code{(compose function1 \#rest functions)}\index{compose} & A function that applies the functions from right to left. &  \\
\code{(complement predicate)}\index{complement} & A function that answers \code{\#t} where \code{predicate} answers \code{\#f}, and \code{\#f} otherwise. &  \\
\code{(disjoin predicate1 \#rest predicates)}\index{disjoin} & A function true when any of the predicates is true of its arguments. &  \\
\code{(conjoin predicate1 \#rest predicates)}\index{conjoin} & A function true when all of the predicates are true of its arguments. &  \\
\code{(curry function \#rest arguments)}\index{curry} & A function that calls \code{function} with these arguments first, then its own. &  \\
\code{(rcurry function \#rest arguments)}\index{rcurry} & A function that calls \code{function} with its own arguments first, then these. &  \\
\code{(always object)}\index{always} & A function that ignores its arguments and answers \code{object}. &  \\
\bottomrule
\end{longtable}}

\section{Numbers}\label{sec:lib-numbers}

{\sloppy The number classes and arithmetic (\dylanbook{chapter 14}). In the compiler to C integers are 63 bits, overflow is an error, and there are no complex numbers: \code{<complex>} has no instances, \code{make-rectangular} and \code{make-polar} signal an error unless the result is real, and so do the functions whose result would be complex, such as \code{(sqrt -4)}. The interactive implementation has Common Lisp's unlimited integers and complex numbers.\par}

{\small\frenchspacing
\begin{longtable}{@{}>{\raggedright\arraybackslash}p{0.43\textwidth}>{\raggedright\arraybackslash}p{0.47\textwidth}>{\centering\arraybackslash}p{0.04\textwidth}@{}}
\toprule
Name & Description & \\
\midrule
\endhead
\code{<number>}\index{<number>} & The class of all numbers. &  \\
\code{<complex>}\index{<complex>} & Complex numbers, the superclass of \code{<real>}. &  \\
\code{<rectangular-complex>}\index{<rectangular-complex>} & Complex numbers with a nonzero imaginary part. & \textsc{l} \\
\code{<real>}\index{<real>} & Real numbers. &  \\
\code{<float>}\index{<float>} & Floating-point numbers. &  \\
\code{<single-float>}\index{<single-float>} & Single-precision floats. &  \\
\code{<double-float>}\index{<double-float>} & Double-precision floats, the class of float literals such as \code{1.5}. &  \\
\code{<extended-float>}\index{<extended-float>} & Extended-precision floats. &  \\
\code{<rational>}\index{<rational>} & Exact rational numbers. &  \\
\code{<ratio>}\index{<ratio>} & Ratios that are not integers, such as \code{3/4}. &  \\
\code{<integer>}\index{<integer>} & Integers. &  \\
\code{(odd? integer)}\index{odd?} & True if \code{integer} is odd. &  \\
\code{(even? integer)}\index{even?} & True if \code{integer} is even. &  \\
\code{(zero? number)}\index{zero?} & True if \code{number} is zero. &  \\
\code{(positive? real)}\index{positive?} & True if \code{real} is greater than zero. &  \\
\code{(negative? real)}\index{negative?} & True if \code{real} is less than zero. &  \\
\code{(integral? number)}\index{integral?} & True if \code{number} has an integer value. &  \\
\code{(+ \#rest numbers)}\index{+} & The sum; \code{(+)} is \code{0}. &  \\
\code{(* \#rest numbers)}\index{*} & The product; \code{(*)} is \code{1}. &  \\
\code{(- number \#rest numbers)}\index{-} & The difference, or the negation of a single argument. &  \\
\code{(/ number \#rest numbers)}\index{/} & The quotient, or the reciprocal of a single argument. Integers divide exactly, giving a ratio when needed. &  \\
\code{(binary+ number1 number2)}\index{binary+} & The generic addition \code{+} is built on; similarly the next five. &  \\
\code{(binary- number1 number2)}\index{binary-} & Generic subtraction. &  \\
\code{(binary* number1 number2)}\index{binary*} & Generic multiplication. &  \\
\code{(binary/ number1 number2)}\index{binary/} & Generic division. &  \\
\code{(unary- number)}\index{unary-} & Generic negation. &  \\
\code{(unary/ number)}\index{unary/} & Generic reciprocal. &  \\
\code{(floor real)}\index{floor} & The greatest integer not above \code{real}, and the remainder. &  \\
\code{(ceiling real)}\index{ceiling} & The least integer not below \code{real}, and the remainder. &  \\
\code{(round real)}\index{round} & The nearest integer (to even on a tie), and the remainder. &  \\
\code{(truncate real)}\index{truncate} & The integer part, toward zero, and the remainder. &  \\
\code{(floor/ real1 real2)}\index{floor/} & The quotient rounded down, and the remainder. &  \\
\code{(ceiling/ real1 real2)}\index{ceiling/} & The quotient rounded up, and the remainder. &  \\
\code{(round/ real1 real2)}\index{round/} & The quotient rounded to nearest, and the remainder. &  \\
\code{(truncate/ real1 real2)}\index{truncate/} & The quotient truncated toward zero, and the remainder. &  \\
\code{(modulo real1 real2)}\index{modulo} & The remainder of \code{floor/}, with the sign of \code{real2}. &  \\
\code{(remainder real1 real2)}\index{remainder} & The remainder of \code{truncate/}, with the sign of \code{real1}. &  \\
\code{(abs number)}\index{abs} & The absolute value. &  \\
\code{(magnitude number)}\index{magnitude} & The magnitude (the absolute value of a real). &  \\
\code{(min real \#rest reals)}\index{min} & The smallest argument. &  \\
\code{(max real \#rest reals)}\index{max} & The largest argument. &  \\
\code{(gcd \#rest integers)}\index{gcd} & The greatest common divisor. &  \\
\code{(lcm \#rest integers)}\index{lcm} & The least common multiple. &  \\
\code{(binary-gcd integer1 integer2)}\index{binary-gcd} & The generic function \code{gcd} is built on. &  \\
\code{(binary-lcm integer1 integer2)}\index{binary-lcm} & The generic function \code{lcm} is built on. &  \\
\code{(numerator rational)}\index{numerator} & The numerator, in lowest terms. &  \\
\code{(denominator rational)}\index{denominator} & The denominator, in lowest terms. &  \\
\code{(rationalize number)}\index{rationalize} & A simple rational close to \code{number}. &  \\
\code{(sqrt number)}\index{sqrt} & The square root, a float (\code{(sqrt 25)} is \code{5.0}). &  \\
\code{(exp number)}\index{exp} & $e$ raised to the power \code{number}. &  \\
\code{(log number)}\index{log} & The natural logarithm. &  \\
\code{(expt number1 number2)}\index{expt} & \code{number1} raised to the power \code{number2}; exact for an integer power of an exact number. &  \\
\code{(sin number) (cos number) (tan number)}\index{sin}\index{cos}\index{tan} & The trigonometric functions, in radians. &  \\
\code{(asin number) (acos number) (atan number)}\index{asin}\index{acos}\index{atan} & Their inverses. &  \\
\code{(atan2 number1 number2)}\index{atan2} & The arc tangent of \code{number1/number2}, using both signs to choose the quadrant. &  \\
\code{(sinh number) (cosh number) (tanh number)}\index{sinh}\index{cosh}\index{tanh} & The hyperbolic functions. &  \\
\code{(asinh number) (acosh number) (atanh number)}\index{asinh}\index{acosh}\index{atanh} & Their inverses. &  \\
\code{(make-rectangular real imaginary)}\index{make-rectangular} & The complex number with these parts. &  \\
\code{(make-polar magnitude angle)}\index{make-polar} & The complex number with this magnitude and angle. &  \\
\code{(real-part number)}\index{real-part} & The real part. &  \\
\code{(imag-part number)}\index{imag-part} & The imaginary part (\code{0} for a real). &  \\
\code{(angle number)}\index{angle} & The angle of a complex number in radians. &  \\
\code{(logior \#rest integers)}\index{logior} & Bitwise inclusive or. &  \\
\code{(logxor \#rest integers)}\index{logxor} & Bitwise exclusive or. &  \\
\code{(logand \#rest integers)}\index{logand} & Bitwise and. &  \\
\code{(logeqv \#rest integers)}\index{logeqv} & Bitwise equivalence (exclusive nor). &  \\
\code{(lognot integer)}\index{lognot} & Bitwise complement. &  \\
\code{(lognand integer1 integer2)}\index{lognand} & The complement of \code{logand}. &  \\
\code{(lognor integer1 integer2)}\index{lognor} & The complement of \code{logior}. &  \\
\code{(logandc1 integer1 integer2)}\index{logandc1} & The and of the complement of \code{integer1} with \code{integer2}. &  \\
\code{(logandc2 integer1 integer2)}\index{logandc2} & The and of \code{integer1} with the complement of \code{integer2}. &  \\
\code{(logorc1 integer1 integer2)}\index{logorc1} & The or of the complement of \code{integer1} with \code{integer2}. &  \\
\code{(logorc2 integer1 integer2)}\index{logorc2} & The or of \code{integer1} with the complement of \code{integer2}. &  \\
\code{(logbit? index integer)}\index{logbit?} & True if bit \code{index} of \code{integer} is 1. &  \\
\code{(ash integer count)}\index{ash} & \code{integer} shifted left \code{count} bits (right if \code{count} is negative). &  \\
\bottomrule
\end{longtable}}

\section{Characters, strings and symbols}\label{sec:lib-chars}

{\sloppy Characters (\dylanbook{chapter 16}), symbols and keywords (\dylanbook{chapter 15}), and the case conversions. The string classes are collections, listed with them below.\par}

{\small\frenchspacing
\begin{longtable}{@{}>{\raggedright\arraybackslash}p{0.43\textwidth}>{\raggedright\arraybackslash}p{0.47\textwidth}>{\centering\arraybackslash}p{0.04\textwidth}@{}}
\toprule
Name & Description & \\
\midrule
\endhead
\code{<character>}\index{<character>} & The class of characters. &  \\
\code{<symbol>}\index{<symbol>} & The class of symbols. &  \\
\code{<keyword>}\index{<keyword>} & The class of keywords such as \code{foo:}. &  \\
\code{(as-uppercase object)}\index{as-uppercase} & The upper-case version of a character, or a new upper-case string. &  \\
\code{(as-lowercase object)}\index{as-lowercase} & The lower-case version of a character, or a new lower-case string. &  \\
\code{(as-uppercase! string)}\index{as-uppercase"!} & Converts \code{string} to upper case in place. &  \\
\code{(as-lowercase! string)}\index{as-lowercase"!} & Converts \code{string} to lower case in place. &  \\
\bottomrule
\end{longtable}}

\section{Collection classes}\label{sec:lib-coll-classes}

{\sloppy The classes of the collection library (\dylanbook{chapter 12}).\par}

{\small\frenchspacing
\begin{longtable}{@{}>{\raggedright\arraybackslash}p{0.43\textwidth}>{\raggedright\arraybackslash}p{0.47\textwidth}>{\centering\arraybackslash}p{0.04\textwidth}@{}}
\toprule
Name & Description & \\
\midrule
\endhead
\code{<collection>}\index{<collection>} & All collections. &  \\
\code{<explicit-key-collection>}\index{<explicit-key-collection>} & Collections whose keys are not simply indices, such as tables. &  \\
\code{<mutable-collection>}\index{<mutable-collection>} & Collections whose elements can be changed. &  \\
\code{<mutable-explicit-key-collection>}\index{<mutable-explicit-key-collection>} & Mutable explicit-key collections. &  \\
\code{<sequence>}\index{<sequence>} & Collections whose keys are the integers from 0. &  \\
\code{<mutable-sequence>}\index{<mutable-sequence>} & Sequences whose elements can be changed. &  \\
\code{<array>}\index{<array>} & Multi-dimensional arrays; make one with \code{dimensions:}. &  \\
\code{<vector>}\index{<vector>} & One-dimensional arrays. &  \\
\code{<simple-object-vector>}\index{<simple-object-vector>} & The vectors \code{vector} makes; they hold any object. &  \\
\code{<stretchy-vector>}\index{<stretchy-vector>} & Vectors that grow as elements are added. &  \\
\code{<string>}\index{<string>} & Strings, the superclass of the two below. &  \\
\code{<byte-string>}\index{<byte-string>} & Strings of 8-bit characters; string literals are byte strings. &  \\
\code{<unicode-string>}\index{<unicode-string>} & Strings of any characters. &  \\
\code{<deque>}\index{<deque>} & Double-ended queues. &  \\
\code{<list>}\index{<list>} & Lists: the empty list and pairs. &  \\
\code{<pair>}\index{<pair>} & Pairs (cons cells). &  \\
\code{<empty-list>}\index{<empty-list>} & The class of the empty list, \code{()}. &  \\
\code{<range>}\index{<range>} & Arithmetic sequences, possibly unbounded. &  \\
\code{<table>}\index{<table>} & Hash tables, keyed by any object, compared with \code{=}. &  \\
\bottomrule
\end{longtable}}

\section{Collection operations}\label{sec:lib-coll-ops}

{\sloppy Operations on every collection, and the iteration protocol that a new collection class implements to take part in them.\par}

{\small\frenchspacing
\begin{longtable}{@{}>{\raggedright\arraybackslash}p{0.43\textwidth}>{\raggedright\arraybackslash}p{0.47\textwidth}>{\centering\arraybackslash}p{0.04\textwidth}@{}}
\toprule
Name & Description & \\
\midrule
\endhead
\code{(element collection key \#key default)}\index{element} & The element at \code{key}; \code{default} if there is none (an error if no default is given). &  \\
\code{((setter element) collection key new-value)} & Sets the element at \code{key}. Usually written \code{(set! (element c k) v)}. &  \\
\code{(size collection)}\index{size} & The number of elements, or \code{\#f} for an unbounded collection. &  \\
\code{((setter size) sequence new-size)} & Changes the size of a stretchy sequence. &  \\
\code{(empty? collection)}\index{empty?} & True if the collection has no elements. &  \\
\code{(do procedure collection \#rest collections)}\index{do} & Calls \code{procedure} on the elements, in parallel across the collections; answers \code{\#f}. &  \\
\code{(map procedure collection \#rest collections)}\index{map} & A new collection of the results of \code{procedure} on the elements. &  \\
\code{(map-as class procedure collection \#rest collections)}\index{map-as} & As \code{map}, but the result is an instance of \code{class}. &  \\
\code{(map-into collection procedure source \#rest sources)}\index{map-into} & Stores the results of \code{procedure} on the sources' elements into \code{collection}; answers it. (Only the sources' elements are passed; see Appendix~\ref{app:differences}.) &  \\
\code{(any? procedure collection \#rest collections)}\index{any?} & The first true result of \code{procedure} on the elements, or \code{\#f}. &  \\
\code{(every? procedure collection \#rest collections)}\index{every?} & True if \code{procedure} is true of every element. &  \\
\code{(reduce procedure initial-value collection)}\index{reduce} & Combines the elements with \code{procedure}, starting from \code{initial-value}. &  \\
\code{(reduce1 procedure collection)}\index{reduce1} & As \code{reduce}, starting from the first element. &  \\
\code{(member? value collection \#key test)}\index{member?} & True if \code{value} is an element (compared with \code{test}, default \code{id?}). &  \\
\code{(find-key collection procedure \#key skip failure)}\index{find-key} & The key of the first element (after skipping \code{skip} matches) for which \code{procedure} is true; \code{failure} (default \code{\#f}) if none. &  \\
\code{(replace-elements! collection predicate new-value-function \#key count)}\index{replace-elements"!} & Replaces each element satisfying \code{predicate} with \code{new-value-function} of it, at most \code{count} times. &  \\
\code{(fill! collection value \#key start end)}\index{fill"!} & Stores \code{value} into every element (of a sequence, from \code{start} to \code{end}). &  \\
\code{(key-sequence collection)}\index{key-sequence} & A sequence of the collection's keys. &  \\
\code{(initial-state collection)}\index{initial-state} & The state of the first element, or \code{\#f} if there is none. &  \\
\code{(next-state collection state)}\index{next-state} & The state of the next element, or \code{\#f} at the end. &  \\
\code{(current-element collection state)}\index{current-element} & The element at \code{state}. &  \\
\code{((setter current-element) collection state new-value)} & Replaces the element at \code{state}. &  \\
\code{(current-key collection state)}\index{current-key} & The key of the element at \code{state}. &  \\
\code{(final-state collection)}\index{final-state} & The state of the last element of a sequence, or \code{\#f}. &  \\
\code{(previous-state collection state)}\index{previous-state} & The state of the element before \code{state}, or \code{\#f}. &  \\
\code{(copy-state collection state)}\index{copy-state} & A copy of \code{state} that can be advanced independently. &  \\
\bottomrule
\end{longtable}}

\section{Sequence operations}\label{sec:lib-sequences}

{\sloppy Operations on sequences. Unless \code{test:} is given, elements are compared with \code{id?}.\par}

{\small\frenchspacing
\begin{longtable}{@{}>{\raggedright\arraybackslash}p{0.43\textwidth}>{\raggedright\arraybackslash}p{0.47\textwidth}>{\centering\arraybackslash}p{0.04\textwidth}@{}}
\toprule
Name & Description & \\
\midrule
\endhead
\code{(add sequence new-element)}\index{add} & A new sequence with \code{new-element} added (at the front of a list). &  \\
\code{(add! sequence new-element)}\index{add"!} & Adds \code{new-element}, altering \code{sequence} if it can; answers the result. &  \\
\code{(add-new sequence new-element \#key test)}\index{add-new} & As \code{add}, unless the element is already present. &  \\
\code{(add-new! sequence new-element \#key test)}\index{add-new"!} & As \code{add!}, unless the element is already present. &  \\
\code{(remove sequence value \#key test count)}\index{remove} & A new sequence without the elements equal to \code{value} (at most \code{count} of them). &  \\
\code{(remove! sequence value \#key test count)}\index{remove"!} & As \code{remove}, altering \code{sequence} if it can. &  \\
\code{(choose predicate sequence)}\index{choose} & A new sequence of the elements satisfying \code{predicate}. &  \\
\code{(choose-by predicate test-sequence value-sequence)}\index{choose-by} & The elements of \code{value-sequence} whose corresponding elements of \code{test-sequence} satisfy \code{predicate}. &  \\
\code{(intersection sequence1 sequence2 \#key test)}\index{intersection} & The elements of \code{sequence1} that are also in \code{sequence2}. &  \\
\code{(union sequence1 sequence2 \#key test)}\index{union} & The elements of either sequence, without duplicates between them. &  \\
\code{(remove-duplicates sequence \#key test)}\index{remove-duplicates} & A new sequence with each element only once. &  \\
\code{(remove-duplicates! sequence \#key test)}\index{remove-duplicates"!} & As \code{remove-duplicates}, altering \code{sequence} if it can. &  \\
\code{(copy-sequence sequence \#key start end)}\index{copy-sequence} & A new sequence of the elements from \code{start} to \code{end}. &  \\
\code{(concatenate sequence \#rest sequences)}\index{concatenate} & A new sequence of all the elements, of the first sequence's kind. &  \\
\code{(concatenate-as class sequence \#rest sequences)}\index{concatenate-as} & As \code{concatenate}, but an instance of \code{class}. &  \\
\code{(replace-subsequence! sequence insert-sequence \#key start)}\index{replace-subsequence"!} & Overwrites elements of \code{sequence}, from \code{start}, with those of \code{insert-sequence}. &  \\
\code{(reverse sequence)}\index{reverse} & A new sequence with the elements in reverse order. &  \\
\code{(reverse! sequence)}\index{reverse"!} & As \code{reverse}, reusing \code{sequence} if it can. &  \\
\code{(sort sequence \#key test stable)}\index{sort} & A new sequence sorted by \code{test} (default \code{<}). &  \\
\code{(sort! sequence \#key test stable)}\index{sort"!} & As \code{sort}, reusing \code{sequence} if it can. &  \\
\code{(first sequence) (second sequence) (third sequence)}\index{first}\index{second}\index{third} & The first, second and third elements. &  \\
\code{(last sequence)}\index{last} & The last element. &  \\
\code{((setter first) sequence value) }\textrm{and likewise}\code{ second, third, last} & Set the first, second, third and last elements. &  \\
\code{(subsequence-position big pattern \#key test count)}\index{subsequence-position} & The index at which \code{pattern} occurs in \code{big} (the \code{count}th occurrence), or \code{\#f}. &  \\
\bottomrule
\end{longtable}}

\section{Lists}\label{sec:lib-lists}

{\sloppy Pairs and lists.\par}

{\small\frenchspacing
\begin{longtable}{@{}>{\raggedright\arraybackslash}p{0.43\textwidth}>{\raggedright\arraybackslash}p{0.47\textwidth}>{\centering\arraybackslash}p{0.04\textwidth}@{}}
\toprule
Name & Description & \\
\midrule
\endhead
\code{(cons car cdr)}\index{cons} & A new pair. &  \\
\code{(list \#rest objects)}\index{list} & A new list of the arguments. &  \\
\code{(list* \#rest objects)}\index{list*} & A new list of the arguments, the last of which is its tail. &  \\
\code{(car list) (cdr list)}\index{car}\index{cdr} & The first element, and the rest, of a pair; both are \code{()} for \code{()}. &  \\
\code{(caar list) }\textrm{\dots}\code{ (cdddr list)}\index{caar}\index{cdddr} & The compositions of \code{car} and \code{cdr} to three levels: \code{caar}, \code{cadr}, \code{cdar}, \code{cddr}, \code{caaar}, \code{caadr}, \code{cadar}, \code{caddr}, \code{cdaar}, \code{cdadr}, \code{cddar}, \code{cdddr}. &  \\
\code{((setter car) pair value) }\textrm{and likewise for the others}\code{} & Store into a pair; every \code{c\dots r} above has a setter. &  \\
\code{(append list \#rest lists)}\index{append} & A new list of the elements of all the lists. &  \\
\code{(find-pair item list \#key test)}\index{find-pair} & The first pair in \code{list} whose \code{car} matches \code{item}, or \code{\#f}. &  \\
\bottomrule
\end{longtable}}

\section{Vectors, arrays, deques, ranges and tables}\label{sec:lib-other-coll}

{\sloppy Operations specific to particular collections.\par}

{\small\frenchspacing
\begin{longtable}{@{}>{\raggedright\arraybackslash}p{0.43\textwidth}>{\raggedright\arraybackslash}p{0.47\textwidth}>{\centering\arraybackslash}p{0.04\textwidth}@{}}
\toprule
Name & Description & \\
\midrule
\endhead
\code{(vector \#rest objects)}\index{vector} & A new simple object vector of the arguments. &  \\
\code{(aref array \#rest indices)}\index{aref} & The element of \code{array} at the indices. &  \\
\code{((setter aref) array \#rest indices-and-new-value)} & Sets an element of an array. &  \\
\code{(dimensions array)}\index{dimensions} & A sequence of the array's dimensions. &  \\
\code{(push deque new-value)}\index{push} & Adds \code{new-value} at the front of the deque. &  \\
\code{(pop deque)}\index{pop} & Removes and answers the first element. &  \\
\code{(push-last deque new-value)}\index{push-last} & Adds \code{new-value} at the end of the deque. &  \\
\code{(pop-last deque)}\index{pop-last} & Removes and answers the last element. &  \\
\code{(range \#key from up-to through by size)}\index{range} & A range from \code{from} (default 0) by \code{by} (default 1), stopping before \code{up-to}, at \code{through}, or after \code{size} elements; unbounded if none is given. &  \\
\code{(remove-key! table key)}\index{remove-key"!} & Removes \code{key} and its element from the table. &  \\
\bottomrule
\end{longtable}}

\section{Conditions}\label{sec:lib-conditions}

{\sloppy Signaling and handling conditions (\dylanbook{chapter 13}).\par}

{\small\frenchspacing
\begin{longtable}{@{}>{\raggedright\arraybackslash}p{0.43\textwidth}>{\raggedright\arraybackslash}p{0.47\textwidth}>{\centering\arraybackslash}p{0.04\textwidth}@{}}
\toprule
Name & Description & \\
\midrule
\endhead
\code{<condition>}\index{<condition>} & The class of all conditions. &  \\
\code{<serious-condition>}\index{<serious-condition>} & Conditions that must be dealt with. &  \\
\code{<error>}\index{<error>} & Errors. &  \\
\code{<simple-error>}\index{<simple-error>} & Errors with a message; init arguments \code{format-string:} and \code{format-arguments:}. &  \\
\code{<type-error>}\index{<type-error>} & A value not of the expected type; init arguments \code{value:} and \code{type:}. &  \\
\code{<warning>}\index{<warning>} & Warnings; their default handler prints them and returns \code{\#f}. &  \\
\code{<simple-warning>}\index{<simple-warning>} & Warnings with a message, like \code{<simple-error>}. &  \\
\code{<restart>}\index{<restart>} & The class of restarts. &  \\
\code{<simple-restart>}\index{<simple-restart>} & A plain restart. &  \\
\code{<abort>}\index{<abort>} & The restart that abandons the current operation. &  \\
\code{(signal condition \#rest arguments)}\index{signal} & Signals a condition (or a \code{<simple-warning>} made from a format string and arguments); answers what a handler returns. &  \\
\code{(error condition \#rest arguments)}\index{error} & Signals an error (or a \code{<simple-error>} made from a format string and arguments); never returns. &  \\
\code{(cerror restart-description condition \#rest arguments)}\index{cerror} & As \code{error}, but with a restart that makes \code{cerror} return \code{\#f}. &  \\
\code{(break \#rest arguments)}\index{break} & Enters the debugger with a message made like \code{signal}'s. The compiler to C, which has no debugger, prints the message and continues. &  \\
\code{(check-type value type)}\index{check-type} & Signals a \code{<type-error>} unless \code{value} is an instance of \code{type}; otherwise answers \code{value}. &  \\
\code{(abort)}\index{abort} & Signals an \code{<abort>}. &  \\
\code{(default-handler condition)}\index{default-handler} & Called when no handler takes a condition: a serious condition enters the debugger (in the compiler to C, is reported and stops the program or thread), a warning is printed, and other conditions answer \code{\#f}. &  \\
\code{(do-handlers function)}\index{do-handlers} & Calls \code{function} with the type, test, function and description of each active handler. &  \\
\code{(restart-query restart)}\index{restart-query} & Asks the user for what a restart needs (in the compiler to C, does nothing and answers \code{\#f}). &  \\
\code{(return-query condition)}\index{return-query} & Asks the user for values to return from \code{signal} (in the compiler to C, answers \code{\#f}). &  \\
\code{(return-allowed? condition)}\index{return-allowed?} & True if a handler may return from this condition. &  \\
\code{(return-description condition)}\index{return-description} & A description of what returning from the condition means. &  \\
\code{(condition-format-string condition)}\index{condition-format-string} & The format string of a simple error, warning or restart. &  \\
\code{(condition-format-arguments condition)}\index{condition-format-arguments} & Its format arguments. &  \\
\code{(type-error-value type-error)}\index{type-error-value} & The value that had the wrong type. &  \\
\code{(type-error-expected-type type-error)}\index{type-error-expected-type} & The type it should have had. &  \\
\bottomrule
\end{longtable}}

\section{Printing and streams}\label{sec:lib-printing}

{\sloppy Output (Chapter~\ref{ch:printing}). The 1992 manual defines none of these; both implementations have all of them.\par}

{\small\frenchspacing
\begin{longtable}{@{}>{\raggedright\arraybackslash}p{0.43\textwidth}>{\raggedright\arraybackslash}p{0.47\textwidth}>{\centering\arraybackslash}p{0.04\textwidth}@{}}
\toprule
Name & Description & \\
\midrule
\endhead
\code{(print object \#key stream verbose)}\index{print} & Prints \code{object} escaped on a line of its own; answers it. &  \\
\code{(write-line object \#key stream)}\index{write-line} & The same as \code{print}. &  \\
\code{(display object \#key stream)}\index{display} & Prints \code{object} plain; answers it. &  \\
\code{(newline \#key stream)}\index{newline} & Ends the line; answers \code{\#f}. &  \\
\code{(pp object \#key stream)}\index{pp} & Prints \code{object} as \code{print} does, laid out over lines as code. &  \\
\code{(write object \#key stream)}\index{write} & Prints \code{object} escaped, with no newline; answers it. &  \\
\code{(print-object object stream escape?)}\index{print-object} & The generic function all printing goes through; add methods to change how a class prints. &  \\
\code{(format stream control-string \#rest arguments)}\index{format} & Formatted output: \code{\textasciitilde a}, \code{\textasciitilde s}, \code{\textasciitilde d}, \code{\textasciitilde\%}, \code{\textasciitilde\&}, \code{\textasciitilde\textasciitilde}. A stream of \code{\#t} is standard output; of \code{\#f}, answer a string. &  \\
\code{<stream>}\index{<stream>} & The class of streams. &  \\
\code{<string-stream>}\index{<string-stream>} & Streams that collect their output in memory. &  \\
\code{*standard-output*}\index{*standard-output*} & The stream to standard output. &  \\
\code{*standard-error*}\index{*standard-error*} & The stream to standard error. &  \\
\code{(stream-contents string-stream)}\index{stream-contents} & A string of everything written to the stream. &  \\
\bottomrule
\end{longtable}}

\section{Threads}\label{sec:lib-threads}

{\sloppy An extension of the compiler to C (Chapter~\ref{ch:extensions}). The special forms \code{with-lock}, \code{define-thread-variable} and \code{dynamic-bind} are in Section~\ref{sec:lib-forms}.\par}

{\small\frenchspacing
\begin{longtable}{@{}>{\raggedright\arraybackslash}p{0.43\textwidth}>{\raggedright\arraybackslash}p{0.47\textwidth}>{\centering\arraybackslash}p{0.04\textwidth}@{}}
\toprule
Name & Description & \\
\midrule
\endhead
\code{<thread>}\index{<thread>} & The class of threads. & \textsc{c} \\
\code{(make-thread function)}\index{make-thread} & Starts a thread that calls \code{function} with no arguments; answers the thread. & \textsc{c} \\
\code{(join-thread thread)}\index{join-thread} & Waits for the thread to end and answers its function's values; signals a \code{<thread-error>} if the thread ended with an unhandled condition. & \textsc{c} \\
\code{(current-thread)}\index{current-thread} & The thread that calls it. & \textsc{c} \\
\code{(thread-yield)}\index{thread-yield} & Lets other threads run. & \textsc{c} \\
\code{<thread-error>}\index{<thread-error>} & The error \code{join-thread} signals for a thread that ended with an unhandled condition. & \textsc{c} \\
\code{(thread-error-thread thread-error)}\index{thread-error-thread} & The thread that failed. & \textsc{c} \\
\code{(thread-error-condition thread-error)}\index{thread-error-condition} & The condition it did not handle. & \textsc{c} \\
\code{<lock>}\index{<lock>} & Mutual-exclusion locks; \code{(make <lock>)}. & \textsc{c} \\
\code{<recursive-lock>}\index{<recursive-lock>} & Locks the holding thread may acquire again. & \textsc{c} \\
\code{(acquire-lock lock)}\index{acquire-lock} & Waits for and takes the lock. & \textsc{c} \\
\code{(release-lock lock)}\index{release-lock} & Releases the lock. & \textsc{c} \\
\code{<condition-variable>}\index{<condition-variable>} & Condition variables; \code{(make <condition-variable>)}. & \textsc{c} \\
\code{(wait-condition-variable condition-variable lock)}\index{wait-condition-variable} & Releases \code{lock}, waits for a notification, and takes \code{lock} again. & \textsc{c} \\
\code{(notify-condition-variable condition-variable)}\index{notify-condition-variable} & Wakes one waiting thread. & \textsc{c} \\
\code{(notify-all-condition-variable condition-variable)}\index{notify-all-condition-variable} & Wakes all waiting threads. & \textsc{c} \\
\code{<atomic-integer>}\index{<atomic-integer>} & Integers updated atomically; \code{(make <atomic-integer> value: n)}. & \textsc{c} \\
\code{(atomic-value atomic-integer)}\index{atomic-value} & The current value. & \textsc{c} \\
\code{(atomic-add! atomic-integer n)}\index{atomic-add"!} & Adds \code{n} atomically and answers the new value. & \textsc{c} \\
\code{(atomic-compare-and-swap! atomic-integer old new)}\index{atomic-compare-and-swap"!} & If the value is \code{old}, makes it \code{new} and answers \code{\#t}; otherwise answers \code{\#f}. & \textsc{c} \\
\bottomrule
\end{longtable}}

\section{Other extensions}\label{sec:lib-other}

{\sloppy Extensions of the library beyond the manual's, in both
implementations unless marked.\par}

{\small\frenchspacing
\begin{longtable}{@{}>{\raggedright\arraybackslash}p{0.43\textwidth}>{\raggedright\arraybackslash}p{0.47\textwidth}>{\centering\arraybackslash}p{0.04\textwidth}@{}}
\toprule
Name & Description & \\
\midrule
\endhead
\code{(current-milliseconds)}\index{current-milliseconds} & The time in milliseconds from a fixed but arbitrary moment, for timing. &  \\
\code{(gensym \#rest prefix)}\index{gensym} & A symbol no program can write, for macros (Chapter~\ref{ch:extensions}); the optional string begins its name. In the interactive implementation, and to macro bodies under the compiler. & \textsc{l} \\
\code{(lisp-variable symbol)}\index{lisp-variable} & The value of the Common Lisp variable named \code{symbol} (Section~\ref{sec:lisp-interop}). & \textsc{l} \\
\code{(lisp-procedure symbol)}\index{lisp-procedure} & A Dylan function that calls the Common Lisp function named \code{symbol}. & \textsc{l} \\
\bottomrule
\end{longtable}}

\section{Special forms and macros}\label{sec:lib-forms}

\index{special form}
These forms are not functions: each has its own rules for which of its
parts are evaluated. They are listed by the shape of a use. The three
marked \textsc{c} are extensions of the compiler to C
(Chapter~\ref{ch:extensions}); \code{define-macro} and
\code{define-for-macros}, with which programs define macros of their
own, are extensions too, in both implementations.




{\small\frenchspacing
\begin{longtable}{@{}>{\raggedright\arraybackslash}p{0.43\textwidth}>{\raggedright\arraybackslash}p{0.47\textwidth}>{\centering\arraybackslash}p{0.04\textwidth}@{}}
\toprule
Name & Description & \\
\midrule
\endhead
\code{(quote object)}\index{quote} & The object itself, unevaluated; also written \code{'object}. &  \\
\code{(begin form \dots)}\index{begin} & Evaluates the forms in order; answers the last one's values. &  \\
\code{(if test consequent alternative)}\index{if} & Evaluates \code{consequent} if \code{test} is true, otherwise \code{alternative}. &  \\
\code{(when test form \dots) (unless test form \dots)}\index{when}\index{unless} & Evaluates the forms if \code{test} is true (false). &  \\
\code{(cond (test form \dots) \dots (else: form \dots))}\index{cond} & Evaluates the forms of the first clause whose test is true. &  \\
\code{(case target ((key \dots) form \dots) \dots (else: form \dots))}\index{case} & Chooses the clause whose keys include the target (compared with \code{id?}). &  \\
\code{(select target test ((match \dots) form \dots) \dots (else: form \dots))}\index{select} & As \code{case}, but the matches are evaluated and compared with the function \code{test}, such as \code{=} or \code{instance?}. &  \\
\code{(and form \dots) (or form \dots)}\index{and}\index{or} & Logical and and or, evaluating only as many forms as needed. &  \\
\code{(bind ((var \dots init) \dots) form \dots)}\index{bind} & Binds local variables (several to the values of one form if more than one is named) and evaluates the forms. &  \\
\code{(set! place new-value)}\index{set"!} & Assigns a variable, or calls the setter of \code{place}: \code{(set! (f x) v)} is \code{((setter f) x v)}. &  \\
\code{(setter name)}\index{setter} & The setter function of \code{name}, \code{(setter name)}. &  \\
\code{(method parameter-list form \dots)}\index{method} & An anonymous method. &  \\
\code{(bind-methods ((name parameter-list form \dots) \dots) form \dots)}\index{bind-methods} & Local methods that may call each other. &  \\
\code{(for ((var init step) \dots) (end-test result \dots) form \dots)}\index{for} & General iteration. &  \\
\code{(for-each ((var collection) \dots) (end-test result \dots) form \dots)}\index{for-each} & Iterates over collections in parallel. &  \\
\code{(dotimes (var count result) form \dots)}\index{dotimes} & Evaluates the forms with \code{var} from 0 below \code{count}. &  \\
\code{(while test form \dots) (until test form \dots)}\index{while}\index{until} & Repeats the forms while (until) \code{test} is true; answers \code{\#f}. &  \\
\code{(bind-exit (var) form \dots)}\index{bind-exit} & Binds \code{var} to an exit procedure that returns from the form. &  \\
\code{(unwind-protect protected-form cleanup-form \dots)}\index{unwind-protect} & Evaluates the cleanup forms however \code{protected-form} is left. &  \\
\code{(handler-bind (type function \#key test description) form \dots)}\index{handler-bind} & Evaluates the forms with a condition handler in effect. &  \\
\code{(handler-case form ((type \#key condition test description) form \dots) \dots)}\index{handler-case} & Evaluates \code{form}; if a condition of a listed type is signaled, exits and evaluates that clause. &  \\
\code{(define name value)}\index{define} & Defines a module variable. &  \\
\code{(define-class name (superclass \dots) slot-spec \dots)}\index{define-class} & Defines a class and its slot accessors; the variable is read-only. &  \\
\code{(define-generic-function name parameter-list)}\index{define-generic-function} & Defines a generic function; the variable is read-only. &  \\
\code{(define-method name parameter-list form \dots)}\index{define-method} & Adds a method to a generic function, creating it if need be. &  \\
\code{(define-slot owner keyword value \dots)}\index{define-slot} & Adds a slot to a class or singleton. &  \\
\code{(define-module name clause \dots)}\index{define-module} & Declares a module (Chapter~\ref{ch:modules}). &  \\
\code{`template ,form ,@form}\index{backquote} & Backquote: a list built from a template, as in Common Lisp. &  \\
\code{(define-macro name parameter-list form \dots)}\index{define-macro} & Defines a macro (Chapter~\ref{ch:extensions}). &  \\
\code{(define-for-macros form \dots)}\index{define-for-macros} & Forms run only while compiling, for macros to use. &  \\
\code{(define-thread-variable name value)}\index{define-thread-variable} & Defines a variable each thread can rebind. & \textsc{c} \\
\code{(dynamic-bind ((name value) \dots) form \dots)}\index{dynamic-bind} & Gives thread variables new values in this thread while the forms run. & \textsc{c} \\
\code{(with-lock (lock) form \dots)}\index{with-lock} & Holds the lock while the forms run. & \textsc{c} \\
\bottomrule
\end{longtable}}

\chapter{Testing}\label{app:testing}
\sloppy

Both implementations come with tests. This appendix describes each
suite: what it tests, how to run it, how long it takes, how to run a
single case, and how to add one. All commands are run from the root of
the source tree.

\begin{notebox}
The expected results of the tests come from the 1992 manual, not from
either implementation. Where the manual gives an example with its
result, the test expects that result; where it specifies behavior
without an example, the test is written from the specification and cites
the chapter and section. Neither implementation is taken as the
reference for the other.
\end{notebox}

\section{The suites at a glance}\label{sec:test-overview}

\begin{center}
\small
\begin{tabular}{>{\raggedright\arraybackslash}p{0.33\textwidth}%
                >{\raggedright\arraybackslash}p{0.58\textwidth}}
\toprule
Command & What it tests \\
\midrule
\cmd{sbcl --script tests/front-end/run.lisp} & the compiler's front end:
  reading, modules, conversion to core forms, diagnostics \\
\cmd{sbcl --script tests/ir/run.lisp} & the IR and its analysis \\
\cmd{sbcl --script tests/codegen/run.lisp} & the compiler's features,
  phase by phase: programs compiled, run, and their output compared \\
\cmd{sbcl --script tests/conformance/run.lisp} & the manual's examples
  and specifications, chapter by chapter \\
\cmd{sbcl --script tests/stress/run.lisp} & 10\textsuperscript{8} tail
  calls; many threads at once \\
\cmd{sh tests/driver/run.sh} & the command line, \cmd{bin/dylanc} \\
\cmd{make -C runtime test} & the C runtime's data representation \\
\code{(dylan:test-examples)} & the interactive implementation, on the
  book's examples \\
\cmd{sbcl --script tests/interpreter/run.lisp} & the compiler's test
  programs, run in the interactive implementation \\
\bottomrule
\end{tabular}
\end{center}

Each command prints one line per failure, beginning \code{FAIL}, and a
summary; the suites of the compiler exit with status 0 when everything
passed and 1 otherwise.

\section{The front end and the IR}\label{sec:test-front-ir}

\file{tests/front-end/cases.lisp} holds 55 cases. Each is a list

\begin{lisp}
(name "dylan source" expected)
\end{lisp}

\noindent where \meta{expected} is either the core forms the source
converts to, as the compiler's \code{unparse} writes them, one per line,
or \code{(:diagnostics "text" ...)}: the diagnostics the source must
produce, in order, each given as a string it must contain. For example:

\begin{lisp}
 (no-complex
  "#c(1 2)"
  (:diagnostics "error: complex numbers are not supported"))
\end{lisp}

After the cases, the runner converts \file{dylan-examples.dyl}, the
book's examples, which must produce no diagnostic but the one error that
its page~68 example is meant to cause.

\file{tests/ir/cases.lisp} holds 18 cases of the form

\begin{lisp}
(name "dylan source" "expected IR" check ...)
\end{lisp}

\noindent The expected IR is the A-normal form of each top-level form;
each \meta{check} names a variable or function and properties it must
have --- whether a variable is captured or assigned, how many references
it has, whether a function escapes, where it is called from and which of
those calls are tail calls, and so on (the file's header lists them).
The IR of \file{dylan-examples.dyl} must then pass the checker of the
IR's invariants.

Neither suite runs a C compiler. Given the argument
\cmd{print}, each shows what every case converts to, which is how to
write the expected text of a new case:

\begin{shell}
$ sbcl --script tests/front-end/run.lisp print
\end{shell}

\section{Programs compiled and run}\label{sec:test-programs}

Three suites --- \file{tests/codegen}, \file{tests/conformance} and
\file{tests/stress} --- consist of whole Dylan programs that are
compiled, built, run, and whose output is compared with what the test
expects. They share one runner, \file{tests/codegen/runner.lisp}.

\begin{description}
\item[\file{tests/codegen/cases*.lisp}] 108 programs testing the
  compiler's features phase by phase: \file{cases.lisp} (calls, tail
  calls, closures, \code{\#rest} and \code{\#key}, multiple values,
  run-time errors), \file{cases-control.lisp} (exits, cleanups,
  conditions), \file{cases-objects.lisp} (classes, slots, dispatch,
  \code{next-method}, reflection), \file{cases-library.lisp},
  \file{cases-threads.lisp}, \file{cases-macros.lisp} and
  \file{cases-optimize.lisp} (type inference, inlining, unboxing,
  escape analysis).
\item[\file{tests/conformance/ch*.lisp}] 237 programs, one file for each
  chapter of the manual from 2 to 18: the manual's examples with the
  results it gives, and its specifications, each test citing its
  section. What a static program cannot do is not tested here; it is
  listed in Appendix~\ref{app:differences}.
\item[\file{tests/stress/cases-*.lisp}] 11 programs that take minutes:
  $10^8$ tail calls through unknown functions, generic functions,
  closures, \code{next-method} and \code{apply}, each in a 256~KB stack;
  and many threads allocating, contending for a bounded buffer,
  interning symbols, signalling conditions, binding thread variables
  and making threads.
\end{description}

\subsection{The case format}\label{sec:test-case-format}

\index{test cases}
A file of cases holds one list of cases, each

\begin{lisp}
(name "dylan source" "expected standard output"
      [:error "text"] [:compile-error "text"] [:compiler-only t])
\end{lisp}

The program is the source string; the strings follow Lisp's rules, so a
double quote inside is written \code{\textbackslash"} and a backslash
\code{\textbackslash\textbackslash}. The output must be exactly the
expected text, newline for newline. With \code{:error}, the program must
instead stop with status~1 and an error message on the standard error
containing the text --- after writing the expected output. With
\code{:compile-error}, the compiler must report an error containing the
text, and the program is not run. \code{:compiler-only t} marks a case
that tests what the interactive implementation does not do by design ---
63-bit integers, a static program's restrictions, the compiler's own
diagnostics, threads --- which \file{tests/interpreter/run.lisp}
(Section~\ref{sec:test-interactive}) skips. From
\file{cases-threads.lisp}:

\begin{lisp}
 (thread-values
  "(define t1 (make-thread (method () (values 1 2))))
   (bind ((a b (join-thread t1))) (print (list a b)))
   (print (join-thread t1))
   (define t2 (make-thread (method () (current-thread))))
   (print (list (id? (join-thread t2) t2) (id? (current-thread) t2)))
   (print (instance? t2 <thread>))"
  "(1 2)
1
(#t #f)
#t
")
\end{lisp}

\noindent and from \file{cases-macros.lisp}, a case that must fail to
compile:

\begin{lisp}
 (macro-not-top-level
  "(define f (method () (define-macro m () 1)))"
  ""
  :compile-error "define-macro must be at top level")
\end{lisp}

Threads run in any order, so a test of threads must print only what does
not depend on the order. Every test program runs with a 256~KB stack
(\cmd{ulimit -s 256}), so a call in tail position that was not compiled
as a jump overflows it quickly and the test fails.

\subsection{Modes}\label{sec:test-modes}

\index{test modes}
Each program is compiled to C once and then built and run in each
\emph{configuration} of a \emph{mode}, named on the command line:

\begin{description}
\item[\cmd{all}] (the default) GCC and Clang, each with \code{musttail}
  and with trampolines: four builds of each program, the first at
  \cmd{-O2} and the others at \cmd{-O0}.
\item[\cmd{quick}] The first of those only: GCC, \code{musttail},
  \cmd{-O2}.
\item[\cmd{whole}] As \cmd{quick}, but with the library compiled with
  each program and tree shaken (\cmd{-{}-whole}), not compiled apart.
\item[\cmd{gc-stress}] GCC, with \code{musttail} and with trampolines,
  each forcing a collection every 1000 allocations
  (\cmd{-{}-gc-stress} with \code{DYLAN\_GC\_STRESS=1000}).
\item[\cmd{address}] AddressSanitizer and UndefinedBehaviorSanitizer,
  GCC and Clang, in debug builds.
\item[\cmd{thread}] ThreadSanitizer, GCC, in a debug build.
\end{description}

Every mode builds with the default collector, the Memory Pool System.
With the environment variable \code{DYLAN\_GC=boehm} set, it builds with
the Boehm-Demers-Weiser collector instead (Section~\ref{sec:dylanc-gc});
the compiler and runtime should pass every suite with either.

\begin{shell}
$ sbcl --script tests/codegen/run.lisp quick
$ sbcl --script tests/conformance/run.lisp            # all
$ sbcl --script tests/stress/run.lisp thread
\end{shell}

The builds go in \file{runtime/build/tests/} (a subdirectory named for
the mode, except in \cmd{all}). The sanitizer and stress modes build
their own variants of the runtime library the first time they need them.

\subsection{Running some of the cases}\label{sec:test-single}

Further arguments name the case files to run, and the cases: any
argument ending in \file{.lisp} is a file to read instead of the suite's
own, and any other (besides the mode) is the name of a case to run,
whichever file holds it.

\begin{shell}
$ sbcl --script tests/codegen/run.lisp quick thread-values
$ sbcl --script tests/codegen/run.lisp quick tests/codegen/cases-macros.lisp
$ sbcl --script tests/conformance/run.lisp quick ch08-true-false
\end{shell}

The summary line counts the programs, the configurations and the
failures, as in \code{1 program in 1 configuration, 0 failures}.

\subsection{How long they take}\label{sec:test-time}

The programs are built and run one after another, so the time grows
with the number of programs and configurations. A build takes one to two
seconds once the compiler and the library's object are cached
(Section~\ref{sec:dylanc-cache}); a configuration whose C compiler and
options are new must first compile the library, which takes some
seconds more. On the 32-thread machine of Chapter~\ref{ch:performance},
the conformance suite in mode \cmd{all} --- 237 programs in four
configurations --- took about a quarter of an hour; \cmd{quick} needs a
fraction of that. The \cmd{address}, \cmd{thread} and \cmd{gc-stress}
modes are much slower, because their programs run slowly, and the stress
suite takes minutes in any mode. The front-end and IR suites run no C
compiler and are quick. Use \cmd{quick} while working, and \cmd{all} and the other
modes before committing a change to the compiler or runtime.

\section{The command line}\label{sec:test-driver}

\cmd{sh tests/driver/run.sh} tests \cmd{bin/dylanc} itself, in a
temporary directory: the default executable name and \cmd{-o}; that the
C is removed unless \cmd{-k} is given, and what \cmd{-k} keeps; the
\code{\#line} directives and \cmd{-{}-no-lines}; \cmd{-C}, and compiling
its C by hand; each \cmd{-O} level, \cmd{-Os} and \cmd{-g};
\cmd{-{}-trampoline}, \cmd{-{}-cc}, \code{CC}, \cmd{-{}-lto} with GCC and
with Clang, and \cmd{-{}-whole}; the library's object cache, and that two
programs' library C is the same text; a program of several modules,
with a file for each; a macro exported from its module, and \cmd{-E};
run-time errors and their exit status; and the exit status of compile
errors and of bad command lines. It prints \code{FAIL} and a description
for each check that fails, and a count of checks and failures.

\section{The runtime}\label{sec:test-runtime}

\cmd{make -C runtime test} builds and runs
\file{runtime/build/test-runtime}, a C program that tests the runtime's
data representation without any Dylan program: the built-in classes as
the compiler emits them (which \file{tests/runtime/generate.lisp}
writes): immediate values, whether each class is a subclass of each
other, how constants print, integer arithmetic at the edges of the
63-bit range and its errors, other numbers, strings and symbols, the
collector, and threads. It runs in moments once built.

\section{The interactive implementation}\label{sec:test-interactive}

\index{test-examples@\code{test-examples}}
The interactive implementation is tested on the examples from the Dylan
book, \file{dylan-examples.dyl}:

\begin{shell}
$ cd src
$ sbcl --eval '(load "load-dylan.lisp")' \
       --eval '(load "full-test.lisp")' \
       --eval '(dylan:test-examples)'
\end{shell}

\code{test-examples} reads the file form by form, echoes each form,
evaluates it, and prints its values --- or, for a form that signals a
condition, \code{;Error:} and the condition's message, after which it
goes on with the next form. It does not compare the results itself. The
expected results are in \file{src/reference/}, the transcripts Thomas
produced from the same file in 1992 (\file{MIT.output} among them);
compare the output with them. Where they and the manual disagree, the
manual decides; the few remaining differences are listed in
\file{src/README}. Run it from the
\file{src} directory, as shown, since \code{test-examples} finds the
examples at \file{../dylan-examples.dyl}; give it another file name as
its argument to run other examples the same way.

\index{tests/interpreter/run.lisp@\file{tests/interpreter/run.lisp}}
The programs the compiler is tested on run in the interactive
implementation too, by \file{tests/interpreter/run.lisp}, which is how
the two implementations are kept in agreement. It takes files of cases in
the form of Section~\ref{sec:test-case-format}, and test names, as the
compiler's runners do; by default all of \file{tests/codegen/cases*.lisp}.
Each program runs in \code{dylan-user}, as a compiled program does, and
what it defined is forgotten before the next; its output must be the
expected output, and the error it signals, if the case expects one, must
have a message containing the expected text --- whether the compiler
would find the error compiling or running makes no difference here, and
the \code{Dylan error:} of a compiled program's report is not expected.
Cases marked \code{:compiler-only} are skipped and counted:

\begin{shell}
$ sbcl --script tests/interpreter/run.lisp
97 programs, 0 failures, 14 compiler-only skipped
$ sbcl --script tests/interpreter/run.lisp tests/codegen/cases-control.lisp
\end{shell}

Every program not so marked must pass in both implementations.

\section{Adding a test}\label{sec:test-adding}

\begin{enumerate}
\item Choose the suite. A feature of the compiler goes in the
  \file{tests/codegen/cases-*.lisp} file for its area; an example or
  specification from the manual goes in the
  \file{tests/conformance/ch}\meta{nn}\file{.lisp} of its chapter, with
  a comment citing the section and saying whether the expected result
  is the manual's example or follows from its text; a test that takes
  long goes in \file{tests/stress/}.
\item Write the case in the format of Section~\ref{sec:test-case-format}:
  a unique name, the program, and the output the manual or the
  specification requires --- not merely what the implementation prints
  today. Print every result with \code{print} or \code{format}, and,
  for threads, print only what does not depend on their order.
\item Run it alone in \cmd{quick} mode until it passes, then in the
  default mode, and in \cmd{gc-stress}, \cmd{address} or \cmd{thread}
  if it exercises memory or threads:
\begin{shell}
$ sbcl --script tests/codegen/run.lisp quick my-new-case
$ sbcl --script tests/codegen/run.lisp my-new-case
\end{shell}
\item For the front end or the IR, add a case to its \file{cases.lisp},
  using the runner's \cmd{print} argument to see what the source
  converts to, and checking that by hand against the manual before
  making it the expected text.
\end{enumerate}

A test that fails because the implementation is wrong is a bug to fix
in the implementation, not in the test; a test is changed only when it
misread the manual.

\fussy


\backmatter
\printindex
\end{document}
